\documentclass[11pt,letterpaper]{article}
\usepackage[T1]{fontenc}
\usepackage[utf8]{inputenc}
\usepackage{lmodern}
\usepackage[margin=1in,headheight=15pt]{geometry}
\usepackage{microtype}
\usepackage{amsmath,amssymb,amsthm,mathtools,mathrsfs}
\usepackage{booktabs,array,longtable}
\usepackage{enumitem}
\usepackage{fancyhdr}
\usepackage{needspace}
\usepackage{xcolor}
\usepackage{xurl}
\usepackage{aliascnt}
\usepackage{hyperref}
\hypersetup{colorlinks=true,linkcolor=blue!45!black,citecolor=blue!45!black,
 urlcolor=blue!45!black,
 pdftitle={Set-Sized Quotients of the Omnific Integers},
 pdfauthor={Merged AI-assisted research report},
 pdfsubject={Omnific integers, universal set-sized quotients, cardinal thresholds}}
\usepackage[nameinlink,noabbrev]{cleveref}
\numberwithin{equation}{section}
\newtheorem{theorem}{Theorem}[section]
\newaliascnt{lemma}{theorem}
\newtheorem{lemma}[lemma]{Lemma}
\aliascntresetthe{lemma}
\newaliascnt{proposition}{theorem}
\newtheorem{proposition}[proposition]{Proposition}
\aliascntresetthe{proposition}
\newaliascnt{corollary}{theorem}
\newtheorem{corollary}[corollary]{Corollary}
\aliascntresetthe{corollary}
\theoremstyle{definition}
\newaliascnt{definition}{theorem}
\newtheorem{definition}[definition]{Definition}
\aliascntresetthe{definition}
\newaliascnt{example}{theorem}
\newtheorem{example}[example]{Example}
\aliascntresetthe{example}
\newaliascnt{question}{theorem}
\newtheorem{question}[question]{Question}
\aliascntresetthe{question}
\theoremstyle{remark}
\newaliascnt{remark}{theorem}
\newtheorem{remark}[remark]{Remark}
\aliascntresetthe{remark}
\theoremstyle{plain}
\newtheorem{maintheorem}{Theorem}
\renewcommand{\themaintheorem}{\Alph{maintheorem}}
\crefname{maintheorem}{Theorem}{Theorems}
\crefname{theorem}{Theorem}{Theorems}
\crefname{lemma}{Lemma}{Lemmas}
\crefname{proposition}{Proposition}{Propositions}
\crefname{corollary}{Corollary}{Corollaries}
\crefname{definition}{Definition}{Definitions}
\crefname{example}{Example}{Examples}
\crefname{question}{Question}{Questions}
\crefname{remark}{Remark}{Remarks}
\crefname{section}{Section}{Sections}
\crefname{subsection}{Section}{Sections}
\crefname{appendix}{Appendix}{Appendices}
\crefname{table}{Table}{Tables}
\crefname{equation}{}{}
\crefname{enumi}{item}{items}
\newcommand{\No}{\mathbf{No}}
\newcommand{\On}{\mathbf{On}}
\newcommand{\Oz}{\mathbf{Oz}}
\newcommand{\SC}{\No[i]}
\newcommand{\GOz}{\Oz[i]}
\newcommand{\Z}{\mathbb Z}
\newcommand{\Q}{\mathbb Q}
\newcommand{\R}{\mathbb R}
\newcommand{\C}{\mathbb C}
\newcommand{\N}{\mathbb N}
\newcommand{\F}{\mathbb F}
\newcommand{\A}{\mathcal A}          % D + Pi_k (positive support side)
\newcommand{\B}{\mathcal B}          % D + m_k  (negative support side)
\newcommand{\mm}{\mathfrak m}        % strictly negative support ideal
\newcommand{\OO}{\mathcal O}         % finite surreal / surcomplex rings
\newcommand{\FF}{\mathscr F}         % Hahn field over the scale subgroup H_a
\newcommand{\EE}{\mathscr E}         % internal algebraically closed field of Q_a
\newcommand{\Qa}{\mathcal Q}         % binomial quotient Oz/(1+omega^a)
\newcommand{\NN}{\mathcal N}         % integral closure (normalization)
\newcommand{\hahn}[2]{#1(\!(X^{#2})\!)}
\newcommand{\Hc}[2]{\mathsf H_{\omega}(#1,#2)}
\newcommand{\ct}{\operatorname{ct}}
\newcommand{\st}{\operatorname{st}}
\newcommand{\supp}{\operatorname{supp}}
\newcommand{\Hom}{\operatorname{Hom}}
\newcommand{\Homone}{\operatorname{Hom}_1}
\newcommand{\End}{\operatorname{End}}
\newcommand{\Ann}{\operatorname{Ann}}
\newcommand{\Frac}{\operatorname{Frac}}
\newcommand{\Der}{\operatorname{Der}}
\newcommand{\Tor}{\operatorname{Tor}}
\newcommand{\Ext}{\operatorname{Ext}}
\newcommand{\Mod}{\operatorname{Mod}}
\newcommand{\Ab}{\mathbf{Ab}}
\newcommand{\pd}{\operatorname{pd}}
\newcommand{\fd}{\operatorname{fd}}
\newcommand{\gen}{\operatorname{gen}}
\newcommand{\cf}{\operatorname{cf}}
\newcommand{\coin}{\operatorname{coin}}
\newcommand{\Span}{\operatorname{span}}
\newcommand{\id}{\operatorname{id}}
\newcommand{\IC}{\operatorname{IC}}
\newcommand{\Int}{\operatorname{Int}}
\newcommand{\Num}{\mathscr N}
\newcommand{\CP}{\operatorname{CP}}
\newcommand{\lcm}{\operatorname{lcm}}
\newcommand{\Jac}{\operatorname{Jac}}
\newcommand{\LC}{\operatorname{LC}}
\newcommand{\Refset}{\operatorname{Ref}_{\mathrm{set}}}
\newcommand{\degw}{\deg_{\omega}}
\newcommand{\Char}{\operatorname{char}}
\DeclareMathOperator{\diag}{diag}
\newcommand{\src}[1]{\textup{\textsf{[#1]}}}
\newcommand{\merge}{\textup{\textsf{[merge]}}}
\setlist[enumerate]{label=(\roman*),leftmargin=2.1em,itemsep=2pt,topsep=4pt}
\setlist[itemize]{leftmargin=1.6em,itemsep=2pt,topsep=4pt}
\setlength{\parskip}{3pt}
\setlength{\parindent}{1.15em}
\emergencystretch=2.5em
\pagestyle{fancy}
\fancyhf{}
\fancyhead[L]{\small Set-sized quotients of the omnific integers}
\fancyhead[R]{\small Research report}
\fancyfoot[C]{\thepage}
\renewcommand{\headrulewidth}{0.3pt}
\setcounter{tocdepth}{2}

\begin{document}
\hypersetup{pageanchor=false}
\begin{titlepage}
\centering
\vspace*{0.3in}
{\large\scshape Surreal algebra\par}
\vspace{0.3in}
{\LARGE\bfseries Set-Sized Quotients of the Omnific Integers\par}
\vspace{0.2in}
{\large The universal constant-term quotient, exact cardinal thresholds,\par
support thresholds, and what survives in large quotients\par}
\vspace{0.3in}
{Merged research report built from eight manuscripts (03, 04, 06, 07, 08, 09, 10, 11)\par}
\vspace{0.08in}
{22 September 2026; merged 23 September 2026\par}
\vspace{0.3in}
\begin{minipage}{0.9\textwidth}
\begin{abstract}
Write an omnific integer uniquely as $n+s$ with $n\in\Z$ and $s$ purely infinite, and let
$\ct(n+s)=n$. Every additive and multiplicative map from $\Oz$ to a set-sized ring, even
a noncommutative, nonreduced or nonunital one, has the form $x\mapsto\ct(x)e$ with $e$
idempotent; every set-sized $\Oz$-module is an abelian group with scalar action through
$\ct$. The same holds for $D+\Pi_K$ over any set field $K$ and unital $D\subseteq K$, for
the Gaussian ring $\Oz[i]$ (with $\Z[i]$), and, on the finite side, for standard part on
the finite surreal and surcomplex rings. Countable supports already force this
collapse while finite supports do not.

Over set-sized exponent groups the theorem becomes a detection threshold. For each of
five explicit set-sized rings (four integer parts inside $\Oz$, of sizes $\kappa^{\aleph_0}$,
$\kappa^{<\kappa}$, $\kappa$ and $\lambda$, and a lexicographic Hahn ring of size
$\kappa^{<\kappa}$), the least size of a ring
or module detecting any nonzero purely infinite element equals the size of the ring;
two of these values, and the exact thresholds of two examples in fixed exponent groups,
are new to this merge. The report also keeps the large-quotient side (a copy of $\No$ and
a Cantor algebra inside $\Oz/(1+\omega^a)$), the integral and complete integral closures
of $\Oz$, the arithmetic of integer-valued polynomials on $\Oz$, and a negative answer to a
posted question on natural ordinal arithmetic. It is an AI-assisted merge of unrefereed
manuscripts; nothing in it is verified in Lean.
\end{abstract}
\end{minipage}
\vfill
\begin{minipage}{0.9\textwidth}\small
\textbf{Status.} Proof-bearing research report assembled from eight AI-assisted
manuscripts written on 22 September 2026. The sources call their main results candidate
original; priority is not certified, no named conjecture is claimed solved, the proofs have
not been refereed, and no statement is formalized. Results added in the merge are marked
\merge; the source of every other statement is tagged \src{NN}.
\end{minipage}
\end{titlepage}
\hypersetup{pageanchor=true}
\begingroup
\small
\setlength{\parskip}{0pt}
\tableofcontents
\endgroup
\clearpage

\section{Introduction}\label{osq:sec:intro}

\subsection{The question and the answer}\label{osq:sub:question}
The omnific integers $\Oz=\Z\oplus\Pi$ form a discretely ordered subring of $\No$ whose
fraction field is all of $\No$. Their purely infinite ideal $\Pi$ is a proper class; each of
its elements has a set-sized normal form, but below any positive exponent there is a proper
class of smaller exponent scales. The eight manuscripts merged here ask the same question:
\emph{what can an ordinary, set-sized ring or module see of an omnific integer?} All eight
give the same answer, by different proofs: exactly its ordinary integer constant term. This
is much stronger than the observation that a proper-class field cannot embed in a set: the
source is a ring, its homomorphisms have large kernels, and the targets may have zero
divisors, nilpotents, noncommutative multiplication and any characteristic (06, 11). Nor is
it a consequence of divisibility of $\Pi$ by every integer, which only disposes of targets
of positive characteristic (04, 07).

The sources then go in different directions: exact cardinal thresholds for set-sized
integer parts (03, 06, 08, 09, 11), the finite side and the support threshold (04, 11), the
large quotients in which a purely infinite element becomes a unit (07), integral closure
(03), polynomial maps (10), and natural ordinal arithmetic (11). This report prints the
common spine once and keeps every result of every source.

\subsection{Main theorems}\label{osq:sub:main}

\begin{maintheorem}[Universal set-sized quotient]\label{osq:main:universal}
Let $K$ be a set field and $D\subseteq K$ a unital subring. Every additive multiplicative map
from $D+\Pi_K$ to a set-sized ring kills $\Pi_K$ and factors uniquely through $\ct$; unital
maps correspond bijectively to unital maps from $D$; every set-sized module is annihilated
by $\Pi_K$. For $\Oz$ the unique unital map to a set-sized ring $S$ is $x\mapsto\ct(x)1_S$, and
unital maps $\GOz\to S$ correspond to $j\in S$ with $j^2=-1$. The kernels admitting set-sized
quotients are exactly the $\ct^{-1}(\mathfrak a)$; for $\Oz$ they are $\Pi$ and $n\Oz$
($n\ge1$). (\Cref{osq:thm:universal,osq:cor:Oz,osq:cor:gaussian,osq:thm:quotients}; all eight
sources and manuscript 05 of the sibling report.)
\end{maintheorem}

\begin{maintheorem}[The finite side and the support threshold]\label{osq:main:support}
Standard part is the universal set-sized quotient of the finite surreal and surcomplex rings,
and more generally of $D+\mm_k$ (04). For every infinite $\lambda$, the rings of normal forms
with supports of size at most $\lambda$ (full surreal exponents) already have universal
set-sized quotient $D$ on both sides, while the finite-support rings admit the augmentation
$\sum c_g\omega^g\mapsto\sum c_g$; a map on the finite-support ring extends to the countable-support
ring if and only if it factors through $\ct$ (the failure for the augmentation is 04's and 11's,
the criterion is 11's). (\Cref{osq:thm:finite,osq:thm:support}.)
\end{maintheorem}

\begin{maintheorem}[Detection threshold equals size]\label{osq:main:thresholds}
For each of five explicit set-sized rings $A$, namely four integer parts inside $\Oz$, of cardinality
$\kappa^{\aleph_0}$ ($\cf\kappa>\aleph_0$; 03), $\kappa^{<\kappa}$ ($\kappa$ regular
uncountable; 06), $\kappa$ (every infinite $\kappa$; 08) and $\lambda$ ($\lambda^{\aleph_0}=\lambda$;
09), and 11's lexicographic Hahn ring of cardinality $\kappa^{<\kappa}$ (coefficient field with at
most $\max(\kappa,2^{\aleph_0})$ elements), the least cardinality of a ring target
or module detecting any nonzero purely infinite element is exactly $|A|$, in ZFC.
(\Cref{osq:thm:thresholds}.) The values for 06's and 11's rings are new here: 06 and 11 prove
exactness only under $\kappa^{<\kappa}=\kappa$, and the unconditional value $\kappa^{<\kappa}$ follows
from 06's own collision lemma (\cref{osq:thm:model06,osq:prop:card11}). The same thresholds hold
for the first and second $\Ext$ groups of the homological package of 06, which extends to all five
rings (\cref{osq:thm:homological,osq:thm:extthreshold}).
\end{maintheorem}

\begin{maintheorem}[The field bound in a fixed exponent group]\label{osq:main:fieldbound}
In a set-sized Hahn integer part, a ring or module detecting an element divisible by $X^{2c}$
has at least as many elements as the field $L_c(F)$ of elements of the ambient field $F$
supported in the scale subgroup $H_c$ (03, generalized through 06's lemma). This bound is exact in 04's example $\Q\oplus_{\mathrm{lex}}\Q$, where it
equals $2^{\aleph_0}$, and in an example with rational coefficients it exceeds by far all the
monomial and weighted counts of 04, 07 and 10, showing that 10's bound need not be attained.
(\Cref{osq:prop:fieldbound,osq:ex:lexQQ,osq:ex:countablecoeff}.)
\end{maintheorem}

\begin{maintheorem}[Large internal fields]\label{osq:main:internal}
If a nonzero purely infinite element becomes a unit in a nonzero ring, the ring contains an
explicitly embedded copy of $\No$. The nonzero quotient $\Oz/(1+\omega^a)$, which has no
nonzero set-sized image, contains an algebraically closed field containing $\SC$, exact
cyclotomic algebras over it, and the Cantor algebra of locally constant functions on $\Z_2$
(07; one proposition of 06). (\Cref{osq:if:thm:transfer,osq:if:thm:binomialkernel,osq:if:thm:cantor}.)
\end{maintheorem}

\begin{maintheorem}[Closures and invisible equations]\label{osq:main:closure}
The complete integral closure of $\Oz$ is $\No$, while its integral closure is a proper,
fine-dense subring with zero conductor that no set generates; its constant slice is the real
algebraic integers, it has no finite quotients, and no nontrivial set-valued valuation of $\No$
contains $\Oz$ (03). For $\gamma>0$ and $m\ge2$ the equation $X^m=\omega^\gamma+1$ has no root in
$\GOz$, yet its image in every set-sized ring has the root $1$ (08).
(\Cref{osq:nm:thm:complete,osq:nm:thm:normalization,osq:nm:thm:valuation,osq:nm:thm:invisible}.)
\end{maintheorem}

\begin{maintheorem}[Polynomial maps]\label{osq:main:polynomial}
The ring of finite-degree surreal polynomials preserving $\Oz^r$ is $\Pi[\mathbf X]\oplus\Int(\Z^r)$
with universal set-sized image $\Int(\Z^r)$; a Newton least-common-multiple condition decides
congruence preservation modulo every omnific integer; the $p$-adic completion is $C(\Z_p^r,\Z_p)$
and most of its characters are not evaluations; rational self-maps of $\Oz$ and $\GOz$ are
polynomial and the bijections are affine (10).
(\Cref{osq:pm:thm:numreal,osq:pm:thm:cp,osq:pm:cor:completion,osq:pm:cor:bijections}.)
\end{maintheorem}

\begin{maintheorem}[Natural ordinal arithmetic]\label{osq:main:ordinal}
The Grothendieck ring of the ordinals under natural operations is not a quotient of $\Oz$; every
unital ring map from $\Oz$ to it has image $\Z$. This answers negatively the quotient question in
MathOverflow question 188430 of Jesse Elliott (11). (\Cref{osq:thm:ordinal}.)
\end{maintheorem}

\subsection{Sources, provenance and where the merge chose}\label{osq:sub:sources}

The report is built from eight manuscripts, all dated 22 September 2026 and placed in commit
\texttt{be06fc8}; none has a reviewer, and none is shipped with this report. Their numbers are
arrival numbers in the batch.

\begin{center}\small
\begin{tabular}{@{}>{\raggedright\arraybackslash}p{0.03\textwidth}>{\raggedright\arraybackslash}p{0.27\textwidth}>{\raggedright\arraybackslash}p{0.08\textwidth}>{\raggedright\arraybackslash}p{0.54\textwidth}@{}}
\toprule
& Manuscript (pages) & Pin & Contribution and placement\\
\midrule
03 & \emph{Cardinal Visibility and Normalization of Omnific Integers} (25) & \texttt{cd80e5e} & field bound (\cref{osq:prop:fieldbound}); model $A^{03}_\kappa$ with threshold $\kappa^{\aleph_0}$ (\cref{osq:sec:model03}); flatness, $\Tor$; integral and complete integral closure, valuations (\cref{osq:sec:closure})\\
04 & \emph{What Set-Sized Algebra Can See of Surreal Arithmetic} (24) & \texttt{cd80e5e} & $D+\Pi_k$ theorem via separation and telescope; finite side (\cref{osq:sec:finite}); support threshold with negative side (\cref{osq:sec:support}); fixed-group bounds; localizations; modules, derivations, algebras\\
06 & \emph{The Universal Set-Sized Quotient of the Omnific Integers} (25) & \texttt{a5c2a97} & \textbf{base}: coefficient-ring theorem for any set field, scaled-field collision lemma, direct module argument; model $A^{06}_\kappa$; homological package (\cref{osq:sec:homological}); matrix representations; polynomial tests; prime-reservoir proposition\\
07 & \emph{Small Quotients and Large Internal Fields of the Omnific Integers} (27) & \texttt{37eefca} & theorem for any set field via weighted certificates; completions; compression $\No\hookrightarrow\FF_a$, transfer, binomial kernel, Cantor algebra (\cref{osq:sec:internal})\\
08 & \emph{What Small Rings Can See of Omnific Integers} (26) & \texttt{37eefca} & Hartogs formulation, nonunital maps; closure formula; invisible equation; class-indexed derivations; criterion; model $A^{08}_\kappa$ for every infinite $\kappa$; cyclic-module lemma\\
09 & \emph{Set-Sized Shadows of Omnific Integers} (20) & \texttt{37eefca} & direct division of $D+\Pi_k$ elements; model $A^{09}_\lambda$, residue fields, non-arithmetic maximal ideals, $\aleph_1$ generators; Gaussian ramification\\
10 & \emph{Set-Sized Shadows of Omnific Arithmetic} (26) & \texttt{a5c2a97} & theorem for any set field incl.\ maps on $\Pi_K$; quantitative bound; endomorphisms; numerical polynomials, $p$-adic completion and characters, lattice rigidity (\cref{osq:sec:polynomial})\\
11 & \emph{What Set-Sized Algebra Can See of the Omnific Integers} (20) & \texttt{37eefca} & pairwise-difference families; localization and presentation calculus; extension criterion; coinitiality theorem; residual targets; ordinal ring (\cref{osq:sec:ordinal})\\
\bottomrule
\end{tabular}
\end{center}

The pins are the repository commits the manuscripts inspected (\texttt{cd80e5e0adc2},
\texttt{a5c2a97df4b5}, \texttt{37eefca1a309}, all of 22 September 2026, Pacific time).
Manuscript~05 of the same batch, \emph{Omnific Integers and Diophantine Rigidity}, is merged
into the sibling report \cite{RepoOdg}. It also proves the unital universal theorem for $\Oz$
(its ``set-target rigidity''), the clearing of set supports and common monomial divisors, and it
asks the ``exact size boundary'' question that 03, 06, 08, 09 and 11 answer; 03 credits the
first three to 05 and 04 credits its finite-quotient background to manuscript~02 of the same
sibling report. Those credits are kept.

\paragraph{Where the merge chose.}
\begin{itemize}
\item \emph{Base.} 06 is the base: its coefficient-ring theorem over any set field (also
proved by 07 and 10) is the most general statement, and its scaled-field collision lemma (with
03's field bound, a special case) is the mechanism that yields the sharp cardinal bounds where
counting monomials does not suffice (03, 06, 11 and the fixed-group examples; 08 and 09 reach
their thresholds by counting monomials). The proof is written in 08's
choice-free Hartogs form, with 08's and 10's nonunital forms and 06's direct module argument.
04's explicit telescope, 09's and 10's direct division, 11's pairwise-difference families and 07's
weighted certificates are kept as second routes (\cref{osq:sub:universalthm}).
\item \emph{Support threshold.} Printed once, in 11's formulation (every infinite support bound,
with the extension criterion) together with 04's negative side (\cref{osq:thm:support}).
\item \emph{Thresholds.} Organized by one theorem (\cref{osq:thm:thresholds}) with the five
constructions of \cref{osq:sec:models}; generator counts are listed per construction and no
general law is stated.
\item \emph{Homological algebra.} 06's package is proved under three hypotheses that all five
constructions satisfy (\cref{osq:thm:homological}).
\item \emph{Foundations.} NBG with choice for sets; global choice, stated by 04, 06, 09 (and 05),
is not used anywhere (\cref{osq:sub:foundations}).
\end{itemize}

\paragraph{Results added in the merge.} Each carries the tag \merge\ and a complete proof:
the field bound for an arbitrary subfield of a fixed-group Hahn field
(\cref{osq:prop:fieldbound}); the exact thresholds $2^{\aleph_0}$ in 04's example and in the
rational-coefficient example (\cref{osq:ex:lexQQ,osq:ex:countablecoeff}), which answer 04's
question for its own example and show that 10's bound need not be attained; the criterion (b$'$) in \cref{osq:thm:criterion};
the exact value $\kappa^{<\kappa}$ of the sizes and thresholds of $A^{06}_\kappa$ and of 11's
example, in ZFC (\cref{osq:lem:card06,osq:thm:model06,osq:prop:card11}), which answers the first
half of 06's question on cardinal arithmetic (08 answers the second half); the organizing
theorem itself (\cref{osq:thm:thresholds}) and the reconciliation of the realized cardinals
(\cref{osq:rem:reconcile}); the homological package and the $\Ext$ thresholds for all five
constructions (\cref{osq:thm:homological,osq:thm:extthreshold}); the negative-side extension
criterion (\cref{osq:thm:support}(iv)); the embedding $\No\hookrightarrow\Frac(\Oz/P)$ for primes
not containing $\Pi$ (\cref{osq:if:rem:composition}); and the explicit elements of the
normalization provided by 08's equation (\cref{osq:nm:rem:rootinN}). Each follows from lemmas
of the sources; none is a new method.

\subsection{Relation to the repository}\label{osq:sub:repository}
The foundations report defines $\Oz=\Pi\oplus\Z$ (\nolinkurl{found:eq:omnific} in
\cite{RepoFound}), and omnific integers appear as kernels and integer parts in the trigonometry,
gamma--zeta, analysis and euclidean-three-space reports; the sibling report \cite{RepoOdg} treats
their Diophantine geometry. The euclidean-three-space report proves the group analogue of the
finite-side theorem for $\mathrm{SO}(3,\No)$ (\cref{osq:sec:finite}). The first-kappa report proves
closedness of bounded-support Hahn fields, which contains the closure step of 06's construction
(06 proves that step itself; \cref{osq:sub:model06}). The
$\kappa$-bounded coefficient rings $\Z+K(\!(G^{<0})\!)_\kappa$ of L'Innocente--Mantova \cite{LM} are
the published precedent for the bounded-support rings of 06 and 11. No statement of this report
has a Lean counterpart; the repository has no omnific module.

\paragraph{Stale statements corrected.} Three sources report a repository code search for
``omnific'' with no hits (06, 08, 11; 09 reports an explicitly incomplete result). At their pins
the repository already defined $\Oz$ at \nolinkurl{found:eq:omnific} and used it in the reports
listed above; what it lacked was a report devoted to omnific arithmetic, and the collection now has
two (this one and \cite{RepoOdg}). 03's companion draft and 04's ``prior manuscript'' were not in the
repository at their pins; they are now manuscripts 05 and 02 of \cite{RepoOdg}. 07's description of
the staging directory (``Gamma--zeta ZIP archives'') was true at its pin and is a snapshot statement.
The shipped audit files keep these statements verbatim; the README records the corrections.

\paragraph{Corrections to the sources.} 03's source lost a relation symbol and prints
``$C$ e $B[i]$'' for $C\ne B[i]$ (\cref{osq:nm:prop:complexdefect}). 09 attributes the separated-scale
division to Proposition~8.2.1 of \cite{LM}; the displayed separated-scale factorization is
Fact~8.0.1 (Gonshor's Theorem~8.6) with its example on p.~46, as 11 cites, while
Proposition~8.2.1 is the truncation criterion for divisibility cited by 08
(\cref{osq:sub:telescope}). 09's hypothesis $\lambda\ge2^{\aleph_0}$ is implied by
$\lambda^{\aleph_0}=\lambda$ (\cref{osq:sub:model09}). 04's remark on its $\Q\oplus_{\mathrm{lex}}\Q$
example is true but misleading (\cref{osq:rem:04misleading}). 06's remark that
$|A^{06}_\kappa/xA^{06}_\kappa|\ge\kappa$ ``under the cardinal-arithmetic hypothesis'' needs no
hypothesis and in fact gives $\kappa^{<\kappa}$ (\cref{osq:sec:homological}). 06, 04 and 09 state
global choice, which is never used (\cref{osq:sub:foundations}). No source contains a false theorem.

\section{Conventions, notation, and the constant-term retraction}
\label{osq:sec:conv}

\subsection{Foundations}\label{osq:sub:foundations}
We work in von Neumann--Bernays--G\"odel class theory (NBG) with the axiom of
choice for sets. Every surreal number, every normal-form support and every
coefficient field or target ring called \emph{set-sized} is a set; the carriers
$\No$, $\Oz$, $\SC$ and most rings of \cref{osq:sec:gaps,osq:sec:universal} are
proper classes. A homomorphism with a class source is a class function obeying
the usual identities, and its restriction to a set is a set.

The sources state different foundations. Manuscript~06 works in G\"odel--Bernays
set theory with global choice, 09 in GBC, and 04 in ``GBC, with ordinary choice
for sets'' (GBC already contains global choice); 07, 08 and 11 work in NBG with
choice for sets, 10 in NBG with choice, and 03 in NBG for the classes and ZFC for
its set-sized constructions; 07 remarks that global choice is unnecessary.
\emph{No argument of any source uses global choice.} The collision step of the
universal theorem (\cref{osq:thm:universal}) is choice-free in the formulation of
08, which we adopt: for a set $S$, its Hartogs ordinal is an ordinal that does not
inject into $S$; it exists in ZF. The same device is used in the rotation-group
theorem \nolinkurl{e3:cut:thm:smallquotient} of \cite{RepoE3}. Choice for sets enters
only where it is visibly needed: maximal ideals of set-sized rings
(\cref{osq:sec:models}), cardinal arithmetic such as
$\kappa^{\aleph_0}=\max(2^{\aleph_0},|[\kappa]^{\le\aleph_0}|)$, and the
existence of residue fields.

\begin{definition}[Set-sized quotients and detection]\label{osq:def:setquotient}
A class ideal $J$ of a class ring $R$ \emph{admits a set-sized quotient} if there is a
surjective class ring homomorphism $R\to S$ onto a set-sized ring with kernel~$J$.
A ring homomorphism $\varphi$ \emph{detects} $x$ if $\varphi(x)\neq0$; a module $M$
detects $x$ if $xM\neq0$.
\end{definition}

This is the convention of 04 (``set-realizable quotient''), 08, 09 (``set-sized
realization''), 10 (``set-sized quotient presentation'') and 11. Manuscripts 06 and 07
alternatively represent a class quotient by Scott codes (the members of a coset of least
set-theoretic rank form a nonempty set); every quotient statement below may be read in
either way, and no set or class whose members are arbitrary proper classes is formed.
Notation such as $\Hom(\Oz,S)$ abbreviates a statement about individual class maps and
their factorizations; after the universal theorem these maps are coded by their
restrictions to a set-sized coefficient ring.

\paragraph{The universe-relative reading.} Manuscripts 04, 10 and 11 add an equivalent
reading: fix a lower universe containing the coefficient fields, targets and supports,
and place the surreal carrier in a higher universe; ``set-sized'' then means small in the
lower universe. The size restriction must be carried in the statement. A set-sized
fragment of surreal arithmetic viewed in a larger universe has its own identity map as a
set-sized target, so it does not satisfy the full-class theorem (11 insists that its main
theorem is relative to one fixed universe, not to fragments viewed in a larger one).

\subsection{Normal forms and the orientation of exponents}\label{osq:sub:orientation}
A nonzero surreal has a unique Conway normal form
$x=\sum_{\gamma\in S}r_\gamma\omega^\gamma$ with nonzero real coefficients and a
\emph{reverse well-ordered set} $S$ of exponents (every nonempty subset has a greatest
element) \cite{Conway,Gonshor,Berarducci}. We use \emph{large monomials}:
$\omega^\gamma$ is infinitely large when $\gamma>0$, and
$\omega^{\gamma+\delta}=\omega^\gamma\omega^\delta$ is Conway's monomial law, not an
exponential. For an abstract ordered abelian group $G$ we write formal monomials $X^g$
with the same orientation, and $\hahn kG$ for the Hahn series over a set field $k$ with
reverse well-ordered set supports in $G$ (for $G=\No$ a proper class of series). The
usual small-monomial notation is recovered by $t^g=X^{-g}$; sources that write $t$
with well-ordered supports reverse all signs. A family is strongly summable if the
union of its supports is reverse well ordered and each exponent receives finitely many
contributions; strong sums are coefficientwise and are \emph{not} limits of partial sums
in the fine order topology.

\emph{Purely infinite} refers to support, not to size or sign: $-\omega$ is purely
infinite, while $\omega-1$ is infinite but not purely infinite. This is the convention of
the foundations report, which defines $\Oz=\Pi\oplus\Z$ with $\Pi$ the purely infinite
normal forms (\nolinkurl{found:eq:omnific} in \cite{RepoFound}); the gamma--zeta report
uses the same symbol \cite{RepoGammaZeta}. The trigonometry report writes the same class
$\Pi$ as the normal forms supported at \emph{negative $t$-exponents}
(\nolinkurl{trigonometry:eq:split} in \cite{RepoTrig}), because it uses $t=\omega^{-1}$; the
two descriptions agree after the sign change $t^{g}=\omega^{-g}$.

\subsection{The rings}\label{osq:sub:rings}
Let $k$ be a set field, $G$ an ordered abelian group (a set, or the class $\No$), and
$D\subseteq k$ a unital subring. Put
\begin{align}
 \Pi_k(G)&=\{f\in\hahn kG:\supp f\subseteq G_{>0}\},&
 \mm_k(G)&=\{f\in\hahn kG:\supp f\subseteq G_{<0}\},\label{osq:eq:ideals}\\
 \A_{D,k}(G)&=D+\Pi_k(G),&
 \B_{D,k}(G)&=D+\mm_k(G).\label{osq:eq:rings}
\end{align}
Both sums are direct as additive groups, and the zero series belongs to both ideals.
We drop $(G)$ when $G=\No$ and write $\Pi=\Pi_\R$. For $k=\R$ and $k=\C$ the
identification $X^g\mapsto\omega^g$ gives $\hahn\R\No=\No$ and $\hahn\C\No=\SC$ (a
union of two reverse well-ordered supports is reverse well ordered, so complexification
is coefficientwise), and
\begin{equation}\label{osq:eq:Ozdef}
 \Oz=\A_{\Z,\R}=\Z\oplus\Pi,\qquad
 \GOz=\Oz+i\Oz=\A_{\Z[i],\C}=\Z[i]\oplus\Pi_\C,\qquad \Pi_\C=\Pi+i\Pi .
\end{equation}
We call $\GOz$ the \emph{Gaussian omnific ring}; it is a subring of $\SC$, not the
surcomplex field, and not the integral closure of $\Z$ in $\SC$ (08, 09). The finite
surreal and surcomplex rings are $\OO_\R=\B_{\R,\R}$ and $\OO_\C=\B_{\C,\C}$, and
$\A_{\R,\R}=\R+\Pi$ is the nonnegative-support ring (03's $P_\R$). For a general set
field $k$ (06, 07, 10) the class field $\hahn k\No$ need not be realized inside the
surreals; nothing below requires such a realization.

\begin{lemma}[Constant term \src{03, 04, 06, 07, 08, 09, 10, 11}]\label{osq:lem:ct}
The map $\ct:\A_{D,k}(G)\to D$ extracting the coefficient of $X^0$ is a surjective
unital ring homomorphism with kernel $\Pi_k(G)$, split by the inclusion of constants.
On $\B_{D,k}(G)$ the same coefficient is a split surjective ring homomorphism with kernel
$\mm_k(G)$; on $\OO_\R$ and $\OO_\C$ it is the standard part, written $\st$.
\end{lemma}
\begin{proof}
All exponents occurring in $\A_{D,k}(G)$ are nonnegative, and a sum of two nonnegative
exponents is zero only when both are zero; so the constant coefficient of a product is the
product of the constant coefficients. Additivity, the kernel and the section are
immediate from uniqueness of the decomposition \eqref{osq:eq:rings}. The negative side is
identical with nonpositive exponents.
\end{proof}

\begin{remark}[What $\ct$ is not]\label{osq:rem:ctnot}
The map $\ct$ is defined on a ring of generally infinite elements; it is not the standard
part, and coefficient extraction at $0$ is not multiplicative on $\No$:
$\omega\cdot\omega^{-1}=1$ although both factors have zero constant coefficient. This is
the warning recorded in \cite{RepoNotation} and \nolinkurl{diff:warn:ct} of
\cite{RepoDiffEq}. Nor is $\ct$ order preserving: $\omega-N>0$ but
$\ct(\omega-N)=-N$ for every ordinary $N\ge1$ (04, 08, 09, 11). On $\A_{D,k}$ the two
maps $\ct$ and $\st$ receive different names because they have different domains
(04), and \eqref{osq:eq:incompatible} below shows that they cannot be combined.
\end{remark}

\begin{lemma}[Degree and units \src{03, 04, 07, 08, 09, 10, 11}]\label{osq:lem:units}
For $x\neq0$ in $\hahn kG$ put $\degw(x)=\max\supp x$. Then
$\degw(xy)=\degw(x)+\degw(y)$ for $x,y\ne0$, so $\A_{D,k}(G)$ and $\B_{D,k}(G)$ are
domains, and
\[
 \A_{D,k}(G)^\times=D^\times,\qquad \Oz^\times=\{\pm1\},\qquad \GOz^\times=\{\pm1,\pm i\}.
\]
In particular every nonconstant element of $\A_{D,k}(G)$, such as $1+\omega^a$ or
$\omega^{\sqrt2}+\omega+1$, is a nonunit.
\end{lemma}
\begin{proof}
The leading coefficient of a product is the product of the leading coefficients, which is
nonzero because $k$ is a field; the degree identity follows. A nonzero element of
$\A_{D,k}(G)$ has nonnegative degree, so $xy=1$ forces both degrees to be $0$, i.e.\
both factors are constants in $D$ with $xy=1$ in $D$. The two unit lists are the units of
$\Z$ and $\Z[i]$.
\end{proof}

Here $\degw$ is the largest growth exponent (04); it is not the ordinal-valued degree of
the order type of a support used in parts of the factorization literature \cite{LM}.

\subsection{Notation and renamings}\label{osq:sub:notation}
The eight manuscripts use five symbols for the purely infinite ideal and three for the
constant term. We use the repository's $\Pi$ and the common $\ct$, and resolve every
collision by renaming in the new material. The renamings, and the tags distinguishing
the different set-sized constructions, are listed below.

\begin{center}\small
\begin{tabular}{@{}>{\raggedright\arraybackslash}p{0.30\textwidth}>{\raggedright\arraybackslash}p{0.40\textwidth}>{\raggedright\arraybackslash}p{0.24\textwidth}@{}}
\toprule
Object & In the sources & Here\\
\midrule
purely infinite ideal & $\mathcal I$ (03, 06, 09), $\mathcal J_k$ (04), $\mathfrak I$ (07, 08, 11), $\mathfrak P$ (10) & $\Pi$, $\Pi_k$, $\Pi_\C$\\
constant term & $c_0$ (06), $\pi$, $\pi_D$ (10), $\varepsilon_D$ (09), $\ct$ (others) & $\ct$\\
coefficientwise constant term on polynomials & $\Pi$ (10), $c_0$ (06) & $\ct_*$\\
$D+\Pi_k$ & $\A_{D,k}$ (04), $A_{D,k}$ (06), $A(D,K)$ (07), $R_D$ (08), $\A_D$ (09), $R(D,K)$ (10), $\mathcal Z_{K,D}$ (11), $A(D,k,G)$ (03) & $\A_{D,k}$, $\A_{D,k}(G)$\\
Gaussian omnific ring & $\mathrm{GOz}$ (06), $\mathbf{Og}$ (09) & $\GOz$\\
surcomplex field & $\No(i)$ (04, 07, 08, 09) & $\SC=\No[i]$ \cite{RepoNotation}\\
Hahn field over $H_a$ & $L_a$ (06), $\mathscr F_a$ (07), $L_h$ (03, set version) & $\FF_a(k)$; $L_h$ for 03's set field\\
localization $\Oz[\omega^{-a}]$ & $L_a$ (07) & $\Oz[\omega^{-a}]$ (never $L_a$)\\
binomial quotient $\Oz/(1+\omega^a)$ & $B_a$ (07) & $\Qa_a$\\
07's compression map $a\Psi$ & $E_a$ & $\Psi_a$\\
integral closure & $B,C,\Int_F(A)$ (03) & $\NN$, $\NN_\C$, $\IC_F(A)$\\
integer-valued polynomials & $\Int(\Z^r)$, $I_r$ (10) & $\Int(\Z^r)$\\
set-sized integer parts & $A_\kappa$ (03, 08), $R_\kappa$ (06), $A_{\lambda,D}$ (09), $R_{\Gamma,<\kappa}$ (11) & $A^{03}_\kappa$, $A^{06}_\kappa$, $A^{08}_\kappa$, $A^{09}_{\lambda,D}$, $R^{11}_{\Gamma,<\kappa}$\\
their ideals & $I_\kappa$, $\mathfrak I_\kappa$, $I_{\lambda,k}$ & $\Pi^{03}_\kappa$, $\Pi^{06}_\kappa$, $\Pi^{08}_\kappa$, $\Pi^{09}_\lambda$\\
small exponents below a support & $\operatorname{Tiny}(S)$ (09), $U(S)$ (10), $G_{\ll h}$ (03) & $U_G(S)$, $H_h(G)$\\
binomial series & $B_m(t)$ (08) & $\mathsf b_m(t)$\\
\bottomrule
\end{tabular}
\end{center}

Throughout, for $a>0$ in an ordered group $G$,
\begin{equation}\label{osq:eq:Ha}
 H_a(G)=\{h\in G: n|h|<a\text{ for every ordinary }n\ge1\},\qquad H_a=H_a(\No),
\end{equation}
and $h\ll a$ means $h>0$ and $h\in H_a(G)$. This is a comparison of \emph{exponents}:
$h\ll a$ does not say that $\omega^h$ is small (it is infinite whenever $h>0$). For a set
$S\subseteq G_{>0}$ we write $U_G(S)=\{a>0: na<s\text{ for all }s\in S,\ n\ge1\}$ (09's
$\operatorname{Tiny}(S)$, 10's $U(S)$). Letters $\kappa,\lambda,\mu,\theta$ denote
cardinals and are restated locally; $\mu=\cf\kappa$ and $\theta=\kappa^{\aleph_0}$ in
the construction of 03.

\section{Support gaps and the purely infinite ideal}\label{osq:sec:gaps}

This section collects the facts that use the full surreal exponent class. All eight
manuscripts prove them, as does manuscript~05 of the sibling report \cite{RepoOdg};
they are background for the new results, not claimed as new. Throughout, $k$ is a set
field and $D\subseteq k$ a unital subring.

\begin{lemma}[Gaps below and above a set \src{03, 04, 06, 07, 08, 09, 10, 11}]\label{osq:lem:gap}
\leavevmode
\begin{enumerate}
\item Every set of surreal numbers has a strict upper and a strict lower bound in $\No$.
\item For every set $S\subseteq\No_{>0}$ there is $h>0$ with $nh<s$ for all $s\in S$
and all ordinary $n\ge1$; in particular $0<h<S$.
\item For such $S$ the class $U_{\No}(S)$ contains the interval $(0,h)$, which is not a
set, and for every ordinal $\Lambda$ the exponents
\begin{equation}\label{osq:eq:scales}
 a_\alpha=\frac{h}{\omega^{\alpha+1}}\qquad(\alpha<\Lambda)
\end{equation}
are distinct elements of $U_{\No}(S)$, each $\ll h$.
\end{enumerate}
\end{lemma}
\begin{proof}
For a set $T$, the cuts $\{T\mid\varnothing\}$ and $\{\varnothing\mid T\}$ are set cuts
and give (i). For (ii), apply the set cut $\{0\mid T\}$ to the set
$T=\{s/n:s\in S,\ n\ge1\}$ of positive surreals; for $S=\varnothing$ take $h=1$. For
(iii), $(0,h)$ is not a set: a positive lower bound of it would have a positive half still
in the interval. The denominators $\omega^{\alpha+1}$ are the Conway monomials of the
ordinals $\alpha+1$ in the \emph{exponent} field; they strictly increase with $\alpha$ and
exceed every ordinary integer, so the $a_\alpha$ are distinct, positive and satisfy
$na_\alpha<h$.
\end{proof}

Only the existence of a separator is used, not its simplicity or birthday (08). The
sources use several equivalent families in place of \eqref{osq:eq:scales}:
$a/\omega^{\alpha+1}$ (04, 05, 06, 08, 09, 03), $a/(\omega+\alpha+1)$ with the
field sum of embedded ordinals (07), $\delta\omega^{-(\alpha+1)}$ (11) and
$\epsilon\xi/\kappa$ (10). Each is a set whenever $\alpha$ ranges over a set.

\begin{example}[Surreal gaps where real gaps fail \src{03, 06, 08, 10, 11}]\label{osq:ex:gap}
The series $f=\sum_{n\ge1}\omega^{1/n}$ lies in $\Pi$; its support has no positive
\emph{real} lower bound but has positive surreal lower bounds, for instance $1/\omega$.
With $a=2/\omega$ and $b=1/\omega$ one has $na<1/m$ for all ordinary $m,n$ (10). In the
fixed exponent group $\R$ the same series is purely infinite but is not divisible by any
$X^h$, $h>0$, inside the nonnegative-support ring (03); this is where the passage from
monomials to the whole ideal fails in a fixed group (\cref{osq:sec:fixed}).
\end{example}

\begin{proposition}[Common monomial divisors \src{03, 04, 06, 07, 08, 09, 10, 11}]\label{osq:prop:common}
Let $A=\A_{D,k}$ and $\Pi_k$ its purely infinite ideal.
\begin{enumerate}
\item Every set-sized family $\mathcal F\subseteq\Pi_k$ satisfies
$\mathcal F\subseteq X^a\Pi_k$ for some $a>0$.
\item Each $f\in\Pi_k$ is a single product of two elements of $\Pi_k$; hence
$\Pi_k^2=\Pi_k$, and $\Pi_k=\bigcup_{a>0}X^aA$ is a directed union of principal ideals.
\item $\Pi_k$ is not generated, as an ideal of $A$, by any set of its elements.
\item $d\Pi_k=\Pi_k$ for every $0\neq d\in D$; in particular $n\Pi_k=\Pi_k$ for every
nonzero ordinary integer $n$.
\end{enumerate}
\end{proposition}
\begin{proof}
(i) The union of the supports of the members of $\mathcal F$ is a set of positive surreals;
choose $0<a<$ that union by \cref{osq:lem:gap}. Translating each support by $-a$ leaves it
positive and reverse well ordered, so $X^{-a}f\in\Pi_k$.
(ii) Apply (i) to $\{f\}$: $f=X^a(X^{-a}f)$ with both factors in $\Pi_k$. If $0<b<a$ then
$X^aA\subseteq X^bA$, so the union is directed.
(iii) If a set generated $\Pi_k$, (i) would place the generated ideal inside $X^aA$; but
$X^{a/2}\in\Pi_k$ is not in $X^aA$, since its unique quotient by $X^a$ in the ambient field
is $X^{-a/2}$, which has a negative exponent.
(iv) Divide every coefficient by $d$, which is invertible in $k$.
\end{proof}

Idempotence and divisibility do not by themselves prove the universal theorem: an
idempotent, rationally divisible ideal can act nontrivially on infinite set-sized
modules, and a homomorphism to an infinite characteristic-zero ring need not kill an
arbitrarily divisible additive subgroup (04, 06, 07, 08). In positive characteristic the
last item already suffices (\cref{osq:rem:poschar}).

\begin{equation}\label{osq:eq:nOz}
 n\Oz=n\Z+\Pi=\ct^{-1}(n\Z)\quad(n\ge1),\qquad
 d\,\GOz=d\Z[i]+\Pi_\C\quad(0\ne d\in\Z[i]).
\end{equation}
Indeed $n\Pi=\Pi$ and $d\Pi_\C=\Pi_\C$ by \cref{osq:prop:common}(iv). Consequently
\begin{equation}\label{osq:eq:intersection}
 \Pi=\bigcap_{n\ge1}n\Oz=\bigcap_{p\text{ prime}}p\Oz,\qquad
 \Pi_\C=\bigcap_{n\ge1}n\GOz,
\end{equation}
since a nonzero ordinary integer (or Gaussian integer) is not divisible by every $n$ or
every prime (08, 10, 11). For a general $D$ the intersection formula can fail: for $D=\Q$
every element of $\A_{\Q,\R}$ is divisible by every ordinary positive integer (11).

\begin{proposition}[Clearing and fraction fields \src{03, 04, 06, 07, 08, 09, 11}]\label{osq:prop:fractions}
\leavevmode
\begin{enumerate}
\item For every set $E\subseteq\hahn k\No$ there is $h>0$ with $X^hE\subseteq\Pi_k$. For
$E\subseteq\No$ or $E\subseteq\SC$ this reads $\omega^hE\subseteq\Pi$, respectively
$\omega^hE\subseteq\Pi_\C$.
\item $\Frac(\A_{D,k})=\hahn k\No$; in particular $\Frac(\Oz)=\No$ and
$\Frac(\GOz)=\SC$.
\item Inverting the positive monomials gives the whole field: if
$\mathcal U=\{X^a:a>0\}\cup\{1\}$ then $\mathcal U^{-1}\A_{D,k}=\hahn k\No$.
\item Neither $\No$ nor $\SC$, nor any localization of $\A_{D,k}$ in which some nonzero
element of $\Pi_k$ becomes a unit, admits a unital homomorphism to a nonzero set-sized
ring or a nonzero set-sized unital module.
\end{enumerate}
\end{proposition}
\begin{proof}
(i) The union of the supports of the members of $E$ is a set; by \cref{osq:lem:gap}(i)
choose $h>0$ larger than $-g$ for each exponent $g$ in it. Then every exponent of $X^hx$,
$x\in E$, is positive. (ii) and (iii) For $x\in\hahn k\No$ apply (i) to $\{x\}$:
$x=(X^hx)/X^h$ with numerator and denominator in $\A_{D,k}$; the reverse inclusions hold
because $\A_{D,k}$ is a subring of the field. (iv) is proved after the universal theorem
(\cref{osq:cor:nofield}).
\end{proof}

Part (i) is manuscript~05's ``Clearing supports'' theorem, now in the sibling report
\cite{RepoOdg}; 03 credits it to 05 and adds the algebra-wide formulation of
\cref{osq:sec:closure}. Clearing one element is weaker than clearing all powers of it
with one denominator, which \cref{osq:sec:closure} also supplies (03).

\section{The collision mechanism}\label{osq:sec:collision}

Every proof of the universal theorem combines two finite multiplicative identities with a
cardinal comparison. This section states the identities once, in the forms used by the
different sources. The first (\cref{osq:lem:collision}, from 06, and in a special form
from 03) gives a uniform proof of the sharp cardinal bounds of
\cref{osq:sec:fixed,osq:sec:models}, including those that counting monomials cannot reach;
the others are genuinely different routes and are kept as such.

\subsection{Scaled-field collision}\label{osq:sub:scaled}

\begin{lemma}[Scaled-field collision \src{06; 03}]\label{osq:lem:collision}
Let $A$ be a commutative subring of a field $F$, let $L\subseteq F$ be a subfield and let
$0\neq m\in F$ satisfy $mL\subseteq A$.
\begin{enumerate}
\item If $\varphi:A\to S$ is a ring homomorphism, $S$ any ring (not necessarily
commutative, unital or reduced), and $\varphi(m^2)\neq0$, then $x\mapsto\varphi(mx)$ is
injective on $L$.
\item If $M$ is a unital $A$-module, $v\in M$ and $m^2v\neq0$, then $x\mapsto(mx)v$ is
injective on $L$.
\end{enumerate}
Consequently, if $|L|>|S|$ then $\varphi(m^2)=0$, and if $|L|>|M|$ then $m^2M=0$.
\end{lemma}
\begin{proof}
Since $1\in L$, $m\in A$. For distinct $x,y\in L$ the element $m/(x-y)$ lies in
$mL\subseteq A$, and in $A$
\begin{equation}\label{osq:eq:collision}
 (mx-my)\cdot\frac{m}{x-y}=m^2 .
\end{equation}
If $\varphi(mx)=\varphi(my)$, applying $\varphi$ gives $\varphi(m^2)=0$; no cancellation
in $S$ is used. For the module, apply $m/(x-y)$ to $(mx)v=(my)v$; commutativity of $A$
turns the result into $m^2v=0$.
\end{proof}

The proof of (i) uses only that $A$ is closed under addition and multiplication and contains
$mL$; so it applies verbatim when $A$ is a nonunital subring such as $\Pi_K$, which is how it is
used in \cref{osq:thm:universal}(iii). The module version bounds $|M|$ itself. Deriving it through the ring $\End_\Z(M)$ would
only bound $|\End_\Z(M)|$, which can be much larger than $|M|$ (06, 03, 09); this matters
for every set-sized threshold below. In 03 the lemma appears as the identity
$\bigl(X^{h/2}(f-f')\bigr)\bigl(X^{h/2}/(f-f')\bigr)=X^h$ for $f\neq f'$ in a
smaller-scale field (03's ``smaller-scale field obstruction''), i.e.\ the case
$m=X^{h/2}$ of \eqref{osq:eq:collision}.

\begin{lemma}[The scale subgroup and its Hahn field \src{03, 06, 07}]\label{osq:lem:Ha}
Let $G$ be an ordered abelian group (a set or $\No$) and $a\in G_{>0}$.
\begin{enumerate}
\item $H_a(G)$ is a convex subgroup of $G$, divisible if $G$ is, and
$H_{c}(G)=H_{2c}(G)$ for every $c\in G_{>0}$.
\item $\hahn k{H_a(G)}$ is a subfield of $\hahn kG$, and $X^c\,\hahn k{H_a(G)}\subseteq
\Pi_k(G)$ for every $c\in G$ with $2c\ge a$; for instance $c=a$, or $c=a/2$ when
$a/2\in G$.
\item For $G=\No$ the field $\FF_a(k)=\hahn k{H_a}$ is a proper class: it contains the
distinct monomials $X^{a_\alpha}$ of \eqref{osq:eq:scales} for every ordinal $\alpha$.
\end{enumerate}
\end{lemma}
\begin{proof}
(i) If $h,h'\in H_a(G)$ then $2n|h|<a$ and $2n|h'|<a$ give $n|h+h'|<a$; negation and
convexity are immediate, and if $G$ is divisible then $h/m$ satisfies the same inequalities.
For every $c>0$, $n|h|<c$ for all $n$ is equivalent to $n|h|<2c$ for all $n$.
(ii) Supports of sums and products of series supported in a subgroup stay in it, and the
Hahn inverse of $f=cX^g(1+u)$, $\supp u<0$, is $c^{-1}X^{-g}\sum_{n\ge0}(-u)^n$, whose
support lies in the subgroup generated by $\supp f$. If $2c\ge a$ and $h\in H_a(G)$ then
$2|h|<a\le 2c$, so $c+h>0$; hence $X^cf$ has positive support for
$f\in\hahn k{H_a(G)}$.
(iii) Each $a_\alpha\ll a$ lies in $H_a$.
\end{proof}

The field $\FF_a=\FF_a(\R)$ is 07's $\mathscr F_a$ and 06's $L_a$ with $k=\R$; it is
real closed and contains a copy of all of $\No$ (\cref{osq:if:prop:closed,osq:if:thm:compression}).

\subsection{The geometric telescope}\label{osq:sub:telescope}

\begin{lemma}[Bounded-exponent telescope \src{03, 04, 05, 06, 07, 08, 09, 10, 11}]\label{osq:lem:telescope}
Let $G$ be an ordered abelian group, and let $0<b<a$ and $a\ll c$ in $G$. Put $d=a-b$.
Then
\begin{equation}\label{osq:eq:telescope}
 q_{c;a,b}=\sum_{n\ge0}X^{c-a-nd}\ \in\ \Pi_k(G),
 \qquad (X^a-X^b)\,q_{c;a,b}=X^c .
\end{equation}
The support of $q_{c;a,b}$ is countable and strictly decreasing. For every ordinary
$N\ge0$,
\begin{equation}\label{osq:eq:finitetelescope}
 (X^a-X^b)\sum_{n=0}^{N}X^{c-a-nd}=X^c-X^{c-(N+1)d},
\end{equation}
and the remainder monomial $X^{c-(N+1)d}$ is still infinitely large.
\end{lemma}
\begin{proof}
For each ordinary $n\ge0$, $a+nd\le(n+1)a<c$, so every exponent is positive; the
exponents strictly decrease with $n$, so every nonempty set of them has a largest element.
Multiplication by $X^a$ gives the exponents $c-nd$ and multiplication by $X^b$ gives
$c-(n+1)d$; at $c$ the coefficient is $1$ and at every $c-md$, $m\ge1$, it is $1-1=0$.
Each coefficient receives at most two contributions, so this is a Hahn identity. The finite
identity is the same computation truncated; its remainder exponent
$c-(N+1)d\ge c-(N+1)a>0$.
\end{proof}

\begin{remark}[No limit is taken and no sum is transported]\label{osq:rem:nolimit}
The remainders in \eqref{osq:eq:finitetelescope} are positive monomials; they do not tend
to $0$ in the order, valuation or fine topology, not even at the ordinary scale~$1$. The
infinite identity is a coefficient identity, not a limit of \eqref{osq:eq:finitetelescope}
(04, 06, 07, 08, 10, 03). A target homomorphism is never applied term by term to
$q_{c;a,b}$: the series is first a single element of the source ring, and only the finite
product identity is mapped (06, 08, 09, 11). This is why no continuity, valuation or
strong-additivity hypothesis on homomorphisms appears anywhere. It is also why a ring of
finite-support series behaves differently (\cref{osq:sec:support}): there
$q_{c;a,b}$ is absent and the nonzero remainder of \eqref{osq:eq:finitetelescope} cannot
be removed.
\end{remark}

\begin{example}[A witness for $\omega$ \src{04, 10}]\label{osq:ex:omegatelescope}
With $c=1$, $b=1/\omega$ and $a=2/\omega$,
$\bigl(\omega^{2/\omega}-\omega^{1/\omega}\bigr)\sum_{n\ge0}\omega^{1-(n+2)/\omega}=\omega$,
while the finite partial products leave the nonzero remainders
$\omega^{1-(N+1)/\omega}$ (04). In the lexicographic group $\Z\times\Z$ with $c=(1,0)$,
$a=(0,2)$, $b=(0,1)$ the quotient is $\sum_{n\ge0}X^{(1,-n-2)}$ and the $N$-th remainder is
$X^{(1,-N-1)}$ (10); under a compatible monomial realization each remainder is still
infinite.
\end{example}

\paragraph{Precedents.} The countable geometric-series product with a monomial result is
already in L'Innocente--Mantova \cite[Remark~3.4.4]{LM}, used there to show that
$\deg\circ\operatorname{ot}$ is not multiplicative when the value group does not embed in
$\R$; 04 names it a direct precursor. The factorization of a series whose support meets two
distinct nonzero Archimedean classes is \cite[Fact~8.0.1]{LM} (Gonshor's
\cite[Thm.~8.6]{Gonshor}), with the geometric factorization displayed on p.~46 of the arXiv
version; 11 cites it as the precedent of the division below. Manuscript~09 credited
``Proposition~8.2.1 on separated-scale division''; that proposition is the truncation
criterion ($b$ divides $c$ if and only if $b$ divides a truncation $T_\Delta(c)$), whose
proof divides the higher-scale part of $c$ by $b$, and 08 cites it as the precedent for
higher-scale divisibility. We credit all three places; none of the sources claims the
geometric mechanism as new.

\begin{lemma}[Weighted certificate \src{07}]\label{osq:lem:weighted}
Let $K$ be a set field, $r,s\in K^\times$, $a>0$ and $\delta>\eta>0$ with $\delta\ll a$.
Then
\begin{equation}\label{osq:eq:weighted}
 \frac{X^a}{rX^\delta-sX^\eta}
 =\frac1r\sum_{k\ge0}\Bigl(\frac sr\Bigr)^kX^{a-\delta-k(\delta-\eta)}\in\Pi_K ,
\end{equation}
and if $\delta=\eta$ and $r\ne s$ the quotient is $(r-s)^{-1}X^{a-\delta}$.
\end{lemma}
\begin{proof}
The exponents are those of \cref{osq:lem:telescope}; multiplying by $rX^\delta-sX^\eta$
gives, at exponent $a-k(\delta-\eta)$ with $k\ge1$, the two contributions
$r\cdot r^{-1}(s/r)^k$ and $-s\cdot r^{-1}(s/r)^{k-1}$, which cancel.
\end{proof}

\subsection{Division of an arbitrary purely infinite element}\label{osq:sub:division}

The telescope divides a monomial. Manuscripts 09, 10 and 11 divide an arbitrary element
directly, which removes the separate reduction to monomials.

\begin{lemma}[Subordinate division \src{09, 10, 11}]\label{osq:lem:division}
Let $G$ be an ordered abelian group, $0\ne s=\sum_{g\in S}c_gX^g\in\Pi_k(G)$ and
$a>b>0$ in $U_G(S)$. Put $d=a-b$. Then
\begin{equation}\label{osq:eq:division}
 \frac{s}{X^a-X^b}=\sum_{g\in S}\sum_{n\ge0}c_g\,X^{g-a-nd}\ \in\ \Pi_k(G),
\end{equation}
so $s=(X^a-X^b)q$ with $q\in\Pi_k(G)$. The support of $q$ lies in
$\{g-a-nd:g\in S,\ n\in\N\}$, of cardinality at most $\max(|S|,\aleph_0)$; it is
countable when $S$ is.
\end{lemma}
\begin{proof}
In the ambient Hahn field $(X^a-X^b)^{-1}=X^{-a}(1-X^{-d})^{-1}=X^{-a}\sum_{n\ge0}X^{-nd}$.
The product with $s$ is a Hahn product: the sum of two reverse well-ordered sets is reverse
well ordered and each exponent has finitely many representations (11 proves this two-support
lemma in full; it is the standard Neumann support lemma). Every exponent that can occur
satisfies $g-a-nd\ge g-(n+1)a>0$ because $(n+1)a<g$; cancellation cannot create new
exponents. Hence the quotient lies in $\Pi_k(G)$.
\end{proof}

The inverse of $X^a-X^b$ is not in the ring; only its product with the larger-scale $s$ is
(11).

\subsection{Separation principles}\label{osq:sub:separation}

\begin{lemma}[Separation \src{04, 09, 11; 03}]\label{osq:lem:separation}
Let $R$ be a ring (possibly a proper class), $e\in R$, and $(x_t)_{t\in T}$ a set-indexed
family with $e\in(x_s-x_t)R$ for all $s\ne t$.
\begin{enumerate}
\item If $\varphi:R\to S$ is a ring homomorphism with $\varphi(e)\ne0$, then
$t\mapsto\varphi(x_t)$ is injective; hence $|T|\le|S|$ if $S$ is a set.
\item If $R$ is commutative, $M$ an $R$-module, $v\in M$ and $ev\ne0$, then
$t\mapsto x_tv$ is injective.
\item If for every cardinal $\lambda$ such a family with $|T|=\lambda$ exists, then every
homomorphism from $R$ to a set-sized ring kills $e$, and $e$ annihilates every set-sized
$R$-module.
\end{enumerate}
\end{lemma}
\begin{proof}
If $\varphi(x_s)=\varphi(x_t)$ with $e=(x_s-x_t)q$, then
$\varphi(e)=(\varphi(x_s)-\varphi(x_t))\varphi(q)=0$; neither commutativity of $S$ nor
preservation of $1$ is used. For (ii), $ev=q(x_s-x_t)v=0$. For (iii), take $\lambda$
larger than the target, or than the Hartogs ordinal of the target, and apply the
pigeonhole principle to a set.
\end{proof}

Item (i) is 04's separation lemma, (ii) is 09's orbit separation and 03's exponent-family
obstruction, and (iii) is 11's collision--divisibility principle.

\begin{theorem}[Divisibility-separated families \src{11}]\label{osq:thm:packing}
Let $\mathcal S\subseteq\Pi_k$ be a set and $\Lambda$ an ordinal. There are distinct
monomials $(u_\alpha)_{\alpha<\Lambda}$ in $\A_{D,k}$ with
\[
 s\in(u_\alpha-u_\beta)\Pi_k\qquad(s\in\mathcal S,\ \alpha\ne\beta).
\]
\end{theorem}
\begin{proof}
Choose $\delta>0$ strictly below the union of the supports of the members of
$\mathcal S$ (\cref{osq:lem:gap}), put $a_\alpha=\delta\omega^{-(\alpha+1)}$ and
$u_\alpha=X^{a_\alpha}$. For $\alpha<\beta$, $0<a_\beta<a_\alpha\le\delta/\omega$, so
$na_\alpha<\delta$ for every $n$, and \cref{osq:lem:division} applies to each $s$;
reversing the indices only changes the sign of the quotient.
\end{proof}

The coefficients of the $u_\alpha$ are $1$; no independence of the coefficients of $s$ is
used, and cancellations inside $s$ or the quotient are harmless (11).

\section{The universal set-sized quotient}\label{osq:sec:universal}

\subsection{The theorem}\label{osq:sub:universalthm}

\begin{theorem}[Universal set-sized quotient \src{04, 06, 07, 08, 09, 10, 11; 03, 05}]\label{osq:thm:universal}
Let $K$ be a set field, $D\subseteq K$ a unital subring, and $A=\A_{D,K}=D+\Pi_K$ with the
full surreal exponent class.
\begin{enumerate}
\item Let $S$ be a set-sized associative ring, not necessarily unital or commutative, and
let $\varphi:A\to S$ be an additive and multiplicative class map, not necessarily
preserving~$1$. Then $\varphi(\Pi_K)=0$, and $\varphi=\psi\circ\ct$ for a unique ring
homomorphism $\psi:D\to S$, namely $\psi=\varphi|_D$.
\item For every set-sized unital ring $S$, restriction to $D$ and composition with $\ct$
are mutually inverse bijections, natural in $S$,
\[
 \Homone(A,S)\ \longleftrightarrow\ \Homone(D,S).
\]
\item Every additive and multiplicative map from the nonunital ring $\Pi_K$ to a set-sized
ring is zero.
\item Every set-sized unital $A$-module $M$ satisfies $\Pi_KM=0$, and its action is
$f\cdot v=\ct(f)v$.
\end{enumerate}
\end{theorem}

\begin{proof}
We follow 06's scaled-field argument in 08's choice-free Hartogs form.

\emph{Positive monomials are killed.} Let $\varphi:A\to S$ be as in (i), or an additive
multiplicative map defined only on $\Pi_K$ as in (iii). Fix $g>0$, put $m=X^{g/2}$ and
$L=\FF_g(K)=\hahn K{H_g}$. By \cref{osq:lem:Ha}(ii) with $c=g/2$, $mL\subseteq\Pi_K$;
also $m/(x-y)\in mL$ for distinct $x,y\in L$, so the identity \eqref{osq:eq:collision}
takes place inside $\Pi_K$. Let $\lambda$ be the Hartogs ordinal of the set $S$, and put
$b_\alpha=g/\omega^{\alpha+1}$ for $\alpha<\lambda$; these are distinct elements of $H_g$
(\cref{osq:lem:gap}(iii)), so $\alpha\mapsto X^{b_\alpha}$ is an injection
$\lambda\to L$. If $\varphi(X^g)=\varphi(m^2)\ne0$, \cref{osq:lem:collision}(i) makes
$x\mapsto\varphi(mx)$ injective on $L$, and composing gives an injection $\lambda\to S$,
which is impossible. Hence $\varphi(X^g)=0$ for every $g>0$.

\emph{The whole ideal is killed.} For $f\in\Pi_K$, \cref{osq:prop:common} gives
$f=X^ah$ with $a>0$ and $h\in\Pi_K$, so $\varphi(f)=\varphi(X^a)\varphi(h)=0$. This proves
(iii), and for (i) it gives $\varphi(d+f)=\varphi(d)$ for $d\in D$, $f\in\Pi_K$. Thus
$\varphi=\varphi|_D\circ\ct$, and $\psi=\varphi|_D$ is forced because $\ct$ is surjective.

\emph{The correspondence.} If $\varphi$ is unital, so is $\psi$. Conversely, for a unital
$\psi:D\to S$ the composite $\psi\circ\ct$ is a unital homomorphism by \cref{osq:lem:ct},
and its restriction to $D$ is $\psi$. Post-composition with a map of targets preserves both
constructions, which proves (ii).

\emph{Modules.} Let $M$ be a set-sized unital $A$-module and $v\in M$. If $X^gv\ne0$,
\cref{osq:lem:collision}(ii) injects $L$, hence the Hartogs ordinal of $M$, into $M$;
so $X^gM=0$ for every $g>0$, and the factorization $f=X^ah$ gives $\Pi_KM=0$. Writing
$f=\ct(f)+(f-\ct(f))$ gives $f\cdot v=\ct(f)v$.
\end{proof}

\paragraph{Sources and routes.} The statement is proved in the following generality by the
sources: any set field $K$ and unital $D\subseteq K$ (06, 07, 10); $k\in\{\R,\C\}$ and any
unital $D\subseteq k$ (04, 09, 11); $\Oz$, $\GOz$ and complex coefficients with any unital
$D\subseteq\C$ (08). Nonunital maps are treated by 08 (as a theorem), 06 (as a remark)
and 10 (including maps defined on $\Pi_K$ alone). Noncommutative targets are allowed in
all eight sources. Manuscript 05 of the sibling report \cite{RepoOdg} proves the unital
case for $\Oz$ with the telescope and a cardinal $\kappa>|A|$, and 03 reproves that case
(credited to 05) through its exponent-family obstruction. The proofs differ, and the
following are genuinely different routes, kept here in brief.
\begin{itemize}
\item \emph{Explicit telescope} (04, 05, 06 appendix, 08). If $\varphi(X^c)\ne0$, then
for distinct $a_\alpha,a_\beta\ll c$ the telescope \eqref{osq:eq:telescope} puts $X^c$ in
$(X^{a_\alpha}-X^{a_\beta})A$, so \cref{osq:lem:separation}(i) makes
$\alpha\mapsto\varphi(X^{a_\alpha})$ injective, which is impossible for $\alpha$ ranging over a
set larger than $S$; \cref{osq:prop:common} then kills all of $\Pi_K$. This is manuscript 05's
proof, which the sibling report \cite{RepoOdg} quotes and refers here for; 05 chooses a cardinal
$\kappa>|S|$ with global choice, where the Hartogs ordinal suffices. The route uses only countable
supports, which is why it survives in the support-bounded rings of \cref{osq:sec:support}.
\item \emph{Direct division of an arbitrary element} (09, 10, 11). If $\varphi(s)\ne0$ for
$s\in\Pi_K$, \cref{osq:lem:division} gives $s\in(X^a-X^b)\Pi_K$ for all $a>b$ in
$U_{\No}(\supp s)$, and \cref{osq:thm:packing} supplies such families of every set size, so
\cref{osq:lem:separation}(iii) applies to $s$ itself; no separate reduction to monomials is
needed. 10 phrases this as an injection $U(\supp s)\to S$, $a\mapsto\varphi(X^a)$, valid also
for nonunital maps on $\Pi_K$; 11 phrases it as divisibility-separated families.
\item \emph{Weighted certificates} (07). With $\delta_\alpha=a/(\omega+\alpha+1)$,
$\alpha<\kappa^+$ for a cardinal $\kappa\ge|S|$, pairwise differences of the monomials
$X^{\delta_\alpha}$ divide $X^a$, and \eqref{osq:eq:weighted} even separates all
$rX^\delta$, $r\in K^\times$.
\item \emph{Modules through the endomorphism ring} (04, 07, 08, 10, 11). A set-sized
module gives a unital ring map $A\to\End_\Z(M)$ into a set-sized, generally
noncommutative ring, and (ii) applies. This is valid for the full class but bounds only
$|\End_\Z(M)|$; the direct argument above (06, 09, 03) and 08's cyclic-module argument
(\cref{osq:lem:cyclic}) are the ones that give set-sized thresholds.
\end{itemize}

\subsection{Consequences for rings and representations}\label{osq:sub:consequences}

\begin{corollary}[Omnific integers \src{03, 04, 05, 06, 07, 08, 09, 10, 11}]\label{osq:cor:Oz}
For every set-sized unital ring $S$ the only unital homomorphism $\Oz\to S$ is
$x\mapsto\ct(x)1_S$. Additive multiplicative maps $\Oz\to S$ into a set-sized ring
correspond bijectively to idempotents $e\in S$, via $x\mapsto\ct(x)e$.
\end{corollary}
\begin{proof}
There is one unital map $\Z\to S$. A ring map $\psi:\Z\to S$ is determined by
$e=\psi(1)$, which satisfies $e^2=e$, and every idempotent defines one. Apply
\cref{osq:thm:universal}(i),(ii).
\end{proof}

\begin{corollary}[Gaussian omnific integers \src{04, 06, 07, 08, 09, 10, 11}]\label{osq:cor:gaussian}
For a set-sized unital ring $S$, unital maps $\GOz\to S$ correspond bijectively to
elements $j\in S$ with $j^2=-1_S$, via
$x+iy\mapsto\ct(x)1_S+\ct(y)j$ $(x,y\in\Oz)$. Every such map discards all
positive-exponent complex coefficients, not only their real parts.
\end{corollary}
\begin{proof}
Apply \cref{osq:thm:universal} with $(D,K)=(\Z[i],\C)$. A unital map from
$\Z[i]=\Z[T]/(T^2+1)$ is a choice of a root of $T^2+1$; in a noncommutative $S$ the
subring generated by $1_S$ and $j$ is commutative, so evaluation is valid.
\end{proof}

The image of $i$ is part of the data: these are ring maps, not maps of complex algebras.
For $(D,K)=(k,k)$ with $k\in\{\R,\C\}$, every unital map from $\A_{\R,\R}$ or $\A_{\C,\C}$
to a set-sized ring factors through the constant coefficient in $k$, and if the maps are
required to fix $k$, the constant term is the unique such map to $k$ (04). For any unital
$D\subseteq\C$, the ring of complex-coefficient normal forms with positive support apart
from a constant in $D$ has universal set-sized quotient $D$ (08); this is the case
$K=\C$ of the theorem.

\begin{corollary}[What all small observations identify \src{04}]\label{osq:cor:observations}
For $x,y\in\A_{D,k}$, $\ct(x)=\ct(y)$ if and only if $\varphi(x)=\varphi(y)$ for every
unital homomorphism $\varphi$ from $\A_{D,k}$ to a set-sized ring. Thus $\Pi_k$ is the
common kernel of all set-sized ring observations.
\end{corollary}
\begin{proof}
The forward direction is the theorem; for the converse use the target $D$ and the map
$\ct$.
\end{proof}

This does not say that the elements of $\Pi_k$ are zero: $\omega\ne0$ in $\Oz$, yet every
set-target ring homomorphism sends it to $0$ (04).

\begin{corollary}[No faithful set-sized representation \src{04, 06, 07, 08, 11}]\label{osq:cor:faithful}
No nonzero element of $\Pi_k$ acts nontrivially on a set-sized unital $\A_{D,k}$-module; in
particular $\A_{D,k}$ has no faithful set-sized module. For every set-sized field $F$ and
set-sized $F$-vector space $V$ (of any dimension), every unital ring map
$\Oz\to\End_F(V)$ is $x\mapsto\ct(x)\id_V$, and every unital ring map $\GOz\to\End_F(V)$
is $x+iy\mapsto\ct(x)\id_V+\ct(y)J$ with $J^2=-\id_V$.
\end{corollary}

\begin{proposition}[Finite matrix representations \src{06, 11}]\label{osq:prop:matrices}
Let $n\ge1$.
\begin{enumerate}
\item The only unital map $\Oz\to M_n(F)$, $F$ a set-sized field, is $x\mapsto\ct(x)I_n$.
\item Up to conjugacy in $\mathrm{GL}_n(\C)$, the unital maps $\GOz\to M_n(\C)$ are
indexed by $p\in\{0,\dots,n\}$, with $i\mapsto\diag(iI_p,-iI_{n-p})$.
\item A unital map $\GOz\to M_n(\R)$ exists exactly when $n$ is even, and then all such
maps are conjugate to the one sending $i$ to $n/2$ diagonal blocks
$\bigl(\begin{smallmatrix}0&-1\\1&0\end{smallmatrix}\bigr)$.
\end{enumerate}
\end{proposition}
\begin{proof}
(i) is \cref{osq:cor:Oz}. By \cref{osq:cor:gaussian} a map in (ii) or (iii) is a matrix
$J$ with $J^2=-I_n$. Over $\C$ its minimal polynomial divides $(T-i)(T+i)$, which has
distinct roots, so $J$ is diagonalizable with eigenvalues $\pm i$, and the multiplicity
$p$ of $i$ determines the conjugacy class. Over $\R$, $J$ makes $\R^n$ a complex vector
space by $(u+iv)w=uw+vJw$, so $n$ is even, and a complex basis gives the block form;
conversely the blocks square to $-I_n$.
\end{proof}

No norm bound, analytic hypothesis or finite-dimensionality is used to kill the purely
infinite part; finite dimension enters only the classification of the possible images of
$i$ (06). For an infinite-dimensional set-sized space the reduction to one operator $J$ with
$J^2=-1$ remains valid.

\paragraph{Algebras.} A set-sized associative $\A_{D,k}$-algebra is a set-sized ring with a
unital map from $\A_{D,k}$ to its center; by the theorem this map factors through $D$, so
set-sized $\A_{D,k}$-algebras and $D$-algebras coincide, with the same underlying rings and
morphisms; noncentral structure maps are covered by the noncommutative form (04).
Enlarging the set-sized algebra in which a calculation takes place therefore never yields
a faithful algebraic model of the full ring, although such calculations can be faithful on
a chosen set-sized subring (04).

\begin{corollary}[No ordered specialization, no small retract, rigid endomorphisms \src{08, 10, 11; 05}]\label{osq:cor:rigid}
\leavevmode
\begin{enumerate}
\item There is no order-preserving unital map from $\Oz$ to a nonzero ordered set-sized
ring (11; 09 records the relevant value $\ct(\omega-1)=-1$).
\item A set-sized unital subring $A\subseteq\Oz$ is the image of a ring retraction
$\Oz\to A$ if and only if $A=\Z$ (08).
\item Every unital class ring endomorphism $\sigma$ of $\Oz$ satisfies
$\ct\circ\sigma=\ct$; hence every automorphism preserves $\Pi$. Every automorphism of
$\GOz$ induces the identity or complex conjugation on $\Z[i]$ (10; the first assertion is
also 05's, see \cite{RepoOdg}).
\end{enumerate}
\end{corollary}
\begin{proof}
(i) The unique map sends the positive element $\omega-1$ to $-1<0$.
(ii) A retraction is $f\mapsto\ct(f)1_A$ by \cref{osq:cor:Oz}, whose image is $\Z$; and
$\ct$ is a retraction onto $\Z$.
(iii) $\ct\circ\sigma$ is a unital map to the set $\Z$, hence equals $\ct$. For an
automorphism apply this to $\sigma$ and $\sigma^{-1}$. In the Gaussian case $\sigma$ and
$\sigma^{-1}$ induce mutually inverse ring automorphisms of $\Z[i]$ fixing $\Z$, and $i$
goes to a root of $T^2+1$ in $\Z[i]$, i.e.\ to $\pm i$.
\end{proof}

Passing to a small workspace is therefore not a homomorphic truncation of the full ring;
the workspace may embed without being a retract (08). The endomorphism statement does not
say that every endomorphism is the identity: $\Z\hookrightarrow\Oz$ composed with $\ct$ is
a noninjective endomorphism.

\begin{corollary}[Fields and unit-forcing localizations \src{04, 06, 09, 11}]\label{osq:cor:nofield}
Let $0\ne u\in\Pi_k$. The class rings $\hahn k\No$ (in particular $\No$ and $\SC$) and
$\A_{D,k}[u^{-1}]$ admit no unital homomorphism to a nonzero set-sized ring and no nonzero
set-sized unital module.
\end{corollary}
\begin{proof}
A unital map from either ring to a set-sized $S$ restricts to $\A_{D,k}$ and so kills $u$;
but $u$ is invertible in the source, so $1_S=\varphi(u)\varphi(u^{-1})=0$. A nonzero
set-sized module would give such a map to the nonzero set-sized ring $\End_\Z(M)$.
\end{proof}

This is not merely the observation that a proper-class field cannot inject into a set: the
localizations need not be fields (04). For example $\Oz[\omega^{-1}]$ is a nonzero class
domain, not a field (\cref{osq:prop:locsupport}), with no nonzero set-sized module. This
completes the proof of \cref{osq:prop:fractions}(iv).

\begin{remark}[Positive characteristic and additive observations]\label{osq:rem:poschar}
If the target has positive characteristic $n$, the conclusion needs no cardinal argument:
$f=n(f/n)$ with $f/n\in\Pi_k$ gives $\varphi(f)=n\varphi(f/n)=0$ (04, 07); more generally
an ideal of $\A_{D,k}$ containing a nonzero element of $D$ contains $\Pi_k$ (07). The
cardinal argument is needed exactly for characteristic-zero targets such as $\Q$, $\R$ or
$\End_\Z(M)$ for infinite torsion-free $M$ (11). The theorem concerns multiplicative
maps: for a fixed $\gamma>0$ the coefficient extraction $x\mapsto[\omega^\gamma]x$ is a
nonzero additive map $\Oz\to\R$ detecting $\omega^\gamma$, but it is not multiplicative,
since $\omega^{\gamma/2}$ has coefficient $0$ at $\gamma$ while its square has coefficient
$1$ (04, 07, 09, 10, 11). Valuations also detect nonconstant parts, but not by unital ring
maps into set-sized rings (04).
\end{remark}

\begin{remark}[What the theorem does not say]\label{osq:rem:notsay}
The theorem is a universal property relative to set-sized targets, not an isomorphism
$\Oz\cong\Z$. The ideal $\Pi$ is nonzero, idempotent, not nilpotent and not generated by
any set, and $\Oz$ is a domain (09, 08). The regular module $\Oz$ and the ideal $\Pi$ are
proper-class modules on which $\Pi$ acts nontrivially, so the module correspondence is not
a Morita equivalence of unrestricted module theories (06). The identity of $\Oz$ and the
inclusion $\Oz\hookrightarrow\No$ are faithful class-valued maps outside the theorem (09).
Individual monomials can be represented faithfully in small rings, for instance in
polynomial rings or Hahn subfields, and every finite calculation lies in a set-sized
subring with a faithful regular representation; the obstruction concerns one global
homomorphism defined on \emph{all} of the class ring, including the collision witnesses at
all lower scales (04, 07, 08). In the language of 11, the set-sized reflection is
$\Refset(\Oz)=\Z$ and $\Refset(\GOz)=\Z[i]$: a surjection $q:R\to D$ onto a set-sized ring
is a universal set-sized quotient if every map from $R$ to a set-sized ring factors uniquely
through $q$; this refers to the explicit universal property, not to a category of all class
rings with set-sized hom-sets.
\end{remark}

\begin{theorem}[Residually set-sized targets \src{11}]\label{osq:thm:residual}
Call a class ring $B$ \emph{residually set-sized} if for each nonzero $b\in B$ there are a
set-sized ring $A_b$ and a unital map $\pi_b:B\to A_b$ with $\pi_b(b)\ne0$ (the quantifier
is elementwise; no set of detecting maps is formed). Every unital map $h:\A_{D,k}\to B$
into a residually set-sized ring kills $\Pi_k$ and factors uniquely through $D$. For a
set-sized commutative ring $E$ and any class $X$ of indeterminates, the ring $E[X]$ of
finite polynomials is residually set-sized.
\end{theorem}
\begin{proof}
If $h(s)\ne0$ for some $s\in\Pi_k$, then $\pi\circ h$ detects $s$ for a detecting map
$\pi$ with set-sized target, contrary to \cref{osq:thm:universal}. A nonzero polynomial
$f\in E[X]$ uses a finite set $Y\subseteq X$ of variables; sending the other variables to
$0$ gives a map $E[X]\to E[Y]$ fixing $f$, with set-sized target.
\end{proof}

\section{Quotients, closures, completions and localizations}\label{osq:sec:quotients}

Throughout, $A=\A_{D,k}$ with $k$ a set field and $D\subseteq k$ a unital subring,
and $\Pi_k=\ker\ct$.

\subsection{Every quotient and its small image}\label{osq:sub:allquotients}

\begin{theorem}[Small reflection of an arbitrary quotient \src{06, 07, 09, 11}]\label{osq:thm:reflection}
Let $J$ be a class ideal of $A$. Then $\ct(J)$ is an ideal of $D$, and for every
set-sized unital ring $S$
\[
 \Homone(A/J,S)\cong\Homone\bigl(D/\ct(J),S\bigr),
\]
the universal map being $[x]\mapsto\ct(x)\bmod\ct(J)$. Its kernel is $(\Pi_k+J)/J$, and
$\ct^{-1}(\ct(J))=\Pi_k+J$.
\end{theorem}
\begin{proof}
$\ct$ is surjective, so the image of an ideal is an ideal. A map from $A/J$ is a map
$\varphi:A\to S$ vanishing on $J$; by \cref{osq:thm:universal},
$\varphi=\psi\circ\ct$ for a unique $\psi$, and $\varphi(J)=0$ if and only if
$\psi(\ct(J))=0$. If $\ct(a)=\ct(b)$ with $b\in J$, then $a-b\in\Pi_k$, which gives the
kernel and the preimage.
\end{proof}

Thus a quotient can be a nonzero ring while all of its set-sized images are zero; a
quotient presentation of a proper-class ring is not automatically a set-sized ring (09).
``Reflection'' is only a name for the displayed mapping property; no category of all proper
classes is used (07).

\begin{theorem}[Ideals with set-sized quotient \src{04, 06, 07, 08, 09, 10, 11}]\label{osq:thm:quotients}
For a class ideal $J$ of $A$ the following are equivalent:
\begin{enumerate}
\item $J$ admits a set-sized quotient (\cref{osq:def:setquotient});
\item $\Pi_k\subseteq J$;
\item $J=\ct^{-1}(\mathfrak a)=\Pi_k+\mathfrak a$ for an ideal $\mathfrak a$ of $D$.
\end{enumerate}
Then $\mathfrak a=J\cap D=\ct(J)$ is unique and $A/J\cong D/\mathfrak a$. For $\Oz$ the
kernels are exactly
\[
 \Pi\ (\text{quotient }\Z)\qquad\text{and}\qquad n\Oz=n\Z+\Pi\ (\text{quotient }\Z/n\Z,\ n\ge1),
\]
the case $n=1$ giving the zero ring. For $\GOz$ they are $\Pi_\C$ (quotient $\Z[i]$) and
$d\GOz=d\Z[i]+\Pi_\C$ for $0\ne d\in\Z[i]$, with finite quotient $\Z[i]/(d)$ of $|d|^2$
elements; associates give the same ideal.
\end{theorem}
\begin{proof}
(i)$\Rightarrow$(ii) by \cref{osq:thm:universal} applied to the quotient map.
(ii)$\Rightarrow$(iii): if $\Pi_k\subseteq J$ then $d+f\in J$ ($f\in\Pi_k$) if and only if
$d\in J\cap D$. (iii)$\Rightarrow$(i): $A\to D\to D/\mathfrak a$ realizes the quotient.
For $\Oz$ the ideals of $\Z$ are $n\Z$, $n\ge0$; $n=0$ gives $\Pi$ and $n\ge1$ gives
$n\Oz$ by \eqref{osq:eq:nOz}. For $\GOz$, $\Z[i]$ is a principal ideal domain (division
with remainder by rounding the real and imaginary parts of a complex quotient), and
\eqref{osq:eq:nOz} gives $d\GOz=d\Z[i]+\Pi_\C$. Multiplication by $d=b+ci$ on
$\Z[i]\cong\Z^2$ has matrix $\bigl(\begin{smallmatrix}b&-c\\c&b\end{smallmatrix}\bigr)$
of determinant $b^2+c^2=|d|^2$, the index of its image; more generally a nonzero ideal
containing $d$ contains the positive integer $d\bar d$, so it has finite index without any
appeal to unique factorization (07, 06, 08).
\end{proof}

The distinction between $n=0$ and $n\ge1$ matters: the characteristic-zero kernel is
$\Pi$, not $0\Oz$; one must not set $n=0$ in $n\Oz=n\Z+\Pi$ (04, 06, 07, 08, 09, 10, 11).
For $D=k$, the only proper kernel with a set-sized quotient is $\Pi_k$, with quotient $k$
(04).

\begin{corollary}[Small primes, maximal ideals and fields \src{06, 07, 08, 10, 11}]\label{osq:cor:smallprimes}
The prime class ideals of $\Oz$ with set-sized quotient are $\Pi$ and $p\Oz$ ($p$ an
ordinary prime), with quotient domains $\Z$ and $\F_p$ and fraction fields $\Q$ and $\F_p$;
the maximal ones are the $p\Oz$. A unital map from $\Oz$ to a set-sized field has kernel
$\Pi$ in characteristic $0$ (its image is $\Z$, not a field) and $p\Oz$ in characteristic
$p$; the only nonzero set-sized field quotients are the $\F_p$. For $\GOz$ replace these by
$\Pi_\C$ and the ideals generated by Gaussian primes, with quotient domains $\Z[i]$ and the
corresponding finite fields.
\end{corollary}

In particular $\Pi$ is prime but not maximal (it lies in $p\Oz$). These results classify
only ideals with set-sized quotients: $(0)$ is prime in $\Oz$ with a proper-class quotient,
and no Zorn argument for a proper class of ideals is used to assert that proper class
ideals lie in maximal class ideals; a maximal ideal outside this list, if one is given,
has a proper-class residue field (07, 08, 11).

\begin{proposition}[Principal quotients \src{04}]\label{osq:prop:principal}
Let $f\in A$ be nonconstant and $\beta=\degw(f)>0$. For distinct $\gamma,\delta\in(0,\beta)$,
$X^\gamma-X^\delta\notin fA$; hence $A/(f)$ has sets of pairwise distinct residues of every
set cardinality and admits no set-sized quotient. Consequently $(f)$ admits a set-sized
quotient exactly when $f\in D\setminus\{0\}$, and then $fA=fD+\Pi_k$ and
$A/(f)\cong D/fD$. Among the ideals with set-sized quotient, $\Pi_k$ is the only one not
generated by a set; every other one is generated by a subset of $D$, and for $\Oz$ these are
the principal ideals $n\Oz$, $n\ge1$.
\end{proposition}
\begin{proof}
The difference has degree $\max(\gamma,\delta)<\beta$, while a nonzero multiple $fq$,
$q\in A$, has degree $\beta+\degw(q)\ge\beta$ (\cref{osq:lem:units}); the interval
$(0,\beta)$ contains sets of every cardinality (\cref{osq:lem:gap}). For $0\ne f\in D$,
$f\Pi_k=\Pi_k$ (\cref{osq:prop:common}); the zero ideal would inject the class $A$ into a
set. For the last assertion write the ideal as $\Pi_k+\mathfrak a$; if
$\mathfrak a=0$ use \cref{osq:prop:common}(iii); otherwise some $0\ne d\in\mathfrak a$ gives
$d\Pi_k=\Pi_k$, so the ideal generated by the set $\mathfrak a$ already contains $\Pi_k$.
\end{proof}

\begin{theorem}[Quotient-size alternative \src{08}]\label{osq:thm:idealsize}
Let $J$ be a class ideal of $\Oz$ with $\Pi\not\subseteq J$. For every set cardinal $\mu$
there is a set of at least $\mu$ elements of $\Oz$ pairwise incongruent modulo $J$.
\end{theorem}
\begin{proof}
Choose $f\in\Pi\setminus J$ and write $f=\omega^ch$ with $c>0$, $h\in\Pi$; then
$\omega^c\notin J$. For an ordinal $\lambda$ of cardinality at least $\mu$ use the
exponents $c/\omega^{\beta+1}$, $\beta<\lambda$: if two of the monomials were congruent
modulo $J$, the telescope \eqref{osq:eq:telescope} would put $\omega^c$ in $J$.
\end{proof}

A statement that a quotient ``is not small'' always has this concrete meaning: for every
set cardinal it contains a larger set of distinct elements; no ``cardinality of $\No$'' is
used (07).

\subsection{Finite congruences, closures and invisible primes}\label{osq:sub:closures}

For a class ideal $J$ of $\Oz$ define its finite-congruence closure
$\overline J^{\,\mathrm{fin}}=\bigcap_{n\ge1}(J+n\Oz)$; the witness in $J$ may depend on
$n$ (08).

\begin{theorem}[Closure of every ideal \src{08, 09}]\label{osq:thm:closure}
If $\ct(J)=d\Z$ with $d\ge0$, then $\overline J^{\,\mathrm{fin}}=J+\Pi=\ct^{-1}(d\Z)$. In
particular $J$ is dense for finite congruences exactly when $\ct(J)=\Z$, and the closure of
$0$ is $\Pi$.
\end{theorem}
\begin{proof}
$\ct(J)$ is an ideal of $\Z$, and $J+n\Oz=\ct^{-1}(d\Z+n\Z)=\Pi+\gcd(d,n)\Z$ by
\eqref{osq:eq:nOz}. If $d>0$ the term $n=d$ is $d\Z$ and all terms contain it; if $d=0$ the
intersection is $\bigcap_nn\Z=0$.
\end{proof}

\begin{corollary}[Dense proper principal ideals \src{08, 09}]\label{osq:cor:dense}
If $u\in\Oz$ is nonconstant with $\ct(u)=\pm1$, then $(u)$ is proper and dense for finite
congruences, and every unital map from $\Oz$ to a nonzero set-sized ring sends $u$ to
$\pm1$. Every nonconstant irreducible $p\in\Oz$ has $\ct(p)=\pm1$; hence $(p)$ is proper,
profinitely dense, and $\Oz/(p)$ has no nonzero set-sized ring image.
\end{corollary}
\begin{proof}
$u$ is a nonunit by \cref{osq:lem:units}; apply \cref{osq:thm:closure} and
\cref{osq:cor:Oz}. For an irreducible $p=c+f$, $0\ne f\in\Pi$: if $c=0$ then
$p=2(p/2)$, and if $|c|>1$ then $p=c(p/c)$ with $p/c=1+f/c$ nonconstant, both
factorizations into nonunits; so $c=\pm1$. A map killing $p$ would kill $\pm1$.
\end{proof}

The converse of the second statement is false: $\omega^2-1=(\omega-1)(\omega+1)$ has
constant term $-1$ and is reducible. The constant term is a complete invariant for
set-sized ring observations, not for factorization inside $\Oz$ (09).

\begin{corollary}[A prime quotient invisible to all set-sized rings \src{06, 08, 09, 11}]\label{osq:cor:primeexample}
Let $q=\omega^{\sqrt2}+\omega+1$. Then $\Oz/(q)$ is a nonzero proper-class integral domain
of characteristic $0$; it admits no unital map to a nonzero set-sized ring and no nonzero
set-sized unital module.
\end{corollary}
\begin{proof}
L'Innocente and Mantova proved that $q$ is prime in $\Oz$ \cite[Theorem~B]{LM}, answering
Gonshor's question \cite[p.~117]{Gonshor}; this primality is imported, not reproved. So
$(q)$ is a proper prime ideal and $\Oz/(q)$ is a nonzero domain. Since $\ct(q)=1$,
\cref{osq:thm:reflection} gives the zero ring as its universal set-sized image. On a
set-sized module, inflated to $\Oz$, $q$ acts both as $\ct(q)=1$ and as $0$, so the module is
zero. If $\Oz/(q)$ were a set, its identity map would contradict the first assertion. No
nonzero integer $n$ lies in $(q)$: $n=qa$ would force $a=n/q$, a nonzero infinitesimal, which
is not omnific; so $\Z\hookrightarrow\Oz/(q)$.
\end{proof}

This is no contradiction with the maximal-ideal theorem for set-sized rings: the source
quotient is not a set, and the corollary does not say that a particular maximal class ideal
fails to exist (06, 09). The same obstruction needs no primality.

\begin{proposition}[Unit-constant ideals \src{06, 07}]\label{osq:prop:unitconstant}
If a class ideal $J$ of $\Oz$ contains an element with constant term $1$, then $\Oz/J$
has no nonzero set-sized unital module and no unital map to a nonzero set-sized ring; if
$J$ is proper, $\Oz/J$ is a nonzero proper-class ring. More generally, if $\ct(f)=c$, the
small reflection of $\Oz/(f)$ is $\Z/c\Z$.
\end{proposition}
\begin{proof}
The element acts as $1$ by \cref{osq:thm:universal} and as $0$ by the quotient relation.
For a principal ideal $\ct((f))=c\Z$; apply \cref{osq:thm:reflection}.
\end{proof}

For example $\omega+1$ and $1+\omega^a$ ($a>0$) are nonunits, so $\Oz/(1+\omega^a)$ is a
nonzero class ring with no nonzero set-sized image; \cref{osq:sec:internal} shows how much
it contains. The benefit of $q$ in \cref{osq:cor:primeexample} is that the quotient is a
domain (06). Under any map to a set-sized ring, $\omega^a$ becomes $0$ and $1+\omega^a$
becomes $1$, which is the arithmetic reason (07).

\begin{theorem}[Arithmetic completions \src{04, 06, 07, 08, 09, 10, 11}]\label{osq:thm:completions}
All finite quotient rings of $\Oz$ are the $\Z/n\Z$, and
\[
 \widehat\Oz:=\varprojlim_{n\ge1}\Oz/n\Oz\cong\widehat\Z,\qquad
 \varprojlim_{r\ge1}\Oz/p^r\Oz\cong\Z_p\quad(p\text{ prime}).
\]
The canonical maps are $\ct$ followed by the usual injections; their kernel is $\Pi$ and
their image is the diagonal copy of $\Z$. For the Gaussian ring,
$\widehat{\GOz}\cong\widehat{\Z[i]}=\widehat\Z[T]/(T^2+1)$ and
$\varprojlim_n\GOz/\ell^n\GOz\cong\Z_\ell[i]=\Z_\ell[T]/(T^2+1)$ for an ordinary prime
$\ell$ (not necessarily a domain), with kernel $\Pi_\C$.
\end{theorem}
\begin{proof}
By \cref{osq:thm:quotients} the finite quotients and their transition maps are those of
$\Z$ (respectively $\Z[i]$), identified by $\ct$; a nonzero integer is not divisible by
every $n$ or every $p^r$. For $\GOz$, every finite quotient has some positive
characteristic $n$, so its kernel contains $n\GOz=n\Z[i]+\Pi_\C$; the ideals $n\Z[i]$ are
cofinal among finite-index ideals, and the limit is free of rank two over $\widehat\Z$
(respectively $\Z_\ell$) with $i^2=-1$.
\end{proof}

The finite quotient system is an ordinary set-indexed system; no inverse limit over a proper
class of targets is formed, and $\Oz$ is not claimed to be an object of the category of small
topological rings (08, 09). The finite-congruence topology on $\Oz$ is not Hausdorff: $\Pi$
is the closure of $0$. It is unrelated to the fine order topology and to Hahn valuation
topologies (07, 10, 11). Every $g\in\Pi$ is divisible in $\Oz$ by every power of $p$, but
this does not represent $g$ as a nonzero $p$-adic integer; for example
$\omega^\omega+\pi\omega+17$ maps to $17$ in every finite quotient and in every $\Z_p$ (07).
The finite-quotient and completion statements were already in the Diophantine manuscripts
of the sibling report \cite{RepoOdg} (04 cites them as prior); the universal theorem explains
why passing from finite targets to arbitrary set-sized targets recovers nothing more.

\subsection{Localizations and presentations}\label{osq:sub:localization}

\begin{proposition}[Adjoining one monomial inverse \src{04, 07}]\label{osq:prop:locsupport}
For $a>0$,
\[
 \A_{D,k}[X^{-a}]=\{x\in\hahn k\No:\ \supp x\subseteq[-Na,\infty)\text{ for some ordinary }N\},
\]
independently of $D$. This is a proper subring of $\hahn k\No$ and not a field; it has no
nonzero set-sized ring target (\cref{osq:cor:nofield}). In particular
$\Oz[\omega^{-a}]$ is a nonzero class domain, not a field, with no nonzero set-sized unital
module.
\end{proposition}
\begin{proof}
An element of the localization is $fX^{-Na}$ with $f\in\A_{D,k}$, whose support is bounded
below by $-Na$. Conversely, if $\supp x\subseteq[-Na,\infty)$, then $X^{(N+1)a}x$ has strictly
positive support, so $x=X^{(N+1)a}x/X^{(N+1)a}$ lies in the localization; the extra factor
$X^a$ removes the restriction on the coefficient at the lower endpoint, which explains the
independence from $D$. The exponent $a\omega$ exceeds every $Na$, so $X^{-a\omega}$ is not in
the ring although $X^{a\omega}$ is; the ring is proper and not a field.
\end{proof}

The same argument shows $\omega^{-\omega}\notin\Oz[1/\omega]$: $\omega^{-\omega}=a/\omega^n$
would force $a=\omega^{n-\omega}\notin\Oz$ (11).

\begin{theorem}[Quotient--localization calculus \src{08, 09, 11}]\label{osq:thm:localization}
Let $R=\A_{D,k}$, $J$ a class ideal and $S\subseteq R$ a multiplicative class containing $1$.
Put $D_0=D/\ct(J)$ and $T=\{\ct(s)+\ct(J):s\in S\}\subseteq D_0$. For every set-sized
commutative ring $B$,
\[
 \Hom\bigl((R/J)[S^{-1}],B\bigr)\cong\Hom\bigl(D_0[T^{-1}],B\bigr);
\]
if $0\in T$ the right side is empty for $B\ne0$. For $R=\Oz$, $J=0$ and $S$ a set-sized
multiplicative subset, the same holds for noncommutative $B$, with universal formula
$f/s\mapsto\ct(f)1_B(\ct(s)1_B)^{-1}$.
\end{theorem}
\begin{proof}
A map from the left ring is a map $h:R\to B$ killing $J$ and sending each $s\in S$ to a unit.
By \cref{osq:thm:universal}, $h=g\circ\ct$ uniquely, and the conditions become
$g(\ct(J))=0$ and $g(\ct(s))\in B^\times$, which are exactly the conditions for a map
$D_0[T^{-1}]\to B$; the sets $\ct(J)$ and $T$ are sets even when $J$ and $S$ are classes.
The natural map to the reflection is surjective on fractions of constants. For the
noncommutative statement, the denominators $\ct(s)1_B$ are central, so the ordinary
universal property of localization applies (08).
\end{proof}

\begin{example}[Three effects \src{08, 11}]\label{osq:ex:localizations}
$\Refset(\Oz[1/(\omega+2)])=\Z[1/2]$; $\Refset(\Oz[1/\omega])=0$;
$\Refset(\Oz/(\omega+1))=0$; and $\Refset(\Oz[1/n])=\Z[1/n]$ for a nonzero ordinary
integer $n$. The second ring is a nonzero domain inside $\No$ and the third is nonzero because
$\omega+1$ is not a unit; both have no nonzero set-sized representation. A nonzero class ring
can therefore have no nonzero set-sized unital module, so set-sized existence arguments for
modules cannot be applied to a proper class of class ideals (11). If $\ct(s)=0$ for some
$s\in S$, the displayed localization of $\Z$ is the zero ring (08). In the language of 09: no
unital map to a nonzero set-sized ring can make a nonzero $f\in\Pi_k$ invertible or make an
element with unit constant term zero.
\end{example}

\begin{theorem}[Algebraic points forget infinite coefficients \src{06, 09, 11}]\label{osq:thm:presentations}
Let $X$ be a set of commuting variables, $(f_j)_{j\in E}$ a set-indexed family in
$\A_{D,k}[X]$, $B=\A_{D,k}[X]/(f_j)$ and $B_0=D[X]/(\ct_*f_j)$, where $\ct_*$ applies
$\ct$ to coefficients. For every set-sized commutative ring $S$ the natural map $B\to B_0$
induces $\Hom(B_0,S)\cong\Hom(B,S)$; thus $\Refset(B)=B_0$.
\end{theorem}
\begin{proof}
A map $\A_{D,k}[X]\to S$ is a map on $\A_{D,k}$, which factors through $\ct$, together with
images of the variables; for such data $f_j$ vanishes exactly when $\ct_*f_j$ does. $B_0$ is a
set because $D$, $X$ and the relations are sets, and $B\to B_0$ is surjective.
\end{proof}

06 states this for finitely many variables and 11 for a set of variables; 09 states the
corresponding reduction of localization and relation problems. The relation $\omega T=1$
reduces to $0=1$, so it has no solution in a nonzero set-sized ring under any representation
of its coefficients, although $T=\omega^{-1}$ solves it in $\No$ (06). For
$B=\Oz[T]/(T^2-\omega)$ one gets $\Refset(B)=\Z[T]/(T^2)$: a set-sized field sends $T$ to
$0$ and a set-sized ring may send it to any square-zero element, while the actual root
$\omega^{1/2}\in\Oz$ gives a map to the class ring $\Oz$ (11). These are universal-property
statements, not quantifier elimination, not a reconstruction of class-valued solutions, and
not an assertion that a scheme with a proper-class coordinate ring is an ordinary scheme
(06, 11).

\subsection{Diophantine observation}\label{osq:sub:dioph}
If $P\in\Z[X_1,\dots,X_m]$, $x\in\Oz^m$ and $\varphi:\Oz\to S$ is a unital set-target map,
then $\varphi(P(x))=P(\ct(x_1)1_S,\dots,\ct(x_m)1_S)$, even for noncommutative $S$, since
the coordinates are integer scalars (04). This does \emph{not} say that $P(x)=0$ whenever
$P(\ct(x))=0$ (take $P(X)=X$ and $x=\omega$), and it does not classify infinite points of
curves; inequalities, supports, valuations and nonmultiplicative coefficient observations
remain available (04).

\begin{proposition}[Conservativity for integer coefficients \src{08}]\label{osq:prop:conservative}
A finite system of polynomial equations with coefficients in $\Z$ has a solution in $\Oz$
if and only if it has one in $\Z$; more generally, positive existential ring formulas with
parameters in $\Z$ have the same truth value in $\Oz$ and $\Z$. With an inequation this
fails: $\exists x\,\exists y\,(x^2=2y^2\wedge y\ne0)$ holds in $\Oz$ (take
$x=\sqrt2\,\omega$, $y=\omega$) but not in $\Z$.
\end{proposition}
\begin{proof}
Apply the homomorphisms $\ct$ and $\Z\hookrightarrow\Oz$, which preserve positive
existential formulas (induction on their construction). A nonzero integer solution of
$x^2=2y^2$ would give a rational square root of $2$; $\ct$ maps the omnific witnesses to
$x=y=0$, preserving the equation but not the inequation.
\end{proof}

The parameter restriction is essential (the coefficient $\omega^\gamma+1$ of
\cref{osq:nm:thm:invisible} is not an integer), and so is positivity (08). The sibling report
\cite{RepoOdg} proves the same constant-term transfer for equations in its Diophantine
setting.

\section{Modules and derivations over the full class ring}\label{osq:sec:modules}

\begin{theorem}[Set-sized modules \src{04, 06, 07, 08, 09, 10, 11}]\label{osq:thm:modules}
Let $A=\A_{D,k}$. Restriction along $D\hookrightarrow A$ and inflation along $\ct$ are mutually
inverse equivalences between set-sized unital $A$-modules and set-sized unital $D$-modules;
the action of $A$ on a set-sized module is $(d+f)m=dm$ ($d\in D$, $f\in\Pi_k$). The
equivalence is the identity on underlying abelian groups and on morphisms, and it is exact.
In particular set-sized $\Oz$-modules are exactly abelian groups with $xm=\ct(x)m$, and
set-sized $\GOz$-modules are exactly $\Z[i]$-modules inflated along $\ct$.
\end{theorem}
\begin{proof}
\Cref{osq:thm:universal}(iv) gives $\Pi_kM=0$ and the action formula. Conversely, a
$D$-module inflates along $\ct$, and restriction recovers the $D$-action because
$\ct|_D=\id$. An additive map between two such modules is $A$-linear exactly when it is
$D$-linear. Kernels, cokernels and exact sequences are computed on underlying groups with
the induced actions.
\end{proof}

A $\GOz$-module structure on an abelian group is the choice of an endomorphism $J$ with
$J^2=-\id$, automatically invertible (08). Category language abbreviates the correspondence
of individual objects and morphisms: each action is coded by its restriction to the set
$D\times M$, and no class of arbitrary class actions is formed (04, 08, 11). The regular
module $A$ is a proper class on which $\Pi_k$ acts nontrivially, so there is no conflict with
recovering a ring from its regular module (04). The equivalence can transport extension
problems between the two small-module categories, but it should not be rephrased as a
statement about homological dimensions computed with arbitrary class-sized modules: the
rank-one free module $A$ is not an object of the small-module category (07).
\Cref{osq:sec:homological} treats the homological algebra over set-sized integer parts.

\begin{proposition}[An additive splitting that is not a module splitting \src{08}]\label{osq:prop:nonsplit}
The exact sequence of class $\Oz$-modules $0\to\Pi\to\Oz\xrightarrow{\ct}\Z\to0$, with
$\Z$ carrying the action $x\cdot n=\ct(x)n$, does not split as a sequence of $\Oz$-modules,
although it splits as a sequence of additive groups and of rings.
\end{proposition}
\begin{proof}
A module section $s$ gives $e=s(1)$ with $\ct(e)=1$ and $ue=s(u\cdot1)=s(0)=0$ for every
$u\in\Pi$. With $u=\omega$ and $\Oz$ a domain, $e=0$, a contradiction.
\end{proof}

\begin{theorem}[Derivations into set-sized modules \src{04, 06, 08, 09, 10, 11}]\label{osq:thm:derivations}
Let $A=\A_{D,k}$ and $M$ a set-sized $A$-module. Every derivation $\partial:A\to M$ (an
additive class map with $\partial(xy)=x\partial(y)+y\partial(x)$) vanishes on $\Pi_k$, and
restriction gives a natural bijection $\Der_\Z(A,M)\cong\Der_\Z(D,M)$ with inverse
$\delta\mapsto\delta\circ\ct$. Consequently:
\begin{enumerate}
\item every derivation $\Oz\to M$ is zero;
\item $\Der_\Z(\GOz,M)\cong M[2]=\{m\in M:2m=0\}$ via $\partial\mapsto\partial(i)$, and
$\Der_{\Z[i]}(\GOz,M)=0$;
\item every $k$-linear derivation of $\A_{k,k}$ or of $\OO_k$ ($k=\R,\C$) into a set-sized
module is zero, and $\Der_D(A,M)=0$.
\end{enumerate}
\end{theorem}
\begin{proof}
$\Pi_kM=0$ by \cref{osq:thm:modules}, and each $x\in\Pi_k$ is a product $uv$ with
$u,v\in\Pi_k$ (\cref{osq:prop:common}), so $\partial(x)=u\partial(v)+v\partial(u)=0$.
Thus $\partial(d+x)=\partial(d)$. Conversely, for a derivation $\delta$ of $D$,
$\delta\circ\ct$ satisfies the Leibniz rule because $A$ acts on $M$ through $\ct$. A
derivation of $\Z$ vanishes ($\partial(1)=2\partial(1)$). A derivation of $\Z[i]$ is
determined by $m=\partial(i)$; differentiating $i^2+1=0$ gives $2im=0$, equivalently $2m=0$
since $i$ acts invertibly; conversely $\partial(a+bi)=bm$ satisfies the Leibniz rule, the
discrepancy for a product being $2bdim=0$. It kills $\Z[i]$ exactly when $m=0$. A $k$-linear
derivation of $k$ vanishes; the finite side uses \cref{osq:thm:finite} and
$\mm_k^2=\mm_k$ (\cref{osq:lem:midempotent}).
\end{proof}

\paragraph{A second route (09).} Since $M$ is a $D$-module, the square-zero extension
$D\ltimes M$ with $(d,m)(e,n)=(de,dn+em)$ is a set-sized ring, and $a\mapsto(\ct(a),\partial(a))$
is a unital ring map $A\to D\ltimes M$; the universal theorem forces it to factor through
$D$.

\begin{example}[The ramification at $2$ survives \src{09, 11}]\label{osq:ex:ramification}
Give $\F_2$ the $\Z[i]$-module structure $i\mapsto1$. Then $\partial(a+bi+f)=b\bmod2$ is a
nonzero $\Z$-derivation $\GOz\to\F_2$. The statement ``all set-valued derivations of the
Gaussian omnific ring vanish'' is false; only the positive-support directions vanish, and
the surviving derivation reflects ordinary ramification of $\Z[i]$ at $2$.
\end{example}

The coefficient-linearity qualification in (iii) cannot be dropped: the theorem reduces
arbitrary derivations of $\A_{\R,\R}$ or $\OO_\R$ to derivations of $\R$ over $\Q$, which
are not zero (04). For a set-sized commutative $D$-algebra $B$, every $D$-algebra map
$A\to B[\varepsilon]/(\varepsilon^2)$ kills $\Pi_k$, so relative first-order deformations
with set-sized coefficients cannot detect $\Pi_k$; no assertion about an absolute
proper-class module of K\"ahler differentials is made (11).

\begin{proposition}[Class-valued derivations \src{06, 08, 10}]\label{osq:prop:classder}
For $\eta\in\No$ let $\ell_\eta:\No\to\R$ extract the coefficient of $\omega^\eta$ from the
normal form of its argument; it is additive. The formula
\[
 D_\eta\Bigl(\sum_x r_x\omega^x\Bigr)=\sum_x \ell_\eta(x)\,r_x\,\omega^x
\]
defines a derivation $\No\to\No$ which restricts to a derivation $\Oz\to\Pi$ into the
class module $\Pi$, preserves strongly summable families, and satisfies
$D_\eta(\omega^{\omega^\theta})=\omega^{\omega^\theta}$ if $\eta=\theta$ and $0$ otherwise.
Every finite family of distinct $D_\eta$ is linearly independent over $\R$, and
$D_0(\omega)=\omega$. Coefficientwise complexification gives the analogous derivations of
$\SC$ and $\GOz$.
\end{proposition}
\begin{proof}
The output support is a subset of the input support, so the formula defines a surreal; the
exponent $0$ is killed because $\ell_\eta(0)=0$, and no negative exponent appears for an
omnific input. Additivity and preservation of strong sums are coefficientwise. At a fixed
exponent $z$ the coefficient of a product is a finite sum over pairs $x+y=z$, and
$\ell_\eta(z)=\ell_\eta(x)+\ell_\eta(y)$; multiplying and summing gives the Leibniz rule.
The exponent $\omega^\theta$ has coefficient $1$ at $\theta$ and $0$ elsewhere, which gives
the displayed values and, by evaluation at the monomials $\omega^{\omega^\theta}$, linear
independence. Finally $\ell_0(1)=1$.
\end{proof}

The proposition combines three sources, none of which states all of it: 08 introduces the
family $D_\eta$ for every $\eta$, as derivations $\Oz\to\Pi$, with the diagonal values and
linear independence; the case $\eta=0$ is 06's Euler-type derivation $\mathcal D$ (06 defines
it on $\No$, with the complex extension) and 10's exponent-weighted derivation $D_h$
($h=\ell_0$; 10 records preservation of strongly summable families). The general-$\eta$
statements on $\No$ and on strong sums follow by the same coefficientwise proof. There is no conflict with
\cref{osq:thm:derivations}: the target $\Pi$ is a proper-class module with its internal
multiplication. Berarducci and Mantova's surreal derivation with $\partial\omega=1$
\cite{BM}, restricted to $\Oz$ with values in $\No$, is another nonzero class-valued
boundary example (09). None of the $D_\eta$ is asserted to be the Berarducci--Mantova
derivation, a fine derivative, or compatible with surreal exponentiation (06, 08, 10), and
no uniqueness is claimed for the established derivation (09). The word ``small'' in
\cref{osq:thm:derivations} cannot be dropped.

\section{The finite side: standard part is universal}\label{osq:sec:finite}

This section is manuscript~04's; no other source treats the negative-support side. Its proof
is shorter because the ambient valuation produces the needed quotients directly, and it is a
different route rather than a sign reversal of the positive side (04). Let $k\in\{\R,\C\}$,
$D\subseteq k$ a unital subring, and for $x\ne0$ let $v(x)=-\degw(x)$, so that
$v(x/y)=v(x)-v(y)$; a nonzero normal form lies in $\mm_k$ exactly when $v(x)>0$.

\begin{lemma}[Division of an infinitesimal by small differences \src{04}]\label{osq:lem:finitediff}
Let $0\ne\varepsilon\in\mm_k$ and $a=v(\varepsilon)>0$. For $0<\gamma<\delta<a$,
$\varepsilon/(\omega^{-\gamma}-\omega^{-\delta})\in\mm_k$.
\end{lemma}
\begin{proof}
The denominator has largest exponent $-\gamma$, so valuation $\gamma$; the quotient has
valuation $a-\gamma>0$, i.e.\ all of its exponents are negative.
\end{proof}

\begin{theorem}[Universal set-sized residue on the finite side \src{04}]\label{osq:thm:finite}
For every set-sized unital ring $S$, restriction gives a bijection
$\Homone(\B_{D,k},S)\to\Homone(D,S)$ with inverse $\psi\mapsto\psi\circ\st$. Every such
homomorphism annihilates $\mm_k$. In particular every unital map from $\OO_\R$ or $\OO_\C$ to
a set-sized ring factors through $\st:\OO_\R\to\R$, respectively $\st:\OO_\C\to\C$.
\end{theorem}
\begin{proof}
Suppose $0\ne\varepsilon\in\mm_k$ and $\varphi(\varepsilon)\ne0$, put $a=v(\varepsilon)$
and choose a set of exponents in $(0,a)$ of cardinality larger than $|S|$
(\cref{osq:lem:gap}). By \cref{osq:lem:finitediff}, $\varepsilon$ is a multiple, by an element
of $\mm_k\subseteq\B_{D,k}$, of the difference of any two of the monomials
$\omega^{-\gamma}$, so \cref{osq:lem:separation}(i) injects the chosen set into $S$, a
contradiction. Hence $\varphi(\mm_k)=0$ and $\varphi(d+\varepsilon)=\varphi(d)$; conversely
$\psi\circ\st$ extends every $\psi$.
\end{proof}

For a target without a specified coefficient-field structure, the induced map on $\R$ or
$\C$ is not presumed to be the identity: the assertion is factorization through the residue
field, not uniqueness of maps from $\OO_\C$ to an arbitrary ring (04). The statement involves
no topology on $S$ and so is stronger than a statement about continuous maps into topological
rings; it must also be distinguished from the existence of standard part, which is classical
and part of the repository's formal infrastructure \cite{RepoBridge}.

\begin{corollary}[Quotients, modules and derivations on the finite side \src{04}]\label{osq:cor:finitequotients}
The kernels of $\B_{D,k}$ admitting a set-sized quotient are exactly $\mm_k+\mathfrak a$
for ideals $\mathfrak a$ of $D$; the only nonzero set-sized quotients of $\OO_\R$ and
$\OO_\C$ are $\R$ and $\C$, up to the residue identification. Set-sized $\B_{D,k}$-modules
are exactly $D$-modules inflated along $\st$ (set-sized $\OO_k$-modules are $k$-vector
spaces), no faithful set-sized module exists, and $\Der_\Z(\B_{D,k},M)\cong\Der_\Z(D,M)$ for
set-sized $M$. For $0\ne\varepsilon\in\mm_k$, the localization $\B_{D,k}[\varepsilon^{-1}]$
has no nonzero set-sized ring target and no nonzero set-sized module.
\end{corollary}
\begin{proof}
Repeat the proofs of \cref{osq:thm:quotients,osq:thm:modules,osq:thm:derivations,osq:cor:nofield}
with \cref{osq:thm:finite} and $\st$ in place of the positive-side theorem and $\ct$; for
$D=k$ the only ideals are $0$ and $k$. Derivations use \cref{osq:lem:midempotent}.
\end{proof}

\begin{lemma}[Idempotence on the finite side \src{04}]\label{osq:lem:midempotent}
$\mm_k^2=\mm_k$ as an ideal of $\B_{D,k}$.
\end{lemma}
\begin{proof}
For $0\ne\varepsilon\in\mm_k$ with $a=v(\varepsilon)$,
$\varepsilon=\omega^{-a/2}(\omega^{a/2}\varepsilon)$ with both factors of valuation $a/2>0$.
\end{proof}

\paragraph{Two residues that cannot be combined (04).} Every element of $\hahn k\No$ has a
unique decomposition along $\Pi_k\oplus k\oplus\mm_k$, obtained by splitting its support at
$0$; so $\A_{k,k}\cap\OO_k=k$ and $\A_{k,k}+\OO_k$ is the whole field additively. The two
constant-coefficient maps agree on the intersection, yet no unital ring map on the field
extends both:
\begin{equation}\label{osq:eq:incompatible}
 \ct(\omega)=0,\qquad\st(\omega^{-1})=0,\qquad\omega\cdot\omega^{-1}=1 .
\end{equation}
The obstruction is multiplicative, not an ambiguity of the additive decomposition, and
\cref{osq:cor:nofield} sharpens it: adjoining one inverse to the positive-support ring
already destroys every nonzero set-sized target.

\paragraph{Relation to the rotation groups.} The euclidean-three-space report proves that
every class homomorphism from $\mathrm{SO}(3,\No)$ to a set-sized group factors through
standard part (\nolinkurl{e3:cut:thm:smallquotient}), that every nontrivial set-sized image is
$\mathrm{SO}(3,\R)$ (\nolinkurl{e3:cut:cor:onlyimage}), and the analogues for $\mathrm{SO}(n,\No)$, $n\ge3$,
and $\mathrm{SU}(n,\SC)$, $n\ge2$ \cite{RepoE3}. \Cref{osq:thm:finite} is the ring analogue on the
finite side; neither result implies the other. The group proof uses exact commutators and
normal generation before a choice-free Hartogs contradiction; the ring proof uses
divisibility by differences of monomials. Manuscripts 06, 07, 08, 09, 10 and 11
acknowledge the rotation theorem as the repository's conceptual predecessor for universal
set-sized quotients, 04 refers to related size obstructions in other algebraic settings,
and none of them uses the rotation theorem as a lemma.

\section{The support threshold: finite supports fail, countable supports suffice}\label{osq:sec:support}

The universal theorem does not need normal forms of every set cardinality. Keep the
\emph{full} exponent class $\No$ and bound only the size of each individual support. For
$k\in\{\R,\C\}$, $D\subseteq k$ unital, and an infinite cardinal $\lambda$, put
\begin{align*}
 \A^{<\lambda}_{D,k}&=D+\{f\in\Pi_k:|\supp f|<\lambda\},&
 \B^{<\lambda}_{D,k}&=D+\{f\in\mm_k:|\supp f|<\lambda\},
\end{align*}
and, in 11's notation, $R_{\le\lambda}=\A^{<\lambda^+}_{D,k}$ and
$R_{\mathrm{fin}}=\A^{<\aleph_0}_{D,k}$. These are rings: the support of a sum lies in a
finite union and that of a product in the image of the product of the two supports under
addition, and for an infinite $\lambda$ finite unions and products of sets of size $<\lambda$
(or $\le\lambda$) keep that bound; regularity of $\lambda$ is not needed. They remain
proper classes, since a one-term series may have any positive exponent (04, 11).

\begin{theorem}[Support threshold \src{04, 11; 03, 07, 09}]\label{osq:thm:support}
\leavevmode
\begin{enumerate}
\item For every uncountable cardinal $\lambda$ and every set-sized unital ring $S$,
restriction induces bijections
\[
 \Homone(\A^{<\lambda}_{D,k},S)\cong\Homone(D,S),\qquad
 \Homone(\B^{<\lambda}_{D,k},S)\cong\Homone(D,S),
\]
with inverses given by $\ct$ and $\st$. In particular countably supported normal forms
($\lambda=\aleph_1$, i.e.\ $R_{\le\aleph_0}$) already have universal set-sized quotient $D$
on both sides, and so does $R_{\le\mu}$ for every infinite $\mu$.
\item For finite supports both assertions fail: the \emph{augmentation}
\[
 E\Bigl(d+\sum_{g\in F}c_g\omega^g\Bigr)=d+\sum_{g\in F}c_g\qquad(F\text{ finite})
\]
is a unital ring homomorphism from each finite-support ring to $k$, sending every monomial,
including $\omega$ and $\omega^{-1}$, to $1$; it does not factor through $\ct$ or $\st$.
\item (Extension criterion.) For a set-sized unital ring $S$ and a unital map
$h:R_{\mathrm{fin}}\to S$ the following are equivalent: (a) $h$ extends to $R_{\le\aleph_0}$;
(b) $h$ extends to $R_{\le\mu}$ for some infinite $\mu$; (c) $h=g\circ\ct$ for a map
$g:D\to S$; (d) $h$ extends to $\A_{D,k}$. Whenever an extension exists it is unique at each
of these support levels. In particular $E$ does not extend to $R_{\le\aleph_0}$.
\item \merge\ The same equivalence holds on the finite side, with $\B^{<\aleph_0}_{D,k}$,
$\B^{<\aleph_1}_{D,k}$, $\B^{<\lambda}_{D,k}$, $\B_{D,k}$ and $\st$ in place of
$R_{\mathrm{fin}}$, $R_{\le\aleph_0}$, $R_{\le\mu}$, $\A_{D,k}$ and $\ct$.
\end{enumerate}
\end{theorem}

\begin{proof}
(i), positive side. Fix $s\in\Pi_k$ with $|\supp s|<\lambda$. For $a>b$ in $U_{\No}(\supp s)$,
the quotient of \cref{osq:lem:division} has support of cardinality at most
$\max(|\supp s|,\aleph_0)<\lambda$, and the monomials $\omega^a$ are in the ring; so the
families of \cref{osq:thm:packing} and all their quotients lie in $\A^{<\lambda}_{D,k}$, and
\cref{osq:lem:separation}(iii) gives $\varphi(s)=0$ (11). Alternatively (04), every
monomial and every telescope $q_{c;a,b}$ has at most countable support, so the explicit-telescope
route (\cref{osq:lem:telescope} with \cref{osq:lem:separation}(i); see the second routes after
\cref{osq:thm:universal}) kills all positive monomials inside the bounded ring (the
field-collision proof of \cref{osq:thm:universal} does not apply here, since the Hahn field over
$H_g$ is not contained in the bounded ring), and the factorization
$s=\omega^a(\omega^{-a}s)$ of \cref{osq:prop:common} only translates the support, keeping
its cardinality.

(i), finite side (04). For $0\ne\varepsilon$ with $|\supp\varepsilon|<\lambda$,
$a=v(\varepsilon)$, $0<\gamma<\delta<a$ and $d=\delta-\gamma$,
\[
 \frac{\varepsilon}{\omega^{-\gamma}-\omega^{-\delta}}=\varepsilon\,\omega^{\gamma}\sum_{n\ge0}\omega^{-nd},
\]
the geometric series being the inverse of $1-\omega^{-d}$. Its support lies in
$\{\beta+\gamma-nd:\beta\in\supp\varepsilon,\ n\in\N\}$, of cardinality at most
$|\supp\varepsilon|\cdot\aleph_0<\lambda$, and its largest exponent is $-a+\gamma<0$. So the
quotient of \cref{osq:lem:finitediff} stays in the bounded ring, and the proof of
\cref{osq:thm:finite} applies. A valuation inequality alone would not control the support
size; this check is needed.

(ii) Additivity and $E(1)=1$ are clear; for finite supports
$E(xy)=\sum_{\alpha,\beta}c_\alpha d_\beta=E(x)E(y)$ by expanding the finite double sum.
$E(\omega)=1\ne0=\ct(\omega)$ and $E(\omega^{-1})=1\ne0=\st(\omega^{-1})$.

(iii) An extension to some $R_{\le\mu}$ factors through $\ct$ by (i), giving (a)$\Rightarrow$(c)
and (b)$\Rightarrow$(c); an extension to $\A_{D,k}$ factors by \cref{osq:thm:universal}, giving
(d)$\Rightarrow$(c). Conversely $g\circ\ct$ is defined on every ring in the list and extends
$h$, and two extensions to one ring both factor through $\ct$ and agree on $D$. Since
$E(\omega)=1\ne\ct(\omega)$, $E$ violates (c).

(iv) The same argument with (i) on the finite side, \cref{osq:thm:finite} and $\st$; the
extension $g\circ\st$ exists on every ring of the list, and uniqueness follows as before.
\end{proof}

\begin{remark}[Why the threshold is about closure, not continuity]\label{osq:rem:closure}
No axiom that target maps preserve countable sums is added: certain countable sums lie
\emph{inside} the source ring, where they satisfy finite identities. An extension of $E$ to
the countable-support ring would send both monomials in \cref{osq:ex:omegatelescope} to $1$
and the product to $1$, forcing $(1-1)E(q)=1$ (04, 11). A map that assigns different values
to some monomials is handled by the full theorem: a large family forces a collision somewhere
(11). The separating families may have cardinality far above $\lambda$; their members are
distinct elements, not terms of one series, and each witness separately has countable
support (04, 07, 11).
\end{remark}

The finite-support augmentation appears independently in 03 (``two necessary warnings''),
04, 09 and 11; 07 describes the same boundary in words (on a finite-support fragment the
certificate \eqref{osq:eq:telescope} need not exist, and a faithful interpretation in a
set-sized target can exist), and manuscript~05 of the sibling report \cite{RepoOdg}
observes that $\Z+\omega\R[\omega]$ has many real evaluation maps that do not extend to
$\Oz$. Proper-class size alone therefore does not prove the universal theorem: the
finite-support rings have a proper class of monomials and arbitrarily small positive
exponents, yet $E$ sees every monomial as $1$ (04). Neither size nor additive divisibility
alone suffices; the proof needs many subordinate scales together with enough support closure
for geometric division (11). For $D=\Z$, finite targets still kill the positive-support
ideal of $R_{\mathrm{fin}}$ ($f=n(f/n)$ in characteristic $n$), so its finite quotient theory is
ordinary congruence theory; having profinite completion $\widehat\Z$ is much weaker than having
$\Z$ universal for all set-sized targets (09).

The quotient and module classifications extend to the bounded rings of (i), and so does the
derivation statement, because the monomial factorizations proving idempotence preserve support
cardinality (04). 04 does \emph{not} infer that these smaller rings have fraction field
$\hahn k\No$: arbitrary normal forms need not have the required support size. The result is
also different from completion or omitted-type questions for bounded-support series over a
\emph{fixed} set-sized exponent group (compare \cite{RepoFKC}): here the exponent class is
still large enough to supply a new separating family for every proposed target (04). The
fixed-group question is the subject of the next section.

\section{Fixed exponent groups: cardinal lower bounds and exact thresholds}\label{osq:sec:fixed}

Over a fixed set-sized exponent group the full-class theorem cannot hold verbatim: the ring
is itself a set, and its identity map is a set-sized target that retains every element
(03, 04, 07, 08, 10, 11). What survives is a lower bound on the size of a target that
detects a given element. This section compares the bounds proved by the sources, proves a
field bound that dominates them, and states the organizing theorem of the report: for each
of the set-sized constructions of \cref{osq:sec:models}, the exact detection threshold is
the size of the ring.

Throughout, $\Gamma$ is an ordered abelian set group, $k$ a set field, $D\subseteq k$ a
unital subring, $F\subseteq\hahn k\Gamma$ a subfield containing $k$ and every monomial
$X^g$, $\Pi(F)=\{f\in F:\supp f\subseteq\Gamma_{>0}\}$ and $A=D+\Pi(F)$. The case
$F=\hahn k\Gamma$ is the ring $\A_{D,k}(\Gamma)$ of \eqref{osq:eq:rings}; the constructions of
\cref{osq:sec:models} take $F$ to be a countable-support, a $<\kappa$-support, or a
controlled subfield.

\subsection{The bounds of the sources}\label{osq:sub:sourcebounds}

\begin{proposition}[Monomial-count bounds \src{03, 04, 07, 10}]\label{osq:prop:monomialbounds}
Let $A=\A_{D,k}(\Gamma)$ and $\varphi:A\to S$ a ring homomorphism.
\begin{enumerate}
\item If $\varphi(X^a)\ne0$ for some $a>0$, then $\gamma\mapsto\varphi(X^\gamma)$ is injective
on $H_a(\Gamma)_{>0}$ (04); more generally $r X^\delta\mapsto\varphi(rX^\delta)$ is injective on
$\{rX^\delta:r\in k^\times,\ 0<\delta\in H_a(\Gamma)\}$, and the same holds if $X^a$ divides
some $g$ with $\varphi(g)\ne0$ (07).
\item If $0\ne x\in\Pi_k(\Gamma)$ and $\varphi(x)\ne0$, then $a\mapsto\varphi(X^a)$ is injective on
$U_\Gamma(\supp x)$; this also holds for a nonunital map defined on $\Pi_k(\Gamma)$ alone (10).
\item On the negative side, if $0\ne\varepsilon\in\mm_k(\Gamma)$, $a=v(\varepsilon)$ and
$\varphi(\varepsilon)\ne0$ for a map $\varphi$ on $\B_{D,k}(\Gamma)$, then
$\gamma\mapsto\varphi(X^{-\gamma})$ is injective on $(0,a)_\Gamma$ (04).
\item For a module $M$ and $v\in M$ with $X^hv\ne0$, $a\mapsto X^av$ is injective on
$\{a\in\Gamma: a\ll h\}$ (03).
\end{enumerate}
No divisibility of $\Gamma$ is needed.
\end{proposition}
\begin{proof}
The telescope, the weighted certificate, the division lemma and the finite-side division
(\cref{osq:lem:telescope,osq:lem:weighted,osq:lem:division,osq:lem:finitediff}) use only
$a$, $\gamma$, $\delta$, subtraction and ordinary integer multiples, so their witnesses lie in
$\A_{D,k}(\Gamma)$ or $\B_{D,k}(\Gamma)$; apply \cref{osq:lem:separation}.
\end{proof}

\begin{example}[Where these bounds are vacuous \src{04, 08, 10}]\label{osq:ex:archimedean}
If $\Gamma$ is Archimedean, $H_a(\Gamma)=\{0\}$ and the telescopes do not exist. For
$\Gamma=\Z$ a reverse well-ordered positive support is finite, $\A_{\Z,\R}(\Z)=\Z+X\R[X]$,
and evaluation at $X=1$ is a map to the set field $\R$ detecting $X$ (10). An arbitrary fixed
Hahn workspace cannot inherit the full theorem merely because its notation resembles $\No$
(08).
\end{example}

\subsection{The field bound}\label{osq:sub:fieldbound}

The monomial counts are far from optimal. The following bound is 03's ``smaller-scale field
obstruction'', obtained there for countable-support fields over divisible groups; the general
form below is a direct application of 06's scaled-field collision lemma and is recorded as a
merge statement.

\begin{proposition}[Field bound \src{03, 06} \merge]\label{osq:prop:fieldbound}
Let $c\in\Gamma_{>0}$ and $L_c(F)=F\cap\hahn k{H_c(\Gamma)}$.
\begin{enumerate}
\item $L_c(F)$ is a subfield of $F$ and $X^cL_c(F)\subseteq\Pi(F)\subseteq A$.
\item If $\varphi:A\to S$ is a ring homomorphism with $\varphi(X^{2c})\neq0$, then
$f\mapsto\varphi(X^cf)$ is injective on $L_c(F)$; if $M$ is an $A$-module and
$X^{2c}v\ne0$, then $f\mapsto(X^cf)v$ is injective on $L_c(F)$.
\item Consequently, if $x\in X^{2c}A$ (for instance $x=X^{2c}$, or $0\ne x\in\Pi(F)$ with
$2c$ a strict lower bound of $\supp x$), every ring homomorphism and every module detecting
$x$ has cardinality at least $|L_c(F)|$.
\end{enumerate}
\end{proposition}
\begin{proof}
(i) $\hahn k{H_c(\Gamma)}$ is a subfield by \cref{osq:lem:Ha}(ii), so $L_c(F)$ is an
intersection of two subfields. For $h\in H_c(\Gamma)$, $|h|<c$, so $c+h>0$ and
$X^cf\in F$ has positive support. (ii) is \cref{osq:lem:collision} with $m=X^c$ and
$L=L_c(F)$, inside the field $\hahn k\Gamma$. (iii) If $x=X^{2c}y$ with $y\in A$ and
$\varphi(x)\ne0$, then $\varphi(X^{2c})\ne0$; if $xv\ne0$ then $X^{2c}(yv)\ne0$. If $2c<\supp x$
then $X^{-2c}x\in F$ has positive support, so it lies in $\Pi(F)\subseteq A$.
\end{proof}

For $F=\Hc{k}{\Gamma}$ and $c=h/2$ this is exactly 03's theorem ($H_{h/2}=H_h$). The bound
contains 04's and 10's monomial counts whenever those monomials lie in $H_c(\Gamma)$, and
07's weighted count, because $L_c(F)$ contains every $rX^\delta$ with $r\in k$ and
$\delta\in H_c(\Gamma)$; it also contains $k$ itself, so with real coefficients every detecting
target has at least $2^{\aleph_0}$ elements (03). The following examples show that it can be
strictly and decisively larger.

\begin{example}[The exact threshold in 04's example \merge]\label{osq:ex:lexQQ}
Let $\Gamma=\Q\oplus_{\mathrm{lex}}\Q$ (first coordinate dominant), $k\in\{\R,\C\}$,
$A=\A_{D,k}(\Gamma)$ with full Hahn supports, and $a=(1,0)$. 04 computes the
positive part $H_a(\Gamma)_{>0}=\{(0,q):q\in\Q_{>0}\}$, so $H_a(\Gamma)=\{0\}\oplus\Q$, and notes that retaining $X^a$ requires only a countable family of
distinct images. With $c=(1/2,0)$, $H_c(\Gamma)=\{0\}\oplus\Q$ and
$L_c=\hahn k{\{0\}\oplus\Q}\cong k(\!(t^\Q)\!)$ has cardinality $2^{\aleph_0}$; and
$|A|\le|\hahn k\Gamma|\le 2^{\aleph_0}\cdot(2^{\aleph_0})^{\aleph_0}=2^{\aleph_0}$, since
$\Gamma$ is countable. Hence the least cardinality of a ring target, and of a module, detecting
any nonzero $x\in X^{(1,0)}A$ is exactly $2^{\aleph_0}$, attained by the identity map and the regular
module. The same holds on the negative side: if $\varepsilon\in\mm_k(\Gamma)$ has
$v(\varepsilon)=(1,0)$, then with $c=(1/4,0)$ the element $X^{-2c}$ divides $\varepsilon$ in
$\B_{D,k}(\Gamma)$ (the cofactor $X^{2c}\varepsilon$ has largest exponent $(-1/2,0)<0$),
$X^{-c}L_c\subseteq\mm_k(\Gamma)$, and \cref{osq:lem:collision} gives the bound $2^{\aleph_0}$,
which is again attained.
\end{example}

In this example 07's weighted bound ($|k^\times|\cdot\aleph_0=2^{\aleph_0}$) is also attained,
while 04's and 10's monomial counts give only $\aleph_0$. With a countable coefficient field
all three source bounds fail:

\begin{example}[None of the monomial or weighted counts is attained \merge]\label{osq:ex:countablecoeff}
Let $k=D=\Q$, $\Gamma=\bigoplus_{n<\omega}\Q e_n$ with the least nonzero index dominant, and
$A=\A_{\Q,\Q}(\Gamma)$ with full Hahn supports. For $x=X^{e_0}$ and $c=e_0/2$,
$H_c(\Gamma)=\bigoplus_{n\ge1}\Q e_n$, and $L_c$ contains the $2^{\aleph_0}$ distinct series
$\sum_{n\in B}X^{e_n}$, $B\subseteq\omega\setminus\{0\}$ (the exponents $e_n$ decrease with $n$,
so each support is reverse well ordered). Since $\Gamma$ is countable,
$|A|\le2^{\aleph_0}\cdot\aleph_0^{\aleph_0}=2^{\aleph_0}$. So the exact detection threshold of
$X^{e_0}$ is $2^{\aleph_0}$, while 04's bound $|H_{e_0}(\Gamma)_{>0}|$, 07's bound
$|\Q^\times|\cdot|H_{e_0}(\Gamma)_{>0}|$ and 10's bound $|U_\Gamma(\{e_0\})|$ are all
$\aleph_0$.
\end{example}

These examples settle the two example-level questions of the sources and leave the general
ones open: 04's question on optimal small targets in a fixed exponent group is answered for
its own example $\Q\oplus_{\mathrm{lex}}\Q$ (the answer is $2^{\aleph_0}$), and for 10's question
(when is the bound of \cref{osq:prop:monomialbounds}(ii) attained?)
\cref{osq:ex:countablecoeff} shows, in ZFC, that the bound need not be attained. The same phenomenon appears in 03's ring
$A^{03}_{\aleph_1}$ of \cref{osq:sec:model03}: there $|U_\Gamma(S)|\le\aleph_1$, while the
exact threshold is $2^{\aleph_0}$, so if the continuum hypothesis fails, 10's bound is not
attained there either. Remark~\ref{osq:rem:04misleading} records the corresponding correction
to a sentence of 04.

\begin{remark}[A correction to 04's commentary]\label{osq:rem:04misleading}
04 remarks that in $\Q\oplus_{\mathrm{lex}}\Q$ ``this is only a countable family'' which
``cannot contradict the existence of a sufficiently large infinite set target, such as the
ring itself''. The statement is true, but the ring itself has $2^{\aleph_0}$ elements and,
by \cref{osq:ex:lexQQ}, no smaller target detects $X^{(1,0)}$; the countable family is not the
obstruction that decides the answer.
\end{remark}

\subsection{Criteria that kill the whole ideal below a cardinal}\label{osq:sub:criteria}

To pass from monomials to the whole ideal one also needs a gap: every element must be
divisible by a positive monomial inside the ring. 03's example $\sum_{n\ge1}X^{1/n}$ over
$\Gamma=\R$ shows that this can fail (\cref{osq:ex:gap}), and 09's uncountable-support series
$f_*$ (\cref{osq:rem:fstar}) shows the same for full supports over a long group.

\begin{theorem}[Quantitative criterion \src{08} \merge]\label{osq:thm:criterion}
Let $\Gamma$ be divisible, $\kappa$ an infinite cardinal, and suppose
\begin{enumerate}
\item[(a)] every nonzero $f\in\Pi(F)$ has a positive exponent $\epsilon\in\Gamma$ strictly
below its support;
\item[(b)] for every $c>0$ in $\Gamma$, the set $\{a\in\Gamma:a\ll c\}$ has at least $\kappa$
elements.
\end{enumerate}
Then every ring homomorphism from $A$ to a ring of cardinality $<\kappa$ factors uniquely
through $\ct:A\to D$. \merge\ The same conclusion holds, with the same proof, if (b) is replaced
by (b$'$): $|L_c(F)|\ge\kappa$ for every $c>0$.
\end{theorem}
\begin{proof}
Fix $c>0$. Under (b), for two distinct $a>b\ll c$ the series $q_{c;a,b}$ equals the rational
expression $X^c/(X^a-X^b)$ in $\hahn k\Gamma$, whose numerator and denominator lie in the field
$F$; so $q_{c;a,b}\in\Pi(F)\subseteq A$, and a nonzero image of $X^c$ would make the images of
$\kappa$ monomials distinct. Under (b$'$) use \cref{osq:prop:fieldbound}(ii) with $c/2$ in
place of $c$. So every positive monomial is killed. For $f\in\Pi(F)$, (a) gives
$f=X^\epsilon g$ with $g=X^{-\epsilon}f\in F$ of positive support, so $f$ is killed.
\end{proof}

Condition (a) is a gap property of the chosen field and (b) a supply of lower Archimedean
scales; neither says that homomorphisms preserve strongly summable families (08). 08 states
the theorem for a subfield $K\subseteq\R$ and constants $\Z$; the proof uses neither
restriction, and it is stated here in the setting of this section.

\begin{theorem}[Coinitiality criterion \src{11}]\label{osq:thm:coinitial}
Let $\Gamma$ be divisible, $\kappa$ uncountable, and $\coin(\Gamma_{>0})\ge\kappa$ (every set of
fewer than $\kappa$ positive elements has a positive strict lower bound). Put
$R_{\Gamma,<\kappa}=D+\{s\in\Pi_k(\Gamma):|\supp s|<\kappa\}$, a ring for every infinite $\kappa$.
Then every unital map $R_{\Gamma,<\kappa}\to S$ with $|S|<\kappa$ factors uniquely through
$\ct$.
\end{theorem}
\begin{proof}
Let $s\ne0$ have $|\supp s|<\kappa$. Choose $\delta>0$ strictly below $\supp s$, and then a
positive strict lower bound $\eta$ of the countable set $\{\delta/n:n\ge1\}$ (possible since
$\kappa>\aleph_0$); every $a\in E=(0,\eta)$ satisfies $na<\delta$. The set $E$ has at least $\kappa$
elements: otherwise it would have a positive strict lower bound $b$, but then $b\in E$ (as
$b<\eta/2\in E$), which is absurd. For distinct $a,b\in E$, \cref{osq:lem:division} gives
$s\in(X^a-X^b)R_{\Gamma,<\kappa}$ with a quotient of support size at most
$\max(|\supp s|,\aleph_0)<\kappa$. Since $|S|<|E|$, two of the monomials collide and $s$ is killed.
\end{proof}

This theorem concerns an ordinary set-sized source and cannot exclude its identity
representation; the target bound is indispensable. Regularity of $\kappa$ is not needed for the
theorem itself, only for the examples of \cref{osq:sub:model11} (11).

\begin{lemma}[Cyclic modules \src{08}]\label{osq:lem:cyclic}
Let $A$ be a commutative ring, $I$ an ideal and $\kappa$ a cardinal such that every ring
homomorphism from $A$ onto a ring of cardinality $<\kappa$ kills $I$. Then $IM=0$ for every
$A$-module $M$ with $|M|<\kappa$.
\end{lemma}
\begin{proof}
For $m\in M$, $A/\Ann(m)\to Am$, $a\mapsto am$, is a bijection, so the quotient ring
$A/\Ann(m)$ has fewer than $\kappa$ elements, and $I\subseteq\Ann(m)$.
\end{proof}

Applying the ring statement to $\End_\Z(M)$ would need an estimate of $|\End_\Z(M)|$, which may
exceed $\kappa$ even when $|M|<\kappa$; the lemma avoids that hidden cardinal-arithmetic
assumption (08). \Cref{osq:prop:fieldbound}(ii) gives the module bound directly (06, 03, 09).

\subsection{The organizing theorem: the threshold is the size of the ring}\label{osq:sub:organizing}

The five constructions of \cref{osq:sec:models} are summarized here; their notation is tagged by
source because they are different rings.

\begin{center}\footnotesize
\setlength{\tabcolsep}{3.5pt}
\begin{tabular}{@{}lllllll@{}}
\toprule
ring & exponent group & supports & coefficients & hypotheses & size & gen.\ of $\Pi_A$\\
\midrule
$A^{03}_\kappa$ & $\bigoplus_{\kappa^{\mathrm{op}}+\kappa}\Q$, 2 arms & countable & $\R$ & $\cf\kappa>\aleph_0$ & $\kappa^{\aleph_0}$ & $\cf\kappa$\\
$A^{06}_\kappa$ & $\Span_\Q\{\omega^{\pm(\alpha+1)}\}$, 2 arms & $<\kappa$ & $\R$ & $\kappa$ reg.\ unctbl. & $\kappa^{<\kappa}$ & $\kappa$\\
$A^{08}_\kappa$ & $\bigoplus_{\xi<\kappa}\Q$, 1 arm & countable, controlled & $\R\cap\overline\Q$ & $\kappa$ infinite & $\kappa$ & $\cf\kappa$\\
$A^{09}_{\lambda,D}$ & $\bigoplus_{\omega_1\times\lambda}\Q$, 1 arm & countable & $k\in\{\R,\C\}$ & $\lambda^{\aleph_0}=\lambda$ & $\lambda$ & $\aleph_1$\\
$R^{11}_{\Gamma,<\kappa}$ & $\bigoplus_{\alpha<\kappa}\Q$, 1 arm & $<\kappa$ & $|K|\le\max(\kappa,2^{\aleph_0})$ & $\kappa$ reg.\ unctbl. & $\kappa^{<\kappa}$ & $\kappa$\\
\bottomrule
\end{tabular}
\end{center}

All are ZFC constructions. The sizes are proved in \cref{osq:sec:models}; for $A^{06}_\kappa$ and
$R^{11}_{\Gamma,<\kappa}$ the size $\kappa^{<\kappa}$ is a merge result (06 and 11 compute the size
only under $\kappa^{<\kappa}=\kappa$).

\begin{theorem}[Detection threshold equals size \src{03, 08, 09} \merge]\label{osq:thm:thresholds}
Let $A$ be any of the five rings of the table, with purely infinite ideal $\Pi_A$, and let
$0\ne x\in\Pi_A$. Then
\[
 \min\{|S|:\varphi:A\to S,\ \varphi(x)\ne0\}
 =\min\{|M|: M\text{ an }A\text{-module},\ xM\ne0\}=|A| .
\]
Every homomorphism from $A$ to a ring of cardinality $<|A|$ factors uniquely through $\ct$, every
$A$-module of cardinality $<|A|$ is annihilated by $\Pi_A$, and every quotient $A/J$ either is
$D/\ct(J)$ with $\Pi_A\subseteq J$ or has cardinality exactly $|A|$. The same holds for the
Gaussian versions of $A^{03}_\kappa$, $A^{06}_\kappa$ and $A^{09}_{\lambda,\Z[i]}$.
\end{theorem}
\begin{proof}
The identity map and the regular module detect $x$ and have cardinality $|A|$. For the lower
bound, each construction has a divisible exponent group and the gap property: $x$ has a positive
strict lower bound $g$ of its support (\cref{osq:lem:smallcommon,osq:lem:kappagap,osq:lem:finitecoord,osq:lem:tail}
and, for $R^{11}_{\Gamma,<\kappa}$, $\coin(\Gamma_{>0})=\kappa$ with \cref{osq:prop:lex}). With
$c=g/2$, \cref{osq:prop:fieldbound}(iii) shows that every ring or module detecting $x$ has at
least $|L_c(F)|$ elements, and in each construction $|L_c(F)|=|A|$ for every $c>0$
(\cref{osq:lem:card03,osq:lem:card06,osq:lem:tail,osq:prop:card11}; for $A^{08}_\kappa$,
$L_c(F)\subseteq F$ has at most $\kappa$ elements and contains the $\kappa$ monomials $X^{e_\xi}$ with
$\xi$ beyond the least index of $c$, by \cref{osq:lem:gammabounds}). A homomorphism
into a smaller ring therefore kills $\Pi_A$ and factors through $\ct$, and a smaller module is
annihilated. If $\Pi_A\not\subseteq J$ the quotient map detects an element of $\Pi_A$, so
$|A/J|\ge|A|$; the reverse inequality is trivial.
\end{proof}

For $A^{03}_\kappa$, $A^{08}_\kappa$ and $A^{09}_{\lambda,D}$ the theorem is proved in the
sources (03 with the field bound; 08 and 09 by counting monomials, which suffices because there
the number of small monomials already equals $|A|$). For $A^{06}_\kappa$ 06 proves the lower
bound $\kappa$ and exactness under $\kappa^{<\kappa}=\kappa$, and for $R^{11}_{\Gamma,<\kappa}$ 11
proves sharpness under $\kappa^{<\kappa}=\kappa$ and $|K|\le\kappa$; the unconditional exact value
$\kappa^{<\kappa}$ is new here (\cref{osq:thm:model06,osq:prop:card11}). The quotient alternative
is 08's and 09's for their rings and a one-line consequence for the others.

\begin{remark}[Reconciling the thresholds \merge]\label{osq:rem:reconcile}
\leavevmode
\begin{itemize}
\item \emph{Realized cardinals.} Since $(\kappa^{\aleph_0})^{\aleph_0}=\kappa^{\aleph_0}$, and
conversely a cardinal $\theta$ with $\theta^{\aleph_0}=\theta$ has $\cf\theta>\aleph_0$ by
K\"onig's theorem (so $\kappa=\theta$ is allowed in 03), the thresholds realized by 03 are
exactly the cardinals $\theta$ with $\theta^{\aleph_0}=\theta$; these are 09's cardinals
$\lambda$. 06 and 11 realize $\kappa^{<\kappa}$ for regular uncountable $\kappa$. 08 realizes every
infinite cardinal, and is the only construction reaching cardinals below $2^{\aleph_0}$ or with
$\kappa^{\aleph_0}>\kappa$: if $rX^c$ lies in the ring for every real $r$, every target detecting $X^{2c}$ has at least
$2^{\aleph_0}$ elements by \cref{osq:lem:collision} with $L=\R$, and 08 escapes this with real
algebraic coefficients and a countable-stage closure.
\item \emph{The same ring, different bounds.} $A^{06}_{\aleph_1}$ and $A^{03}_{\aleph_1}$ are
isomorphic: both are the countable-support real Hahn integer parts over the lexicographic sum
$\bigoplus_J\Q$ with $J$ of order type $\omega_1^{\mathrm{op}}+\omega_1$ (06's generators satisfy
$\epsilon_\beta\ll\epsilon_\alpha\ll h_\alpha\ll h_\beta$ for $\alpha<\beta$). 06 proves the
threshold is at least $\aleph_1$ and exactly $\aleph_1$ under CH; 03 proves it is exactly
$2^{\aleph_0}$ in ZFC; 06's own example already notes that without CH its ring need not have size
$\aleph_1$. These agree.
\item \emph{Generators.} The generator counts are $\cf\kappa$ (03, 08), $\kappa$ (06, and for
11's ring by the argument of (H3) in \cref{osq:sec:homological}) and $\aleph_1$ (09). In each case the proof exhibits a coinitial chain of monomials in the positive
cone of the exponent group of that length; the counts differ because the constructions differ,
and no general formula is asserted for other exponent groups or support restrictions.
\item \emph{Fraction fields.} $\Frac A=F$ is proved for 03 and 06, which use a second
(reversed) arm to clear negative exponents; 08 explicitly does not claim it for its one-armed
controlled field, and 09 and 11 do not address it.
\item \emph{Hypotheses.} Regularity (06) and $\cf\kappa>\aleph_0$ (03) are used for simultaneous
gaps below fewer than $\kappa$ (respectively countably many) full supports; 08 avoids both through
its finite-coordinate invariant, which confines each element to finitely many coordinates.
\end{itemize}
\end{remark}

\section{Set-sized integer parts with exact thresholds}\label{osq:sec:models}

Five constructions replace the proper-class supply of scales by a cardinal. They are
different rings and are tagged by source. Each is an explicit subring of $\Oz$ (or of
$\GOz$), and each is a ZFC construction of sets. For every one of them
\cref{osq:thm:thresholds} gives the exact detection threshold; this section supplies the
constructions and the size computations it uses. Manuscript 05 (sibling report
\cite{RepoOdg}) asked which set-sized exponent groups admit analogues of the common-divisor
and set-target theorems (its ``exact size boundary'' question); 03, 06, 08, 09 and 11 each
answer it with a sharp family, and none claims a classification of all exponent groups or
support restrictions.

\subsection{Countable supports over a two-armed group \texorpdfstring{\src{03}}{[03]}}\label{osq:sec:model03}

For an ordered abelian group $G$ and $k\in\{\R,\C\}$ let $\Hc kG$ be the Hahn series over $k$
with \emph{at most countable} reverse well-ordered support (any countable reverse order type is
allowed; the subscript means countable support, not order type $\omega$). It is the
$\aleph_1$-bounded Hahn field in the terminology of \cite{LM}.

\begin{lemma}[Countable-support Hahn fields \src{03}]\label{osq:lem:hcfield}
$\Hc kG$ is a field. If $G$ is divisible, $\Hc\R G$ is real closed, $\Hc\C G$ is algebraically
closed, and $\Hc\C G=\Hc\R G[i]$.
\end{lemma}
\begin{proof}
Finitely many series lie in the full Hahn field over the countable subgroup $H$ generated by their
supports, all of whose supports are countable; so sums, products and inverses stay in $\Hc kG$.
For closedness, place the coefficient supports of a polynomial in their divisible hull, still
countable; the full Hahn field over a divisible group with real closed (algebraically closed)
coefficients is real closed (algebraically closed) \cite[Fact~2.1.1]{LM}, and it lies inside
$\Hc kG$. Real and imaginary parts split a countable support into two subsets.
\end{proof}

With $A(D,k,G)=D+\{f\in\Hc kG:\supp f\subseteq G_{>0}\}$, the constant term is a split surjection,
the units are $D^\times$ (\cref{osq:lem:units}), and for $(D,k)=(\Z,\R)$ the ring is a discretely
ordered integer part of $\Hc\R G$: write $x=P+r+\epsilon$ by positive, zero and negative exponents and
take $P+\lfloor r\rfloor$, except $P+r-1$ when $r\in\Z$ and $\epsilon<0$ (discarding all negative
exponents would send $-X^{-g}$ to $0$, which is not its floor). The same floor rule gives the omnific
integer part of $\No$ (03, 06, 08; compare \cite{RepoTrig}).

\paragraph{The construction.} Let $\kappa$ be an infinite cardinal with $\mu=\cf\kappa>\aleph_0$,
$\theta=\kappa^{\aleph_0}$, $J_\kappa=\kappa^{\mathrm{op}}+\kappa$ with indices $p_\alpha$ (reversed arm)
and $q_\alpha$ (forward arm), and
\[
 G_\kappa=\bigoplus_{j\in J_\kappa}\Q e_j,
\]
ordered lexicographically with the \emph{least} index of a nonzero vector deciding its sign; earlier
indices are larger scales ($ne_{j'}<e_j$ for $j<j'$). $\iota(e_{p_\alpha})=\omega^\alpha$,
$\iota(e_{q_\alpha})=\omega^{-(\alpha+1)}$ extends to an ordered additive embedding
$G_\kappa\hookrightarrow\No$, and $\sum a_gX^g\mapsto\sum a_g\omega^{\iota(g)}$ is an ordered field
embedding. Put
\[
 F^{03}_\kappa=\Hc\R{G_\kappa},\qquad A^{03}_\kappa=A(\Z,\R,G_\kappa)=F^{03}_\kappa\cap\Oz,\qquad
 \Pi^{03}_\kappa=\ker(\ct|_{A^{03}_\kappa}),
\]
and use the superscript $\C$ for the parallel construction over $\C$ and $\Z[i]$:
$F^{03,\C}_\kappa=F^{03}_\kappa[i]$, $A^{03,\C}_\kappa=A^{03}_\kappa[i]$,
$\Pi^{03,\C}_\kappa=\Pi^{03}_\kappa+i\Pi^{03}_\kappa$.

\begin{lemma}[Room on both sides \src{03}]\label{osq:lem:twosided}
\leavevmode
\begin{enumerate}
\item Every family of fewer than $\mu$ elements of $G_\kappa$ is bounded above and below in $G_\kappa$.
\item Every family of fewer than $\mu$ positive elements has a positive strict lower bound.
\item For every $h>0$, a final segment of the forward arm, of cardinality $\kappa$, consists of basis
vectors $e_{q_\alpha}\ll h$.
\end{enumerate}
\end{lemma}
\begin{proof}
Fewer than $\mu$ vectors have fewer than $\mu$ support indices, bounded below $\kappa$ in each arm. If
$\beta$ exceeds all reversed-arm ordinals occurring, $p_\beta$ precedes all occurring indices, so
$e_{p_\beta}$ exceeds the absolute value of each vector; this gives (i). If $\gamma$ exceeds all
forward-arm ordinals occurring, $q_\gamma$ follows every occurring index, and each positive vector,
having a positive leading coefficient at an earlier index, exceeds $e_{q_\gamma}>0$; this gives (ii),
and (iii) is the case of one vector.
\end{proof}

The reversed arm supplies a denominator large enough to shift a countable support upward; the forward
arm supplies a positive divisor below a countable positive support and many indistinguishable smaller
scales (03).

\begin{lemma}[Small families have common divisors \src{03}]\label{osq:lem:smallcommon}
Every family of fewer than $\mu$ elements of $\Pi^{03}_\kappa$ (or $\Pi^{03,\C}_\kappa$) has a common
divisor $X^g$, $g>0$, with all quotients in the ideal.
\end{lemma}
\begin{proof}
The union of fewer than $\mu$ countable supports has fewer than $\mu$ elements, $\mu$ being regular
and uncountable; apply \cref{osq:lem:twosided}(ii) and translate by $-g$.
\end{proof}

\begin{lemma}[Every smaller-scale field has size $\theta$ \src{03}]\label{osq:lem:card03}
For either coefficient pair, the field, the ring, its purely infinite ideal and every
$L_h=\Hc k{H_h(G_\kappa)}$ ($h>0$) have cardinality $\theta=\kappa^{\aleph_0}$.
\end{lemma}
\begin{proof}
There are at most $\kappa^{\aleph_0}$ countable supports and at most
$(2^{\aleph_0})^{\aleph_0}=2^{\aleph_0}\le\theta$ coefficient choices for each. Conversely, for a
forward tail $T\subseteq\kappa$ with $e_{q_\alpha}\ll h$ ($\alpha\in T$), each countable $S\subseteq T$
gives $f_S=\sum_{\alpha\in S}X^{e_{q_\alpha}}\in L_h$ (increasing ordinals give decreasing exponents),
and $L_h\supseteq k$; since $\kappa^{\aleph_0}=\max(2^{\aleph_0},|[\kappa]^{\le\aleph_0}|)$
(\cref{osq:lem:cardinalfacts}), $|L_h|\ge\theta$. Multiplication by $X^{h/2}$ embeds $L_h$ in the
ideal.
\end{proof}

\begin{theorem}[Cardinal realization \src{03}]\label{osq:thm:model03}
For $A=A^{03}_\kappa$ (or $A^{03,\C}_\kappa$) with ideal $\Pi_A$:
\begin{enumerate}
\item $|F^{03}_\kappa|=|A|=\theta$ and $\Frac(A)=F^{03}_\kappa$;
\item every unital map $A\to R$ with $|R|<\theta$ factors uniquely through $\ct$ (for the Gaussian
ring the factorization through the Gaussian constant term is unique for a given map; different
square roots of $-1$ give different maps);
\item every $A$-module of cardinality $<\theta$ is annihilated by $\Pi_A$, and for every
$0\ne f\in\Pi_A$ the least cardinality of a ring target detecting $f$, and of a module on which $f$
acts nontrivially, is exactly $\theta$;
\item $\Pi_A$ needs exactly $\mu$ generators; it is flat with $\Pi_A^2=\Pi_A$, the module $D=A/\Pi_A$
is not flat, and $\Tor^A_n(D,N)=0$ for every $D$-module $N$ and $n>0$.
\end{enumerate}
No inaccessible cardinal, CH or other additional axiom is used.
\end{theorem}
\begin{proof}
(i) Sizes by \cref{osq:lem:card03}. For $x\in F^{03}_\kappa$, \cref{osq:lem:twosided}(i) gives
$h\in G_\kappa$ larger than $-g$ for every exponent $g$ of $x$, so $x=X^hx/X^h$ with numerator and
denominator in $A$. (ii) and (iii) are \cref{osq:thm:thresholds}: by \cref{osq:lem:smallcommon} each
$f\in\Pi_A$ is $X^gb$ with $b\in\Pi_A$, and \cref{osq:prop:fieldbound} with \cref{osq:lem:card03}
gives the lower bound $\theta$ (03's own proof is exactly this). (iv) Generators: choose a cofinal
sequence $(\alpha_\xi)_{\xi<\mu}$ in $\kappa$; the countably many index supports of the exponents of
$f\in\Pi_A$ are bounded on the forward arm, so $X^{e_{q_{\alpha_\xi}}}$ divides $f$ in $A$ for large
$\xi$, and these $\mu$ monomials generate; a family of fewer than $\mu$ elements lies in some $X^hA$ by
\cref{osq:lem:smallcommon}, which misses $X^{h/2}$. Flatness, idempotence and Tor are
\cref{osq:thm:homological}.
\end{proof}

For $\kappa=\aleph_1$: $2^{\aleph_0}\le\aleph_1^{\aleph_0}\le(2^{\aleph_0})^{\aleph_0}=2^{\aleph_0}$,
so $A^{03}_{\aleph_1}$ has the size of the continuum, every target of size below the continuum sees only
the constant term, whether or not CH holds, and the ideal needs $\aleph_1$ generators. For every
infinite $\rho$, $\kappa=2^\rho$ gives $\theta=2^\rho$ with $\cf(2^\rho)>\rho$; so arbitrarily large
thresholds are realized in ZFC. This is a realization theorem for this family, not a classification
of ordered exponent groups. The case $\cf\kappa=\aleph_0$ is excluded because countable supports can
then run cofinally through an arm (03).

\begin{lemma}[Cardinal facts \src{03}]\label{osq:lem:cardinalfacts}
For every infinite $\kappa$, $\kappa^{\aleph_0}=\max(2^{\aleph_0},|[\kappa]^{\le\aleph_0}|)$; for every
infinite $\rho$, $(2^\rho)^{\aleph_0}=2^\rho$ and $\cf(2^\rho)>\rho$.
\end{lemma}
\begin{proof}
A countable sequence is determined by its countable range and a sequence in that range; conversely
binary sequences embed in $\kappa^\N$ and each countable subset is the range of an enumeration (choice).
$(2^\rho)^{\aleph_0}=2^{\rho\aleph_0}=2^\rho$. If $2^\rho$ were a union of $\tau\le\rho$ families
$E_\xi$ of binary functions on $\rho$, each of size $<2^\rho$, split $\rho$ into $\tau$ blocks of size
$\rho$ and on block $\xi$ choose a pattern that is not the restriction of any member of $E_\xi$; the
assembled function lies in no $E_\xi$ (K\"onig).
\end{proof}

\subsection{Supports of size \texorpdfstring{$<\kappa$}{< kappa} over a two-sided group \texorpdfstring{\src{06}}{[06]}}\label{osq:sub:model06}

Fix an uncountable regular cardinal $\kappa$. In the additive group of $\No$ put
\[
 h_\alpha=\omega^{\alpha+1},\qquad\epsilon_\alpha=\omega^{-(\alpha+1)}\quad(\alpha<\kappa),\qquad
 \Gamma^{06}_\kappa=\Span_\Q\{h_\alpha,\epsilon_\alpha:\alpha<\kappa\}
\]
(finite rational combinations of distinct Conway monomials, hence linearly independent). This is a
divisible ordered group of cardinality $\kappa$, and $\epsilon_\beta\ll\epsilon_\alpha\ll h_\alpha\ll
h_\beta$ for $\alpha<\beta$; here $\ll$ refers to the outer exponents.

\begin{lemma}[Simultaneous upper and lower scales \src{06}]\label{osq:lem:kappagap}
For every $T\subseteq\Gamma^{06}_\kappa$ with $|T|<\kappa$ there is $\beta<\kappa$ with
$\epsilon_\beta\ll|t|\ll h_\beta$ for all $t\in T\setminus\{0\}$.
\end{lemma}
\begin{proof}
The generators occurring in the finite expressions of the members of $T$ have fewer than $\kappa$
indices, bounded below some $\beta$ by regularity. The leading generator of a nonzero expression
determines its Archimedean class, which lies strictly between those of $\epsilon_\beta$ and
$h_\beta$.
\end{proof}

For $k\in\{\R,\C\}$ put $K^{06}_{\kappa,k}=\{f\in\hahn k{\Gamma^{06}_\kappa}:|\supp f|<\kappa\}$,
$K^{06}_\kappa=K^{06}_{\kappa,\R}$, $\Pi^{06}_\kappa=\{f\in K^{06}_\kappa:\supp f>0\}$ and
$A^{06}_\kappa=\Z+\Pi^{06}_\kappa$ (06's $R_\kappa$); $X^g\mapsto\omega^g$ makes them subsets of $\No$
and $\SC$.

\begin{proposition}[Field closure and integer part \src{06}]\label{osq:prop:fieldclosure}
$K^{06}_{\kappa,\R}$ is real closed, $K^{06}_{\kappa,\C}=K^{06}_{\kappa,\R}[i]$ is algebraically
closed, $A^{06}_\kappa=K^{06}_\kappa\cap\Oz$, $\Frac(A^{06}_\kappa)=K^{06}_\kappa$, and $A^{06}_\kappa$
is a discretely ordered integer part of $K^{06}_\kappa$.
\end{proposition}
\begin{proof}
Inverting a series of support size $\mu<\kappa$ (factor the leading term and use the geometric
series) gives support in a translate of the monoid generated by at most $\max(\mu,\aleph_0)<\kappa$
elements; sums and products also keep the bound. For closedness, the coefficients of a polynomial use
fewer than $\kappa$ exponents; the rational span $\Delta$ of these has $|\Delta|<\kappa$, the full field
$\hahn\C\Delta$ is algebraically closed \cite[Corollary~4]{Poonen} and lies in $K^{06}_{\kappa,\C}$; the
real field with its leading-term order has this algebraically closed complexification, so it is real
closed. For $f\in K^{06}_\kappa$, \cref{osq:lem:kappagap} gives $h_\beta$ above $|g|$ for every
exponent $g$, so $X^{h_\beta}f\in\Pi^{06}_\kappa$ and $f$ is a quotient. The floor rule of
\cref{osq:sec:model03} gives the integer part, and a nonzero element with positive leading exponent
exceeds every real, so there is nothing strictly between $0$ and $1$.
\end{proof}

Field closure uses only that $\kappa$ is uncountable; regularity is used for simultaneous gaps and
generation (06). The first-kappa report proves the closedness of $\{f\in k(\!(t^\Gamma)\!):|\supp f|<\kappa\}$
for every uncountable $\kappa$, possibly singular (\nolinkurl{fkc:thm:main} in \cite{RepoFKC}); 06's
closure argument is not a novelty claim, and singular support bounds may preserve field closure while
failing the simultaneous gap property used here. The two families $h_\alpha$ and $\epsilon_\alpha$ are
deliberate: the $\epsilon_\alpha$ provide smaller positive gaps and the $h_\alpha$ shift supports upward,
and a one-sided family would not automatically give the fraction-field assertion (06).

\begin{proposition}[A chain of principal ideals \src{06}]\label{osq:prop:kappachain}
$\Pi^{06}_\kappa=\bigcup_{\alpha<\kappa}X^{\epsilon_\alpha}A^{06}_\kappa$ is a strictly increasing union;
the ideal is idempotent and its least number of generators is $\kappa$.
\end{proposition}
\begin{proof}
For $\alpha<\beta$, $\epsilon_\alpha-\epsilon_\beta>0$ gives the inclusion, strict because the reverse
quotient has a negative exponent. A gap $\epsilon_\alpha<\supp f$ (\cref{osq:lem:kappagap}) gives
$f\in X^{\epsilon_\alpha}\Pi^{06}_\kappa$, hence the union and idempotence. Fewer than $\kappa$ elements
have fewer than $\kappa$ support exponents in total (regularity), so they lie in some $X^aA^{06}_\kappa$,
which misses $X^{a/2}$.
\end{proof}

\begin{lemma}[Size of the ring and of every smaller-scale field \merge]\label{osq:lem:card06}
For every $c>0$ in $\Gamma^{06}_\kappa$, the field $L_c=\{f\in K^{06}_\kappa:\supp f\subseteq
H_c(\Gamma^{06}_\kappa)\}$ has cardinality at least $\kappa^{<\kappa}$, and
\[
 |K^{06}_{\kappa,\R}|=|K^{06}_{\kappa,\C}|=|A^{06}_\kappa|=|\Pi^{06}_\kappa|=\kappa^{<\kappa}.
\]
\end{lemma}
\begin{proof}
The finite expression of $c$ uses indices below some $\alpha_0<\kappa$, and its Archimedean class is
that of its leading generator; so $\epsilon_\alpha\ll c$ for every $\alpha>\alpha_0$. For
$B\subseteq(\alpha_0,\kappa)$ with $|B|<\kappa$, the series $\sum_{\alpha\in B}X^{\epsilon_\alpha}$ has
support $\{\epsilon_\alpha:\alpha\in B\}$, reverse well ordered because $\epsilon_\alpha$ decreases as
$\alpha$ increases, of size $|B|<\kappa$, and inside $H_c$; so it lies in $L_c$, which is a field by
\cref{osq:prop:fieldbound}(i). Distinct $B$ give distinct series, and the number of subsets of size
$<\kappa$ of a set of size $\kappa$ is $\sum_{\nu<\kappa}\kappa^{\nu}=\kappa^{<\kappa}$. Hence
$|L_c|\ge\kappa^{<\kappa}$, and $X^cL_c\subseteq\Pi^{06}_\kappa$ gives the same lower bound for the ideal,
ring and fields. Conversely a series is given by a support of size $\nu<\kappa$ and a coefficient function:
at most $\kappa^\nu\cdot(2^{\aleph_0})^\nu$ choices, which is $\le\kappa^\nu$ for infinite $\nu$ and
$\le\kappa\cdot2^{\aleph_0}\le\kappa^{\aleph_0}\le\kappa^{<\kappa}$ for finite $\nu$ (as $\kappa>\aleph_0$). Summing over
the $\le\kappa$ values of $\nu$ gives at most $\kappa^{<\kappa}$.
\end{proof}

\begin{theorem}[The exact threshold of $A^{06}_\kappa$ \src{06} \merge]\label{osq:thm:model06}
Let $\kappa$ be uncountable and regular. Every homomorphism $A^{06}_\kappa\to S$ with
$|S|<\kappa^{<\kappa}$ factors through $\ct$, every $A^{06}_\kappa$-module $M$ with $|M|<\kappa^{<\kappa}$
satisfies $\Pi^{06}_\kappa M=0$, and for each $0\ne x\in\Pi^{06}_\kappa$ the least cardinality of a ring
target detecting $x$, and of a module on which $x$ acts nontrivially, is exactly $\kappa^{<\kappa}$. The
same holds for $A^{06}_\kappa[i]$ with $\Z[i]$.
\end{theorem}
\begin{proof}
\cref{osq:thm:thresholds} with \cref{osq:lem:kappagap,osq:lem:card06}: each $0\ne x\in\Pi^{06}_\kappa$ has a
gap $\epsilon_\beta<\supp x$, and with $c=\epsilon_\beta/2$ the field $L_c$ has cardinality at least
$\kappa^{<\kappa}=|A^{06}_\kappa|$.
\end{proof}

06 proves the lower bound $\kappa$ (with the explicit telescope family of the $\kappa$ monomials
$X^{a+\epsilon_\alpha}$, each witness having countable support) and, under the extra hypothesis
$\kappa^{<\kappa}=\kappa$, that $|A^{06}_\kappa|=\kappa$ and the threshold is exactly $\kappa$; that is
the special case of the theorem where $\kappa^{<\kappa}=\kappa$. The lower bounds of 06 need no cardinal
arithmetic, and neither does the exact value $\kappa^{<\kappa}$, which comes from 06's own collision lemma
(compare 03's field bound). This answers the first half of 06's question ``Cardinal arithmetic without the
matching hypothesis'': the least size of a target detecting a fixed nonzero $x\in\Pi^{06}_\kappa$ is
$\kappa^{<\kappa}$ in ZFC. The second half (``a different integer-part construction of size exactly
$\kappa$ with the same lower-bound property under weaker hypotheses'') is answered by 08 in ZFC for every
infinite $\kappa$ (\cref{osq:thm:model08}), for a ring of size exactly $\kappa$ whose fraction field,
unlike $\Frac A^{06}_\kappa=K^{06}_\kappa$, is not claimed to be its ambient field.

\begin{example}[The first uncountable cardinal \src{06}]\label{osq:ex:CH}
Under CH, $\aleph_1^{<\aleph_1}=\aleph_1^{\aleph_0}=2^{\aleph_0}=\aleph_1$, so $A^{06}_{\aleph_1}$ has
size $\aleph_1$: every action on a countable abelian group is an ordinary integer action, while the regular
module of size $\aleph_1$ sees all purely infinite elements. Without CH the lower bound at $\aleph_1$ holds,
the ring has size $\aleph_1^{\aleph_0}=2^{\aleph_0}$, and by \cref{osq:thm:model06} the exact threshold is
$2^{\aleph_0}$; this ring is isomorphic to $A^{03}_{\aleph_1}$ (\cref{osq:rem:reconcile}).
\end{example}

\subsection{A controlled field of size exactly \texorpdfstring{$\kappa$}{kappa} \texorpdfstring{\src{08}}{[08]}}\label{osq:sub:model08}

Let $\kappa$ be any infinite cardinal and $\Gamma^{08}_\kappa=\bigoplus_{\xi<\kappa}\Q e_\xi$, ordered
lexicographically with the least nonzero index deciding the sign, so $0<ne_\eta<e_\xi$ for $\xi<\eta$;
it is divisible of cardinality $\kappa$, and each nonzero element has a finite coordinate support.

\begin{lemma}[Scale and support bounds \src{08}]\label{osq:lem:gammabounds}
For every $c>0$ there are $\kappa$ positive $a\ll c$ (namely $e_\xi$ for $\xi$ beyond the least nonzero
index of $c$). If $S\subseteq(\Gamma^{08}_\kappa)_{>0}$ has fewer than $\cf\kappa$ elements, some
$\beta<\kappa$ has $e_\beta\ll s$ for all $s\in S$.
\end{lemma}
\begin{proof}
A union of fewer than $\cf\kappa$ finite subsets of $\kappa$ is bounded; take $\beta$ above it.
\end{proof}

Let $K=\R\cap\overline\Q$, $H=\hahn K{\Gamma^{08}_\kappa}$ and $H_{\mathrm{ctbl}}\subseteq H$ the series
with at most countable support. $H_{\mathrm{ctbl}}$ is a real closed subfield: finitely many series lie in
$\hahn K\Delta$ for the countable divisible subgroup $\Delta$ generated by their supports, which is real
closed \cite[Fact~2.1.1]{LM} and has only countable supports. A \emph{truncation} of a series is its
restriction to an initial segment of its support in decreasing order; a countable support has countably
many. Let $F_0$ be the real closure in $H_{\mathrm{ctbl}}$ of the field generated by $K$ and all
$X^g$; given $F_n$, adjoin all truncations of its elements and let $F_{n+1}$ be the real closure in
$H_{\mathrm{ctbl}}$ of the resulting field; put $F=\bigcup_nF_n$.

\begin{lemma}[The controlled field \src{08}]\label{osq:lem:controlled}
$F$ is real closed and truncation-closed, has cardinality $\kappa$, consists of series with countable
support, and contains $K$ and every monomial.
\end{lemma}
\begin{proof}
An algebraic extension of an infinite field of size $\kappa$ has size $\le\kappa$ ($\kappa$ polynomials,
finitely many roots each), and adjoining countably many truncations of each of $\kappa$ elements adds at
most $\kappa$ elements; so every $F_n$, and $F$, has size $\kappa$, and at least $\kappa$ because of the
monomials. Truncations of an element of $F_n$ lie in $F_{n+1}$, and finitely many coefficients lie in a
common real closed stage.
\end{proof}

\begin{lemma}[Finite-coordinate invariant \src{08}]\label{osq:lem:finitecoord}
Every $f\in F$ has $\supp f\subseteq\Gamma_E=\bigoplus_{\xi\in E}\Q e_\xi$ for some finite $E\subseteq\kappa$
(its support need not be finite).
\end{lemma}
\begin{proof}
For finite $E$, $H_E=\hahn K{\Gamma_E}$ is real closed and lies in $H_{\mathrm{ctbl}}$ ($\Gamma_E$ is
countable); the $H_E$ form a directed family. An element of the field generated by the monomials is a
rational expression in finitely many of them, so lies in some $H_E$. An element of $F_0$ is algebraic over
that field, hence over some $H_E$ (finitely many coefficients); as $H_E$ is real closed and
$H_{\mathrm{ctbl}}$ is an ordered field containing it, it has no proper ordered algebraic extension there,
so the element lies in $H_E$. Truncations of elements of $H_E$ stay in $H_E$, finite rational expressions
stay in $H_{E\cup E'}$, and real closure is handled by the same algebraicity argument; induct over $n$.
\end{proof}

This invariant removes regularity: even at a singular $\kappa$ of countable cofinality, the coordinates in
one element are confined to a finite set. The proof uses countable supports to control the number of
truncations and finite coordinate sets to find smaller positive scales (08).

The map $\iota(\sum q_\xi e_\xi)=\sum q_\xi\omega^{-\xi}$ is an increasing additive embedding of
$\Gamma^{08}_\kappa$ into $\No$ (the least nonzero coordinate decides the sign on both sides), and
$\Phi(\sum a_gX^g)=\sum a_g\omega^{\iota(g)}$ is an ordered field embedding $H\hookrightarrow\No$
preserving truncations and constant terms. Put $F^{08}_\kappa=\Phi(F)$,
$A^{08}_\kappa=F^{08}_\kappa\cap\Oz=\{f\in F:\supp f\ge0,\ \ct(f)\in\Z\}$ and
$\Pi^{08}_\kappa=\ker(\ct|_{A^{08}_\kappa})$. Since $\iota(e_0)=1$, the monomial $X^{e_0}$ is the actual
surreal $\omega$. $A^{08}_\kappa$ is a discretely ordered integer part of $F$ of cardinality $\kappa$:
truncate $x$ to its positive part $P\in F$ and apply the floor rule of \cref{osq:sec:model03} to
$x=P+r+\epsilon$ with $r\in K$.

\begin{theorem}[Sharp models at every infinite cardinal \src{08}]\label{osq:thm:model08}
For every infinite cardinal $\kappa$:
\begin{enumerate}
\item $F^{08}_\kappa$ is a truncation-closed real closed subfield of $\No$ with countable supports,
$|F^{08}_\kappa|=|A^{08}_\kappa|=\kappa$, and $A^{08}_\kappa$ is an integer part of it;
\item every homomorphism from $A^{08}_\kappa$ to a ring of cardinality $<\kappa$ factors through $\ct$;
\item every quotient ring of $A^{08}_\kappa$ is a quotient $\Z/n\Z$ ($n\ge0$) under $\ct$ or has cardinality
exactly $\kappa$;
\item $\Pi^{08}_\kappa$ needs exactly $\cf\kappa$ generators;
\item every unital $A^{08}_\kappa$-module of cardinality $<\kappa$ has action $f\cdot m=\ct(f)m$;
\item the positive-characteristic maximal residue fields are the $\F_p$, with maximal ideals $pA^{08}_\kappa$;
characteristic-zero maximal residue fields exist and all have cardinality exactly $\kappa$.
\end{enumerate}
No hypothesis such as CH, GCH, $\kappa^{\aleph_0}=\kappa$ or an
inaccessible cardinal is used.
\end{theorem}
\begin{proof}
(i) is \cref{osq:lem:controlled} and the paragraph above. Each $f\in\Pi^{08}_\kappa$ has support in some
$\Gamma_E$ (\cref{osq:lem:finitecoord}), so $e_\beta$, $\beta>\max E$, is a positive strict lower bound; and
\cref{osq:lem:gammabounds} gives $\kappa$ scales below every $c$. These are the hypotheses of
\cref{osq:thm:criterion}, which gives (ii).
(iii) If $|A/J|<\kappa$, (ii) applied to $A\to A/J$ gives $\Pi\subseteq J$; otherwise $|A/J|\le|A|=\kappa$.
(iv) With a cofinal $C\subseteq\kappa$ of size $\cf\kappa$, for $f\in\Pi^{08}_\kappa$ choose $\beta\in C$ above
its coordinate set; then $X^{-e_\beta}f\in\Pi^{08}_\kappa$, so $\{X^{e_\beta}:\beta\in C\}$ generates.
Fewer than $\cf\kappa$ elements have coordinate sets bounded by some $\beta$, so they lie in $(X^{e_\beta})$,
which misses $X^{e_\beta/2}$. (v) is \cref{osq:lem:cyclic} with (ii). (vi) Coefficient division gives
$n\Pi^{08}_\kappa=\Pi^{08}_\kappa$, so $nA=n\Z+\Pi^{08}_\kappa$ and $A/pA\cong\F_p$ is already a field; a maximal
ideal of characteristic $p$ contains $pA$ and equals it. The element $u=1+X^{e_0}$ is a nonunit
(\cref{osq:lem:units}), so the set-sized ring $A$ has a maximal ideal $M\ni u$ (ordinary maximal-ideal
theorem); if a nonzero integer $n$ lay in $M$, then $X^{e_0}\in n\Pi^{08}_\kappa\subseteq M$ and $1\in M$. A
characteristic-zero residue field of size $<\kappa$ would receive a map through $\Z$ by (ii), with image $\Z$,
not a field; and it has size $\le\kappa$ as a quotient.
\end{proof}

At $\kappa=\aleph_0$ the theorem restricts only finite targets, and the characteristic-zero residue fields are
countable and detect infinite monomials. For regular $\kappa$ the generator count is $\kappa$; at a singular
$\kappa$ it is smaller, so a countably generated ideal can be invisible to all ring targets below a much larger
singular $\kappa$ of countable cofinality. Two resources are separated: monomial generators coinitial in the
support scales ($\cf\kappa$) and pairwise collision witnesses below a fixed scale ($\kappa$) (08).

In the residue field of a maximal ideal containing $1+\omega$ the image of $\omega$ is $-1$, not $0$, and that
field has size $\kappa$; so for every infinite $\kappa$ the subring $A^{08}_\kappa\subseteq\Oz$ has a
characteristic-zero field quotient of cardinality $\kappa$ in which $\omega\mapsto-1$. This quotient map cannot
extend to $\Oz$, where every set-sized image sends $\omega$ to $0$: an explicit extension obstruction between a
set-sized integer part and the class ring (08). 08 does not identify $\Frac(A^{08}_\kappa)$ with
$F^{08}_\kappa$, imposes no initial-segment condition or birthday cutoff on the embedding, and designs the
construction for truncation, support control and cardinality, not for closure under a derivation or an
exponential; the isomorphism types of the residue fields are not classified.

\subsection{Countable supports over a long one-armed group \texorpdfstring{\src{09}}{[09]}}\label{osq:sub:model09}

Fix a cardinal $\lambda$ with $\lambda^{\aleph_0}=\lambda$; then automatically $\lambda\ge2^{\aleph_0}$,
since $\lambda^{\aleph_0}\ge2^{\aleph_0}$ (09 states $\lambda\ge2^{\aleph_0}$ as a separate hypothesis; it is
redundant). For example $\lambda=2^\nu$ for any infinite $\nu$. Neither regularity nor CH is assumed. Let
$L_\lambda=\omega_1\times\lambda$ with the lexicographic order and
$\Gamma^{09}_\lambda=\bigoplus_{\ell\in L_\lambda}\Q e_\ell$, the least nonzero index deciding the sign.

\begin{lemma}[A large small-exponent tail \src{09}]\label{osq:lem:tail}
For every countable $S\subseteq(\Gamma^{09}_\lambda)_{>0}$ there is $\alpha_0<\omega_1$ with
$e_{(\alpha,\beta)}\in U_{\Gamma}(S)$ for all $\alpha_0\le\alpha<\omega_1$ and $\beta<\lambda$; in particular
$U_\Gamma(S)$ contains $\lambda$ elements, and every countable positive set has a positive strict lower bound.
\end{lemma}
\begin{proof}
The countable union of the finite index supports has first coordinates bounded below some $\alpha_0<\omega_1$;
each $s\in S$ has its least index before every $(\alpha,\beta)$ with $\alpha\ge\alpha_0$, with positive
coefficient there.
\end{proof}

For $k\in\{\R,\C\}$ and unital $D\subseteq k$ let $H_{\lambda,k}$ be the countable-support Hahn field over
$\Gamma^{09}_\lambda$, $\Pi^{09}_{\lambda,k}$ its positive-support part and
$A^{09}_{\lambda,D}=D+\Pi^{09}_{\lambda,k}$. It is a domain with units $D^\times$, the quotients of
\cref{osq:lem:division} of its elements stay in it (countable supports), and $|A^{09}_{\lambda,D}|=\lambda$:
there are at most $\lambda^{\aleph_0}=\lambda$ countable supports and $|k|^{\aleph_0}=2^{\aleph_0}\le\lambda$
coefficient choices, and $\lambda$ positive basis monomials. With $j(\alpha,\beta)=\lambda\cdot\alpha+\beta$
(ordinal operations), $\iota(\sum q_\ell e_\ell)=\sum q_\ell\omega^{-j(\ell)}$ is an ordered embedding and
$\sum c_\gamma X^\gamma\mapsto\sum c_\gamma\omega^{\iota(\gamma)}$ embeds $H_{\lambda,\R}$ in $\No$ and
$H_{\lambda,\C}$ in $\SC$, sending $A^{09}_{\lambda,\Z}$ into $\Oz$ and $A^{09}_{\lambda,\Z[i]}$ into $\GOz$;
here $X^{e_\ell}\mapsto\omega^{\omega^{-j(\ell)}}$, the group coordinates describing exponents.

\begin{theorem}[Cardinal barrier \src{09}]\label{osq:thm:barrier}
Let $A=A^{09}_{\lambda,D}$ and $\Pi_A=\Pi^{09}_{\lambda,k}$.
\begin{enumerate}
\item If $\varphi:A\to B$ is a unital ring map with $\varphi(\Pi_A)\ne0$, then $|\varphi(A)|\ge\lambda$.
\item If $M$ is an $A$-module with $\Pi_AM\ne0$, then $|M|\ge\lambda$.
\item All maps to rings of size $<\lambda$, and all modules of size $<\lambda$, factor through $\ct$.
\item For every ideal $J$, either $\Pi_A\subseteq J$ and $A/J\cong D/\ct(J)$, or $|A/J|=\lambda$.
\end{enumerate}
In particular, in ZFC there is an explicit $A\subseteq\Oz$ of cardinality $2^{\aleph_0}$ (take
$\lambda=2^{\aleph_0}$, $D=\Z$) whose maps to countable rings factor through $\ct$, whose quotients are finite,
countably infinite, or of cardinality $2^{\aleph_0}$ (the last for every quotient not induced from $\Z$), and
whose countable modules are abelian groups with constant-term action.
\end{theorem}
\begin{proof}
For $f\in\Pi_A$ detected by $\varphi$ or by a vector, $\supp f$ is countable, \cref{osq:lem:tail} supplies
$\lambda$ small exponents, and \cref{osq:lem:separation} with \cref{osq:lem:division} injects them into
$\varphi(A)$ (with the module structure given by $\varphi$ and the vector $1$) or into $M$. (iii) follows,
and (iv) as in \cref{osq:thm:thresholds}.
\end{proof}

When the continuum exceeds $\aleph_1$ the last statement also excludes all intermediate quotient cardinalities;
this is conditional only on that cardinal inequality (09).

\begin{lemma}[Recovering the coefficient field \src{09}]\label{osq:lem:coefffield}
Let $\varphi:A\to F$ be a unital map to a field with $\varphi(\Pi_A)\ne0$. There is a unique embedding
$\sigma:k\hookrightarrow F$ with $\varphi(rg)=\sigma(r)\varphi(g)$ for $r\in k$, $g\in\Pi_A$, and
$\sigma|_D=\varphi|_D$; in particular $F$ has characteristic $0$.
\end{lemma}
\begin{proof}
Take $f\in\Pi_A$ with $\varphi(f)\ne0$ and $\sigma(r)=\varphi(rf)/\varphi(f)$ (note $rf\in\Pi_A$ even if
$r\notin A$). It is additive and unital, $(rf)(sf)=(rsf)f$ gives multiplicativity, and it is injective as a
map of fields. The identity $(rg)f=(rf)g$ gives the displayed formula and independence of $f$.
\end{proof}

\begin{theorem}[Residue fields at the threshold \src{09}]\label{osq:thm:residue}
Let $A=A^{09}_{\lambda,D}$.
\begin{enumerate}
\item For every $0\ne f\in\Pi_A$ there is a maximal ideal $\mathfrak m$ with $f\equiv1\pmod{\mathfrak m}$ and
$|A/\mathfrak m|=\lambda$.
\item Every maximal ideal not containing $\Pi_A$ has a residue field of cardinality $\lambda$ with a canonical
copy of $k$.
\item For $D=\Z$, the maximal ideals containing $\Pi_A$ are the $pA$, with residue fields $\F_p$, and every other
maximal ideal has a characteristic-zero residue field of cardinality $\lambda$.
\item There are at least $\lambda$ distinct maximal ideals not containing $\Pi_A$.
\end{enumerate}
Consequently the least cardinality of a field (and of a module) detecting a given $0\ne f\in\Pi_A$ is exactly
$\lambda$; and for $D=\Z$ the intersection of all finite-index ideals is $\Pi_A\ne0$ while $\Jac(A)=0$.
\end{theorem}
\begin{proof}
(i) $1-f$ is a nonconstant nonunit; a maximal ideal of the set-sized ring $A$ containing it sends $f$ to $1$, so
the residue field has size $\ge\lambda$ by \cref{osq:thm:barrier} and $\le\lambda$ as a quotient. (ii) is
\cref{osq:thm:barrier} with \cref{osq:lem:coefffield}; (iii) uses $pA=p\Z+\Pi_A$. (iv) For $\ell<\ell'$ in
$L_\lambda$ put $\gamma=e_\ell$, $a=e_{\ell'}$; then $na<\gamma$ for all $n$, and
$q=-\sum_{n\ge1}X^{\gamma-na}\in\Pi_A$ satisfies $(1-X^a)q=X^\gamma$, so
$1=(1-X^\gamma)+(1-X^a)q$: the $\lambda$ proper ideals $(1-X^{e_\ell})$ are pairwise comaximal, and maximal
ideals above them (set-sized choice) are distinct and do not contain $\Pi_A$. For the consequences, the
residue field of (i) is a field target and a module on which $f$ acts as the identity. Every finite target has
size $<\lambda$, so the finite residual is $\Pi_A$; if $0\ne x=n+f$, a prime not dividing $n\ne0$, or (i) when
$n=0$, gives a maximal ideal avoiding $x$.
\end{proof}

All of this concerns set-sized rings and is not an application of the maximal-ideal theorem to $\Oz$. The
distinction between finite ring probes and field probes is sharp: the extra information is carried by residue
fields of the full workspace cardinality (09). The number of maximal ideals and the isomorphism types of the
residue fields are not determined; distinct maximal ideals need not give nonisomorphic fields.
$\Pi^{09}_{\lambda,k}=\bigcup_{\alpha<\omega_1}X^{e_{(\alpha,0)}}A$ is an increasing chain, and the least number of
generators of the ideal is exactly $\aleph_1$ (the chain gives $\aleph_1$ generators; a countable family lies in
$X^{2h}A$ for a small $h$ by \cref{osq:lem:tail}, which misses $X^h$); flatness is in
\cref{osq:thm:homological}.

\begin{remark}[Why the support restriction is real \src{09}]\label{osq:rem:fstar}
In the unrestricted Hahn field over $\Gamma^{09}_\lambda$ the series $f_*=\sum_{\alpha<\omega_1}X^{e_{(\alpha,0)}}$
has reverse well-ordered but uncountable support with no positive lower bound in $\Gamma^{09}_\lambda$ (a late
enough $e_{(\alpha,0)}$ lies below any given $h>0$); it is not divisible by any positive monomial within the
nonnegative-support ring and is not in $A^{09}_{\lambda,D}$. No assertion is made about the quotient
classification of that larger unrestricted ring.
\end{remark}

\subsection{Supports of size \texorpdfstring{$<\kappa$}{< kappa} over a one-armed group \texorpdfstring{\src{11}}{[11]}}\label{osq:sub:model11}

\Cref{osq:thm:coinitial} (11) applies to any divisible set group $\Gamma$ with $\coin(\Gamma_{>0})\ge\kappa$. The
coinitiality condition says that every set of fewer than $\kappa$ positive elements has a positive strict lower
bound.

\begin{proposition}[Long lexicographic groups \src{11}]\label{osq:prop:lex}
For $\kappa$ uncountable and regular and $\Gamma=\bigoplus_{\alpha<\kappa}\Q e_\alpha$ with the least nonzero
index dominant, $|\Gamma|=\kappa$ and $\coin(\Gamma_{>0})=\kappa$. The group $\Gamma^{06}_\kappa$ also has
coinitiality $\kappa$, so $A^{06}_\kappa=R^{11}_{\Gamma^{06}_\kappa,<\kappa}$ with $(D,K)=(\Z,\R)$ is an instance
of \cref{osq:thm:coinitial}.
\end{proposition}
\begin{proof}
The $e_\beta$ with $\beta$ beyond the least index of a positive vector lie below it, so $(e_\beta)_{\beta<\kappa}$
is coinitial; fewer than $\kappa$ positive vectors have boundedly many leading indices (regularity), and a later
basis vector lies below all of them. For $\Gamma^{06}_\kappa$ use \cref{osq:lem:kappagap}.
\end{proof}

11 proves sharpness under the displayed extra hypotheses $|K|\le\kappa$ and $\kappa^{<\kappa}=\kappa$: then
$|R^{11}_{\Gamma,<\kappa}|=\kappa$ and its identity map is a target of size $\kappa$ not factoring through $\ct$.
11 states explicitly that these hypotheses are assumptions for the sharpness example, not assertions about all
uncountable regular cardinals, and that regularity enters only the lexicographic example.

\begin{proposition}[Exact size and threshold \merge]\label{osq:prop:card11}
In \cref{osq:prop:lex}, for every set field $K$ and every $c>0$ the field
$L_c=\{f\in\hahn K\Gamma:|\supp f|<\kappa,\ \supp f\subseteq H_c(\Gamma)\}$ has at least $\kappa^{<\kappa}$
elements. If $|K|\le\max(\kappa,2^{\aleph_0})$ (which includes 11's hypothesis $|K|\le\kappa$), then
$|R^{11}_{\Gamma,<\kappa}|=\kappa^{<\kappa}$, and the least cardinality of a
ring target or module detecting any nonzero element of the positive-support ideal is exactly $\kappa^{<\kappa}$.
\end{proposition}
\begin{proof}
Exactly as in \cref{osq:lem:card06}: the series $\sum_{\alpha\in B}X^{e_\alpha}$, $B$ a subset of size $<\kappa$
of the indices beyond the leading index of $c$, lie in $L_c$ ($e_\alpha$ decreases with $\alpha$), and the
$<\kappa$-support series form a field because inverses have support in a monoid generated by fewer than $\kappa$
elements. The upper bound counts supports of size $\nu<\kappa$ and coefficient functions: for infinite $\nu$,
$|K|^\nu\le\max(\kappa^\nu,2^\nu)=\kappa^\nu$, and for finite $\nu$,
$|K|^\nu\le\max(\kappa,2^{\aleph_0})\le\kappa^{\aleph_0}\le\kappa^{<\kappa}$; also $|D|\le|K|$, and the
support choices number $\kappa^\nu$.
The threshold follows from \cref{osq:prop:fieldbound}, the coinitiality gap, and the identity map.
\end{proof}

So the sharpness of \cref{osq:thm:coinitial} in the lexicographic example holds without the hypothesis
$\kappa^{<\kappa}=\kappa$ once the bound is stated as $\kappa^{<\kappa}$, which equals $\kappa$ exactly when
$\kappa^{<\kappa}=\kappa$; the coefficient field is still assumed to satisfy $|K|\le\max(\kappa,2^{\aleph_0})$.
With $(D,K)=(\Z,\R)$ this ring is a Hahn integer part of the kind treated in this section; 11 does not
embed it in $\No$, and no embedding is used here.

\section{Flat ideals and homological reflection}\label{osq:sec:homological}

All rings, modules, resolutions, tensor products and $\Tor$, $\Ext$ groups in this section are
ordinary set-sized objects over the set-sized rings of \cref{osq:sec:models}. No class-sized
resolution is formed, and no directed-colimit theorem is applied to the proper-class union
$\Pi=\bigcup_{a>0}\omega^a\Oz$ (03, 06, 09). The package is 06's, proved there for $A^{06}_\kappa$;
03 and 09 prove its flatness and $\Tor$ parts for their rings. The proofs use only the following
properties, which all five constructions have:

\begin{enumerate}
\item[(H1)] $A$ is a commutative domain, $A=D\oplus I$ with $I=\ker\ct$ an ideal and $D=A/I$;
\item[(H2)] $A$ is built over a divisible exponent group $\Gamma$, and every $f\in I$ has the form
$f=X^gb$ with $g\in\Gamma_{>0}$ and $b\in I$; hence $I=\bigcup_{h\in\Gamma_{>0}}X^hA$, a directed
set-indexed union of principal ideals;
\item[(H3)] $I$ is not finitely generated.
\end{enumerate}
(H2) is the gap property used in \cref{osq:thm:thresholds} (a positive strict lower bound $g$ of
$\supp f$ gives $b=X^{-g}f\in I$). (H3) holds because the least number of generators is $\cf\kappa$,
$\kappa$, $\cf\kappa$, $\aleph_1$, respectively $\kappa$ (\cref{osq:thm:model03}(iv),
\cref{osq:prop:kappachain}, \cref{osq:thm:model08}(iv), the chain after \cref{osq:thm:residue}, and for
$R^{11}_{\Gamma,<\kappa}$ the coinitial chain $(X^{e_\alpha})_{\alpha<\kappa}$ of \cref{osq:prop:lex}:
fewer than $\kappa$ elements have a common monomial divisor $X^h$, and $X^hA$ misses $X^{h/2}$); these
counts are infinite.

\begin{lemma}[Projective ideals of a domain \src{06}]\label{osq:lem:projideal}
A nonzero projective ideal of a commutative domain is finitely generated.
\end{lemma}
\begin{proof}
Fix $0\ne a\in J$. For $A$-linear $f:J\to A$ and $x\in J$, $af(x)=f(ax)=xf(a)$, so $f(a)=0$ forces
$f=0$. A dual basis $(u_s,f_s)$ has $x=\sum_sf_s(x)u_s$ with finitely many nonzero terms for each $x$;
only finitely many $f_s(a)$ are nonzero, all other $f_s$ vanish identically, and the finitely many
corresponding $u_s$ generate $J$.
\end{proof}

\begin{theorem}[Homological package \src{03, 06, 09} \merge]\label{osq:thm:homological}
Let $(A,I,D)$ satisfy (H1)--(H3). Then:
\begin{enumerate}
\item $I$ is flat, $I^2=I$ (each element is a product of two elements of $I$), and $I$ is not projective;
\item for every $D$-module $N$ (inflated along $\ct$), $I\otimes_AN=0$ and $\Tor^A_n(D,N)=0$ for
$n>0$; in particular $\Tor^A_n(D,D)=0$;
\item $D$ is not flat: $\fd_AD=1$, and $\Tor^A_1(D,A/xA)\cong D$ for every $0\ne x\in I$;
\item for all $D$-modules $N,M$ and all $n\ge0$, inflation induces natural isomorphisms
$\Ext^n_A(N,M)\cong\Ext^n_D(N,M)$;
\item $\pd_AD\ge2$;
\item the sequence $0\to I\to A\xrightarrow{\ct}D\to0$ does not split, so $\Ext^1_A(D,I)\ne0$;
\item if $(x_\alpha)_{\alpha\in\Lambda}$ generates $I$, $P=\bigoplus_{\alpha}Ae_\alpha\to I$,
$e_\alpha\mapsto x_\alpha$, and $J$ is its kernel, then $\Ext^2_A(D,J)\ne0$.
\end{enumerate}
\end{theorem}
\begin{proof}
(i) Each $X^hA$ is free of rank one because $A$ is a domain, the union in (H2) is directed
($X^hA\subseteq X^{h'}A$ for $h'<h$), and a filtered colimit of flat modules is flat
\cite[Lemma~10.39.3]{StacksFlat}. The factorization $f=X^gb$ writes $f$ as a product of two elements
of $I$, so $I^2=I$. Projectivity would contradict (H3) by \cref{osq:lem:projideal}.
(ii) Since $IN=0$, $I\otimes_AN\cong(I/I^2)\otimes_DN=0$ (the map $x\otimes n\mapsto(x+I^2)\otimes n$ is
well defined in reverse because $(ab)\otimes n=b\otimes an=0$ for $a,b\in I$); tensoring the flat
resolution $0\to I\to A\to D\to0$ with $N$ computes all positive $\Tor$ groups as $0$.
(iii) The same resolution gives $\fd_AD\le1$. For $0\ne x\in I$, $0\to A\xrightarrow{x}A\to A/xA\to0$ is a
free resolution; after $\otimes_AD$ the first map is $0$ because $x\in I$, so $\Tor_1^A(D,A/xA)\cong D\ne0$
and $D$ is not flat. Equivalently, the injection $xA\hookrightarrow A$ becomes the zero map $D\to D$ after
tensoring (03, 09).
(iv) Take a projective $A$-resolution $P_\bullet\to N$. By (ii), $D\otimes_AP_\bullet$ has homology $N$ in
degree $0$ only, and its terms are projective $D$-modules (direct summands of free ones), so it is a
projective $D$-resolution; the adjunction $\Hom_A(P_\bullet,M)\cong\Hom_D(D\otimes_AP_\bullet,M)$ is an
isomorphism of complexes, and comparison of resolutions makes the result natural.
(v) If $\pd_AD\le1$, the kernel $I$ of the free surjection $A\to D$ would be projective (Schanuel), against
(i).
(vi) A section $s$ gives $r=s(1)$ with $\ct(r)=1$ and $xr=s(x\cdot1)=0$ for $0\ne x\in I$, so $r=0$ in the
domain $A$.
(vii) $P\to I$ is surjective; $0\to J\to P\to I\to0$ does not split (else $I$ would be a summand of a free
module, hence projective), giving $0\ne\Ext^1_A(I,J)$; applying $\Hom_A(-,J)$ to
$0\to I\to A\to D\to0$ with $A$ projective gives $\Ext^1_A(I,J)\cong\Ext^2_A(D,J)$, represented by the
spliced sequence $0\to J\to P\to A\to D\to0$.
\end{proof}

In standard terms, $A\to D$ is a homological ring epimorphism ($D\otimes_AD=D$, all higher self-$\Tor$
vanish) but not a flat one (03, 06); this terminology is optional. The exact value of $\pd_AD$ is not
determined: only $\pd_AD\ge2$ and $\fd_AD=1$ are proved (06). The nonzero $\Tor_1^A(D,A/xA)$ does not
contradict (ii), because $A/xA$ is not a $D$-module: $I\not\subseteq xA$, since a gap below $\supp x$
exhibits a smaller positive monomial of $I$ outside $xA$ (06). $I/I^2=0$ is not a finite-presentation or
smoothness statement (09).

\begin{theorem}[Exact size of cohomological witnesses \src{06} \merge]\label{osq:thm:extthreshold}
Let $A$ be one of the five rings of \cref{osq:sub:organizing}, with $D=\Z$ (or any admissible $D$), and
$\tau=|A|$. For every $A$-module $M$ with $|M|<\tau$ and every $n>0$, $\Ext^n_A(D,M)=0$, and for any two
$A$-modules of size $<\tau$ the $\Ext$ groups over $A$ are those of the underlying $D$-modules. For
$n\in\{1,2\}$,
\[
 \min\{|M|:\Ext^n_A(D,M)\ne0\}=\tau .
\]
For $A^{06}_\kappa$ this minimum is $\kappa^{<\kappa}$ in ZFC; under $\kappa^{<\kappa}=\kappa$ it is $\kappa$,
which is 06's statement. For $A^{03}_\kappa$, $A^{08}_\kappa$ and $A^{09}_{\lambda,D}$ it is $\kappa^{\aleph_0}$,
$\kappa$ and $\lambda$.
\end{theorem}
\begin{proof}
Modules of size $<\tau$ are annihilated by $I$ (\cref{osq:thm:thresholds}), so they are inflated $D$-modules
and \cref{osq:thm:homological}(iv) applies; $D$ is free over itself. For the upper bounds, $|I|=\tau$ and
$\Ext^1_A(D,I)\ne0$. For $n=2$ choose a generating family of cardinality at most $\tau$ (the generator counts of (H3) are all $\le\tau$); then $|P|\le\tau$, so
$|J|\le\tau$, and $\Ext^2_A(D,J)\ne0$ with the lower bound forces $|J|=\tau$.
\end{proof}

The $\Ext$ groups are computed in the full category of $A$-modules, not only among small modules; this is why
\cref{osq:thm:homological}(iv) is an ambient derived-functor statement and not merely an equivalence of small
module categories (06). The two witnesses show that a quotient can look homologically projective against every
small coefficient module while failing projectivity in the ambient category. The whole argument works for the
Gaussian versions with $D=\Z[i]$, where $\Ext_D$ means extensions of $\Z[i]$-modules (06). Finally, for
$0\ne x\in\Pi^{06}_\kappa$ the module $A^{06}_\kappa/xA^{06}_\kappa$ is detected by $I$ (it is not a $D$-module), so
its carrier has at least $\kappa^{<\kappa}$ elements; 06 stated this with the qualifier ``under the
cardinal-arithmetic hypothesis'', which is not needed.

\section{Large internal fields and the binomial quotient}
\label{osq:sec:internal}

The universal theorem (\cref{osq:thm:universal}) describes what a set-sized target
can see of $\Oz$: every purely infinite element dies, and only the constant term
survives. This section, taken from manuscript~07 together with one proposition of
manuscript~06, shows how much survives in a large target once a single purely infinite
element becomes a \emph{unit}. The answer is a great deal: the target then contains a
real closed Hahn field into which all of $\No$ embeds as an ordered field. In the
nonzero binomial quotient $\Qa_a=\Oz/(1+\omega^a)$, which has no nonzero set-sized ring
image at all, the same mechanism produces an algebraically closed class field containing
a copy of $\SC$, exact finite cyclotomic algebras over it, and an atomless Boolean
algebra of idempotents.

The dependence of the results is as follows (07). An ordered additive compression of
$\No$ and the closedness of set-sized Hahn fields give a copy of $\No$ inside the
compressed field $\FF_a$ (\cref{osq:if:prop:closed,osq:if:thm:compression});
localization transfers $\FF_a$ into every quotient in which a suitable purely infinite
element is invertible (\cref{osq:if:thm:transfer}); a support argument gives exact
binomial kernels (\cref{osq:if:thm:binomialkernel}); and finite Chinese remainder
decompositions give the internal field $\EE_a$ and the Cantor algebra
(\cref{osq:if:cor:Ea,osq:if:thm:cantor}). The absence of small images of $\Qa_a$
comes from the other branch, \cref{osq:thm:universal,osq:thm:reflection}. No
factorization conjecture enters either chain. What manuscript~07 calls its Main
Theorem~B is \cref{osq:if:thm:transfer} together with \cref{osq:if:thm:compression},
and its Main Theorem~C is the combination of
\cref{osq:if:lem:localization,osq:if:cor:Ea,osq:if:thm:cantor} with the absence of
small images of $\Qa_a$ recorded before \cref{osq:if:lem:localization}.

\subsection{Compressed exponent fields}\label{osq:if:sub:compressed}

Fix $a>0$. Recall from \eqref{osq:eq:Ha} and \cref{osq:lem:Ha} that
$H_a=\{h\in\No:n|h|<a\text{ for all ordinary }n\ge1\}$ is a convex divisible
subgroup of $(\No,+)$ containing a proper class of positive elements, and that
$\FF_a(k)=\hahn{k}{H_a}$ is a field for every set field $k$. For $k=\R$ and $k=\C$ we
realize these fields inside the surreals and surcomplex numbers,
\begin{equation}\label{osq:if:eq:Fa}
\begin{aligned}
 \FF_a&=\FF_a(\R)=\{f\in\No:\supp f\subseteq H_a\},\\
 \FF_a(i)&=\FF_a(\C)=\{f\in\SC:\supp f\subseteq H_a\}.
\end{aligned}
\end{equation}
Taking real and imaginary parts of coefficients shows $\FF_a(i)=\FF_a+i\FF_a$.
Manuscript~07 calls them \emph{compressed exponent fields}; the phrase refers to the
restriction on exponents, not to finiteness of supports. Both are proper classes.
Manuscript~07 checks directly that $H_a$ is a subgroup by the inequality
$n|h+h'|\le n|h|+n|h'|<a/2+a/2$, using the defining inequalities at $2n$; this is
part of \cref{osq:lem:Ha}.

\begin{proposition}[Closedness of the compressed fields \src{07}]\label{osq:if:prop:closed}
The class $\FF_a$ is a real closed subfield of $\No$, and $\FF_a(i)$ is an algebraically
closed subfield of $\SC$.
\end{proposition}

\begin{proof}
Since $H_a$ is a subgroup, sums and products of series supported in $H_a$ are supported
in $H_a$; the inverse of a nonzero Hahn series has support in the subgroup generated by
its own support (factor out the leading monomial and use the Hahn geometric series), so
it stays in $H_a$ as well. Hence $\FF_a$ and $\FF_a(i)$ are subfields
(\cref{osq:lem:Ha}).

Let $P$ be a nonconstant polynomial with coefficients in $\FF_a(i)$. The union of the
supports of its finitely many coefficients is a set; its rational span is a set-sized
divisible ordered subgroup $G\subseteq H_a$, because $H_a$ is a divisible subgroup. The
series with support in $G$ form the set-sized Hahn fields $\R(\!(\omega^G)\!)\subseteq
\FF_a$ and $\C(\!(\omega^G)\!)\subseteq\FF_a(i)$, whose operations are those of $\No$ and
$\SC$, and $P$ has coefficients in $\C(\!(\omega^G)\!)$. For divisible $G$ this field is
algebraically closed by the Hahn-field theorem \cite[Corollary~4]{Poonen} (stated there
for well-ordered supports; the substitution $t^g=\omega^{-g}$ converts). So $P$ has a root
in $\C(\!(\omega^G)\!)\subseteq\FF_a(i)$, and $\FF_a(i)$ is algebraically closed.

For the real field, note first that $\R(\!(\omega^G)\!)$ is real closed. If
$0<f\in\R(\!(\omega^G)\!)$, write a square root of $f$ in $\C(\!(\omega^G)\!)$ as $u+iv$
with $u,v\in\R(\!(\omega^G)\!)$; then $u^2-v^2=f$ and $2uv=0$. If $u=0$ then
$f=-v^2\le0$, which is impossible; so $v=0$ and $u^2=f$. If $Q$ has odd degree and
coefficients in $\R(\!(\omega^G)\!)$, coefficientwise conjugation is an automorphism of
$\C(\!(\omega^G)\!)$ fixing $Q$, so the nonreal roots of $Q$ occur in conjugate pairs
with equal multiplicities, and an odd number of roots counted with multiplicity must be
real. An ordered field in which positive elements are squares and odd-degree polynomials
have roots is real closed. Now let $0<f\in\FF_a$, or let $Q$ be of odd degree over
$\FF_a$; choosing $G$ as above for $f$ or for the coefficients of $Q$, the required square
root or root lies in $\R(\!(\omega^G)\!)\subseteq\FF_a$. Hence $\FF_a$ is real closed.
Only set-sized instances of the classical theorem have been used.
\end{proof}

\begin{remark}[Bounded in the ambient field, yet a proper class \src{07}]
\label{osq:if:rem:bounded}
Every positive $f\in\FF_a$ satisfies $\omega^{-a}<f<\omega^a$: its leading exponent $h$
satisfies $|h|<a/2$ and its leading coefficient is a positive real. A field contained in
such an ambient interval need not be small, since the interval is itself a proper class;
$\FF_a$ is a proper class, and by \cref{osq:if:thm:compression} it even contains an
ordered-field copy of $\No$. Bounds in the surrounding field do not give bounds inside an
isomorphic copy of $\No$.
\end{remark}

\subsection{An ordered additive compression of the surreals}\label{osq:if:sub:compression}

Knowing that $\FF_a$ contains Hahn fields of every set size (\cref{osq:lem:Ha}) would
already show considerable largeness. Manuscript~07 does more: it embeds all of $\No$ into
$\FF_a$ as an ordered field. Define
\begin{equation}\label{osq:if:eq:theta}
 \vartheta(x)=\frac12\Bigl(\frac{x}{1+|x|}-3\Bigr),\qquad
 \vartheta:\No\longrightarrow(-2,-1).
\end{equation}
For $x\ge0$ the inner fraction is $x/(1+x)=1-1/(1+x)\in[0,1)$, and for $x<0$ it is
$x/(1-x)=-1+1/(1-x)\in(-1,0)$. Each branch is strictly increasing and they agree at $0$,
so $\vartheta$ is strictly increasing with values in $(-2,-1)$. For a real-coefficient
normal form put
\begin{equation}\label{osq:if:eq:Psi}
 \Psi\Bigl(\sum_\gamma r_\gamma\omega^\gamma\Bigr)=\sum_\gamma r_\gamma\omega^{\vartheta(\gamma)},
 \qquad \Psi_a(x)=a\Psi(x).
\end{equation}
(Manuscript~07 writes $E_a$ for $\Psi_a$; the letter $\EE_a$ is reserved here for the
internal field of \cref{osq:if:cor:Ea}.)

\begin{lemma}[Ordered additive compression \src{07}]\label{osq:if:lem:compression}
The map $\Psi_a:(\No,+)\to(H_a,+)$ is an injective, order-preserving, $\R$-linear
homomorphism of additive groups.
\end{lemma}

\begin{proof}
A strictly increasing map sends a reverse well-ordered set to a reverse well-ordered set,
so the support $\vartheta(\supp x)$ is a reverse well-ordered set and $\Psi(x)$ exists.
Since $\vartheta$ is injective, distinct exponents stay distinct, and the coefficient of
$\Psi(x)$ at $\vartheta(\gamma)$ is $x_\gamma$; hence coefficientwise addition, with all
cancellations, and multiplication by real scalars are preserved. If $x\ne0$, the greatest
exponent $\gamma_0$ of $\supp x$ is sent to the greatest exponent $\vartheta(\gamma_0)$ of
$\supp\Psi(x)$ with the same coefficient; since the sign of a nonzero normal form is the
sign of its leading coefficient, $\Psi$ maps positive elements to positive elements, and
being additive it is injective and order preserving. The leading exponent of a nonzero
$\Psi(x)$ lies in $(-2,-1)$, so $\Psi(x)$ is infinitesimal and $n|\Psi(x)|<1$ for every
ordinary $n\ge1$. Multiplying by $a>0$ gives $n|\Psi_a(x)|<a$, that is,
$\Psi_a(x)\in H_a$, and multiplication by the positive surreal $a$ preserves additivity,
$\R$-linearity and order.
\end{proof}

The map $\vartheta$ is not additive, and $\Psi$ is not asserted to be multiplicative. The
construction deliberately has two levels: first an additive exponent group is compressed,
and only then is the additive map used as a substitution on monomials.

\begin{theorem}[A full surreal field in every compressed field \src{07}]
\label{osq:if:thm:compression}
The formula
\begin{equation}\label{osq:if:eq:Theta}
 \Theta_a\Bigl(\sum_\gamma r_\gamma\omega^\gamma\Bigr)=\sum_\gamma r_\gamma\omega^{\Psi_a(\gamma)}
\end{equation}
defines an injective ordered-field homomorphism $\Theta_a:\No\hookrightarrow\FF_a$ fixing
$\R$. It preserves every set-indexed strongly summable family. Extending coefficients,
$x+iy\mapsto\Theta_a(x)+i\Theta_a(y)$, gives an injective field homomorphism
$\SC\hookrightarrow\FF_a(i)$ fixing $\C$.
\end{theorem}

\begin{proof}
The support of the right side of \eqref{osq:if:eq:Theta} is $\Psi_a(\supp x)$, the image of
a reverse well-ordered set under a strictly increasing map into $H_a$
(\cref{osq:if:lem:compression}); so the output exists and lies in $\FF_a$. If
$(x_\lambda)$ is a strongly summable set-indexed family, the union of the supports is
reverse well ordered and each exponent receives finitely many contributions; since
$\Psi_a$ is injective and strictly increasing, the same holds for the family
$(\Theta_a(x_\lambda))$ with the same contributions at $\Psi_a(\gamma)$ as at $\gamma$,
so $\Theta_a(\sum_\lambda x_\lambda)=\sum_\lambda\Theta_a(x_\lambda)$.

For multiplication, write $x=\sum r_\gamma\omega^\gamma$ and $y=\sum s_\delta\omega^\delta$.
Because $\Psi_a$ is additive and injective,
$\Psi_a(\gamma)+\Psi_a(\delta)=\Psi_a(\varepsilon)$ holds exactly when
$\gamma+\delta=\varepsilon$, and every sum $\Psi_a(\gamma)+\Psi_a(\delta)$ lies in the image
of $\Psi_a$. Hence the coefficient of $\Theta_a(x)\Theta_a(y)$ at $\Psi_a(\varepsilon)$ is
$\sum_{\gamma+\delta=\varepsilon}r_\gamma s_\delta=(xy)_\varepsilon$, and its coefficients
at exponents outside the image vanish; so $\Theta_a(xy)=\Theta_a(x)\Theta_a(y)$.
Additivity is clear, and constants are fixed because $\Psi_a(0)=0$. A nonzero $x$ has a
nonzero image whose leading coefficient is that of $x$, so $\Theta_a$ is injective and
preserves signs, hence order. It is an injective unital ring map between fields, hence a
field embedding. The complex extension is the same coefficient-and-support argument with
complex coefficients.
\end{proof}

\begin{example}[\src{07}]\label{osq:if:ex:Theta}
Since $\vartheta(0)=-3/2$, one has $\Psi_a(1)=a\omega^{-3/2}$ and
\[
 \Theta_a(\omega)=\omega^{a\omega^{-3/2}} .
\]
\end{example}

No preservation of the ordinal embedding, of birthdays or simplicity, or of the surreal
exponential is asserted or needed: $\Theta_a$ is an explicit ordered-field embedding with a
controlled normal-form support image. Constructions of self-similar surreal subclasses and
strongly linear maps have a substantial prior literature, notably the surreal
substructures of Bagayoko and van der Hoeven \cite{BagayokoHoeven}. Manuscript~07 uses the
elementary construction above to make the transfer theorem independent of a general
classification of surreal substructures, and does not claim that surreal self-similarity
itself is new.

\subsection{Transfer by inverting a purely infinite element}\label{osq:if:sub:transfer}

\begin{lemma}[Multiplication by a purely infinite denominator \src{07}]\label{osq:if:lem:gf}
Let $0\ne g\in\Pi$ and choose $a>0$ with $a<\gamma$ for every $\gamma\in\supp g$ (possible
by \cref{osq:lem:gap}). If $f\in\FF_a$, then $gf\in\Pi$. The same holds for
$0\ne g\in\Pi_\C$ and $f\in\FF_a(i)$.
\end{lemma}

\begin{proof}
Every $h\in\supp f$ satisfies $|h|<a/2$. Every exponent occurring in the product before
cancellation has the form $\gamma+h$ with $\gamma\in\supp g$, and
$\gamma+h>a-a/2>0$. Hahn multiplication is defined and cancellation cannot create an
exponent outside the sum of the supports, so $gf$ has strictly positive support and real
(respectively complex) coefficients.
\end{proof}

\begin{theorem}[Localization and quotient transfer \src{07}]\label{osq:if:thm:transfer}
Under the hypotheses of \cref{osq:if:lem:gf}, the subfield $\FF_a$ of $\No$ lies in the
localization $\Oz[1/g]\subseteq\No$. If $q:\Oz\to B$ is a unital homomorphism to a nonzero
class ring and $q(g)$ is invertible in $B$, then
\begin{equation}\label{osq:if:eq:transfer}
 \iota_g(f)=q(gf)\,q(g)^{-1}\qquad(f\in\FF_a)
\end{equation}
is an injective field homomorphism $\FF_a\hookrightarrow B$, and $\iota_g\circ\Theta_a$ is
an injective field homomorphism $\No\hookrightarrow B$. For $q:\GOz\to B$ and
$0\ne g\in\Pi_\C$ with $q(g)$ invertible, the same formula gives an injective field
homomorphism $\FF_a(i)\hookrightarrow B$, and hence a copy of $\SC$ in $B$. In either case
$B$ is a proper class.
\end{theorem}

\begin{proof}
Inside $\No$, $f=(gf)/g$ with $gf\in\Pi\subseteq\Oz$ by \cref{osq:if:lem:gf}, so
$f\in\Oz[1/g]$. The elements of $q(\Oz)$ commute pairwise, and $q(g)^{-1}$ commutes with
every element $y$ commuting with $q(g)$ (multiply $q(g)y=yq(g)$ on both sides by
$q(g)^{-1}$). Hence $q(\Oz)$ and $q(g)^{-1}$ generate a commutative subring of $B$, and all
computations below take place there; no commutativity of $B$ is needed. The map
\eqref{osq:if:eq:transfer} is additive and $\iota_g(1)=1$. For $f,f'\in\FF_a$ the
elements $gf$, $gf'$, $gff'$ lie in $\Oz$ and $(gf)(gf')=g\cdot(gff')$, so
\[
 \iota_g(f)\iota_g(f')=q(gf)q(gf')q(g)^{-2}=q(g)q(gff')q(g)^{-2}=\iota_g(ff').
\]
(Equivalently, $\iota_g$ is the restriction to $\FF_a$ of the extension of $q$ to
$\Oz[1/g]$ given by the universal property of localization.) If $f\ne0$ and
$\iota_g(f)=0$, then $1_B=\iota_g(f)\iota_g(f^{-1})=0$, contradicting $B\ne0$; so
$\iota_g$ is injective. Composing with \cref{osq:if:thm:compression} gives the copy of
$\No$. The complex case is identical, with \cref{osq:if:lem:gf} for $\Pi_\C$ and the
complex extension of $\Theta_a$. Finally, if $B$ were set-sized,
\cref{osq:thm:universal} would give $q(g)=0$, which is not invertible in a nonzero ring.
\end{proof}

\begin{remark}[\src{07}]\label{osq:if:rem:transfer}
If the target $B$ is noncommutative, the image of the localization still lies in a
commutative subring, as the proof shows; no commutativity condition on all of $B$ is
needed. The hypothesis on $a$ is always available in the full surreal exponent class, even
when the exponents of $g$ approach zero inside a smaller ordered group
(\cref{osq:lem:gap}). No inference of this kind is valid for an arbitrary fixed
set-sized exponent group without an additional lower-bound hypothesis.
\end{remark}

\begin{corollary}[Quotients by a unit-constant perturbation \src{07}]
\label{osq:if:cor:unitconstant}
If $0\ne g\in\Pi$, then $\Oz/(1+g)$ is nonzero and contains a class copy of $\No$. If
$0\ne g\in\Pi_\C$, then $\GOz/(1+g)$ is nonzero and contains a class copy of $\SC$. Neither
quotient has a unital homomorphism to a nonzero set-sized ring.
\end{corollary}

\begin{proof}
The element $1+g$ is nonconstant, hence a nonunit by \cref{osq:lem:units}, so the
quotients are nonzero. In the quotient $g=-1$ is a unit, so
\cref{osq:if:thm:transfer} applies. Since $\ct(1+g)=1$, the ideal $(1+g)$ has constant-term
image all of $\Z$ (respectively $\Z[i]$), and \cref{osq:thm:reflection} makes the
universal set-sized image the zero ring; a unital map from the zero ring exists only to the
zero ring.
\end{proof}

\begin{corollary}[A conditional statement about nonarithmetic primes \src{07}]
\label{osq:if:cor:primefield}
Let $P$ be a prime ideal of $\Oz$ such that $P\cap\Z=0$ and $P\not\subseteq\Pi$. Then
$\Oz/P$ contains a class copy of $\No$ and is a $\Q$-algebra. Likewise, if $P$ is a prime
ideal of $\GOz$ with $P\cap\Z[i]=0$ and $P\not\subseteq\Pi_\C$, then $\GOz/P$ contains a class
copy of $\SC$ and is a $\Q(i)$-algebra.
\end{corollary}

\begin{proof}
Choose $f\in P$ with nonzero constant term $c$. The series $f/c$ still lies in the omnific
ring: its constant term is $1$ and its other coefficients are real (respectively complex).
Since $f=c\cdot(f/c)$, $P$ is prime and $c\notin P$, we get $f/c\in P$. Write $f/c=1+g$
with $g$ purely infinite; $g\ne0$ because $c\notin P$. In the quotient $g=-1$ is a unit, and
\cref{osq:if:thm:transfer} applies. For every nonzero ordinary coefficient $d$ ($d\in\Z$,
respectively $d\in\Z[i]$), $g/d$ lies in the ring, and in the quotient
$d\cdot(-g/d)=-g=1$; so $d$ is invertible, which gives the $\Q$-, respectively
$\Q(i)$-, algebra structure.
\end{proof}

The corollary is deliberately conditional. Manuscript~07 does not assert that prime ideals
with these properties exist; see \cref{osq:if:sub:boundaries}.

\begin{example}[Invertibility cannot be replaced by nonvanishing \src{07}]
\label{osq:if:ex:nilpotent}
The quotient $\Oz/(\omega^a)$ is nonzero, and the class of $\omega^{a/2}$ is a nonzero
element of square zero: it is nonzero because $\omega^{a/2}/\omega^a=\omega^{-a/2}\notin\Oz$.
Since $(\omega^a)\subseteq\Pi$, this quotient still maps onto $\Z$ by constant terms, and so
it cannot contain a unital copy of $\Q$: composing such a copy with that map would give a
unital map $\Q\to\Z$. A surviving nilpotent and a surviving unit are therefore
substantively different.
\end{example}

The same mechanism, with a fraction field in place of the quotient, applies to every prime
ideal not containing the purely infinite ideal. Here $X^g$ is the monomial of
$\hahn{k}{\No}$, which is $\omega^g$ when $k=\R$ or $k=\C$.

\begin{proposition}[A scaled field in every non-arithmetic residue field \src{06}]
\label{osq:if:prop:reservoir}
Let $k$ be a set field, $D\subseteq k$ a unital subring, $A=\A_{D,k}$, and let $P$ be a
prime class ideal of $A$ with $\Pi_k\not\subseteq P$. There exists $a>0$ for which the
class field $\FF_a(k)=\hahn{k}{H_a}$ embeds in $\Frac(A/P)$.
\end{proposition}

\begin{proof}
Choose $f\in\Pi_k\setminus P$ and write $f=X^gh$ with $g>0$ and $h\in\Pi_k$
(\cref{osq:lem:gap} applied to $\supp f$). Then $X^g\notin P$, since otherwise $f\in P$.
Put $a=g/2$ and $m=X^a$. As $m^2=X^g\notin P$ and $P$ is prime, $\bar m\ne0$ in the domain
$A/P$. Since $m\FF_a(k)\subseteq\Pi_k\subseteq A$ (\cref{osq:lem:Ha}), we may define
\[
 \iota(x)=\frac{\overline{mx}}{\bar m}\in\Frac(A/P)\qquad(x\in\FF_a(k)).
\]
Additivity and $\iota(1)=1$ are immediate. Multiplicativity follows by reducing the
identity $(mx)(my)=m(mxy)$ in $A$ and dividing by $\bar m^2$. For $x\ne0$, the identity
$(mx)(m/x)=m^2$ in $A$ shows $\overline{mx}\ne0$, because $\overline{m^2}\ne0$. Thus the
unital field homomorphism $\iota$ is injective.
\end{proof}

This is the scaled-field collision mechanism of \cref{osq:lem:collision} in a fraction
field. The conclusion says more than that the residue field is large: it contains a
specific Hahn field built on all exponent scales infinitesimal relative to one positive
$a$. It is a necessary structural condition on such residue fields; manuscript~06 states
that it does not classify those prime ideals, assert that they exist, or assert that
every field of this form is a residue field.

\begin{remark}[The composed embedding \merge]\label{osq:if:rem:composition}
Let $P$ be a prime class ideal of $\Oz$ with $\Pi\not\subseteq P$. Then there is an
injective field homomorphism $\No\hookrightarrow\Frac(\Oz/P)$ whose first factor is an
ordered-field embedding,
\[
 \No\xrightarrow{\ \Theta_a\ }\FF_a\xrightarrow{\ \iota\ }\Frac(\Oz/P).
\]
Indeed, \cref{osq:if:prop:reservoir} with $(D,k)=(\Z,\R)$ gives $a>0$ and an injective field
homomorphism $\iota:\FF_a\to\Frac(\Oz/P)$, and \cref{osq:if:thm:compression} gives
$\Theta_a$; a composite of injective field homomorphisms is one. With $(D,k)=(\Z[i],\C)$ and
the complex extension of $\Theta_a$, every prime $P$ of $\GOz$ with $\Pi_\C\not\subseteq P$
gives $\SC\hookrightarrow\Frac(\GOz/P)$. Neither source states this composition.

The hypotheses of \cref{osq:if:prop:reservoir} and \cref{osq:if:cor:primefield} are
different. If $P\cap\Z=0$ and $P\not\subseteq\Pi$, then $\Pi\not\subseteq P$: otherwise,
for $x\in P$ we would have $x-\ct(x)\in\Pi\subseteq P$, hence $\ct(x)\in P\cap\Z=0$, so
$P\subseteq\Pi$. Conversely, $\Pi\not\subseteq P$ forces $P\cap\Z=0$, since a nonzero
$n\in P$ would give $\Pi=n\Pi\subseteq P$ (\cref{osq:prop:common}), but it allows
$P\subseteq\Pi$, for instance $P=0$, where the statement is the trivial
$\Frac(\Oz)=\No$ (\cref{osq:prop:fractions}). Thus 07's hypothesis is the stronger one and
yields the stronger conclusion (a copy of $\No$ inside $\Oz/P$ itself, and a
$\Q$-algebra structure), while 06's hypothesis also covers primes contained in $\Pi$, with
the conclusion only in the fraction field.
\end{remark}

\subsection{The binomial quotient and its localization}\label{osq:if:sub:binomial}

Fix $a>0$ and write
\[
 \Qa_a=\Oz/(1+\omega^a),\qquad q_a:\Oz\longrightarrow\Qa_a
\]
for the quotient map (07's $B_a$). The element $1+\omega^a$ is a nonunit
(\cref{osq:lem:units}), so $\Qa_a\ne0$; and since $\ct(1+\omega^a)=1$,
\cref{osq:thm:reflection} shows that $\Qa_a$ has no unital homomorphism to a nonzero
set-sized ring. A purely arithmetic calculation shows why: in a set-sized target
$\omega^a$ becomes $0$ (\cref{osq:thm:universal}), so $1+\omega^a$ becomes $1$, and a
target in which it vanishes has $1=0$. Its failure to have small ring images is therefore
not caused by its being zero, a field, or a domain; it is none of these by
\cref{osq:if:cor:boolean}. By \cref{osq:if:cor:unitconstant} it contains a copy of $\No$;
the rest of this section identifies much more of its internal structure.

The support description of the localization
\[
 \Oz[\omega^{-a}]=\{f\in\No:\supp f\subseteq[-Na,\infty)\text{ for some ordinary }N\ge0\}
\]
is \cref{osq:prop:locsupport}; manuscripts~04 and~07 both prove it. (07 calls this ring
$L_a$.)

\begin{lemma}[Localization of the binomial quotient \src{07}]\label{osq:if:lem:localization}
\leavevmode
\begin{enumerate}
\item $\omega^a\FF_a\subseteq\Pi$; in particular $\FF_a\subseteq\Oz[\omega^{-a}]$.
\item There is a unique ring homomorphism $\widetilde q_a:\Oz[\omega^{-a}]\to\Qa_a$
extending $q_a$. It is surjective with kernel $(1+\omega^a)\Oz[\omega^{-a}]$, so
\begin{equation}\label{osq:if:eq:localquotient}
 \Qa_a\cong\Oz[\omega^{-a}]/(1+\omega^a)\Oz[\omega^{-a}] .
\end{equation}
\item The restriction of $\widetilde q_a$ to $\FF_a$ is the injective field homomorphism
\begin{equation}\label{osq:if:eq:iota}
 \iota_a(f)=-q_a(\omega^af)\qquad(f\in\FF_a).
\end{equation}
\end{enumerate}
\end{lemma}

\begin{proof}
(i) Every $h\in H_a$ satisfies $a+h>0$, so translating a support in $H_a$ by $a$ gives a
strictly positive support (\cref{osq:lem:Ha}); then $f=\omega^{-a}(\omega^af)$.

(ii) The subring of $\No$ generated by $\Oz$ and $\omega^{-a}$ consists of the finite sums
$\sum_{j\le N}h_j\omega^{-ja}=\omega^{-Na}\sum_{j\le N}h_j\omega^{(N-j)a}$, so every element
has the form $\omega^{-Na}h$ with $N\ge0$ and $h\in\Oz$. Any extension of $q_a$ sends
$\omega^{-a}$ to $q_a(\omega^a)^{-1}=-1$, which proves uniqueness; define
$\widetilde q_a(\omega^{-Na}h)=(-1)^Nq_a(h)$. This is well defined: if
$\omega^{-Na}h=\omega^{-Ma}h'$ with $N\le M$, then $h'=\omega^{(M-N)a}h$ and
$(-1)^Mq_a(h')=(-1)^M(-1)^{M-N}q_a(h)=(-1)^Nq_a(h)$. Writing two elements over a common
power $\omega^{-Na}$ shows that $\widetilde q_a$ is a ring homomorphism; it extends $q_a$
and is surjective. If $\widetilde q_a(\omega^{-Na}h)=0$, then $q_a(h)=0$, so
$h=(1+\omega^a)y$ with $y\in\Oz$ and
$\omega^{-Na}h=(1+\omega^a)\omega^{-Na}y\in(1+\omega^a)\Oz[\omega^{-a}]$. Conversely
$\widetilde q_a((1+\omega^a)\omega^{-Na}y)=(-1)^Nq_a((1+\omega^a)y)=0$ for $y\in\Oz$.

(iii) By (i) and (ii), $\widetilde q_a(f)=\widetilde q_a(\omega^{-a}\cdot\omega^af)
=-q_a(\omega^af)$. A unital ring homomorphism from a field to a nonzero ring is injective.
\end{proof}

Henceforth $\FF_a$ is identified with its image $\iota_a(\FF_a)\subseteq\Qa_a$ whenever we
work inside $\Qa_a$; all coefficient fields below are embedded with the identity of
$\Qa_a$. This identification is not constant-term extraction: a nonintegral real scalar
$r$ is represented in $\Qa_a$ by the omnific integer $-r\omega^a$.

\subsection{Exact binomial kernels}\label{osq:if:sub:kernel}

For an ordinary integer $n\ge1$ put $x_n=q_a(\omega^{a/n})\in\Qa_a$. Certainly
$x_n^n=q_a(\omega^a)=-1$. The substantive question is whether the full ambient omnific
ring, with its arbitrary surreal exponents, creates additional polynomial relations over
$\FF_a$.

\begin{theorem}[Exact binomial kernel \src{07}]\label{osq:if:thm:binomialkernel}
For every $n\ge1$, evaluation at $x_n$ induces an injective ring homomorphism
\begin{equation}\label{osq:if:eq:binomialinjection}
 \FF_a[T]/(T^n+1)\hookrightarrow\Qa_a .
\end{equation}
Equivalently, the kernel of $\FF_a[T]\to\Qa_a$, $T\mapsto x_n$, is exactly $(T^n+1)$.
\end{theorem}

\begin{proof}
The polynomial $T^n+1$ lies in the kernel. Dividing by this monic polynomial over the
field $\FF_a$, it suffices to show that no nonzero remainder
\[
 R(T)=\sum_{j=0}^{n-1}r_jT^j\qquad(r_j\in\FF_a)
\]
vanishes at $x_n$. Put $\xi=\omega^{a/n}\in\Oz$; then $R(\xi)\in\Oz[\omega^{-a}]$ by
\cref{osq:if:lem:localization}(i), and $\widetilde q_a(R(\xi))=R(x_n)$.

\emph{The quotient in $\No$.} The series $U=\sum_{k\ge0}(-1)^k\omega^{-(k+1)a}$ has the
decreasing support $\{-(k+1)a\}$, and $(1+\omega^a)U=\sum_{k\ge0}(-1)^k\omega^{-(k+1)a}
+\sum_{k\ge0}(-1)^k\omega^{-ka}=1$, each exponent receiving at most two contributions that
cancel except at exponent $0$. Hence $U=(1+\omega^a)^{-1}$ in $\No$, and multiplying the
finite sum $R(\xi)=\sum_jr_j\omega^{ja/n}$ by $U$ gives
\begin{equation}\label{osq:if:eq:remainderdivision}
 \frac{R(\xi)}{1+\omega^a}=\sum_{k\ge0}\sum_{j=0}^{n-1}(-1)^kr_j\,\omega^{(j/n-k-1)a},
\end{equation}
provided the right side is a strongly summable family, which we now check together with
the absence of cancellation between blocks.

\emph{Block cosets.} The support of the block indexed by $(j,k)$ lies in the coset
$c_{j,k}+H_a$ with centre $c_{j,k}=(j/n-k-1)a$, because $\supp r_j\subseteq H_a$. Distinct
pairs give distinct cosets. Indeed, $c_{j,k}-c_{j',k'}=(m/n)a$ with
$m=j-j'-n(k-k')\in\Z$, and $(m/n)a\in H_a$ only for $m=0$, since for $m\ne0$ one has
$n\cdot|(m/n)a|=|m|a\ge a$. From $j-j'=n(k-k')$ and $|j-j'|<n$ we get $j=j'$ and $k=k'$.
Cosets of $H_a$ are disjoint, so distinct blocks have disjoint supports, and each exponent
receives a contribution from at most one block.

\emph{Reverse well-ordering.} As $(j,k)$ runs over $\{0,\dots,n-1\}\times\N$, the integer
$j-n(k+1)$ runs bijectively over the integers $m\le-1$, so the centres are the points
$(m/n)a$, $m\le-1$, a strictly descending discrete progression with step $a/n$. Every
offset $h\in H_a$ satisfies $4n|h|<a$, that is, $|h|<a/(4n)$, so the block with centre
$(m/n)a$ lies in the open interval of radius $a/(4n)$ about its centre; these intervals are
pairwise disjoint and ordered as their centres. Given a nonempty subset $S$ of the union of
the supports, the set of $m$ whose block meets $S$ is a nonempty set of integers bounded
above by $-1$, so it has a largest element $m^*$; the elements of $S$ in block $m^*$
exceed all other elements of $S$, and since that block has reverse well-ordered support,
$S$ has a greatest element. All supports and the index family are sets. Hence the right
side of \eqref{osq:if:eq:remainderdivision} is a legitimate strongly summable Hahn
expression, and \eqref{osq:if:eq:remainderdivision} holds because Hahn multiplication
distributes over strongly summable families. (Equivalently, multiplying the right side by
$1+\omega^a$ cancels consecutive $k$-blocks and leaves exactly $R(\xi)$.)

\emph{Unbounded-below support.} Choose $j$ with $r_j\ne0$ and $h\in\supp r_j$. For every
$k\ge0$ the exponent $h+(j/n-k-1)a$ lies only in block $(j,k)$, so it occurs in
\eqref{osq:if:eq:remainderdivision} with the nonzero coefficient $(-1)^k(r_j)_h$. Since
$h<a$ and $j/n<1$, this exponent is less than $a+a-(k+1)a=(1-k)a$, which is below $-Na$
as soon as $k\ge N+1$. Thus the support of \eqref{osq:if:eq:remainderdivision} is not
contained in $[-Na,\infty)$ for any ordinary $N$, and by \cref{osq:prop:locsupport} the
quotient $R(\xi)/(1+\omega^a)$ does not belong to $\Oz[\omega^{-a}]$.

\emph{Conclusion.} If $R(x_n)=0$, then $R(\xi)\in\ker\widetilde q_a
=(1+\omega^a)\Oz[\omega^{-a}]$ by \cref{osq:if:lem:localization}(ii), so
$R(\xi)=(1+\omega^a)u$ with $u\in\Oz[\omega^{-a}]$. Division in the field $\No$ identifies
$u$ with \eqref{osq:if:eq:remainderdivision}, a contradiction. Hence $R(x_n)\ne0$.
\end{proof}

\begin{remark}[\src{07}]\label{osq:if:rem:hiddenfactors}
The proof does not restrict the exponents of an alleged ambient factor $u$ to rational
multiples of $a$; the full support characterization of $\Oz[\omega^{-a}]$ is exactly what
rules out such hidden factors. A calculation solely in a finite group algebra (the span of
the monomials $\omega^{ja/n}$) would not establish the theorem.
\end{remark}

\begin{corollary}[Internal algebraically closed field \src{07}]\label{osq:if:cor:Ea}
Let $u=x_2=q_a(\omega^{a/2})$. The subring $\EE_a=\FF_a[u]\subseteq\Qa_a$ is a field,
isomorphic over $\FF_a$ to $\FF_a(i)$ by $i\mapsto u$. It is algebraically closed and
contains a class copy of $\SC$.
\end{corollary}

\begin{proof}
By \cref{osq:if:thm:binomialkernel} with $n=2$, $\EE_a\cong\FF_a[T]/(T^2+1)$. Since $\FF_a$
is real closed (\cref{osq:if:prop:closed}), $-1$ is not a square in it, $T^2+1$ is
irreducible, and the quotient is the field $\FF_a+i\FF_a=\FF_a(i)$, which is algebraically
closed by \cref{osq:if:prop:closed}. The complex extension of $\Theta_a$
(\cref{osq:if:thm:compression}) embeds $\SC$ in it.
\end{proof}

There is no order contradiction. Although $\omega^{a/2}$ is positive in $\Oz$, its square
becomes $-1$ in $\Qa_a$; the order of $\Oz$ does not descend to $\Qa_a$. The corollary is a
statement about rings and fields, not about ordered quotients.

More generally, finite real factorization and the Chinese remainder theorem give
\begin{equation}\label{osq:if:eq:finiteCRT}
 \FF_a[T]/(T^n+1)\cong
 \begin{cases}
 \EE_a^{\,m}, & n=2m,\\[2pt]
 \FF_a\times\EE_a^{\,m}, & n=2m+1.
 \end{cases}
\end{equation}
The roots of $T^n+1$ are the $n$ distinct ordinary complex numbers
$\exp(i\pi(2k+1)/n)$, which lie in $\C\subseteq\FF_a(i)$. For even $n$ none is real, and
the $m$ conjugate pairs give $m$ distinct monic quadratic factors with real coefficients,
irreducible over $\FF_a$ because their roots are not in $\FF_a$. For odd $n$ the only real
root is $-1$, and the remaining roots give $m$ such quadratics. All roots are simple, so
the factors are pairwise coprime, and each quadratic quotient is isomorphic to
$\FF_a(i)\cong\EE_a$. The isomorphisms in \eqref{osq:if:eq:finiteCRT} need not identify a
chosen imaginary unit in the same way at different stages; the dyadic tower below makes
coherent choices.

\subsection{The dyadic tower and its Cantor algebra}\label{osq:if:sub:cantor}

For $r\ge1$ define
\[
 \xi_r=q_a(\omega^{a/2^r})=x_{2^r},\qquad \mathcal S_r=\FF_a[\xi_r]\subseteq\Qa_a
\]
(07's $S_r$). Then $\xi_r^2=\xi_{r-1}$ for $r\ge2$, $\mathcal S_r\subseteq\mathcal S_{r+1}$,
and $\mathcal S_1=\EE_a$, so each $\mathcal S_r$ is an $\EE_a$-algebra. In this subsection
$i$ denotes the element $u=\xi_1$ of $\EE_a$; under $\EE_a\cong\FF_a(i)$ the ordinary
complex field $\C\subseteq\FF_a(i)$ corresponds to $\R+\R u\subseteq\EE_a$, and ordinary
complex numbers are read in $\EE_a$ through this identification.

\begin{lemma}[A single stage over the fixed complex field \src{07}]\label{osq:if:lem:stage}
Let $m=2^{r-1}$. As $\EE_a$-algebras,
\begin{equation}\label{osq:if:eq:stage}
 \mathcal S_r\cong\EE_a[T]/(T^m-i),\qquad T\mapsto\xi_r .
\end{equation}
\end{lemma}

\begin{proof}
The $\EE_a$-algebra map $\EE_a[T]\to\mathcal S_r$, $T\mapsto\xi_r$, is surjective, since
$\EE_a[\xi_r]=\FF_a[\xi_r^m,\xi_r]=\mathcal S_r$, and it kills $T^m-i$ because
$\xi_r^m=\xi_1=i$. For injectivity on the quotient, suppose
$\sum_{j<m}(a_j+b_ji)\xi_r^j=0$ with $a_j,b_j\in\FF_a$. Substituting $i=\xi_r^m$ gives
$\sum_{j<m}a_j\xi_r^j+\sum_{j<m}b_j\xi_r^{j+m}=0$, a polynomial relation over $\FF_a$ of
degree less than $2m=2^r$ satisfied by $\xi_r=x_{2^r}$. By
\cref{osq:if:thm:binomialkernel} with $n=2^r$ all its coefficients vanish, so all $a_j$ and
$b_j$ are zero.
\end{proof}

The roots of $T^m-i$ can be chosen in the ordinary complex subfield of $\EE_a$:
\begin{equation}\label{osq:if:eq:zetas}
 \zeta_{r,k}=\exp\Bigl(i\,\frac{\pi/2+2\pi k}{2^{r-1}}\Bigr)\qquad(k\in\Z/2^{r-1}\Z).
\end{equation}
The expression is evaluated in $\C$ and then embedded in $\EE_a$; no exponential function
on surcomplex arguments is assumed. These $m$ roots are distinct, and
\[
 \zeta_{r+1,k}^{\,2}=\zeta_{r,\,k\bmod 2^{r-1}}\qquad(k\in\Z/2^r\Z).
\]
Since the roots are distinct, the Chinese remainder theorem and \eqref{osq:if:eq:stage}
give an isomorphism of $\EE_a$-algebras
\begin{equation}\label{osq:if:eq:stageCRT}
 \Phi_r:\mathcal S_r\longrightarrow\EE_a^{\Z/2^{r-1}\Z},\qquad
 P(\xi_r)\longmapsto\bigl(P(\zeta_{r,k})\bigr)_k\qquad(P\in\EE_a[T]).
\end{equation}
The inclusion $\mathcal S_r\subseteq\mathcal S_{r+1}$ replaces $\xi_r$ by $\xi_{r+1}^2$, so
$P(\xi_r)=P(\xi_{r+1}^2)$ has stage-$(r+1)$ coordinates
$P(\zeta_{r+1,k}^2)=P(\zeta_{r,k\bmod m})$. In coordinates the transition map is therefore
exactly
\[
 (c_k)_{k\bmod m}\longmapsto(c_{k\bmod m})_{k\bmod 2m},
\]
which repeats each coordinate on its two lifts. This compatibility, and not merely the
existence of separate finite product decompositions, is what identifies the direct limit.

\begin{theorem}[Cantor algebra theorem \src{07}]\label{osq:if:thm:cantor}
Let $\LC(\Z_2,\EE_a)$ be the class of functions from the $2$-adic integers $\Z_2$ to $\EE_a$
that factor through some finite quotient $\Z/2^N\Z$, with pointwise operations; thus
$\LC(\Z_2,\EE_a)=\bigcup_{N\ge0}\EE_a^{\Z/2^N\Z}$ with the transition maps repeating each
coordinate on its two lifts. There is an explicit $\EE_a$-algebra isomorphism
\begin{equation}\label{osq:if:eq:cantor}
 \bigcup_{r\ge1}\mathcal S_r\;\cong\;\LC(\Z_2,\EE_a),
\end{equation}
under which $\xi_r$ corresponds to the function $z\mapsto\zeta_{r,\,z\bmod 2^{r-1}}$. The
left side is the subalgebra of $\Qa_a$ generated by $\FF_a$ and the classes of
$\omega^{a/2^r}$, $r\ge1$.
\end{theorem}

\begin{proof}
Identify $\EE_a^{\Z/2^{r-1}\Z}$ with the functions $\Z_2\to\EE_a$ factoring through
$\Z/2^{r-1}\Z$, via $c\mapsto(z\mapsto c_{z\bmod 2^{r-1}})$. By the computation of the
transition maps, the function attached to an element of $\mathcal S_r$ by $\Phi_r$ is the
same as the one attached by $\Phi_{r+1}$, so the $\Phi_r$ combine into a map $\Phi$ on the
increasing union; it is an $\EE_a$-algebra homomorphism, and it is injective because each
$\Phi_r$ is. Conversely, every function factoring through $\Z/2^N\Z$ lies in the image of
the bijection $\Phi_{N+1}$, so $\Phi$ is surjective. Evaluating $T$ at the roots gives the
function attached to $\xi_r$. Finally, $\xi_{r'}=\xi_r^{2^{r-r'}}\in\mathcal S_r$ for
$r'\le r$, so the union of the $\mathcal S_r$ is the subalgebra generated by $\FF_a$ and all
$\xi_r$.
\end{proof}

The isomorphism is an algebraic statement about a specified subalgebra; it is not a
classification of all of $\Qa_a$.

\begin{remark}[Locally constant functions \src{07}]\label{osq:if:rem:lc}
The topological terminology is justified: a locally constant function on the ordinary
compact space $\Z_2$ has a finite subcover of neighbourhoods of constancy and hence factors
through one finite quotient, and conversely such a factorization is locally constant. This
uses only the ordinary Cantor space; no topology is imposed on the proper-class field
$\EE_a$, and one can equally take the finite-factorization condition as the definition.
\end{remark}

\begin{remark}[Evaluation at a point \src{07}]\label{osq:if:rem:evaluation}
Evaluation at a point $z\in\Z_2$ is a homomorphism from the subalgebra
\eqref{osq:if:eq:cantor} onto $\EE_a$. Its target is a proper-class field, not a set-sized
ring, and no assertion is made that it extends to all of $\Qa_a$. The construction is
therefore fully compatible with the absence of small ring images of $\Qa_a$.
\end{remark}

For $m=2^{r-1}$ and $k\in\Z/m\Z$ define, in $\mathcal S_r$,
\begin{equation}\label{osq:if:eq:idempotent}
 e_{r,k}=\frac1m\sum_{\ell=0}^{m-1}\Bigl(\frac{\xi_r}{\zeta_{r,k}}\Bigr)^{\ell}.
\end{equation}
At the root $\zeta_{r,j}$ the sum has value $1$ for $j=k$ and $0$ otherwise, by the finite
geometric-sum identity for the $m$th root of unity $\zeta_{r,j}/\zeta_{r,k}$. Thus
$\Phi_r(e_{r,k})$ is the indicator of the coordinate $k$, and
\begin{align}
 e_{r,k}^2&=e_{r,k},\qquad e_{r,k}e_{r,j}=0\ (k\ne j),\qquad
 \sum_{k\bmod m}e_{r,k}=1,\label{osq:if:eq:idempotentlaws}\\
 e_{r,k}&=e_{r+1,k}+e_{r+1,k+m}.\label{osq:if:eq:refine}
\end{align}
Each $e_{r,k}$ is nonzero because $\Phi_r$ is injective, which rests on
\cref{osq:if:thm:binomialkernel}. The refinement formula \eqref{osq:if:eq:refine} follows
from the coordinate transition, or by checking values at all roots of stage $r+1$.

The first nontrivial instance can be written without trigonometric notation. Put
\[
 y=\xi_2,\qquad u=\xi_1=y^2,\qquad \zeta=\frac{1+u}{\sqrt2}\in\EE_a .
\]
Then $u^2=-1$ and $\zeta^2=(1+2u+u^2)/2=u=y^2$, so $(y/\zeta)^2=1$, and
\begin{equation}\label{osq:if:eq:firstidempotents}
 e_+=\frac12\Bigl(1+\frac y\zeta\Bigr),\qquad e_-=\frac12\Bigl(1-\frac y\zeta\Bigr)
\end{equation}
satisfy $e_\pm^2=e_\pm$, $e_+e_-=\frac14(1-(y/\zeta)^2)=0$ and $e_++e_-=1$. Nonvanishing is
not inferred from the suggestive formula: the roots at stage $2$ are
$\zeta_{2,0}=\zeta$ and $\zeta_{2,1}=-\zeta$, so $e_+=e_{2,0}$, $e_-=e_{2,1}$, and
$\Phi_2$ identifies $\mathcal S_2\cong\EE_a\times\EE_a$ and sends $e_+,e_-$ to $(1,0)$ and
$(0,1)$. This is why the exact-kernel theorem is an essential part of the argument.

\begin{corollary}[Boolean algebra and failure of finiteness conditions \src{07}]
\label{osq:if:cor:boolean}
The idempotents of $\bigcup_r\mathcal S_r$ form precisely the Boolean algebra of clopen
subsets of $\Z_2$, and this Boolean algebra is atomless. The ring $\Qa_a$ has an infinite
sequence of pairwise orthogonal nonzero idempotents, is neither a domain nor a local ring,
and fails the ascending chain condition on ideals.
\end{corollary}

\begin{proof}
A field has only the idempotents $0$ and $1$, so an idempotent of $\LC(\Z_2,\EE_a)$ is the
indicator of a set that is a union of residue classes modulo some $2^N$, that is, of a
clopen subset of $\Z_2$; conversely, by compactness every clopen subset is a finite union
of such residue classes, and its indicator is idempotent. Every nonempty clopen set
contains a residue class modulo some $2^N$, which splits into two nonempty residue classes
modulo $2^{N+1}$; so there are no atoms.

For $j\ge0$ the sets $C_j=\{z\in\Z_2:z\equiv2^j\pmod{2^{j+1}}\}$ (the $z$ of $2$-adic
valuation exactly $j$) are nonempty, clopen and pairwise disjoint. Their indicators give
pairwise orthogonal nonzero idempotents $f_j\in\Qa_a$. The ideals
\[
 (f_0)\subsetneq(f_0,f_1)\subsetneq(f_0,f_1,f_2)\subsetneq\cdots
\]
of $\Qa_a$ strictly increase: if $f_{n+1}=\sum_{j\le n}c_jf_j$ with $c_j\in\Qa_a$,
multiplication by $f_{n+1}$ would give $f_{n+1}=0$. This proves the failure of the
ascending chain condition without forming any class-sized generating family. A nontrivial
idempotent $e$ satisfies $e(1-e)=0$, so $\Qa_a$ is not a domain; and in a local ring (for
every $x$, one of $x$ and $1-x$ is a unit) either $e$ or $1-e$ is a unit, which forces
$e=1$ or $e=0$.
\end{proof}

Grouping a sufficiently fine dyadic partition of $\Z_2$ into any prescribed number $d\ge1$
of nonempty blocks yields a unital subalgebra of $\Qa_a$ isomorphic to $\EE_a^d$.
Manuscript~07 claims neither that all idempotents of $\Qa_a$ lie in the Cantor algebra
\eqref{osq:if:eq:cantor} nor that $\Qa_a$ itself is a ring of locally constant functions.
No infinite sum of the orthogonal sequence $(f_j)$ is used.

\subsection{Gaussian binomial quotients}\label{osq:if:sub:gaussian}

Let
\[
 \Qa_a^\C=\GOz/(1+\omega^a),\qquad q_a^\C:\GOz\to\Qa_a^\C .
\]
Since $\GOz=\Oz+i\Oz$ and $\omega^{-a}$ is real, $\GOz[\omega^{-a}]=\Oz[\omega^{-a}]
+i\,\Oz[\omega^{-a}]$, so by \cref{osq:prop:locsupport} it consists of the surcomplex normal
forms whose supports are bounded below by some $-Na$. It contains $\FF_a(i)$, and the proof of
\cref{osq:if:lem:localization} applies verbatim with $\GOz$ in place of $\Oz$: the
extension $\widetilde q_a^\C$ of $q_a^\C$ has kernel $(1+\omega^a)\GOz[\omega^{-a}]$, and
it embeds the coefficient field by
\[
 f\longmapsto-q_a^\C(\omega^af)\qquad(f\in\FF_a(i)).
\]
For $f=i$ this is the Gaussian imaginary unit of the quotient, since
$-q_a^\C(\omega^ai)=-q_a^\C(\omega^a)q_a^\C(i)=q_a^\C(i)$.

\begin{theorem}[Complex exact kernels and dyadic algebra \src{07}]
\label{osq:if:thm:complexkernel}
For every $n\ge1$ there is an injective ring homomorphism
\[
 \FF_a(i)[T]/(T^n+1)\hookrightarrow\Qa_a^\C,\qquad T\mapsto q_a^\C(\omega^{a/n}),
\]
and its source is isomorphic to $\FF_a(i)^n$. The increasing union of the dyadic stages
$n=2^r$, $r\ge0$, is isomorphic to $\LC(\Z_2,\FF_a(i))$. The ring $\Qa_a^\C$ is nonzero,
has no unital homomorphism to a nonzero set-sized ring, and contains a class copy of $\SC$.
\end{theorem}

\begin{proof}
The proof of \cref{osq:if:thm:binomialkernel} uses only supports and the nonvanishing of
coefficients, not their reality, so it applies verbatim over $\FF_a(i)$ and
$\GOz[\omega^{-a}]$: a nonzero remainder $R$ with coefficients in $\FF_a(i)$ gives the
expansion \eqref{osq:if:eq:remainderdivision}, whose blocks lie in the distinct cosets
$(j/n-k-1)a+H_a$, whose union of supports is reverse well ordered, and whose support is
not bounded below by any $-Na$; so $R(\omega^{a/n})\notin(1+\omega^a)\GOz[\omega^{-a}]$.
The polynomial $T^n+1$ has $n$ distinct roots $\exp(i\pi(2k+1)/n)$ in $\C\subseteq\FF_a(i)$,
so the Chinese remainder theorem gives $\FF_a(i)[T]/(T^n+1)\cong\FF_a(i)^n$ by evaluation
at the roots.

For $n=2^r$ label the roots $\eta_{r,k}=\exp(i(\pi+2\pi k)/2^r)$, $k\in\Z/2^r\Z$. Then
$\eta_{r,k}^2=\eta_{r-1,\,k\bmod 2^{r-1}}$, and the stage-$(r-1)$ generator
$q_a^\C(\omega^{a/2^{r-1}})$ is the square of the stage-$r$ generator; so, exactly as in
\cref{osq:if:thm:cantor}, the coordinate transitions repeat each entry on its two lifts,
and the union is $\LC(\Z_2,\FF_a(i))$. The ring is nonzero because $1+\omega^a$ is not
among the units $\pm1,\pm i$ of $\GOz$ (\cref{osq:lem:units}); it has no nonzero set-sized
image by \cref{osq:thm:reflection}, since $\ct(1+\omega^a)=1$ generates $\Z[i]$; and it
contains a copy of $\SC$ by the complex extension of $\Theta_a$
(\cref{osq:if:thm:compression}).
\end{proof}

Coefficientwise complex conjugation preserves $\GOz$ and, because $1+\omega^a$ is real,
the ideal $(1+\omega^a)$; hence it descends to an involution of $\Qa_a^\C$. It commutes
with the coefficient-field embedding $f\mapsto-q_a^\C(\omega^af)$ because the multiplier
$-\omega^a$ is real. This supplies a compatible involution without requiring an order or an
analytic interpretation of the quotient.

\subsection{Boundaries of the large-field results}\label{osq:if:sub:boundaries}

The small-image theorem and the large-field results concern different hypotheses and
reinforce one another (07). A homomorphism to a set-sized ring kills every purely infinite
element (\cref{osq:thm:universal}); a homomorphism to a nonzero ring that makes even one
purely infinite element invertible has a huge field in its target
(\cref{osq:if:thm:transfer}). That target is necessarily a proper class, but the explicit
field embedding gives considerably more than this size conclusion, and the Cantor algebra
describes zero-divisor structure that no cardinality argument captures. Manuscript~07
insists on the following distinctions.

\paragraph{A nonzero element versus a unit.} A nonzero nilpotent, as in
\cref{osq:if:ex:nilpotent}, does not produce a field. Invertibility is what allows the
formula $f=(gf)/g$ to be transported to a quotient. No cancellation in a ring with zero
divisors has been used.

\paragraph{A specified subalgebra versus the entire quotient.} The exact kernel theorem
identifies the subalgebras $\FF_a[x_n]$, and their dyadic union is a Cantor algebra.
Surreal normal forms with other supports can still create further elements of $\Qa_a$.
Neither a classification of all idempotents of $\Qa_a$ nor a description of its prime
ideals has been proved.

\paragraph{Class ideals versus set-theoretic Zorn arguments.} Zorn's lemma applies to sets
of candidates. The ring $\Qa_a$ is a proper class, so the usual existence argument for
maximal ideals cannot simply be repeated over its class of ideals. Statements such as
\cref{osq:if:cor:primefield} are deliberately conditional: it is not concluded that a
proper-class nonarithmetic maximal ideal exists merely because $(1+\omega^a)$ is proper.
Likewise (06) \cref{osq:if:prop:reservoir} is a necessary condition on residue fields, not an
existence theorem.

\paragraph{Strong sums versus limits.} Every infinite sum used above, namely the geometric
series for $(1+\omega^a)^{-1}$ and the block expansion \eqref{osq:if:eq:remainderdivision},
has a set-sized reverse well-ordered support and finitely many contributions at each
exponent, and is verified coefficientwise. Finite geometric remainders of the kind in
\cref{osq:lem:telescope} do not tend to zero, and no limit argument is used.

\paragraph{Finite checks.} The program accompanying manuscript~07 checks exact finite
geometric identities, finite binomial remainders, the first nontrivial idempotents
\eqref{osq:if:eq:firstidempotents} and finite coordinate refinement maps (07 records that
all $621$ named checks passed). It does not represent arbitrary surreal normal forms, and
it does not verify the ambient support argument, namely that the infinite expansion
\eqref{osq:if:eq:remainderdivision} lies outside $\Oz[\omega^{-a}]$; that is a
mathematical proof in the text. For a formalization, 07 recommends keeping the support
characterization of $\Oz[\omega^{-a}]$ and the exact binomial kernel separate from the
routine finite Chinese remainder calculations; no corresponding Lean declarations are
claimed.

\section{Integral closure, complete integral closure, and invisible equations}
\label{osq:sec:closure}

The universal theorem uses exponent scales \emph{below} the supports of purely infinite
elements. This section uses a scale \emph{above} all negative exponents, which yields one
denominator for a whole set-generated algebra (03). The consequences separate three
notions that coincide for Noetherian domains: membership in $\Oz$, integrality over
$\Oz$, and almost integrality over $\Oz$. The complete integral closure of $\Oz$ is all
of $\No$, while its integral closure $\NN$ is a proper, fine-dense subring that no set of
generators produces (03). The section ends with a monic equation of 08 that has no root in
$\GOz$ although its image in every set-sized ring has the root $1$; its positive root is an
explicit element of $\NN\setminus\Oz$.

\paragraph{Standing notation.}
As in \cref{osq:sec:model03}, $\kappa$ is an infinite cardinal with
$\mu=\cf(\kappa)>\aleph_0$, $\theta=\kappa^{\aleph_0}$, $G_\kappa$ is the two-armed
lexicographic group of 03, $F^{03}_\kappa=\Hc{\R}{G_\kappa}$ is the field of
countably supported Hahn series over $G_\kappa$, identified with its image in $\No$, and
\[
 A^{03}_\kappa=F^{03}_\kappa\cap\Oz=\Z+\Pi^{03}_\kappa,
\]
where $\Pi^{03}_\kappa$ consists of the elements of $F^{03}_\kappa$ with strictly positive
support. We write $X^g$ ($g\in G_\kappa$) for the monomials of $F^{03}_\kappa$; under the
identification they are surreal monomials. The Gaussian versions are
$F^{03}_\kappa[i]=\Hc{\C}{G_\kappa}$, $A^{03}_\kappa[i]=\Z[i]+\Pi^{03}_\kappa+i\Pi^{03}_\kappa$
and the ideal $\Pi^{03}_\kappa+i\Pi^{03}_\kappa$. The field $F^{03}_\kappa$ is real closed
and $F^{03}_\kappa[i]$ is algebraically closed (\cref{osq:sec:model03}). We write
$\overline\Z$ for the ring of complex algebraic integers and
$\overline\Z_\R=\overline\Z\cap\R$ for the real ones. For an inclusion of rings
$A\subseteq R$ the \emph{conductor} is $(A:R)=\{a\in A:aR\subseteq A\}$.

\subsection{Uniform denominators}\label{osq:nm:sub:denominators}

The clearing statement \cref{osq:prop:fractions}(i) clears one set of surreals. Applied to
a well-chosen set it clears a whole algebra.

\begin{theorem}[One denominator for a set-generated algebra \src{03}]\label{osq:nm:thm:onedenominator}
For every set $S\subseteq\No$ there is a positive monomial $d=\omega^h\in\Pi$ with
\begin{equation}\label{osq:nm:eq:fullabsorption}
 d\,\Oz[S]\subseteq\Pi .
\end{equation}
For every set $S\subseteq\SC$ there is such a $d$ with
\begin{equation}\label{osq:nm:eq:complexabsorption}
 d\,\GOz[S]\subseteq\Pi_\C .
\end{equation}
In particular $d$ is a nonzero element of the conductor $(\Oz:\Oz[S])$, respectively
$(\GOz:\GOz[S])$.
\end{theorem}

\begin{proof}
Although $\Oz[S]$ is in general a proper class, the ring $\Z[S]$ is a set: it consists of
the values of ordinary integer polynomials in finitely many members of $S$. By
\cref{osq:prop:fractions}(i) there is $h>0$ with $\omega^h\Z[S]\subseteq\Pi$. Every element
of $\Oz[S]$ is a finite sum $\sum_ja_jb_j$ with $a_j\in\Oz$ and $b_j\in\Z[S]$, and
\[
 \omega^h\sum_ja_jb_j=\sum_ja_j(\omega^hb_j)\in\Pi
\]
because $\Pi$ is an ideal of $\Oz$ (\cref{osq:lem:ct}). The complex statement is the same
argument with the set $\Z[i][S]$, the complex clause of \cref{osq:prop:fractions}(i) and
the ideal $\Pi_\C$ of $\GOz$. Finally $\Pi\subseteq\Oz$ and $\Pi_\C\subseteq\GOz$.
\end{proof}

\begin{remark}[Conductors \src{03}]\label{osq:nm:rem:conductor}
The theorem gives more than a nonzero element of the conductor $(\Oz:\Oz[S])$: it gives
a monomial whose products with the whole algebra land in the strictly positive-support
ideal $\Pi$. The real clearing statement and its proof are due to manuscript~05 (now in
the sibling report \cite{RepoOdg}); the algebra formulation and the closure consequences
below are 03's.
\end{remark}

\begin{theorem}[The bounded-size version \src{03}]\label{osq:nm:thm:smallabsorption}
Let $A$ be $A^{03}_\kappa$ or $A^{03}_\kappa[i]$, with ambient field $F=F^{03}_\kappa$,
respectively $F^{03}_\kappa[i]$, constant ring $D=\Z$, respectively $\Z[i]$, and ideal
$P=\Pi^{03}_\kappa$, respectively $\Pi^{03}_\kappa+i\Pi^{03}_\kappa$. For every set
$S\subseteq F$ with $|S|<\mu=\cf(\kappa)$ there is $h>0$ in $G_\kappa$ with
\begin{equation}\label{osq:nm:eq:smallabsorption}
 X^h A[S]\subseteq P .
\end{equation}
In particular $\Frac(A^{03}_\kappa)=F^{03}_\kappa$ and
$\Frac(A^{03}_\kappa[i])=F^{03}_\kappa[i]$.
\end{theorem}

\begin{proof}
The ring $D[S]$ has cardinality at most $\max(|S|,\aleph_0)<\mu$, since $\mu$ is
uncountable. Each of its elements has countable support (in the complex case, take the
union of the supports of the real and imaginary coordinates). The union $U$ of these
supports is a union of fewer than $\mu$ countable sets, so $|U|<\mu$ because $\mu$ is
regular and uncountable. By \cref{osq:lem:twosided}(i) the family
$\{-g:g\in U\}\cup\{0\}$, which has fewer than $\mu$ members, has an upper bound
$b\in G_\kappa$; put $h=b+e$ for any $e>0$ in $G_\kappa$. Then $h>0$ and $h>-g$ for every
$g\in U$. Multiplication by $X^h$ translates each support by $h$, so it maps every element
of $D[S]$ to a countably supported series with strictly positive support, that is, into
$P$. Every element of $A[S]$ is a finite sum $\sum_ja_jb_j$ with $a_j\in A$ and
$b_j\in D[S]$, and $X^h\sum_ja_jb_j=\sum_ja_j(X^hb_j)\in P$ because $P$ is an ideal of $A$.

For $x\in F$ apply this to $S=\{x\}$: then $x=(X^hx)/X^h$ is a quotient of elements of
$A$, and $X^h\in P\subseteq A$ is nonzero.
\end{proof}

This supplies the fraction-field clause $\Frac(A^{03}_\kappa)=F^{03}_\kappa$ of 03's
cardinal realization (compare \cref{osq:thm:model03}). Clearing one element is weaker than
clearing \emph{all powers} of that element by the same denominator; both theorems supply
the latter as well, since $x^n\in A[x]$ for every $n$.

\subsection{Almost integrality and complete integral closure}\label{osq:nm:sub:complete}

\begin{definition}[\src{03}]\label{osq:nm:def:almost}
Let $A$ be a domain contained in a field $F$. An element $x\in F$ is \emph{almost integral}
over $A$ if there is $0\ne d\in A$ with $dx^n\in A$ for every ordinary $n\ge0$. When
$F=\Frac(A)$, the set $A^*$ of these elements is the \emph{complete integral closure} of
$A$. The \emph{integral closure} $\IC_F(A)$ consists instead of the roots in $F$ of monic
polynomials in $A[T]$.
\end{definition}

These are the standard notions of commutative algebra; they differ in non-Noetherian
settings, and we take the terminology from \cite{BGR} without importing any result of
that paper as a proof. The quantifier over powers ranges over ordinary finite powers, not
over ordinal iterations or surreal exponentiation. Each individual witness $d$ and each
monic polynomial is set-sized data, so the definitions apply verbatim to the class rings
$\Oz$ and $\GOz$.

\begin{theorem}[Every ambient element is almost integral \src{03}]\label{osq:nm:thm:complete}
Relative to their fraction fields,
\begin{align}
 \Oz^*&=\No,& \GOz^*&=\SC,\label{osq:nm:eq:fullcic}\\
 (A^{03}_\kappa)^*&=F^{03}_\kappa,& (A^{03}_\kappa[i])^*&=F^{03}_\kappa[i].\label{osq:nm:eq:localcic}
\end{align}
More generally, every ring $R$ with $A\subseteq R\subseteq F$, where $(A,F)$ is one of
the four pairs $(\Oz,\No)$, $(\GOz,\SC)$, $(A^{03}_\kappa,F^{03}_\kappa)$,
$(A^{03}_\kappa[i],F^{03}_\kappa[i])$, has $\Frac(R)=F$ and $R^*=F$.
\end{theorem}

\begin{proof}
The fraction fields are the indicated ambient fields by \cref{osq:prop:fractions}(ii) and
\cref{osq:nm:thm:smallabsorption}, so there are no further elements to consider. For $x$
in the ambient field apply \cref{osq:nm:thm:onedenominator} or
\cref{osq:nm:thm:smallabsorption} to $S=\{x\}$. The single nonzero denominator $d$
supplied there satisfies $dx^n\in dA[x]\subseteq A$ for every ordinary $n\ge0$, so $x$ is
almost integral. If $A\subseteq R\subseteq F$, then $F=\Frac(A)\subseteq\Frac(R)\subseteq
F$, and the same $d\in A\subseteq R$ gives $dx^n\in A\subseteq R$.
\end{proof}

\begin{example}[Two explicit denominators \src{03}]\label{osq:nm:ex:almost}
The rational $1/2$ is almost integral over $\Oz$, witnessed by $d=\omega$: indeed
$\omega/2^n\in\Pi$ for all $n\ge0$. The monomial $\omega^{-1}$ is almost integral,
witnessed by $d=\omega^\omega$: indeed $\omega^{\omega-n}\in\Pi$ for every ordinary $n$,
since $\omega-n>0$. Neither claim uses a limit as $n$ tends to an infinite ordinal.
Neither element is integral (\cref{osq:nm:lem:reciprocal} below).
\end{example}

\subsection{Integral closure: reciprocals, small units and density}\label{osq:nm:sub:integral}

The geometric statements are proved first for an arbitrary integer part of a real closed
field; this isolates what is a consequence of integer-part geometry rather than of the
proper-class size of $\No$ (03).

\begin{lemma}[An integral reciprocal was already present \src{03}]\label{osq:nm:lem:reciprocal}
Let $A$ be a domain in a field $F$ and $0\ne a\in A$. If $a^{-1}$ is integral over $A$,
then $a^{-1}\in A$. Hence the integral closure of $A$ in $F$ contains the inverse of no
nonunit of $A$.
\end{lemma}

\begin{proof}
Suppose $(a^{-1})^n+c_{n-1}(a^{-1})^{n-1}+\dots+c_0=0$ with $c_j\in A$ and $n\ge1$.
Multiplying by $a^{n-1}$ gives
\[
 a^{-1}=-c_{n-1}-c_{n-2}a-\dots-c_0a^{n-1}\in A. \qedhere
\]
\end{proof}

Consequently $1/2$ is not integral over $\Oz$, $A^{03}_\kappa$, $\GOz$ or
$A^{03}_\kappa[i]$, despite \cref{osq:nm:ex:almost}: $2$ is a nonunit of each of these
rings by \cref{osq:lem:units} (the units of $A^{03}_\kappa\subseteq\Oz$ are among
$\Oz^\times=\{\pm1\}$, and similarly in the Gaussian case). More generally, for $h>0$ and
$0\ne c\in\R$ the negative monomial $c\,\omega^{-h}$ is not integral over $\Oz$, and for
$0\ne c\in\C$ it is not integral over $\GOz$: its reciprocal $c^{-1}\omega^h$ lies in the
purely infinite ideal $\Pi$, respectively $\Pi_\C$, and is a nonunit. The same holds for
$cX^{-h}$ over $A^{03}_\kappa$ and $A^{03}_\kappa[i]$ when $h>0$ in $G_\kappa$.

\begin{lemma}[Integer parts \src{03}]\label{osq:nm:lem:intpart}
Recall that a subring $A$ of an ordered field $F$ is an \emph{integer part} if for every
$x\in F$ there is a unique $\lfloor x\rfloor\in A$ with $\lfloor x\rfloor\le x<\lfloor
x\rfloor+1$. Then $\Oz$ is an integer part of $\No$ and $A^{03}_\kappa$ is an integer part
of $F^{03}_\kappa$. Every nonzero element of $\Oz$ has absolute value at least $1$, and
every nonzero element of $\GOz$ has surreal modulus $|a+ib|=\sqrt{a^2+b^2}$ at least $1$.
\end{lemma}

\begin{proof}
A nonzero element of $\Oz$ is either a nonzero ordinary integer or has positive degree
(\cref{osq:lem:units}); in the second case its absolute value exceeds every real number.
Hence $|a|\ge1$ for $0\ne a\in\Oz$, and for $a+ib\in\GOz$ with $a,b\in\Oz$ not both zero,
$\sqrt{a^2+b^2}\ge\max(|a|,|b|)\ge1$. This lower bound also gives uniqueness of floors:
two elements of $\Oz$ in a half-open interval of length $1$ differ by an element of
absolute value less than $1$. For existence, write the normal form of $x\in\No$ uniquely
as $x=P+r+\epsilon$ with $P\in\Pi$, $r\in\R$ and $\supp\epsilon$ strictly negative, so
that $\epsilon$ is zero or infinitesimal. Take $\lfloor x\rfloor=P+\lfloor r\rfloor$,
except when $r\in\Z$ and $\epsilon<0$, where $\lfloor x\rfloor=P+r-1$. In every case
$\lfloor x\rfloor\le x<\lfloor x\rfloor+1$. For $x\in F^{03}_\kappa$ the same decomposition
has $P\in\Pi^{03}_\kappa$, so the floor lies in $A^{03}_\kappa$.
\end{proof}

The exceptional case in the floor formula matters: discarding all negative exponents sends
$-\omega^{-1}$ to $0$, which is not its floor.

\paragraph{Why integral elements form a ring.}
That the integral elements over a commutative ring $A$ form a ring, and that integrality is
transitive, is finite algebra, which we recall because it is applied to class rings. If
$x$ satisfies a monic equation of degree $n$ over $A$, then $A[x]=A+Ax+\dots+Ax^{n-1}$ is
a finitely generated $A$-module, by reducing powers with that equation. If $x$ and $y$ are
integral, $A[x,y]$ is generated by finitely many monomials $w_1=1,w_2,\dots,w_N$ in $x$
and $y$. For $z\in A[x,y]$ write $zw_k=\sum_la_{kl}w_l$ with $a_{kl}\in A$; the determinant
trick shows that $\det(zI-(a_{kl}))$ annihilates every $w_l$, in particular $w_1=1$, so it
is zero, and this is a monic equation for $z$. The same argument, applied to the finitely
many coefficients of a monic equation and to the element, proves transitivity. Only
finitely many elements are involved at each step, so the arguments apply unchanged to
$\Oz$, $\GOz$ and their overrings in $\No$ and $\SC$ (03).

\begin{theorem}[Small integral units \src{03}]\label{osq:nm:thm:smallunits}
Let $F$ be a real closed ordered field, $A\subseteq F$ an integer part, and $N$ the
integral closure of $A$ in $F$. For every $H\in A_{>0}$ the element
$u_H=\sqrt{H^2+1}-H$ is a positive unit of $N$ satisfying
\begin{equation}\label{osq:nm:eq:smallunit}
 u_H^2+2Hu_H-1=0,\qquad u_H^{-1}=u_H+2H,\qquad 0<u_H<\frac1{2H}.
\end{equation}
For every $\epsilon\in F_{>0}$ there is such a unit with $0<u_H<\epsilon$.
\end{theorem}

\begin{proof}
Since $\sqrt{H^2+1}>H>0$, the element $u_H$ is positive. Squaring
$u_H+H=\sqrt{H^2+1}$ gives the first equation, which is monic over $A$; so $u_H\in N$.
The same equation says $u_H(u_H+2H)=1$, and $u_H+2H\in N$ because $N$ is a ring containing
$A$; so $u_H$ is a unit of $N$. Rationalizing,
\[
 u_H=\frac{1}{\sqrt{H^2+1}+H}<\frac1{2H}.
\]
An integer part is cofinal in $F$: $\lfloor 1/(2\epsilon)\rfloor+1\in A$ is positive and
exceeds $1/(2\epsilon)$. Taking $H$ to be this element gives $u_H<1/(2H)<\epsilon$.
\end{proof}

\begin{theorem}[Dense normalization of an integer part \src{03}]\label{osq:nm:thm:density}
Under the hypotheses of \cref{osq:nm:thm:smallunits}, $N$ is dense in $F$ for the order
topology. More precisely, for every $x\in F$ and every $\epsilon\in F_{>0}$ there is
$y\in N$ with
\begin{equation}\label{osq:nm:eq:density}
 0\le x-y<\epsilon .
\end{equation}
\end{theorem}

\begin{proof}
Choose a positive unit $u<\epsilon$ of $N$ by \cref{osq:nm:thm:smallunits} and let
$n=\lfloor x/u\rfloor\in A$. Multiplying $n\le x/u<n+1$ by $u>0$ gives $0\le x-un<u<\epsilon$;
put $y=un\in N$. To meet a prescribed open interval $(a,b)$, apply this with
$x=(a+b)/2$ and $\epsilon<(b-a)/2$; then $y\in(x-\epsilon,x]\subseteq(a,b)$.
\end{proof}

No completion and no sequential argument is used. The same proof is valid for the order
topology of $\No$ (the fine topology), with all positive surreal radii. Results on dense
subfields and integer parts in other settings provide context \cite{Kuhlmann}; the object
here is the integral closure of a specified integer part, not an arbitrary dense subfield,
and the normalization theorem is not attributed to that work. The small-unit argument is
a general elementary lemma about integer parts, and 03 makes no priority claim for it.

\paragraph{The normalizations.}
We write
\begin{equation}\label{osq:nm:eq:normalizations}
 \NN=\IC_{\No}(\Oz),\qquad \NN_\kappa=\IC_{F^{03}_\kappa}(A^{03}_\kappa).
\end{equation}
These are 03's $B$ and $B_\kappa$. Manuscript~03 wrote $\Int_F(A)$ for the integral
closure of $A$ in $F$; we write $\IC_F(A)$ because $\Int$ denotes rings of integer-valued
polynomials, such as $\Int(\Z^r)$, elsewhere in this report. Note that $\NN_\kappa$ is
defined relative to $F^{03}_\kappa$; only $\NN_\kappa\subseteq\NN\cap F^{03}_\kappa$ is
automatic, and no descent equality is claimed.

\begin{theorem}[Set-generated intermediate algebras are uniformly discrete \src{03}]\label{osq:nm:thm:discrete}
For every set $S\subseteq\No$ the algebra $C=\Oz[S]$ has a positive surreal separation
bound $r$:
\[
 |x-y|\ge r\qquad\text{for all distinct }x,y\in C .
\]
It is a closed, uniformly discrete subclass of $\No$ in the fine order topology. The same
holds for $A^{03}_\kappa[S]\subseteq F^{03}_\kappa$, in the order topology of
$F^{03}_\kappa$, whenever $|S|<\mu$.
\end{theorem}

\begin{proof}
Choose the positive monomial $d$ of \cref{osq:nm:thm:onedenominator}, respectively
$d=X^h$ of \cref{osq:nm:thm:smallabsorption}. For distinct $x,y\in C$ the product
$d(x-y)$ is a nonzero element of $\Pi$ (respectively $\Pi^{03}_\kappa$), hence of the
integer part, so $|d(x-y)|\ge1$ by \cref{osq:nm:lem:intpart}. Thus $|x-y|\ge1/d=:r$.

Uniform discreteness gives an isolating interval of radius $r/3$ around each point. For
closedness, let $z\notin C$. If no point of $C$ lies within $r/3$ of $z$, that
neighbourhood misses $C$. Otherwise let $c\in C$ be such a point; every other point of $C$
is at distance greater than $2r/3$ from $z$ by the triangle inequality, so a radius
smaller than both $r/3$ and $|z-c|$ gives a neighbourhood of $z$ disjoint from $C$.
\end{proof}

This is a class-size statement about a proper-class base ring; it is not the familiar
assertion that a finite extension of a discrete ring is discrete in an Archimedean
topology (03).

\begin{theorem}[Proper dense normalization, zero conductor, no set of generators \src{03}]\label{osq:nm:thm:normalization}
For the full omnific ring,
\begin{equation}\label{osq:nm:eq:normalization}
 \Oz\subsetneq\NN\subsetneq\No,\qquad \overline{\NN}^{\,\mathrm{fine}}=\No,\qquad
 (\Oz:\NN)=0,
\end{equation}
and $\NN$ is not generated as an $\Oz$-algebra by any set of its elements. For the
set-sized ring $A^{03}_\kappa$, the normalization $\NN_\kappa$ satisfies
$A^{03}_\kappa\subsetneq\NN_\kappa\subsetneq F^{03}_\kappa$, is dense in $F^{03}_\kappa$,
has conductor $(A^{03}_\kappa:\NN_\kappa)=0$, and requires at least $\mu$ algebra
generators over $A^{03}_\kappa$.
\end{theorem}

\begin{proof}
The element $\sqrt2\in\R\subseteq F^{03}_\kappa$ is integral, being a root of $T^2-2$,
and it is in neither integer part, whose real elements are the ordinary integers; so the
first inclusions are strict. The element $1/2$ is not integral by
\cref{osq:nm:lem:reciprocal}, so the second inclusions are strict. The fields $\No$ and
$F^{03}_\kappa$ are real closed and $\Oz$, $A^{03}_\kappa$ are integer parts
(\cref{osq:nm:lem:intpart}); density is \cref{osq:nm:thm:density}.

For the conductor, fix $a\ne0$ in the integer part. Choose a positive unit $u$ of the
normalization with $u<1/|a|$ (\cref{osq:nm:thm:smallunits}). Then $au$ lies in the
normalization and $0<|au|<1$, so $au$ is not in the integer part by
\cref{osq:nm:lem:intpart}. Hence $a$ is not in the conductor, and this holds for every
$a\ne0$.

If $\NN=\Oz[S]$ for a set $S$, then \cref{osq:nm:thm:onedenominator} would give a nonzero
element $d$ of the conductor, a contradiction. (Equivalently, $\NN$ would be uniformly
discrete by \cref{osq:nm:thm:discrete}, contradicting density.) The same argument with
\cref{osq:nm:thm:smallabsorption} shows that $\NN_\kappa\ne A^{03}_\kappa[S]$ whenever
$|S|<\mu$.
\end{proof}

Integral closure therefore changes the fine geometry of the ring, although $\NN$ still
contains neither $1/2$ nor any nonzero negative monomial $c\,\omega^{-h}$. For
$\kappa=\aleph_1$ the theorem says, without the continuum hypothesis, that the
normalization cannot be obtained by adjoining countably many elements to the
continuum-sized ring $A^{03}_{\aleph_1}$ (here $\theta=\aleph_1^{\aleph_0}=2^{\aleph_0}$).
For general $\kappa$ the theorem gives only the \emph{lower bound} $\mu$ on the number of
algebra generators of $\NN_\kappa$, not an exact generation number; the exact count $\mu$
of 03's cardinal realization (\cref{osq:thm:model03}) concerns the ideal
$\Pi^{03}_\kappa$, not the normalization.

\begin{remark}[Dense does not mean sequentially approximable \src{03}]\label{osq:nm:rem:dense}
The fine topology of $\No$ has the size obstructions discussed in the foundations report
\cite{RepoFound}. The density statement chooses a point separately for each positive
surreal radius; it does not produce a nonconstant convergent set-indexed net. Indeed, every
fixed set of elements of $\NN$ lies in a single uniformly discrete algebra $\Oz[S]$ by
\cref{osq:nm:thm:discrete}. The dense normalization requires access to more than any
fixed set of its additional generators.
\end{remark}

\subsection{Constant slices, truncation, and a unit equation}\label{osq:nm:sub:slices}

The results above do not give a complete normal-form test for integrality. They do give
exact tests on the nonnegative-support slice, and they show why leading-term tests fail
(03). Recall $\A_{\R,\R}=\R+\Pi$ (03's $P_\R$), on which $\ct$ is a ring homomorphism onto
$\R$ by \cref{osq:lem:ct}, although $\A_{\R,\R}$ is not an integer part; its countable
version is $\R+\Pi^{03}_\kappa$.

\begin{theorem}[The real constant and support slices \src{03}]\label{osq:nm:thm:realslice}
\begin{equation}\label{osq:nm:eq:realslice}
 \NN\cap\A_{\R,\R}=\overline\Z_\R+\Pi=\A_{\overline\Z_\R,\R},\qquad
 \NN\cap\R=\overline\Z_\R,\qquad \NN\cap\Q=\Z .
\end{equation}
The same formulas hold for $\NN_\kappa$, with $\A_{\R,\R}$ and $\Pi$ replaced by
$\R+\Pi^{03}_\kappa$ and $\Pi^{03}_\kappa$ in the first formula.
\end{theorem}

\begin{proof}
Let $x=r+P$ with $r\in\R$ and $P\in\Pi$, and suppose
$x^n+a_{n-1}x^{n-1}+\dots+a_0=0$ with $a_j\in\Oz$. All terms lie in $\A_{\R,\R}$, and
applying the ring homomorphism $\ct$ gives $r^n+\ct(a_{n-1})r^{n-1}+\dots+\ct(a_0)=0$, a
monic equation over $\Z$. So $r\in\overline\Z_\R$. Conversely, a real algebraic integer $r$
is integral over $\Z\subseteq\Oz$, and $P\in\Pi\subseteq\Oz$, so $r+P$ is integral over
$\Oz$ because integral elements form a ring. This is the first equality; intersecting with
$\R$ gives the second, and the third holds because a rational root of a monic integer
polynomial is an integer. For $\NN_\kappa$ the coefficients lie in
$A^{03}_\kappa\subseteq\Oz$, the element lies in $\R+\Pi^{03}_\kappa$, and the same
argument applies; conversely $\overline\Z_\R\subseteq\R\subseteq F^{03}_\kappa$ and
$\Pi^{03}_\kappa\subseteq A^{03}_\kappa$.
\end{proof}

\begin{example}[\src{03}]\label{osq:nm:ex:slice}
The element $\pi\omega+\sqrt2$ is integral over $\Oz$, whereas $\pi\omega+\pi$ is not. The
nature of a coefficient at a \emph{positive} exponent is irrelevant to the test; only the
ordinary constant coefficient is constrained. The element $\omega+1/2$ is almost integral
(\cref{osq:nm:thm:complete}) but not integral, since $1/2\notin\overline\Z_\R$.
\end{example}

\begin{proposition}[Normalization is not truncation closed \src{03}]\label{osq:nm:prop:truncation}
Let $h>0$ and $H=\omega^h$ (or $H=X^h$ with $h>0$ in $G_\kappa$). The integral unit $u_H$
of \cref{osq:nm:thm:smallunits} has the normal form
\begin{equation}\label{osq:nm:eq:binomialunit}
 u_H=\sum_{n\ge1}\binom{1/2}{n}\omega^{-(2n-1)h}
 =\tfrac12\omega^{-h}-\tfrac18\omega^{-3h}+\tfrac1{16}\omega^{-5h}
 -\tfrac5{128}\omega^{-7h}+\tfrac7{256}\omega^{-9h}-\cdots .
\end{equation}
Its leading monomial $\frac12\omega^{-h}$ is not integral over the integer part.
Therefore neither $\NN$ nor $\NN_\kappa$ is closed under leading normal-form truncation.
\end{proposition}

\begin{proof}
By \cref{osq:nm:lem:binomial} below with $m=2$ and $\gamma=2h$, the series
$s=\omega^h\mathsf b_2(\omega^{-2h})$ satisfies $s^2=\omega^{2h}(1+\omega^{-2h})=H^2+1$.
Its leading term is $\omega^h>0$, so $s$ is the positive square root $\sqrt{H^2+1}$ (an
ordered field has at most one positive square root of a given element). Hence
\[
 u_H=s-\omega^h=\omega^h\bigl(\mathsf b_2(\omega^{-2h})-1\bigr)
 =\sum_{n\ge1}\binom{1/2}{n}\omega^{h-2nh},
\]
a series with countable, reverse well-ordered support $\{-(2n-1)h:n\ge1\}$. The first five
binomial coefficients are $\binom{1/2}{1}=\frac12$, $\binom{1/2}{2}=-\frac18$,
$\binom{1/2}{3}=\frac1{16}$, $\binom{1/2}{4}=-\frac5{128}$ and
$\binom{1/2}{5}=\frac7{256}$. In the countable model the same computation takes place in
$F^{03}_\kappa$. The leading monomial $\frac12\omega^{-h}$ is the reciprocal of the nonunit
$2\omega^h$ of the integer part, so it is not integral by \cref{osq:nm:lem:reciprocal}.
\end{proof}

This is stronger than the statement that some integral elements have negative exponents:
their negative tails can be essential, since deleting all but the leading term destroys
integrality. In the opposite direction, \cref{osq:nm:thm:realslice} gives a complete
constant-term criterion when negative exponents are absent (03).

\begin{proposition}[A unit equation acquires solutions only after normalization \src{03}]\label{osq:nm:prop:unitdioph}
Let $A$ be $\Oz$ or $A^{03}_\kappa$ and $0\ne H\in A$. The equation
\begin{equation}\label{osq:nm:eq:unitdioph}
 U^2+2HU-1=0
\end{equation}
has no solution in $A$. For $H>0$ its two roots in the real closed ambient field are
$u_H$ and $-2H-u_H$; both are integral units and belong to the normalization.
\end{proposition}

\begin{proof}
A solution in $A$ would satisfy $U(U+2H)=1$, so $U$ would be a unit of $A$, hence
$U=\pm1$ (\cref{osq:lem:units}; $A^{03}_\kappa\subseteq\Oz$). Substituting $U=1$ gives
$2H=0$ and substituting $U=-1$ gives $-2H=0$, both contrary to $H\ne0$. The two roots are
$u_H$ and $-2H-u_H$ because their sum is $-2H$. Each is integral by the monic equation, and
their product is $-1$, so each is a unit of the normalization.
\end{proof}

Thus membership, integrality and almost integrality give three different arithmetic
environments even for a single quadratic equation with a parameter. This does not
contradict constant-term transfer for systems with ordinary integer coefficients: here a
possibly infinite coefficient $H$ is fixed, and changing the ambient ring changes the
allowed witnesses (03).

\subsection{The surcomplex and Gaussian versions}\label{osq:nm:sub:complex}

The Gaussian omnific ring $\GOz=\Oz+i\Oz=\Z[i]\oplus\Pi_\C$ is the natural Gaussian
analogue of $\Oz$; it is not discretely ordered. It has the split Gaussian constant term
$\ct:\GOz\to\Z[i]$ (\cref{osq:lem:ct}), its units are $\pm1,\pm i$ (\cref{osq:lem:units}),
and each nonzero element has surreal modulus at least $1$ (\cref{osq:nm:lem:intpart}). Put
\begin{equation}\label{osq:nm:eq:complexnormalization}
 \NN_\C=\IC_{\SC}(\GOz),\qquad
 \NN_{\kappa,\C}=\IC_{F^{03}_\kappa[i]}(A^{03}_\kappa[i]),
\end{equation}
03's $C$ and $C_\kappa$. Since $i$ is integral over $\Oz$, transitivity shows that these are
also the integral closures of $\Oz$ in $\SC$ and of $A^{03}_\kappa$ in $F^{03}_\kappa[i]$.
Complex conjugation is a ring automorphism of $\SC$ preserving $\GOz$ (and of
$F^{03}_\kappa[i]$ preserving $A^{03}_\kappa[i]$), so it preserves $\NN_\C$ and
$\NN_{\kappa,\C}$. Consequently
\begin{equation}\label{osq:nm:eq:realintersection}
 \NN_\C\cap\No=\NN,\qquad \NN_{\kappa,\C}\cap F^{03}_\kappa=\NN_\kappa .
\end{equation}
Complexification nevertheless does not commute with normalization, and the failure is
located at the prime $2$.

\begin{proposition}[The complexification defect is killed by $2$ \src{03}]\label{osq:nm:prop:complexdefect}
There are strict inclusions of additive subclasses
\begin{equation}\label{osq:nm:eq:complexdefect}
 \NN[i]\subsetneq\NN_\C\subsetneq\tfrac12\NN[i],\qquad \NN_\C[1/2]=\NN[1/2][i].
\end{equation}
The last object of the chain is a fractional module, not asserted to be a ring. In
particular $2\NN_\C\subseteq\NN[i]$ although $\NN_\C\ne\NN[i]$. The same inclusions hold
for $\NN_\kappa$ and $\NN_{\kappa,\C}$, and the set-sized quotient
$\NN_{\kappa,\C}/\NN_\kappa[i]$ is a nonzero $\NN_\kappa$-module annihilated by $2$.
\end{proposition}

\begin{proof}
Every element of $\NN[i]$ is integral over $\GOz$, so $\NN[i]\subseteq\NN_\C$. If
$z=x+iy\in\NN_\C$ with $x,y\in\No$, then $\bar z\in\NN_\C$, so
\[
 2x=z+\bar z\in\NN_\C\cap\No=\NN,\qquad 2y=-i(z-\bar z)\in\NN_\C\cap\No=\NN
\]
by \eqref{osq:nm:eq:realintersection}. Hence $2\NN_\C\subseteq\NN[i]$, that is,
$\NN_\C\subseteq\frac12\NN[i]$. Inverting $2$ in $2\NN_\C\subseteq\NN[i]\subseteq\NN_\C$
gives $\NN_\C[1/2]=\NN[i][1/2]=\NN[1/2][i]$.

The root of unity $\zeta=(-1+i\sqrt3)/2$ satisfies $\zeta^2+\zeta+1=0$, so
$\zeta\in\NN_\C$; its real coordinate $-1/2$ is not in $\NN$ by
\cref{osq:nm:lem:reciprocal}, so $\zeta\notin\NN[i]$. Finally $1/2\in\frac12\NN[i]$ is not
in $\NN_\C$, again by the reciprocal lemma, since $2$ is a nonunit of $\GOz$. The proof is
unchanged in $F^{03}_\kappa[i]$, which contains $\zeta$; there $\zeta$ represents a nonzero
class of $\NN_{\kappa,\C}/\NN_\kappa[i]$, and this quotient of $\NN_\kappa$-modules is
annihilated by $2$ because $2\NN_{\kappa,\C}\subseteq\NN_\kappa[i]$.
\end{proof}

\begin{theorem}[Dense Gaussian normalization and its constant slice \src{03}]\label{osq:nm:thm:complexnormal}
The ring $\NN_\C$ is a proper dense subring of $\SC$ for the surreal-modulus topology, its
conductor $(\GOz:\NN_\C)$ is zero, and it cannot be generated over $\GOz$ by a set.
Moreover, with $\A_{\C,\C}=\C+\Pi_\C$,
\begin{equation}\label{osq:nm:eq:complexslice}
 \NN_\C\cap\A_{\C,\C}=\overline\Z+\Pi_\C=\A_{\overline\Z,\C},\qquad
 \NN_\C\cap\C=\overline\Z .
\end{equation}
For $\NN_{\kappa,\C}$ the analogous properness, density, slice (with
$\Pi^{03}_\kappa+i\Pi^{03}_\kappa$ in place of $\Pi_\C$) and zero-conductor statements
hold, and fewer than $\mu$ algebra generators over $A^{03}_\kappa[i]$ do not suffice.
\end{theorem}

\begin{proof}
Let $z=x+iy\in\SC$ and $\epsilon>0$. Choose a positive real integral unit
$u<\epsilon/2$ of $\NN$ (\cref{osq:nm:thm:smallunits}) and put
\[
 n=\lfloor x/u\rfloor,\qquad m=\lfloor y/u\rfloor,\qquad w=u(n+im).
\]
The floors lie in $\Oz$ and $u\in\NN\subseteq\NN_\C$, so $w\in\NN_\C$. Each coordinate of
$z-w$ lies in $[0,u)$, hence $|z-w|<\sqrt2\,u<\epsilon$. This proves density. The element
$1/2$ is excluded by \cref{osq:nm:lem:reciprocal}, which proves properness.

For $0\ne a\in\GOz$ choose such a $u$ with $0<u<1/|a|$. Then $au\in\NN_\C$ and
$0<|au|<1$, so $au\notin\GOz$ by \cref{osq:nm:lem:intpart}; thus the conductor is zero.
If $\NN_\C=\GOz[S]$ for a set $S$, \cref{osq:nm:thm:onedenominator} would give a nonzero
element of the conductor; for $\NN_{\kappa,\C}$ and fewer than $\mu$ generators,
\cref{osq:nm:thm:smallabsorption} does the same. The two absorption theorems also show
directly that every such small-generated algebra has a uniform modulus separation bound,
as in \cref{osq:nm:thm:discrete}.

Finally, let $z=c+P$ with $c\in\C$ and $P\in\Pi_\C$, integral over $\GOz$. All terms of its
monic equation lie in $\A_{\C,\C}$, and applying the ring homomorphism
$\ct:\A_{\C,\C}\to\C$ shows that $c$ is integral over $\Z[i]$; since $\Z[i]$ is integral
over $\Z$, $c\in\overline\Z$. Conversely such a $c$ is integral over the base, and $P$
already belongs to it. Restricting to constants gives the second equality. The same
arguments work in $F^{03}_\kappa[i]$, using the real integer part $A^{03}_\kappa$ of
$F^{03}_\kappa$ for the floors.
\end{proof}

The original integer-part rings have many finite quotients:
$\Oz/n\Oz\cong\Z/n\Z$ for ordinary $n\ge1$ by \eqref{osq:eq:nOz}, and likewise
$A^{03}_\kappa/nA^{03}_\kappa\cong\Z/n\Z$, because the purely infinite ideal is divisible
by every nonzero ordinary integer (compare \cref{osq:thm:quotients}). Their
normalizations behave oppositely.

\begin{theorem}[No nonzero finite ring targets \src{03}]\label{osq:nm:thm:nofinite}
None of $\NN$, $\NN_\kappa$, $\NN_\C$, $\NN_{\kappa,\C}$ admits a unital homomorphism to a
nonzero finite ring. In particular none of them has a proper ideal of finite index.
\end{theorem}

\begin{proof}
Suppose such a homomorphism exists. Its image is a nonzero finite commutative ring, even
if the target is not commutative, and a maximal ideal of the image gives a unital
homomorphism onto a finite field $\F_q$. The ordinary monic polynomial $T^q-T-1$ is
negative at $0$ and positive at $2$, so it has a real root $\alpha\in(0,2)$. This $\alpha$
is a real algebraic integer and belongs to all four normalizations by
\cref{osq:nm:thm:realslice,osq:nm:thm:complexnormal}. Its image $\bar\alpha$ would satisfy
$\bar\alpha^q-\bar\alpha-1=0$ in $\F_q$, whereas every element of $\F_q$ satisfies
$\bar\alpha^q-\bar\alpha=0$; so $1=0$ in $\F_q$, which is impossible. A proper ideal of
finite index would give a unital homomorphism onto a nonzero finite ring.
\end{proof}

The proof uses only the real algebraic integer constants. Accordingly this finite-quotient
obstruction is a general consequence of containing those constants, not a phenomenon
specific to surreal supports. It says nothing about infinite set-sized homomorphic images
of the normalizations; none are classified here (03).

\subsection{Set-valued valuations cannot contain the omnific ring}\label{osq:nm:sub:valuations}

A \emph{valuation} of a field $F$ (possibly a proper class, such as $\No$ or $\SC$) is a
map $v:F^\times\to\Gamma$ to an ordered abelian group with $v(xy)=v(x)+v(y)$ and
$v(x+y)\ge\min(v(x),v(y))$ whenever $x+y\ne0$; its valuation ring consists of $0$ and the
elements of nonnegative value. No compatibility with the ordering of $F$ is assumed. The
value group is \emph{set-sized} if $\Gamma$ is a set.

\begin{lemma}[No nonzero set-valued ordered quotient of $(\No,+)$ \src{03}]\label{osq:nm:lem:orderedmap}
Every order-preserving additive class map $w:(\No,+)\to\Gamma$ into a set-sized ordered
abelian group is zero.
\end{lemma}

\begin{proof}
Suppose $w(x)\ne0$ for some $x$. Replacing $x$ by $-x$ if necessary, $w(x)>0$; then $x\ne0$,
and $x<0$ would give $w(x)\le w(0)=0$, so $h:=x>0$ and $w(h)>0$. Choose a cardinal
$\lambda>|\Gamma|$ and consider the set of surreals $\alpha h$ (surreal products) for
ordinals $\alpha<\lambda$. If $\alpha<\beta$, then $\beta\ge\alpha+1$ in the surreal
order, the ordinal successor being the surreal sum with $1$. Hence
$\beta h-\alpha h\ge h$ and, by additivity and monotonicity,
\[
 w(\beta h)-w(\alpha h)=w(\beta h-\alpha h)\ge w(h)>0 .
\]
So the $\lambda$ elements $\alpha h$ have pairwise distinct images in $\Gamma$, contradicting
$\lambda>|\Gamma|$. Only a set of ordinals was used.
\end{proof}

\begin{theorem}[Full-class valuation obstruction \src{03}]\label{osq:nm:thm:valuation}
Let $v:\No^\times\to\Gamma$ be a valuation with set-sized value group. If $v(a)\ge0$ for
every $0\ne a\in\Oz$, then $v$ is trivial. Likewise, a valuation of $\SC$ with set-sized
value group whose valuation ring contains $\GOz$ is trivial.
\end{theorem}

\begin{proof}
Define $w(g)=v(\omega^g)$ for $g\in\No$. The monomial law $\omega^{g+g'}=\omega^g\omega^{g'}$
makes $w$ additive. If $g>0$ then $\omega^g\in\Pi\subseteq\Oz$, so $w(g)\ge0$; hence
$g\le g'$ implies $w(g')-w(g)=w(g'-g)\ge0$, and $w$ is order preserving. By
\cref{osq:nm:lem:orderedmap}, $w=0$.

For $x\ne0$, clear the set $\{x,x^{-1}\}$ by \cref{osq:prop:fractions}(i): there is a
monomial $d=\omega^h$ with $dx,dx^{-1}\in\Pi\subseteq\Oz$. Since $v(d)=w(h)=0$,
\[
 v(x)=v(dx)\ge0,\qquad -v(x)=v(x^{-1})=v(dx^{-1})\ge0,
\]
so $v(x)=0$. The complex proof uses the same real monomials and the complex clause of
\cref{osq:prop:fractions}(i), with $dx,dx^{-1}\in\Pi_\C\subseteq\GOz$.
\end{proof}

\begin{remark}[Why the size restriction is essential \src{03}]\label{osq:nm:rem:valsize}
The theorem does not say that $\No$ has no interesting valuations. Its normal-form
valuation $x\mapsto-\max\supp x$ takes values in the proper-class group $(\No,+)$ and is
negative on positive-exponent monomials, so it does not satisfy the containment hypothesis
anyway. Nor is the theorem a consequence of complete integral closure alone: in the
lexicographically ordered group $\Z\oplus\Z$, the values $v(d)=(1,0)$ and $v(x)=(0,-1)$
would give $v(dx^n)=(1,-n)>0$ for every ordinary $n$ although $v(x)<0$. Almost integrality
by itself cannot eliminate valuations of higher rank.

For ordinary set-sized domains the integral closure is characterized by valuation
overrings with set-sized value groups \cite[Lemmas 10.50.3 and 10.50.11]{StacksVal}. That
valuative criterion of integrality cannot be imported to a class ring while retaining only
set-sized value groups: $1/2$ is not integral over $\Oz$, yet
\cref{osq:nm:thm:valuation} rules out every nontrivial valuation of that restricted kind
whose ring contains $\Oz$. No assertion is made about the existence or classification of
\emph{proper-class-valued} valuation overrings separating nonintegral surreal elements.
\end{remark}

\subsection{An equation invisible to every set-sized image}\label{osq:nm:sub:invisible}

Manuscript~08 writes $B_m(t)$ for the binomial series; we write $\mathsf b_m(t)$.

\begin{lemma}[The binomial identity \src{08}]\label{osq:nm:lem:binomial}
For an ordinary integer $m\ge2$ the formal power series
\[
 \mathsf b_m(t)=\sum_{n\ge0}\binom{1/m}{n}t^n\in\Q[[t]]
\]
satisfies $\mathsf b_m(t)^m=1+t$. For every $\gamma\in\No_{>0}$ the substitution
$t=\omega^{-\gamma}$ is a valid Hahn substitution, and
$\mathsf b_m(\omega^{-\gamma})^m=1+\omega^{-\gamma}$ in $\No$. The same holds for
$t=X^{-\gamma}$ with $\gamma>0$ in any ordered abelian group, in particular in
$F^{03}_\kappa$ (this last sentence is 03's use of the binomial expansion in its Hahn fields, not 08's).
\end{lemma}

\begin{proof}
Write $b_n=\binom{1/m}{n}$; then $b_0=1$ and $(n+1)b_{n+1}=(1/m-n)b_n$. The coefficient of
$t^n$ in $(1+t)\mathsf b_m'$ is $(n+1)b_{n+1}+nb_n=b_n/m$, so
$(1+t)\mathsf b_m'=\frac1m\mathsf b_m$. Put $c=\mathsf b_m^m=\sum_nc_nt^n$. Then
$(1+t)c'=m\,\mathsf b_m^{m-1}(1+t)\mathsf b_m'=c$ and $c_0=1$. Comparing coefficients gives
$(n+1)c_{n+1}=(1-n)c_n$ for $n\ge0$, so $c_1=1$, $c_2=0$ and hence $c_n=0$ for all $n\ge2$.
Thus $c=1+t$.

After substitution the support lies in $\{-n\gamma:n\ge0\}$, which is reverse well
ordered. In any fixed finite product each exponent $-n\gamma$ receives finitely many
contributions, because an ordinary integer $n\ge0$ has only finitely many ordered
decompositions into a fixed number of nonnegative integers. Hence
$\sum_nc_nt^n\mapsto\sum_nc_n\omega^{-n\gamma}$ is a ring homomorphism
$\Q[[t]]\to\No$, and the formal identity remains a Hahn identity.
\end{proof}

\begin{theorem}[Universal small-target failure of root detection \src{08}]\label{osq:nm:thm:invisible}
Fix $\gamma\in\No_{>0}$ and an ordinary integer $m\ge2$. The equation
\begin{equation}\label{osq:nm:eq:invisible}
 T^m=\omega^\gamma+1
\end{equation}
has no root in $\GOz$, and therefore none in $\Oz$. Nevertheless, for every unital ring
homomorphism $\varphi:\Oz\to R$ with $R$ set-sized, the image equation
$T^m=\varphi(\omega^\gamma)+1_R$ has the solution $T=1_R$. The same holds for unital
homomorphisms with domain $\GOz$.
\end{theorem}

\begin{proof}
Put $r=\omega^{\gamma/m}\mathsf b_m(\omega^{-\gamma})\in\No$. By
\cref{osq:nm:lem:binomial}, $r^m=\omega^\gamma(1+\omega^{-\gamma})=\omega^\gamma+1$. For
every ordinary complex $m$th root of unity $\zeta$, the element $\zeta r\in\SC$ is a
root. These are $m$ distinct roots (as $r\ne0$), so they exhaust the roots of the
degree-$m$ polynomial $T^m-\omega^\gamma-1$ in the field $\SC$; no algebraic-closedness
argument is needed. The complex normal form of each root begins
\begin{equation}\label{osq:nm:eq:rootexpansion}
 \zeta r=\zeta\omega^{\gamma/m}+\frac{\zeta}{m}\,\omega^{\gamma/m-\gamma}
 +\sum_{n\ge2}\zeta\binom{1/m}{n}\omega^{\gamma/m-n\gamma}.
\end{equation}
The exponent $\gamma/m-\gamma$ is strictly negative because $m\ge2$, its coefficient
$\zeta/m$ is nonzero, and it differs from every other exponent $\gamma/m-n\gamma$
($n\ne1$) of the expansion. So this negative-exponent term cannot cancel, and no root lies
in $\GOz=\Z[i]\oplus\Pi_\C$, whose elements have nonnegative support.

On the other hand, by \cref{osq:thm:universal} every unital homomorphism $\varphi$ from
$\Oz$ to a set-sized ring kills $\Pi$, so $\varphi(\omega^\gamma)=0$ and the image equation
is $T^m=1_R$, solved by $1_R$. For a unital homomorphism with domain $\GOz$, the universal
theorem with $D=\Z[i]$ and $k=\C$ (see \cref{osq:cor:gaussian} for the description of these
maps) again gives $\varphi(\omega^\gamma)=0$.
\end{proof}

This is not merely a failure of a local--global test by congruences modulo ordinary primes.
Even allowing \emph{every set-sized ring} as a target does not reveal the missing negative
tail in \eqref{osq:nm:eq:rootexpansion}. The obstruction belongs to the full normal-form
arithmetic and survives in the proper-class ring itself (08).

\begin{corollary}[\src{08; 03}]\label{osq:nm:cor:notintclosed}
The ring $\Oz$ is not integrally closed in its fraction field $\No$.
\end{corollary}

\begin{proof}
The positive square root of $\omega+1$ lies in $\No=\Frac(\Oz)$
(\cref{osq:prop:fractions}(ii)). It is integral over $\Oz$, being a root of the monic
polynomial $T^2-(\omega+1)$, but it is not in $\Oz$ by \cref{osq:nm:thm:invisible} with
$m=2$ and $\gamma=1$.
\end{proof}

The same corollary is contained in 03's \cref{osq:nm:thm:normalization}, with the witness
$\sqrt2$, and manuscript~05 of the sibling report \cite{RepoOdg} proves it with the same
witness $\sqrt2$. This elementary corollary is included to clarify the distinction between
an integer part and an integral closure; it is not offered as a new factorization theorem
(08).

\begin{remark}[Explicit elements of $\NN\setminus\Oz$ \merge]\label{osq:nm:rem:rootinN}
Let $r$ be the positive root of $T^m-(\omega^\gamma+1)$ in $\No$, for $\gamma>0$ and
$m\ge2$; it is the element $r=\omega^{\gamma/m}\mathsf b_m(\omega^{-\gamma})$ of the proof
of \cref{osq:nm:thm:invisible}. Then
\[
 r\in\NN\setminus\Oz,\qquad \zeta r\in\NN_\C\setminus\GOz\quad(\zeta^m=1).
\]
Indeed $r$ is a root of a monic polynomial with coefficients in $\Oz$, so it lies in $\NN$
by definition, and $\zeta r$ lies in $\NN_\C$ for the same reason; neither lies in the
respective integer ring by \cref{osq:nm:thm:invisible}. Thus 08's equation exhibits
explicit elements of the normalization whose absence from $\Oz$ is invisible to every
set-sized image: under every unital homomorphism from $\Oz$ to a set-sized ring the
defining equation of $r$ becomes $T^m=1$, which already has the root $1$. Equivalently, if
$\psi$ is a unital homomorphism from $\NN$ (or from $\Oz[r]$) to a set-sized ring, then its
restriction to $\Oz$ kills $\omega^\gamma$ by \cref{osq:thm:universal}, so $\psi(r)^m=1$.
Unlike $\sqrt2$, the elements $r$ are not constants: by
\eqref{osq:nm:eq:rootexpansion} they have infinite leading term $\omega^{\gamma/m}$ and a
nonzero negative tail. Nothing is claimed here about the set-sized images of $\NN$
themselves.
\end{remark}

\section{Arithmetic of polynomial maps on the omnific lattices}\label{osq:sec:polynomial}

This section is manuscript~10's arithmetic of finite-degree polynomial and rational maps on
the real lattice $\Oz\subset\No$ and the Gaussian lattice $\GOz\subset\SC$. It uses from the
core only the universal theorem \cref{osq:thm:universal}, the quotient list
\cref{osq:thm:quotients} with the identity \eqref{osq:eq:nOz}, the constant-term lemma
\cref{osq:lem:ct} and the unit lists of \cref{osq:lem:units}.

\paragraph{Renamings.} Manuscript~10's symbols are renamed as follows. Its ideals
$\mathfrak P$ and $\mathfrak P_\C$ are $\Pi$ and $\Pi_\C$; its constant-term map $\pi$ is
$\ct$; its projection $\Pi$ on polynomial rings (the \emph{shadow}) is $\ct_*$, the
coefficientwise constant term; its Gaussian lattice is $\GOz$ inside $\SC$; its ring
$R(D,K)$ is $\A_{D,K}$; its binomial polynomials $B_k(X)$ are written $\binom Xk$; its
$L_k=\lcm(1,\dots,k)$ is written $\lcm_k$; and its rings $S_r$ and $I_r$ are written
$\Num_r(\Oz)$ and $\Int(\Z^r)$ (the letter $I_r$ would collide with ideals). The numbered
results keep 10's content.

\paragraph{Notation.} Throughout, $r\ge1$ is an ordinary integer,
$\mathbf X=(X_1,\dots,X_r)$, and bold letters denote $r$-tuples. For an ordinary $k\ge0$,
\[
 \binom Xk=\frac{X(X-1)\cdots(X-k+1)}{k!},\qquad \binom X0=1,
\]
and for a multi-index $\mathbf k=(k_1,\dots,k_r)$ we put
$\binom{\mathbf X}{\mathbf k}=\prod_{j=1}^r\binom{X_j}{k_j}$. We write $\Delta_jF(\mathbf
X)=F(\mathbf X+\mathbf e_j)-F(\mathbf X)$, and $\Delta=\Delta_1$ when $r=1$. Put
$\lcm_0=1$ and $\lcm_k=\lcm(1,\dots,k)$ for $k\ge1$. Recall
$\A_{\Q,\R}=\Q+\Pi$ and $\A_{\Q(i),\C}=\Q(i)+\Pi_\C$; by \cref{osq:lem:ct} the constant
term $\ct$ is a ring homomorphism on both, with kernels $\Pi$ and $\Pi_\C$, so
$\Pi$ is an ideal of $\Q+\Pi$ and $\Pi_\C$ is an ideal of $\Q(i)+\Pi_\C$. Its
coefficientwise extension
\[
 \ct_*:(\Q+\Pi)[\mathbf X]\to\Q[\mathbf X],\qquad
 \ct_*:(\Q(i)+\Pi_\C)[\mathbf X]\to\Q(i)[\mathbf X]
\]
is therefore a ring homomorphism with kernel $\Pi[\mathbf X]$, respectively
$\Pi_\C[\mathbf X]$. Define
\begin{align*}
 \Num_r(\Oz)&=\{F\in\No[\mathbf X]:F(\Oz^r)\subseteq\Oz\},&
 \Int(\Z^r)&=\{f\in\Q[\mathbf X]:f(\Z^r)\subseteq\Z\},\\
 \Num_r(\GOz)&=\{F\in\SC[\mathbf X]:F(\GOz^r)\subseteq\GOz\},&
 \Int(\Z[i]^r)&=\{f\in\Q(i)[\mathbf X]:f(\Z[i]^r)\subseteq\Z[i]\}.
\end{align*}
These are rings under pointwise operations. All degrees, finite sums and factorials are
ordinary.

\subsection{Lattice separation and an omitted congruence type}\label{osq:pm:sub:lattice}

\begin{lemma}[Lattice separation and cofinality \src{10}]\label{osq:pm:lem:lattice}
For $\Lambda=\Oz$ and for $\Lambda=\GOz$, every nonzero $x\in\Lambda$ satisfies
$|x|\ge1$. The finite elements of $\Oz$ are exactly the ordinary integers. The positive
elements of $\Oz$ are cofinal in $\No$. The units are $\Oz^\times=\{\pm1\}$ and
$\GOz^\times=\{\pm1,\pm i\}$ \textup{(}\cref{osq:lem:units}\textup{)}.
\end{lemma}
\begin{proof}
A nonzero real purely infinite series has a positive leading exponent, so its absolute
value exceeds every positive real. Hence $n+u$ with $n\in\Z$ and $0\ne u\in\Pi$ is
infinite, the finite elements of $\Oz$ are the ordinary integers, and every nonzero
element of $\Oz$ has absolute value at least $1$. In the Gaussian case write
$x=a+ib$ with $a,b\in\Oz$, not both zero; the surcomplex modulus dominates either
coordinate, $|x|=\sqrt{a^2+b^2}\ge\max(|a|,|b|)\ge1$.

Given $y\in\No$, choose an exponent $g>0$ larger than the leading exponent of $y$ (any
$g>0$ if $y=0$). Then $\omega^g\in\Oz$ is positive and exceeds $n|y|$ for every ordinary
$n$, in particular $\omega^g>|y|$. This proves cofinality. The unit lists are
\cref{osq:lem:units}.
\end{proof}

\begin{remark}[Modulus and standard part \src{10}]\label{osq:pm:rem:modulus}
The surcomplex modulus needs only real closedness of $\No$: $|a+ib|=\sqrt{a^2+b^2}$ for
$a,b\in\No$, and it is multiplicative and satisfies the triangle inequality. The
standard part of a finite surcomplex number is taken coordinatewise in $\C$; on the
finite rings $\OO_\R$ and $\OO_\C$ it is the map $\st$ of \cref{osq:lem:ct}.
\end{remark}

\begin{proposition}[A finitely satisfiable countable congruence type \src{10}]
\label{osq:pm:prop:type}
Fix an ordinary prime $p$ and $\alpha\in\Z_p\setminus\Z$. For $e\ge1$ let $a_e$ be the
representative of $\alpha\bmod p^e$ in $\{0,\dots,p^e-1\}$. The countable family of ring
formulas
\begin{equation}\label{osq:pm:eq:type}
 \exists y\ (x-a_e=p^e y),\qquad e=1,2,\dots,
\end{equation}
is finitely satisfiable in $\Z\subset\Oz$ but has no realization in $\Oz$. An explicit
choice is
\begin{equation}\label{osq:pm:eq:lacunary}
 \alpha=\sum_{j\ge0}p^{2^j}\in\Z_p .
\end{equation}
\end{proposition}
\begin{proof}
For a finite subfamily let $E$ be its largest index and take $x=a_E$. Since the residues
are compatible, $a_E\equiv a_e\pmod{p^e}$ for $e\le E$, which gives witnesses $y\in\Z$ for
every formula of the subfamily. If $x\in\Oz$ realized all formulas, applying the ring
homomorphism $\ct$ would give an ordinary integer $\ct(x)$ with
$\ct(x)\equiv a_e\equiv\alpha\pmod{p^e}$ for every $e$ (by \eqref{osq:eq:nOz},
$x-a_e\in p^e\Oz$ means $\ct(x)-a_e\in p^e\Z$). In $\Z_p$ this forces
$\ct(x)=\alpha$, contradicting $\alpha\notin\Z$.

For \eqref{osq:pm:eq:lacunary}: the powers $p^{2^j}$ are distinct, so the base-$p$ digits
of $\alpha$ are $1$ at the positions $2^j$ and $0$ elsewhere. They are not eventually
zero, so $\alpha$ is not a nonnegative integer; and they are not eventually $p-1$ (for
$p=2$ the positions $2^j+1<2^{j+1}$ with $j\ge1$ carry the digit $0$), so $\alpha$ is not a
negative integer.
\end{proof}

\begin{remark}[Ring types, not saturation \src{10}]\label{osq:pm:rem:saturation}
This elementary omitted type concerns the ring structure. It shows that the
order-theoretic richness of the full surreal class does not give a corresponding
realization property for countable ring types. No use of the term ``saturated model'' for
a proper-class structure is needed, and none is made.
\end{remark}

\subsection{Polynomials preserving the omnific lattices}\label{osq:pm:sub:numerical}

\begin{lemma}[Omnific binomial closure \src{10}]\label{osq:pm:lem:binomial}
For every $x\in\Oz$ and every ordinary $k\ge0$, $\binom xk\in\Oz$. More generally, if
$f\in\Q[\mathbf X]$ and $f(\Z^r)\subseteq\Z$, then $f(\Oz^r)\subseteq\Oz$.
\end{lemma}
\begin{proof}
Write $x_j=n_j+u_j$ with $n_j\in\Z$ and $u_j\in\Pi$. Expanding each monomial of $f$ shows
that $f(\mathbf n+\mathbf u)-f(\mathbf n)$ is a finite sum of terms, each a rational
multiple of a product of ordinary integers and of powers of the $u_j$ containing at least
one factor $u_j$. Since $\Pi$ is an ideal of $\Q+\Pi$, every such term lies in $\Pi$; thus
\[
 f(\mathbf n+\mathbf u)-f(\mathbf n)\in\Pi .
\]
As $f(\mathbf n)\in\Z$, we get $f(\mathbf x)\in\Z+\Pi=\Oz$. The binomial case is the
case $f=\binom Xk$, which takes integer values at all ordinary integers, including
negative ones: $\binom{-n}{k}=(-1)^k\binom{n+k-1}{k}$.
\end{proof}

Thus $\Oz$ is a binomial ring in the usual torsion-free sense: it is $\Z$-torsion free and
closed under $x\mapsto\binom xk$ \cite{ElliottBinomial}. The terminology is classical; the
support argument above explains its application here without importing an abstract
free-ring theorem.

\begin{theorem}[Exact finite-grid test \src{10}]\label{osq:pm:thm:grid}
Let $F\in\No[\mathbf X]$ with $\deg_{X_j}F\le d_j$ for $j=1,\dots,r$. The following are
equivalent:
\begin{enumerate}
 \item $F(\Oz^r)\subseteq\Oz$;
 \item $F(n_1,\dots,n_r)\in\Oz$ for all ordinary integers $0\le n_j\le d_j$;
 \item in the unique expansion
 \begin{equation}\label{osq:pm:eq:newton}
  F(\mathbf X)=\sum_{\mathbf k\le\mathbf d}c_{\mathbf k}\binom{\mathbf X}{\mathbf k}
 \end{equation}
 all coefficients $c_{\mathbf k}$ belong to $\Oz$.
\end{enumerate}
Explicitly,
\begin{equation}\label{osq:pm:eq:newtoncoef}
 c_{\mathbf k}=(\Delta_1^{k_1}\cdots\Delta_r^{k_r}F)(\mathbf 0)
 =\sum_{\boldsymbol\ell\le\mathbf k}
   (-1)^{\sum_j(k_j-\ell_j)}
   \Bigl(\prod_{j=1}^r\binom{k_j}{\ell_j}\Bigr)F(\boldsymbol\ell).
\end{equation}
\end{theorem}
\begin{proof}
The polynomials $\binom X0,\dots,\binom Xd$ form a basis of the univariate polynomials of
degree at most $d$ over the field $\No$: their degrees are $0,\dots,d$ and their leading
coefficients are nonzero. Taking products, the $\binom{\mathbf X}{\mathbf k}$ with
$\mathbf k\le\mathbf d$ form a basis of the polynomials with $\deg_{X_j}\le d_j$, which
gives existence and uniqueness of \eqref{osq:pm:eq:newton}. Since
$\Delta\binom Xm=\binom X{m-1}$ for $m\ge1$ and $\Delta\binom X0=0$, one has
$\Delta_j^{k_j}\binom{X_j}{m_j}=\binom{X_j}{m_j-k_j}$ for $m_j\ge k_j$ and $0$ otherwise;
as $\binom 0m=0$ for $m>0$, applying $\Delta_1^{k_1}\cdots\Delta_r^{k_r}$ to
\eqref{osq:pm:eq:newton} and evaluating at $\mathbf 0$ leaves exactly $c_{\mathbf k}$.
Writing $\Delta_j=\tau_j-1$ with the shift $\tau_j$ and expanding by the binomial theorem
gives the last expression in \eqref{osq:pm:eq:newtoncoef}.

(i)$\Rightarrow$(ii) is restriction to ordinary grid points. Under (ii), the last
expression in \eqref{osq:pm:eq:newtoncoef} is an integral linear combination of grid values
in $\Oz$, which proves (iii). Under (iii), each factor $\binom{x_j}{k_j}$ with $x_j\in\Oz$
lies in $\Oz$ by \cref{osq:pm:lem:binomial}, and evaluating \eqref{osq:pm:eq:newton}
proves (i).
\end{proof}

The test is finite for a given finite-degree polynomial, although the preservation
property quantifies over the entire class $\Oz^r$. It is a mathematical test with exact
surreal input; no claim is made that arbitrary surreal coefficients have an effective
finite encoding.

\begin{theorem}[The real numerical-polynomial ring \src{10}]\label{osq:pm:thm:numreal}
Every $F\in\Num_r(\Oz)$ has a unique decomposition
\begin{equation}\label{osq:pm:eq:decomp}
 F=U+f,\qquad U\in\Pi[\mathbf X],\quad f\in\Int(\Z^r),
\end{equation}
so that $\Num_r(\Oz)=\Pi[\mathbf X]\oplus\Int(\Z^r)$ additively. The coefficientwise
constant term $\ct_*(F)=f$ is a surjective ring homomorphism
$\Num_r(\Oz)\to\Int(\Z^r)$ with kernel $\Pi[\mathbf X]$, split by the inclusion
$\Int(\Z^r)\subseteq\Num_r(\Oz)$. For $\mathbf x\in\Oz^r$,
\begin{equation}\label{osq:pm:eq:evshadow}
 \ct\bigl(F(\mathbf x)\bigr)=\ct_*(F)\bigl(\ct(x_1),\dots,\ct(x_r)\bigr).
\end{equation}
\end{theorem}
\begin{proof}
Let $F\in\Num_r(\Oz)$. By \cref{osq:pm:thm:grid} its Newton coefficients $c_{\mathbf k}$
lie in $\Oz$; write $c_{\mathbf k}=m_{\mathbf k}+u_{\mathbf k}$ with $m_{\mathbf k}\in\Z$
and $u_{\mathbf k}\in\Pi$, and put
\[
 U=\sum_{\mathbf k}u_{\mathbf k}\binom{\mathbf X}{\mathbf k},\qquad
 f=\sum_{\mathbf k}m_{\mathbf k}\binom{\mathbf X}{\mathbf k}.
\]
The binomials have rational coefficients and $\Pi$ is a $\Q$-vector space, so
$U\in\Pi[\mathbf X]$; products of binomials are integer valued on $\Z^r$, so
$f\in\Int(\Z^r)$. Conversely, a polynomial in $\Pi[\mathbf X]$ takes values in $\Pi$ on
$\Oz^r$, because $\Pi$ is an ideal of $\Q+\Pi\supseteq\Oz$, and $\Int(\Z^r)\subseteq
\Num_r(\Oz)$ by \cref{osq:pm:lem:binomial}; hence $\Pi[\mathbf X]+\Int(\Z^r)
\subseteq\Num_r(\Oz)$. The sum is direct because a coefficient that is both rational and
purely infinite is zero.

By \eqref{osq:pm:eq:decomp} all coefficients of $F$ lie in $\Q+\Pi$, so $\Num_r(\Oz)$ is a
subring of $(\Q+\Pi)[\mathbf X]$ and $\ct_*$ restricts to a ring homomorphism on it. It
sends $U$ to $0$ and fixes $f$, so its restriction is the projection onto the second
summand; its kernel is $\Pi[\mathbf X]$ and the inclusion of $\Int(\Z^r)$ is a ring
section. Finally write $x_j=n_j+u_j$ with $n_j=\ct(x_j)$. Then $U(\mathbf x)\in\Pi$ and
$f(\mathbf x)-f(\mathbf n)\in\Pi$ by the proof of \cref{osq:pm:lem:binomial}, while
$f(\mathbf n)\in\Z$. Hence $\ct(F(\mathbf x))=f(\mathbf n)$, which is
\eqref{osq:pm:eq:evshadow}.
\end{proof}

The decomposition is an additive direct sum whose first summand is an ideal; it is
\emph{not} a direct product of rings.

\begin{example}[\src{10}]\label{osq:pm:ex:numreal}
The polynomial
\[
 F(X)=3\binom X2+\omega^{\sqrt2}X^3+\Bigl(\sum_{m\ge1}\omega^{1/m}\Bigr)(X-1)
\]
is omnific-integer valued and has shadow $\ct_*(F)=3\binom X2$. By contrast
$\omega^{-1}X$ fails already at $X=1$. No negative-exponent coefficient can be hidden in a
finite-degree lattice-preserving polynomial: by \cref{osq:pm:thm:numreal} every
coefficient of an element of $\Num_r(\Oz)$ lies in $\Q+\Pi$.
\end{example}

\begin{theorem}[Gaussian numerical polynomials \src{10}]\label{osq:pm:thm:numgaussian}
For every ordinary $r\ge1$,
\begin{equation}\label{osq:pm:eq:gaussiandecomp}
 \Num_r(\GOz)=\Pi_\C[\mathbf X]\oplus\Int(\Z[i]^r)
\end{equation}
additively. The coefficientwise constant term $\ct_*$ is a surjective ring homomorphism
$\Num_r(\GOz)\to\Int(\Z[i]^r)$ with kernel $\Pi_\C[\mathbf X]$, split by inclusion, and
for $\mathbf x\in\GOz^r$, with $\ct:\GOz\to\Z[i]$ the Gaussian constant term,
\begin{equation}\label{osq:pm:eq:evshadowC}
 \ct\bigl(F(\mathbf x)\bigr)=\ct_*(F)\bigl(\ct(x_1),\dots,\ct(x_r)\bigr).
\end{equation}
\end{theorem}
\begin{proof}
Let $F\in\Num_r(\GOz)$ with $\deg_{X_j}F\le d_j$. The Newton expansion
\eqref{osq:pm:eq:newton} and formula \eqref{osq:pm:eq:newtoncoef} hold over the field
$\SC$ by the same basis argument. The ordinary grid points lie in $\GOz^r$, so
\eqref{osq:pm:eq:newtoncoef} shows $c_{\mathbf k}\in\GOz=\Z[i]+\Pi_\C$. Expanding the
binomials, each monomial coefficient of $F$ is a $\Q$-linear combination of the
$c_{\mathbf k}$, hence lies in $\Q(i)+\Pi_\C$. Put $f=\ct_*(F)\in\Q(i)[\mathbf X]$ and
$U=F-f\in\Pi_\C[\mathbf X]$. For an ordinary Gaussian tuple $\mathbf z\in\Z[i]^r$ we have
$U(\mathbf z)\in\Pi_\C$, because $\Pi_\C$ is an ideal of $\Q(i)+\Pi_\C$, and
$F(\mathbf z)=w+v$ with $w\in\Z[i]$, $v\in\Pi_\C$. Then $f(\mathbf z)-w=v-U(\mathbf z)$
lies in $\Q(i)\cap\Pi_\C=0$, so $f(\mathbf z)=w\in\Z[i]$. Thus $f\in\Int(\Z[i]^r)$.

Conversely write a Gaussian omnific input as $\mathbf z+\mathbf u$ with
$\mathbf z\in\Z[i]^r$ and $u_j\in\Pi_\C$. For $f\in\Int(\Z[i]^r)$ the expansion argument
of \cref{osq:pm:lem:binomial} gives $f(\mathbf z+\mathbf u)-f(\mathbf z)\in\Pi_\C$, so
$f(\mathbf z+\mathbf u)\in\Z[i]+\Pi_\C=\GOz$; and a polynomial in $\Pi_\C[\mathbf X]$
takes values in $\Pi_\C$ on $\GOz^r$. This proves \eqref{osq:pm:eq:gaussiandecomp};
the sum is direct since $\Q(i)\cap\Pi_\C=0$. As in \cref{osq:pm:thm:numreal},
$\Num_r(\GOz)$ is a subring of $(\Q(i)+\Pi_\C)[\mathbf X]$, on which $\ct_*$ is a ring
homomorphism; this gives multiplicativity, the kernel and the section. The same
computation, $\ct(F(\mathbf z+\mathbf u))=f(\mathbf z)$, proves
\eqref{osq:pm:eq:evshadowC}.
\end{proof}

\begin{remark}[Why the real Newton test cannot simply be copied \src{10}]
\label{osq:pm:rem:gaussnewton}
The ordinary binomial polynomials need not be Gaussian-integer valued:
\[
 \binom i2=\frac{i(i-1)}2=\frac{-1-i}{2}\notin\Z[i].
\]
Consequently values on an ordinary integer grid do not by themselves certify
preservation of $\GOz$ (the grid values are necessary, not sufficient). The Gaussian
theorem retains the ordinary ring $\Int(\Z[i]^r)$ rather than imposing a false binomial
basis on it; the arithmetic of $\Int(\Z[i]^r)$ is not replaced by the simpler real
Newton basis.
\end{remark}

\begin{theorem}[Small-target reflection for numerical polynomials \src{10}]
\label{osq:pm:thm:numreflection}
Every class ring homomorphism $\varphi$ from $\Num_r(\Oz)$ to a set-sized ring $A$
factors uniquely through
\[
 \ct_*:\Num_r(\Oz)\longrightarrow\Int(\Z^r).
\]
The analogous assertion holds for $\Num_r(\GOz)$ with universal image $\Int(\Z[i]^r)$.
Therefore the set-sized quotient presentations of either numerical-polynomial ring are
precisely the quotient presentations of its ordinary shadow pulled back along $\ct_*$:
their kernels are the class ideals $\ct_*^{-1}(\mathfrak a)$ for ideals $\mathfrak a$ of
$\Int(\Z^r)$, respectively $\Int(\Z[i]^r)$, with quotients $\Int(\Z^r)/\mathfrak a$,
respectively $\Int(\Z[i]^r)/\mathfrak a$.
\end{theorem}
\begin{proof}
The restriction of $\varphi$ to the ideal $\Pi$ of constant polynomials is a (possibly
nonunital) class ring homomorphism $\Pi\to A$, hence zero by \cref{osq:thm:universal}.
Every element of $\Pi[\mathbf X]$ is a finite sum of products $u\,\mathbf X^{\mathbf m}$
with $u\in\Pi$ and $\mathbf X^{\mathbf m}\in\Num_r(\Oz)$, so
$\varphi(u\mathbf X^{\mathbf m})=\varphi(u)\varphi(\mathbf X^{\mathbf m})=0$ and $\varphi$
kills $\Pi[\mathbf X]=\ker\ct_*$ (\cref{osq:pm:thm:numreal}). Hence
$\varphi=\psi\circ\ct_*$ with $\psi$ the restriction of $\varphi$ to the section
$\Int(\Z^r)$; $\psi$ is unique because $\ct_*$ is surjective. The Gaussian case is
identical, with $\Pi_\C$ and \cref{osq:pm:thm:numgaussian}.

If $\varphi$ is surjective onto a set-sized ring with kernel $J$, then $\psi$ is surjective
and $J=\ct_*^{-1}(\ker\psi)$. Conversely, for an ideal $\mathfrak a$ of the shadow ring,
the composite of $\ct_*$ with the quotient map onto the set $\Int(\Z^r)/\mathfrak a$ is a
set-sized quotient presentation with kernel $\ct_*^{-1}(\mathfrak a)$. This is the
kernel argument of \cref{osq:thm:quotients}.
\end{proof}

\subsection{Congruences modulo every omnific integer}\label{osq:pm:sub:congruence}

\begin{definition}[\src{10}]\label{osq:pm:def:cp}
A function $F:\Oz\to\Oz$ is \emph{congruence preserving} if $a-b$ divides $F(a)-F(b)$ in
$\Oz$ for all $a,b\in\Oz$; equivalently, for every $m\in\Oz$, $a\equiv b\pmod{m\Oz}$
implies $F(a)\equiv F(b)\pmod{m\Oz}$. We write $\CP(\Oz)$ for the congruence-preserving
elements of $\Num_1(\Oz)$ and $\CP(\Z)$ for the elements $f\in\Int(\Z)$ such that $a-b$
divides $f(a)-f(b)$ in $\Z$ for all $a,b\in\Z$.
\end{definition}

The two formulations agree: the forward implication is transitivity of divisibility, and
the reverse one takes $m=a-b$. When $a=b$ the condition reads $0\mid0$, which holds.
Differences and finite sums of congruence-preserving functions are congruence preserving.

\begin{lemma}[Integral divided differences of weighted binomials \src{10}]
\label{osq:pm:lem:divdiff}
For every ordinary $k\ge1$, the identity
\begin{equation}\label{osq:pm:eq:divdiff}
 \frac{\binom{X+T}{k}-\binom Xk}{T}
   =\sum_{j=1}^k\binom X{k-j}\,\frac1j\binom{T-1}{j-1}
\end{equation}
holds in $\Q[X,T]$. In particular $\lcm_k\binom Xk$ is congruence preserving on $\Oz$.
\end{lemma}
\begin{proof}
Vandermonde's identity gives $\binom{X+T}k=\sum_{j=0}^k\binom X{k-j}\binom Tj$. For
$j\ge1$ the falling-factorial formula gives $\binom Tj=\frac Tj\binom{T-1}{j-1}$.
Subtracting the term $j=0$ and cancelling the polynomial factor $T$ gives
\eqref{osq:pm:eq:divdiff}.

After multiplication by $\lcm_k$, every scalar $\lcm_k/j$ with $1\le j\le k$ is an
ordinary integer, so by \cref{osq:pm:lem:binomial} the right side evaluates into $\Oz$ for
all $X,T\in\Oz$. Taking $X=b$ and $T=a-b$ gives
$\lcm_k\binom ak-\lcm_k\binom bk=(a-b)w$ with $w\in\Oz$, including for nonzero infinite
$a-b$; the case $a=b$ is trivial.
\end{proof}

\begin{theorem}[All-modulus Newton criterion \src{10}]\label{osq:pm:thm:cp}
Let $F\in\Num_1(\Oz)$ have Newton expansion $F=\sum_{k=0}^dc_k\binom Xk$, $c_k\in\Oz$.
The following are equivalent:
\begin{enumerate}
 \item $F$ is congruence preserving on $\Oz$;
 \item $c_k\in\lcm_k\Oz$ for every $0\le k\le d$;
 \item $\ct(c_k)\in\lcm_k\Z$ for every $0\le k\le d$;
 \item the ordinary numerical polynomial $\ct_*(F)$ is congruence preserving on $\Z$.
\end{enumerate}
Consequently, within the respective numerical-polynomial rings,
\begin{equation}\label{osq:pm:eq:CPdecomp}
 \CP(\Oz)=\Pi[X]\oplus\CP(\Z).
\end{equation}
\end{theorem}
\begin{proof}
(ii)$\Rightarrow$(i). Write $c_k=\lcm_kd_k$ with $d_k\in\Oz$. Each summand
$d_k\lcm_k\binom Xk$ is congruence preserving by \cref{osq:pm:lem:divdiff} (multiply the
divisibility by $d_k$), and a finite sum of congruence-preserving functions is congruence
preserving.

(i)$\Rightarrow$(ii). We show $c_k\in\lcm_k\Oz$ by induction on $k$. For $k=0$ it is
automatic since $\lcm_0=1$. Suppose it holds for all $j<k$. Then
$\sum_{j<k}c_j\binom Xj$ is congruence preserving by the previous paragraph, so
\[
 G(X)=F(X)-\sum_{j<k}c_j\binom Xj=\sum_{j\ge k}c_j\binom Xj
\]
is congruence preserving. Since $\binom nj=0$ for ordinary $0\le n<j$ and
$\binom kk=1$, we get $G(0)=\dots=G(k-1)=0$ and $G(k)=c_k$. For each $1\le m\le k$,
comparing the inputs $k$ and $k-m$ gives $m\mid G(k)-G(k-m)=c_k$ in $\Oz$. By
\eqref{osq:eq:nOz}, the ordinary integer $\ct(c_k)$ is divisible by every $m\le k$, hence
by $\lcm_k$, and again by \eqref{osq:eq:nOz}, $c_k\in\lcm_k\Oz$.

(ii)$\Leftrightarrow$(iii) is \eqref{osq:eq:nOz}. For (iii)$\Leftrightarrow$(iv), note
that $\ct_*(F)=\sum_k\ct(c_k)\binom Xk$ is the Newton expansion of $\ct_*(F)$, since the
binomials have rational coefficients. The same two arguments over $\Z$ (the identity
\eqref{osq:pm:eq:divdiff} with $X,T\in\Z$, and the same induction with divisibility in
$\Z$) show that an $f\in\Int(\Z)$ with Newton coefficients $m_k\in\Z$ is congruence
preserving on $\Z$ if and only if $m_k\in\lcm_k\Z$ for all $k$. Applied to $m_k=\ct(c_k)$,
this is (iii)$\Leftrightarrow$(iv).

For \eqref{osq:pm:eq:CPdecomp}, decompose $F=U+f$ by \cref{osq:pm:thm:numreal}. The Newton
coefficients of $U\in\Pi[X]$ are integral combinations of values $U(\ell)\in\Pi$, and
$\Pi\subseteq\lcm_k\Oz$ by \eqref{osq:eq:nOz}; so $U\in\CP(\Oz)$ by (ii). An
$f\in\Int(\Z)$ lies in $\CP(\Oz)$ exactly when it lies in $\CP(\Z)$, by
(i)$\Leftrightarrow$(iv) applied to $f=\ct_*(f)$. Since $\ct_*(F)=f$, the equivalence
(i)$\Leftrightarrow$(iv) for $F$ now gives $F\in\CP(\Oz)$ if and only if $f\in\CP(\Z)$.
\end{proof}

The divisibility condition in (i) includes infinite differences and therefore every
infinite omnific modulus; no extra coefficient condition appears at those moduli. The
ordinary least-common-multiple pattern is classical, and related finite-ring formulations
are proved in \cite{CGG}; Newton bases and integer-valued polynomials are treated in
\cite{CC}. The content specific to the present domain is that the pattern suffices for
differences at every surreal scale, not only for congruences modulo ordinary integers.

\begin{example}[Two kinds of denominators \src{10}]\label{osq:pm:ex:cp}
The polynomial $12\binom X4$ is congruence preserving ($\lcm_4=12$), although its leading
coefficient $1/2$ is not an ordinary integer. On the other hand
\[
 F(X)=\binom X2+\omega^{\sqrt2}X^3
\]
is integer valued on all of $\Oz$ but not congruence preserving:
$F(2)-F(0)=1+8\omega^{\sqrt2}$ has odd constant term and so is not divisible by $2$ in
$\Oz$, by \eqref{osq:eq:nOz}. The purely infinite term does not repair the odd constant
term.
\end{example}

\begin{remark}[What \cref{osq:pm:thm:cp} does not say \src{10}]\label{osq:pm:rem:cpnot}
The theorem does not say that reduction modulo an infinite omnific integer is an ordinary
finite congruence; such a quotient need not have any set-sized presentation (by
\cref{osq:thm:quotients} it has one only when it contains $\Pi$). It says that the
polynomial \emph{operation} respects those congruences, with an explicit quotient
supplied by \eqref{osq:pm:eq:divdiff}. It also does not classify all congruence-preserving
class functions $\Oz\to\Oz$: the finite Newton expansion is essential to the statement and
the proof, and no infinite interpolation formula or full-class summation is implicit.
\end{remark}

\subsection{\texorpdfstring{$p$}{p}-adic characters and completions}\label{osq:pm:sub:padic}

In this subsection $p$ is an ordinary prime, $e\ge1$ and $r\ge1$ are ordinary integers,
and $C(\Z_p^r,M)$ denotes the continuous functions from $\Z_p^r$ with its $p$-adic
topology to $M$ ($M=\Z/p^e\Z$ or $\F_p$ discrete, or $M=\Z_p$). All $p$-adic fields,
rings and function spaces here are ordinary set-sized structures.

\begin{lemma}[Continuous extension \src{10}]\label{osq:pm:lem:padicext}
Every $f\in\Int(\Z^r)$ extends, by its rational-coefficient formula, to a continuous
function $\Z_p^r\to\Z_p$, and this continuous extension is unique.
\end{lemma}
\begin{proof}
A rational polynomial defines a continuous function $\Q_p^r\to\Q_p$. Ordinary integer
tuples are dense in $\Z_p^r$ and take values in $\Z\subseteq\Z_p$; since $\Z_p$ is closed
in $\Q_p$, all values on $\Z_p^r$ lie in $\Z_p$. Density also gives uniqueness.
\end{proof}

Reduction of values modulo $p^e$ is therefore a ring homomorphism
$\Int(\Z^r)\to C(\Z_p^r,\Z/p^e\Z)$. Its surjectivity needs an approximation argument; we
prove it explicitly by indicator polynomials rather than assume it, so no unproved
approximation assertion is used. For $0\le a<p^k$ write $a=a_0+a_1p+\dots+a_{k-1}p^{k-1}$
with digits $0\le a_j<p$, and put
\begin{equation}\label{osq:pm:eq:indicator}
 E_{a,k}(X)=\prod_{j=0}^{k-1}\Bigl[1-\Bigl(\binom X{p^j}-a_j\Bigr)^{p-1}\Bigr]
 \in\Int(\Z),
 \qquad
 I_{a,k,e}(X)=E_{a,k}(X)^{p^{e-1}}.
\end{equation}

\begin{lemma}[Digit and indicator congruences \src{10}]\label{osq:pm:lem:indicator}
For every $x\in\Z_p$, with $\mathbf 1_{a,k}(x)=1$ if $x\equiv a\pmod{p^k}$ and
$\mathbf 1_{a,k}(x)=0$ otherwise,
\[
 E_{a,k}(x)\equiv\mathbf 1_{a,k}(x)\pmod p,
 \qquad
 I_{a,k,e}(x)\equiv\mathbf 1_{a,k}(x)\pmod{p^e}.
\]
\end{lemma}
\begin{proof}
First let $x$ be a nonnegative ordinary integer with base-$p$ digits $x_j$. In
$\F_p[T]$ the Frobenius identity gives
\[
 (1+T)^x=\prod_{j\ge0}(1+T)^{x_jp^j}=\prod_{j\ge0}\bigl(1+T^{p^j}\bigr)^{x_j}.
\]
The coefficient of $T^{p^j}$ on the right is $x_j$: the factors of index below $j$ have
combined degree at most $\sum_{i<j}(p-1)p^i=p^j-1$, and factors of index above $j$
contribute only exponents $0$ or at least $p^{j+1}$. Hence $\binom x{p^j}\equiv
x_j\pmod p$. Both sides are continuous functions $\Z_p\to\F_p$ of $x$
(\cref{osq:pm:lem:padicext}; the $j$th digit depends only on $x\bmod p^{j+1}$), and the
nonnegative integers are dense in $\Z_p$, so the congruence holds for all $x\in\Z_p$ with
$x_j$ the $j$th $p$-adic digit. By Fermat, $t^{p-1}=1$ for $0\ne t\in\F_p$, so the $j$th
factor of $E_{a,k}(x)$ is $\equiv1$ if $x_j=a_j$ and $\equiv0$ otherwise. The product is
the indicator of $x_j=a_j$ for all $j<k$, that is, of $x\equiv a\pmod{p^k}$.

For the lift, let $z\in\Z_p$. If $z\equiv0\pmod p$ then $z^{p^{e-1}}\equiv0\pmod{p^e}$,
because $p^{e-1}\ge e$. If $z=1+pw$, then $(1+pw)^{p^{e-1}}\equiv1\pmod{p^e}$ by
induction from $(1+p^sv)^p\equiv1\pmod{p^{s+1}}$ for $s\ge1$: in
$(1+p^sv)^p=1+p^{s+1}v+\sum_{i=2}^{p}\binom pi p^{si}v^i$ the middle binomial coefficients
are divisible by $p$ and $sp\ge s+1$; this is also valid for $p=2$. Applying these two
cases to $z=E_{a,k}(x)$ gives the second congruence.
\end{proof}

\begin{theorem}[Exact finite-level function rings \src{10}]\label{osq:pm:thm:finitelevel}
The coefficientwise constant term and $p$-adic evaluation induce natural ring
isomorphisms
\begin{equation}\label{osq:pm:eq:finitelevel}
 \Num_r(\Oz)/p^e\Num_r(\Oz)\ \simeq\ \Int(\Z^r)/p^e\Int(\Z^r)
 \ \simeq\ C(\Z_p^r,\Z/p^e\Z).
\end{equation}
\end{theorem}
\begin{proof}
By \cref{osq:pm:thm:numreal}, $\Num_r(\Oz)=\Pi[\mathbf X]\oplus\Int(\Z^r)$ additively,
and $p^e\Pi[\mathbf X]=\Pi[\mathbf X]$ because $\Pi$ is a $\Q$-vector space. Hence
$p^e\Num_r(\Oz)=\Pi[\mathbf X]\oplus p^e\Int(\Z^r)$, and the ring homomorphism $\ct_*$
induces the first isomorphism.

For the second, the map sends $f$ to $x\mapsto f(x)\bmod p^e$
(\cref{osq:pm:lem:padicext}). If $f=p^eg$ with $g\in\Int(\Z^r)$, its reduction is zero.
Conversely, if $f$ vanishes modulo $p^e$ at every $p$-adic tuple, its values at ordinary
integer tuples are divisible by $p^e$, so $f/p^e\in\Int(\Z^r)$. The kernel is exactly
$p^e\Int(\Z^r)$.

Let $g\in C(\Z_p^r,\Z/p^e\Z)$. The target is discrete, so every point has a residue box
$\mathbf y+p^{k}\Z_p^r$ on which $g$ is constant; by compactness of $\Z_p^r$ finitely many
of these boxes cover it. Let $k$ be the largest of their levels. Every level-$k$ box
$\mathbf y+p^k\Z_p^r$ lies in the covering box that contains $\mathbf y$, so $g$ is
constant on every residue box modulo $p^k$. For
$\mathbf a\in\{0,\dots,p^k-1\}^r$ choose an ordinary integer $c_{\mathbf a}$ lifting the
value of $g$ on the box of $\mathbf a$, and put
\[
 f(\mathbf X)=\sum_{\mathbf a}c_{\mathbf a}\prod_{j=1}^rI_{a_j,k,e}(X_j)\in\Int(\Z^r).
\]
By \cref{osq:pm:lem:indicator}, at $\mathbf x$ in the box of $\mathbf a$ the product
indexed by $\mathbf a$ is $\equiv1$ and all others are $\equiv0$ modulo $p^e$, so
$f(\mathbf x)\equiv g(\mathbf x)\pmod{p^e}$. This proves surjectivity.
\end{proof}

\begin{corollary}[The $p$-adic completion \src{10}]\label{osq:pm:cor:completion}
There is a canonical isomorphism of topological rings
\begin{equation}\label{osq:pm:eq:completion}
 \varprojlim_{e\ge1}\Num_r(\Oz)/p^e\Num_r(\Oz)\ \simeq\ C(\Z_p^r,\Z_p),
\end{equation}
where the left side carries the inverse-limit topology of the discrete quotients and the
right side the uniform $p$-adic topology. The canonical map
$\Num_r(\Oz)\to C(\Z_p^r,\Z_p)$ sends $F$ to the continuous extension of $\ct_*(F)$, has
dense image, and has kernel $\Pi[\mathbf X]$.
\end{corollary}
\begin{proof}
The isomorphisms of \cref{osq:pm:thm:finitelevel} commute with reduction of values from
level $e+1$ to level $e$. A coherent sequence of functions $g_e:\Z_p^r\to\Z/p^e\Z$
determines a unique function $g:\Z_p^r\to\Z_p$ pointwise, and $g$ is continuous exactly
when every $g_e=g\bmod p^e$ is continuous, since the sets $y+p^e\Z_p$ form a base of the
topology of $\Z_p$. This identifies the inverse limit with $C(\Z_p^r,\Z_p)$ as rings.
The kernel of the projection to level $e$ corresponds to the functions that vanish modulo
$p^e$ everywhere, which are the basic neighbourhoods of $0$ in the uniform $p$-adic
topology; so the topologies agree. Surjectivity at every finite level gives density of
the image.

If the continuous extension of $f\in\Int(\Z^r)$ is zero, then $f$ vanishes on $\Z^r$, and
a rational polynomial vanishing on $\Z^r$ is zero (write it as a polynomial in $X_r$ with
coefficients in $\Q[X_1,\dots,X_{r-1}]$; for each ordinary tuple
$(n_1,\dots,n_{r-1})$ the resulting univariate polynomial has infinitely many roots, so
all coefficients vanish on $\Z^{r-1}$; induct on $r$). With \cref{osq:pm:thm:numreal} the
kernel on $\Num_r(\Oz)$ is exactly $\Pi[\mathbf X]$. Equivalently
$\bigcap_ep^e\Int(\Z^r)=0$, because an ordinary integer divisible by every $p^e$ is
zero.
\end{proof}

\begin{remark}[Not a profinite ring completion \src{10}]\label{osq:pm:rem:notprofinite}
For $r\ge1$ the characteristic functions of residue boxes at increasing depths are
infinitely many distinct elements of $C(\Z_p^r,\F_p)$. Therefore
$\Num_r(\Oz)/p\Num_r(\Oz)$ is infinite (and discrete in its quotient topology), as are all
the rings in \eqref{osq:pm:eq:finitelevel}. The completion \eqref{osq:pm:eq:completion} is
a $p$-adic completion, not an inverse limit of finite rings. This differs from the
genuine profinite completion of $\Oz$: its finite quotient presentations are the rings
$\Z/n\Z$ with kernels $n\Oz$ (\cref{osq:thm:quotients}), so
$\varprojlim_n\Oz/n\Oz\simeq\widehat\Z$.
\end{remark}

The ordinary polynomial-approximation content is part of the classical theory of
integer-valued polynomials \cite{CC}. The explicit proof above also identifies exactly
which omnific coefficients disappear, rather than only invoking density for the ordinary
subring. The compactness of $\Z_p^r$ used here is ordinary $p$-adic compactness, not
compactness of an omnific interval.

\begin{theorem}[$p$-adic classification of characters \src{10}]
\label{osq:pm:thm:characters}
Every unital class ring homomorphism $\chi:\Num_r(\Oz)\to\F_p$ has the form
\begin{equation}\label{osq:pm:eq:character}
 \chi_{\boldsymbol\alpha}(F)=\ct_*(F)(\boldsymbol\alpha)\bmod p
\end{equation}
for a unique $\boldsymbol\alpha\in\Z_p^r$, and every $\boldsymbol\alpha\in\Z_p^r$ defines
such a character. The prime ideals of $\Num_r(\Oz)$ containing $p\Num_r(\Oz)$ are exactly
the kernels $\ker\chi_{\boldsymbol\alpha}$, with distinct $\boldsymbol\alpha$ giving
distinct kernels; every such prime ideal is maximal with quotient $\F_p$.
\end{theorem}
\begin{proof}
The map $\chi_{\boldsymbol\alpha}$ is a composite of ring homomorphisms
(\cref{osq:pm:thm:numreal,osq:pm:lem:padicext}), so it is a character. Conversely, a
unital $\chi$ kills $p\Num_r(\Oz)$, so by \cref{osq:pm:thm:finitelevel} with $e=1$ it is
$\psi\circ\rho$, where $\rho(F)=\ct_*(F)\bmod p$ and $\psi:C(\Z_p^r,\F_p)\to\F_p$ is a
unital ring homomorphism. At level $k$, the indicators of the $p^{kr}$ residue boxes modulo
$p^k$ are pairwise orthogonal idempotents with sum $1$. Their images are idempotents of
$\F_p$, hence $0$ or $1$, pairwise orthogonal with sum $1$, so exactly one box of level
$k$ is selected. A box indicator is the sum of the indicators of its children at level
$k+1$, and a child indicator $c$ satisfies $c=c\cdot b$ for its parent $b$; so the
selected boxes are nested. They are nonempty compact sets with diameters tending to $0$,
so by completeness of $\Z_p^r$ they have exactly one common point $\boldsymbol\alpha$.
Every $h\in C(\Z_p^r,\F_p)$ is constant on the residue boxes of some level $k$
(as in the proof of \cref{osq:pm:thm:finitelevel}), so $h$ is an $\F_p$-linear
combination of box indicators and $\psi(h)$ is the value of $h$ on the selected box, that
is, $h(\boldsymbol\alpha)$. Hence $\chi=\chi_{\boldsymbol\alpha}$. Distinct points are
separated by a residue-box indicator, which is $\rho(F)$ for some $F$ by surjectivity; so
$\boldsymbol\alpha$ is unique.

The prime ideals of $\Num_r(\Oz)$ containing $p\Num_r(\Oz)$ correspond, through the
explicit set-sized presentation $\rho$, to the prime ideals of $C(\Z_p^r,\F_p)$. Every
element $z$ of this ring satisfies $z^p=z$ pointwise. The quotient by a prime ideal is a
commutative domain in which every element is a root of $X^p-X$; such a domain has at most
$p$ elements, and it contains the $p$ distinct images of the constants, so it is $\F_p$.
Thus every such prime is maximal with quotient $\F_p$, and the quotient map is a character,
hence some $\chi_{\boldsymbol\alpha}$. A unital surjection onto $\F_p$ is determined by its
kernel, because $\F_p$ has only the identity automorphism; with the uniqueness of
$\boldsymbol\alpha$ this shows that distinct parameters give distinct kernels.
\end{proof}

\begin{corollary}[Most characters are not omnific evaluations \src{10}]
\label{osq:pm:cor:nonevaluation}
A character $\Num_r(\Oz)\to\F_p$ obtained by evaluating at a tuple $\mathbf x\in\Oz^r$
and then applying a unital ring homomorphism $\Oz\to\F_p$ has parameter
$\boldsymbol\alpha=(\ct(x_1),\dots,\ct(x_r))\in\Z^r\subset\Z_p^r$. Conversely every
parameter in $\Z^r$ arises in this way. Parameters in $\Z_p^r\setminus\Z^r$ give
characters that are not evaluation at any omnific tuple.
\end{corollary}
\begin{proof}
By \cref{osq:thm:universal} the unital homomorphism $\Oz\to\F_p$ is uniquely $\ct$
followed by reduction modulo $p$. By \eqref{osq:pm:eq:evshadow}, its composite with
evaluation at $\mathbf x$ is $F\mapsto\ct_*(F)(\ct(\mathbf x))\bmod p$, the character with
parameter $\ct(\mathbf x)$. An ordinary integer tuple is itself an allowed input, which
gives the converse. The last assertion follows from the uniqueness of the parameter in
\cref{osq:pm:thm:characters}.
\end{proof}

For $r=1$ the lacunary $p$-adic integer \eqref{osq:pm:eq:lacunary} of the omitted type in
\cref{osq:pm:prop:type} gives a concrete nonevaluation character. There are uncountably
many parameters and only countably many integer ones. The enormous class of omnific inputs
does not enlarge the evaluation parameters beyond the ordinary integers, because the
constant term of every omnific input is an ordinary integer.

\subsection{Rational self-maps and gaps at arbitrary surreal scales}
\label{osq:pm:sub:rigidity}

Throughout this subsection $(E,\Lambda)$ is one of the pairs $(\No,\Oz)$ or $(\SC,\GOz)$;
in the second the absolute value is the surreal-valued modulus. The input and output
lattices are the full classes, not a fixed set-sized truncation. We write $\Num_1(\Lambda)$
for $\Num_1(\Oz)$, respectively $\Num_1(\GOz)$, and $E(X)$ for the rational functions in
one variable over the class field $E$. We use repeatedly the following estimate, valid in
$E$ by the properties of the modulus: if $Q(X)=\sum_{j\le n}q_jX^j\in E[X]$ has
$q_n\ne0$ and $|y|\ge R_Q:=1+2\sum_{j<n}|q_j/q_n|$, then
\begin{equation}\label{osq:pm:eq:polyest}
 \tfrac12|q_n||y|^n\le|Q(y)|\le\tfrac32|q_n||y|^n,
\end{equation}
because $|\sum_{j<n}q_jy^j|\le\sum_{j<n}|q_j||y|^{n-1}\le\frac12|q_n||y|^n$.

\begin{lemma}[Finite differences of a proper rational function \src{10}]
\label{osq:pm:lem:rationaldiff}
Let $r\in E(X)$ be nonzero and proper, and let $\Delta r(X)=r(X+1)-r(X)$. For every
ordinary $k\ge0$, $\Delta^kr$ is nonzero and proper. If the leading Laurent term of $r$ at
infinity is $cX^{-m}$ with $m\ge1$, then the leading term of $\Delta^kr$ is
\begin{equation}\label{osq:pm:eq:rationalleading}
 (-1)^kc\,m(m+1)\cdots(m+k-1)X^{-m-k},
\end{equation}
the empty product for $k=0$ being $1$.
\end{lemma}
\begin{proof}
Compute in the field of formal Laurent series in $X^{-1}$ over a set-sized subfield of $E$
containing the finitely many coefficients involved; the rational functions over that
subfield embed in it, a rational function is proper exactly when its expansion has only
negative exponents, and the shift $X\mapsto X+1$ becomes the substitution
$X^{-1}\mapsto X^{-1}(1+X^{-1})^{-1}$. The binomial series gives
\[
 (X+1)^{-n}-X^{-n}=X^{-n}\bigl((1+X^{-1})^{-n}-1\bigr)
 =-nX^{-n-1}+\text{terms of lower degree in }X .
\]
Since $r=cX^{-m}+(\text{terms }X^{-n},\ n>m)$, and $\Delta$ sends each such lower term to
terms of degree at most $-n-1<-m-1$, the leading term of $\Delta r$ is $-cmX^{-m-1}$.
Induction on $k$ gives \eqref{osq:pm:eq:rationalleading}. The coefficient is nonzero in
characteristic zero and the exponent is negative, which proves both assertions.
\end{proof}

\begin{theorem}[Polynomiality of rational lattice self-maps \src{10}]
\label{osq:pm:thm:rational}
Let $f\in E(X)$ satisfy $f(x)\in\Lambda$ for every $x\in\Lambda$ at which $f$ is defined.
Then $f\in E[X]$. In particular its reduced rational presentation has no finite poles
anywhere in $E$.
\end{theorem}
\begin{proof}
Divide the numerator by the denominator to write $f=P+r$ with $P\in E[X]$ and $r$ proper,
and suppose $r\ne0$. Choose an ordinary $k\ge1$ larger than $\deg P$ (any $k\ge1$ if
$P=0$). Then $\Delta^kP=0$, so $g:=\Delta^kf=\Delta^kr$ is nonzero and proper by
\cref{osq:pm:lem:rationaldiff}. Write $f$ and $g$ in reduced form, $f=N_f/D_f$ and
$g=N_g/D_g$, with $\deg N_g<\deg D_g$.

We find $B\in\No$ such that for every $x\in\No$ with $x>B$ all of $f(x),\dots,f(x+k)$ are
defined, $g(x)$ is defined, and
\begin{equation}\label{osq:pm:eq:smallrational}
 0<|g(x)|<1 .
\end{equation}
By \eqref{osq:pm:eq:polyest}, $D_f(y)\ne0$ for $|y|\ge R_{D_f}$, $N_g(y)\ne0$ for
$|y|\ge R_{N_g}$, and for $|y|\ge\max(1,R_{N_g},R_{D_g})$,
\[
 |g(y)|\le\frac{\frac32|\ell(N_g)||y|^{\deg N_g}}{\frac12|\ell(D_g)||y|^{\deg D_g}}
 \le3\,\frac{|\ell(N_g)|}{|\ell(D_g)|}\,|y|^{-1},
\]
where $\ell(\cdot)$ is the leading coefficient. Take $B$ larger than $1$, than
$R_{D_f},R_{N_g},R_{D_g}$ and than $3|\ell(N_g)/\ell(D_g)|$; for real $x>B$ the points
$x,\dots,x+k$ are real and larger than $B$, so every requirement holds. These are finitely
many inequalities in the ordered field $\No$.

By \cref{osq:pm:lem:lattice} choose $x\in\Oz$ with $x>B$. Then $x+j\in\Oz\subseteq\Lambda$
for $0\le j\le k$, $f$ is defined at these points, and
\[
 g(x)=\sum_{j=0}^k(-1)^{k-j}\binom kjf(x+j)\in\Lambda .
\]
This contradicts \eqref{osq:pm:eq:smallrational} and the lattice separation of
\cref{osq:pm:lem:lattice}. Therefore $r=0$ and $f=P$.
\end{proof}

The finite-difference argument is a general discrete-lattice method. It is included as a
supporting result, not asserted to be a new ordinary theorem about rational functions on
$\Z$.

Recall that $y\in E$ is finite if $|y|$ is bounded by an ordinary integer and
infinitesimal if $|y|<1/n$ for every ordinary $n\ge1$; the standard part of a finite
element is the unique real (respectively complex) number at infinitesimal distance, taken
coordinatewise in $\SC$ (\cref{osq:pm:rem:modulus}).

\begin{lemma}[Finite roots and reduction \src{10}]\label{osq:pm:lem:finiteroot}
A root in $E$ of a monic polynomial with finite coefficients is finite. The standard part
is a ring homomorphism on the finite elements of $E$ and carries such a root to a root of
the coefficientwise standard-part polynomial. For finite $z\in\SC$,
\[
 \st(|z|)=|\st(z)|.
\]
\end{lemma}
\begin{proof}
Let $z^d+\sum_{j<d}c_jz^j=0$ with $|c_j|\le C$ for an ordinary integer $C$. If
$|z|>1+dC$, then
\[
 \Bigl|\sum_{j<d}c_jz^j\Bigr|\le dC|z|^{d-1}<|z|^d,
\]
which excludes a root; so every root is finite. Sums and products of finite elements with
infinitesimal differences have infinitesimal differences, so $\st$ is a ring homomorphism
on the finite elements (\cref{osq:lem:ct} for $\OO_\R$ and $\OO_\C$) and commutes with
evaluation of polynomials; applying it to the equation gives the second assertion. For
finite $z\in\SC$, $|z|\le|\Re z|+|\Im z|$ is finite; taking standard parts in
$|z|^2=(\Re z)^2+(\Im z)^2$ gives $\st(|z|)^2=|\st(z)|^2$, and $\st(|z|)\ge0$, so
uniqueness of nonnegative real square roots gives the identity.
\end{proof}

For an ordinary $d\ge2$ define the ordinary integer
\begin{equation}\label{osq:pm:eq:qd}
 q_{E,d}=
 \begin{cases}
  1,&E=\No\text{ and }d\text{ odd},\\
  2,&E=\No\text{ and }d\text{ even},\\
  d,&E=\SC.
 \end{cases}
\end{equation}

\begin{theorem}[A bounded number of lattice preimages at every scale \src{10}]
\label{osq:pm:thm:localcount}
Let $F\in\Num_1(\Lambda)$ have degree $d\ge2$, and let $H\in\No_{>0}$. There is a positive
surreal $T_0$ such that, for every $T\in\Oz$ with $T>T_0$, the class
\begin{equation}\label{osq:pm:eq:preimages}
 \{u\in\Lambda:|F(u)-F(T)|\le H\}
\end{equation}
has at most $q_{E,d}$ elements. In particular $F(\Lambda)$ has at most $q_{E,d}$ points in
the closed ball of radius $H$ about $F(T)$.
\end{theorem}
\begin{proof}
Write $F(X)=a_dX^d+\sum_{j<d}a_jX^j$ with $a_d\ne0$, and put
$T_0=\omega\cdot\max\bigl(1,H/|a_d|,|a_j/a_d|\ (j<d)\bigr)$. For $T>T_0$ we have $T>1$
and, for every ordinary $n\ge1$, $T>n|a_j/a_d|$ and $T>nH/|a_d|$. Since $d-j\ge1$ and
$d-1\ge1$, it follows that
\begin{equation}\label{osq:pm:eq:scaling}
 b_j:=\frac{a_j}{a_dT^{d-j}}\ (j<d)\quad\text{and}\quad\frac{H}{|a_d|T^{d-1}}
 \quad\text{are infinitesimal.}
\end{equation}
Fix such a $T$ and an element $u$ of \eqref{osq:pm:eq:preimages}, and put $z=u/T$. Dividing
$F(u)-F(T)$ by $a_dT^d$ gives
\begin{equation}\label{osq:pm:eq:normalizedroot}
 z^d+\sum_{j<d}b_jz^j=1+\sum_{j<d}b_j+\delta,\qquad
 |\delta|\le\frac{H}{|a_d|T^d}\le\frac{H}{|a_d|T^{d-1}},
\end{equation}
so $\delta$ is infinitesimal. Thus $z$ is a root of a monic polynomial with finite
coefficients, hence finite by \cref{osq:pm:lem:finiteroot}, and taking standard parts
gives $\st(z)^d=1$. There are at most $d$ such residues in $\C$; among real residues only
$1$ is possible for odd $d$, and only $\pm1$ for even $d$. So there are at most $q_{E,d}$
possible residues $\st(u/T)$.

It remains to show that each residue belongs to at most one lattice input. Let $u,v$ be
elements of \eqref{osq:pm:eq:preimages} with $\st(u/T)=\st(v/T)=\zeta$, and put $z=u/T$,
$w=v/T$. Write $F(u)-F(v)=(u-v)Q_F(u,v)$ with
$Q_F(u,v)=\sum_{j=1}^da_j\sum_{\ell=0}^{j-1}u^{j-1-\ell}v^\ell$. Direct expansion gives
\begin{equation}\label{osq:pm:eq:normalizeddifference}
 \frac{Q_F(u,v)}{a_dT^{d-1}}
 =\sum_{\ell=0}^{d-1}z^{d-1-\ell}w^\ell
 +\sum_{j=1}^{d-1}b_j\sum_{\ell=0}^{j-1}z^{j-1-\ell}w^\ell .
\end{equation}
Its standard part is $d\zeta^{d-1}$, of modulus $d$ because $\zeta^d=1$; by
\cref{osq:pm:lem:finiteroot} the standard part of its modulus is $d$, so
\[
 |Q_F(u,v)|>\frac d2|a_d|T^{d-1}.
\]
If $u\ne v$, the triangle inequality gives $|F(u)-F(v)|\le2H$, hence
\[
 |u-v|=\frac{|F(u)-F(v)|}{|Q_F(u,v)|}<\frac{4H}{d|a_d|T^{d-1}}<1,
\]
the middle quantity being infinitesimal by \eqref{osq:pm:eq:scaling}. This contradicts
lattice separation for $0\ne u-v\in\Lambda$ (\cref{osq:pm:lem:lattice}). So $u=v$, and
\eqref{osq:pm:eq:preimages} has at most $q_{E,d}$ elements. Every image point in the ball
is the image of an element of \eqref{osq:pm:eq:preimages}, which gives the last assertion.
\end{proof}

The proof uses the finite residue polynomial $Z^d-1$ only as a label for possible clusters.
It assumes no convergence of an iterative root-finding scheme, no compactness of an
omnific interval, and no existence theorem for new roots in the full field.

\begin{theorem}[Arbitrary-scale holes in nonlinear polynomial images \src{10}]
\label{osq:pm:thm:holes}
Let $F\in\Num_1(\Lambda)$ have degree $d\ge2$. For every $H>0$ and every $B\ge0$ in $\No$
there is $c\in\Lambda$ with $|c|>B$ such that
\begin{equation}\label{osq:pm:eq:hole}
 \{y\in E:|y-c|\le H\}\cap F(\Lambda)=\varnothing .
\end{equation}
For the real lattice this is an empty closed interval; for the Gaussian lattice it is an
empty closed surcomplex disk.
\end{theorem}
\begin{proof}
Put $q=q_{E,d}$ and choose, by \cref{osq:pm:lem:lattice}, a positive $M\in\Oz$ with $M>H$.
Let $L=M+1$ and $R=3qL+M$. Apply \cref{osq:pm:thm:localcount} with radius $R$ to get
$T_0$, and choose $T\in\Oz$ with $T>T_0$ and also
\[
 |F(T)|>B+3qL .
\]
This is possible: by \eqref{osq:pm:eq:polyest}, for $T\ge\max(1,R_F)$ one has
$|F(T)|\ge\frac12|a_d|T^d\ge\frac12|a_d|T$, so any $T\in\Oz$ exceeding $T_0$, $1$, $R_F$
and $2(B+3qL)/|a_d|$ works (\cref{osq:pm:lem:lattice}).

Set $y=F(T)\in\Lambda$ and consider the $q+1$ lattice points $c_j=y+3jL$,
$j=0,\dots,q$. The closed balls of radius $M$ about them are pairwise disjoint, since
distinct centres are at distance at least $3L>2M$. Each lies in the closed ball of radius
$R$ about $y$, since $|x-y|\le|x-c_j|+3jL\le M+3qL$. That ball contains at most $q$ points
of $F(\Lambda)$ by \cref{osq:pm:thm:localcount}, so one of the $q+1$ smaller balls
contains none. Its centre $c=c_j$ satisfies $|c|\ge|y|-3qL>B$, and since $H<M$ the ball
of radius $H$ about $c$ is disjoint from $F(\Lambda)$ as well.
\end{proof}

Failure of surjectivity is therefore not a single small defect near zero: the gaps exceed
any prescribed infinite surreal radius and occur beyond any prescribed bound in modulus,
that is, at arbitrarily large moduli.

\begin{corollary}[All rational bijections of the two lattices \src{10}]
\label{osq:pm:cor:bijections}
A rational function $F\in E(X)$ defining a bijection $\Lambda\to\Lambda$ has the form
\[
 F(X)=uX+b,\qquad b\in\Lambda,\quad u\in\Lambda^\times,
\]
and every such map is a bijection. Thus the rational bijections of $\Oz$ are
$X\mapsto\pm X+b$ with $b\in\Oz$, and those of $\GOz$ are $X\mapsto uX+b$ with
$b\in\GOz$ and $u\in\{\pm1,\pm i\}$.
\end{corollary}
\begin{proof}
By \cref{osq:pm:thm:rational}, $F$ is a polynomial, so $F\in\Num_1(\Lambda)$. A constant is
not a bijection of $\Lambda$. Degree at least two is excluded by \cref{osq:pm:thm:holes},
because the centre $c$ there is itself a lattice point missing from the image. Hence
$F(X)=aX+b$ with $b=F(0)\in\Lambda$ and $a=F(1)-F(0)\in\Lambda$. Surjectivity gives
$x\in\Lambda$ with $F(x)=b+1$, so $ax=1$ and $a$ is a unit of $\Lambda$. Conversely, an
affine map with slope $u\in\Lambda^\times$ has the lattice-preserving inverse
$X\mapsto u^{-1}(X-b)$. The unit lists are those of \cref{osq:lem:units}.
\end{proof}

Here ``rational'' means a rational function in one variable over the ambient field $E$.
Arbitrary functions $\Lambda\to\Lambda$, and maps such as coordinatewise conjugation on
$\GOz$, are not being classified.

\subsection{Scope of the polynomial arithmetic}\label{osq:pm:sub:scope}

Two mechanisms act in this section and have different hypotheses. The reflection theorem
\cref{osq:pm:thm:numreflection} transports the size-sensitive divisibility mechanism of
\cref{osq:thm:universal} through finite polynomial algebra; that mechanism applies to
arbitrary set coefficient fields in the class Hahn construction. The results of
\cref{osq:pm:sub:rigidity} use instead \emph{finite reduction plus lattice separation}: at
a large scale $T$, polynomial preimages of a prescribed output ball reduce to roots of
$Z^d-1$, and a nonzero reduced divided difference would force two inputs of one cluster to
lie less than $1$ apart, which the lattice forbids. This gives
\cref{osq:pm:thm:localcount,osq:pm:thm:holes} without a count of points in a surreal
interval or a real-valued asymptotic density; it uses the real or surcomplex modulus, the
standard part and the unit-spaced lattices. Neither mechanism should be used as a
substitute for the other.

Manuscript~10 proposes \cref{osq:pm:thm:numreflection} and the arbitrary-scale bound
\cref{osq:pm:thm:localcount} as new; the other results of this section are explicit
consequences or transfers of classical arithmetic facts (Newton interpolation,
binomial-ring facts \cite{ElliottBinomial}, the least-common-multiple pattern \cite{CGG},
$p$-adic approximation by integer-valued polynomials \cite{CC}), proved here to make the
scope complete. The results do not classify all class ideals of $\Oz$ or of
$\Num_r(\Oz)$, all automorphisms of $\Oz$, or the factorization of arbitrary omnific
integers, and they neither construct nor use a global surcomplex exponential. They do not
assert a general transfer principle from $\Z$ to $\Oz$; the omitted type of
\cref{osq:pm:prop:type} is a concrete warning against such a principle. The Gaussian
theorem \cref{osq:pm:thm:numgaussian} does not replace the arithmetic of $\Int(\Z[i]^r)$
by a binomial basis, and no Gaussian analogue of the congruence criterion
\cref{osq:pm:thm:cp} is given. For $r\ge1$ the completion of
\cref{osq:pm:cor:completion} is a $p$-adic, not a profinite, completion. The
classification of \cref{osq:pm:cor:bijections} concerns rational functions only.
Manuscript~10 records a question suggested by this material: for the full omnific
integers, determine how much an automorphism can do inside the ideal $\Pi$ while fixing
the ordinary shadow $\ct$ (every ring endomorphism of $\Oz$ preserves $\ct$, by the
uniqueness in \cref{osq:thm:universal}).

\section{Natural ordinal arithmetic is not a quotient of the omnific integers}\label{osq:sec:ordinal}

This section is manuscript~11's. Let $\mathcal G_{\mathrm{nat}}$ be the Grothendieck ring of the
ordinals under Hessenberg's natural sum and natural product. Every ordinal has a finite Cantor normal
form $\sum_jn_j\omega^{\gamma_j}$, each $\gamma_j$ has its own finite Cantor normal form
$\sum_\alpha m_{j,\alpha}\omega^\alpha$, and under natural multiplication
$\omega^{\gamma_j}=\prod_\alpha(\omega^{\omega^\alpha})^{m_{j,\alpha}}$. Uniqueness of these forms gives
the free commutative semiring $\N[T_\alpha:\alpha\in\On]$, and adjoining additive inverses gives
\begin{equation}\label{osq:eq:ordinalring}
 \mathcal G_{\mathrm{nat}}\cong\Z[T_\alpha:\alpha\in\On],
\end{equation}
with $T_\alpha$ represented by $\omega^{\omega^\alpha}$ under the natural embedding in $\No$. Natural
operations, not the usual noncommutative ordinal operations, are essential. The description
\eqref{osq:eq:ordinalring} appears in the question \cite{MO188430}; the argument above fixes the operations.

\begin{theorem}[The ordinal quotient question \src{11}]\label{osq:thm:ordinal}
Every unital ring map $h:\Oz\to\mathcal G_{\mathrm{nat}}$ has image exactly the ordinary integral
constants. In particular $\mathcal G_{\mathrm{nat}}$ is not isomorphic to a quotient ring of $\Oz$.
More generally, no ring map $\Oz\to\Z[X]$ with $X\ne\varnothing$ is surjective.
\end{theorem}
\begin{proof}
\emph{First proof.} $\Z[X]$ is residually set-sized by \cref{osq:thm:residual} (for any class $X$ of
indeterminates), so $h$ kills $\Pi$ and factors through $\ct$; its image is $\Z$, a proper subring
since $T_0\notin\Z$.

\emph{Second proof, specific to this target.} The additive group of $\Pi$ is divisible
($s/n\in\Pi$ for every ordinary $n\ge1$), while no nonzero element of $\Z[T_\alpha]$ is divisible by
every $n$ (each of its finitely many integral coefficients would be). So every additive map from $\Pi$
to $\Z[X]$ is zero, and a unital ring map from $\Oz=\Z+\Pi$ has image $\Z$.
\end{proof}

The question was posted by Jesse Elliott as the final addition to MathOverflow question~188430,
``Transcendence degree of the surreals over the subfield generated by the ordinals'' (30 November 2014)
\cite{MO188430}; the displayed answer, by Eric Wofsey, concerns the transcendence degree, a different
question. The answer to the quotient question is negative. 11 gives the short additive proof precisely so
that the answer is not mistaken for a difficult new factorization theorem: no claim is made that the
elementary obstruction was unknown to specialists or absent from every other source. The first proof
illustrates the general residual theorem, and the second isolates the weaker mechanism that suffices for
this target. The obstruction is not a cardinal comparison, since both rings may be proper classes; it is
their divisibility and representation properties (11).

\section{Boundaries, non-claims and open questions}\label{osq:sec:boundaries}

\subsection{What the results do not say}\label{osq:sub:notsay}
The universal theorem is a statement about homomorphisms into set-sized rings and about
set-sized modules. It is not an isomorphism $\Oz\cong\Z$; it does not classify class ideals,
class modules, automorphisms or factorizations of $\Oz$; it concerns multiplicative maps and
says nothing against additive or valuation-theoretic observations; and it fails verbatim for every
fixed set-sized exponent group, whose identity map is a set-sized target. Size alone does not prove
it (finite supports have proper-class many monomials and an augmentation), nor does divisibility
alone (only targets of positive characteristic are handled). The set-sized thresholds are exact
for the five constructions and the two fixed-group examples treated; no classification of exponent
groups or support restrictions is claimed, and generator counts follow no general law stated here.
The large-quotient results describe specified subalgebras of $\Oz/(1+\omega^a)$ and conditional
statements about primes; no existence theorem for nonarithmetic class primes is claimed.

As an internal ring, $\Oz$ has all surreal scales available through its fraction field, a nonzero
proper-class ideal $\Pi$, nonzero class-valued derivations and nonarithmetic prime quotients; as a
source of representations on sets it carries no more scalar information than $\Z$. The gap is caused
not by ignoring limits but by finite identities that become available once the source contains the
appropriate Hahn sums (06). Adapting 06's list of hypotheses that must remain visible to all the
constructions, the proofs use four different size restrictions, each for its own reason: individual supports are sets, so that a common positive gap exists; the full class of smaller
scales supplies arbitrarily large collision families; in the set-sized models an uncountable bound
admits the countable witnesses, and regularity (or cofinality, or the finite-coordinate invariant)
bounds unions of fewer than the relevant number of supports; and cardinal arithmetic enters only in
computing the size of the ring, which is the exact threshold. Dropping any of them needs a new
argument: finite supports lack the witnesses (\cref{osq:sec:support}), full supports over a fixed set
group need not have gaps (\cref{osq:rem:fstar}), and singular support bounds may preserve field closure
but not the simultaneous gap property.

\subsection{Non-claims of the sources}\label{osq:sub:nonclaims}
Every limitation stated by a source survives. They are listed and numbered by source;
several recur in the text where the corresponding result is proved.

\paragraph{03.} (1) Not peer reviewed; no Lean declarations. (2) ``Developed here'' is relative to the
inspected material; no certified priority. (3) Neither Conway's general refinement problem nor the general factorization problems for omnific
integers are settled.
(4) No classification of exponent groups or support restrictions; the cardinal theorem is a
realization theorem for its family. (5) $\cf\kappa=\aleph_0$ is excluded. (6) Only the lower bound
$\mu$ on the number of algebra generators of $\NN_\kappa$ is proved. (7) No complete integrality test
for mixed supports. (8) No descent equality $\NN\cap F^{03}_\kappa=\NN_\kappa$. (9) No statement on
proper-class-valued valuations. (10) Infinite set-sized images of the normalizations are not
classified, and the absence of finite images is a general consequence of containing the real
algebraic integers. (11) No class tensor products and no class prime-ideal theorem; all $\Tor$ is over
set-sized rings. (12) The finite checks do not verify the cardinal or class arguments. (13) Restricted
Hahn fields, truncation integer parts, complete-integral-closure terminology and the small-unit lemma
are standard or elementary, with no priority claim; Kuhlmann is context only. (14) The full-class
theorem, common divisors and clearing are credited to the companion manuscript (05). (15) Density produces no convergent set-indexed net: every set-generated intermediate algebra is
uniformly discrete, so the dense normalization needs more than any set of generators. (16)
$\tfrac12\NN[i]$ is not asserted to be a ring. (17) The flatness result is an application of standard
module theory, not a new general flatness theorem. (18) Density does not impose continuity on abstract
ring homomorphisms. (19) No claim about unrestricted surcomplex analytic continuation is made.

\paragraph{04.} (1) Unrefereed; no Lean formalization; the universal theorems are candidate
contributions without certified priority. (2) No classification of all class ideals, no class Zorn
argument, no claim that proper class ideals extend to maximal ones. (3) The module equivalence covers
set-sized modules only, not class modules, group representations of $\mathrm{GL}_n(\Oz)$, or
additive homomorphisms. (4) It does not say that $P(\ct(x))=0$ implies $P(x)=0$. (5) It does not
classify infinite points of curves; the earlier project's curve questions are not declared solved.
(6) Nothing is inferred about fraction fields of the bounded-support rings. (7) The fixed-workspace
results are lower bounds only. (8) No named open conjecture is solved. (9) Additive coefficient maps
escape the theorem and are not counterexamples. (10) The telescope has a direct precursor
\cite[Remark~3.4.4]{LM}; the constant term, finite quotients and common shifts come from the prior
manuscript (02 of \cite{RepoOdg}). (11) The existence of standard part is classical and in the
repository's infrastructure. (12) The support-threshold theorem differs from fixed-group completion or
omitted-type questions. (13) The finite checks do not prove normal forms, infinite identities, class
factorizations or novelty. (14) The quotient, module, derivation and localization results are a
connected theorem package deduced from the universal theorems, not unrelated discoveries.

\paragraph{06.} (1) Unrefereed, not Lean-verified; priority not established by the targeted searches.
(2) $\pd\Z$ is shown only to be at least $2$. (3) Residue fields of nonarithmetic primes are not
classified; the prime-reservoir proposition is a necessary condition, not an existence theorem.
(4) The primality of $\omega^{\sqrt2}+\omega+1$ is imported; no factorization, refinement, GCD or
unique-factorization conjecture is used or claimed. (5) No automorphisms, class ideals or factorization
patterns are described. (6) No elementarity claim for $A^{06}_\kappa\subseteq\Oz$; no compactness
argument for a proper-class theory. (7) The Euler derivation is not asserted to be the
Berarducci--Mantova derivation or compatible with exponentiation. (8) Real closedness of
bounded-support Hahn fields is not a novelty claim, and regularity is not claimed necessary for field
closure. (9) The idea of a universal set-sized quotient is not new to the repository. (10) The module
correspondence is not a Morita equivalence of unrestricted module theories. (11) The invisible prime
quotient does not say that a particular maximal class ideal fails to exist. (12) The polynomial-test
theorem does not assert that a proper-class coordinate ring defines an ordinary scheme. (13) The
finite-congruence observation is elementary and not a strong novelty claim. (14) Numerical tests and a
successful build are not proofs; the formalization route was not executed.

\paragraph{07.} (1) Conway's refinement conjecture is neither settled nor used (the Abramov
development \cite{Abramov} is cited for its claim and caveats only); no historical priority; no referee;
no Lean. (2) The Cantor algebra is a specified subalgebra; neither the idempotents nor the prime ideals
of $\Qa_a$ are classified, and $\Qa_a$ is not claimed to be a function ring. (3) The statement on
nonarithmetic primes is conditional; no existence is asserted. (4) Surreal self-similarity is credited
to Bagayoko--van der Hoeven. (5) $\Theta_a$ is not claimed to preserve the ordinal embedding,
birthdays, simplicity or the exponential; $\vartheta$ is not additive and $\Psi$ not multiplicative.
(6) Quotient, completion and module classifications are consequences, not separate major discoveries.
(7) The small-image theorem does not say that no set-valued invariant sees positive exponents. (8) The
module equivalence is not a homological-dimension statement over class modules. (9) The global theorem
is not a theorem about small approximations of $\Oz$. (10) The finite checks do not verify the
transfinite or ambient-support arguments. (11) The CRT isomorphisms at different stages need not be
coherent; evaluation at a point of $\Z_2$ is not claimed to extend to $\Qa_a$; no topology on $\EE_a$
is used. (12) Class prime ideals are not classified, and no Zorn argument for a proper class of
ideals is used. (13) No exponential function on surcomplex arguments is assumed.

\paragraph{08.} (1) Candidate-original, not certified priority, not refereed, no Lean; the broad
factorization questions for $\Oz$ are not addressed. (2) $\Frac(A^{08}_\kappa)=F^{08}_\kappa$ is not
claimed. (3) Residue-field isomorphism types are not classified. (4) No closure under derivations or
exponentials, no initial-segment or birthday condition on the embedding. (5) The collision mechanism
has published precedents (\cite[Proposition~8.2.1]{LM}); the constant retraction and finite reductions
are not principal novelty claims. (6) Fixed Hahn workspaces lie outside the theorem (an Archimedean group has no
small scales at all). (7) The non-integral-closure corollary is not a factorization theorem. (8) Class-sized modules
are not classified; the theorem is not an isomorphism with $\Z$, and individual monomials can be
represented faithfully in small rings. (9) The Hahn derivations $D_\eta$ are not an
exponential-compatible calculus. (10) No maximal-ideal theorem for proper-class rings; no assertion that
proper class ideals lie in maximal ones. (11) The quantitative criterion gives sufficient conditions
only. (12) The finite checks prove neither the class theorem, nor summability, nor real closedness.
(13) The primality of $\omega^{\sqrt2}+\omega+1$ is an imported published theorem.

\paragraph{09.} (1) Novelty provisional after a targeted comparison; no referee; no proof-assistant
verification. (2) No named factorization conjecture solved; primality imported. (3) The number of
maximal ideals and the isomorphism types of the residue fields of $A^{09}_{\lambda,D}$ are not
determined. (4) No assertion about the unrestricted (uncountable-support) ring over
$\Gamma^{09}_\lambda$. (5) No class Zorn argument; the Jacobson statements concern set-sized rings.
(6) $\GOz$ is a subring of the surcomplex field, not the field. (7) The checks do not verify infinite
supports, cardinal arithmetic or the existence of maximal ideals. (8) The geometric identity is
elementary with precedents; profinite, derivation and flatness statements are consequences; the
Gaussian $2$-torsion computation is standard. (9) No global surcomplex exponential, sine, integration or
topology is used; the arithmetic quotient is not an ordered quotient. (10) The rotation quotient is a
conceptual predecessor, not imported. (11) No new exponential derivation is constructed and no
uniqueness of the Berarducci--Mantova derivation is claimed. (12) $I/I^2=0$ is not a
finite-presentation or smoothness claim.

\paragraph{10.} (1) Unrefereed and not machine-checked; targeted review only; no named factorization
conjecture solved. (2) For $r\ge1$ the $p$-adic completion is not a profinite completion. (3) No
classification of all class ideals, automorphisms or factorizations. (4) No transfer principle from
$\Z$ to $\Oz$. (5) $\Int(\Z[i]^r)$ is not replaced by a binomial basis. (6) The derivation $D_h$ is not the
Berarducci--Mantova derivation. (7) The universe-relative reading must carry the size restriction in
its statement. (8) The rational classification is not a classification of arbitrary functions or of
coordinatewise conjugation. (9) The congruence theorem does not make reduction modulo an infinite
omnific integer an ordinary finite congruence, and does not classify congruence-preserving class
functions. (10) Classical ingredients (Newton interpolation, binomial rings, the least-common-multiple
pattern, $p$-adic approximation) are credited, not claimed. (11) The grid test assumes exact surreal
input; no effective encoding is claimed. (12) The finite checks do not test class functions, support
lemmas or priority.

\paragraph{11.} (1) The main theorem is relative to one fixed universe, not to fragments viewed in a
larger universe. (2) No classification of prime or maximal class ideals. (3) No class module of
K\"ahler differentials. (4) The presentation theorem is not quantifier elimination or a
reconstruction of class-valued solutions. (5) Sharpness of the coinitiality theorem is stated under
$\kappa^{<\kappa}=\kappa$ and $|K|\le\kappa$ as explicit extra assumptions (the unconditional value is a
merge result, \cref{osq:prop:card11}). (6) Necessary closure conditions for intermediate rings are not
given. (7) The ordinal answer is not claimed hard or unknown to specialists. (8) Bounded-support
coefficient rings and lower-scale geometric division have precedents in \cite[Section~8]{LM}.
(9) Targeted novelty assessment only; the finite checks prove no infinite-support or class-size
statement. (10) The rotation-group theme is acknowledged, not originated. (11) The only answer to
the MathOverflow question concerns transcendence degree. (12) No primality or refinement conjecture is
newly settled; the ideals, class-valued representations and factorization theory of quotients invisible
to set-sized rings are not determined, and ordinary factorization questions about $\Oz$ remain
outside the conclusions.

\paragraph{The merge.} (1) The results marked \merge\ are consequences of the sources' lemmas; they
are no more refereed than the sources and are not formalized. (2) No general law is asserted for the
generator counts. (3) For full-support Hahn rings over a prescribed set group, exact thresholds are computed only in
\cref{osq:ex:lexQQ,osq:ex:countablecoeff}; the general questions of 04 and 10 stay open. (4) The
support-threshold theorem is stated for $k\in\{\R,\C\}$ as in 04 and 11, not for every set field.
(5) No Gaussian analogue of 10's congruence-preservation criterion is given.

\subsection{Open questions and their status}\label{osq:sub:questions}
The sources pose the following questions, each as a question suggested by the proofs rather than as
an established open problem. Several coincide and are merged.

\begin{question}[Least detecting targets over a fixed exponent group; 03, 04, 07, 08, 09, 10, 11]\label{osq:q:fixed}
For a set-sized ordered group $\Gamma$, a family of admissible supports and a specified element (or
the whole positive-support ideal), what is the least cardinality of a ring target detecting it, and
which intrinsic support conditions characterize it?
\end{question}
\emph{Status.} Open in general. The field bound (\cref{osq:prop:fieldbound}) is a lower bound that is
exact for the five constructions and for the two examples of \cref{osq:sub:fieldbound}; there it answers
04's ``optimal small targets'' question for its own example, and for 10's question ``when is the
bound attained'' it shows that the bound need not be attained (in the rational example, and for
$A^{03}_{\aleph_1}$ if CH fails). 08's criterion
and 11's coinitiality theorem give sufficient conditions (07's ``sharp set-sized analogues'' and 09's
``quantitative problem'' are the same question).

\begin{question}[Cardinal arithmetic without the matching hypothesis; 06]\label{osq:q:kappa}
For $\kappa$ uncountable regular with $\kappa^{<\kappa}>\kappa$, what is the least size of a target
detecting a nonzero element of $\Pi^{06}_\kappa$, and is there an integer part of size exactly $\kappa$
with the same lower bound?
\end{question}
\emph{Status.} Answered. The least size is $\kappa^{<\kappa}$ in ZFC (\cref{osq:thm:model06}, a merge
result from 06's collision lemma); 08 constructs, for every infinite $\kappa$ in ZFC, an integer part of
size exactly $\kappa$ with threshold $\kappa$ (\cref{osq:thm:model08}), without claiming its fraction field.

\begin{question}[Higher projective dimension; 06]\label{osq:q:pd}
What is $\pd_A\Z$ for the set-sized integer parts, and how does it depend on the cofinality of the
principal-ideal chain?
\end{question}
\emph{Status.} Open; only $\fd_AD=1$ and $\pd_AD\ge2$ are known, now for all five constructions.

\begin{question}[Nonarithmetic primes and their residue fields; 06, 07]\label{osq:q:primes}
Which prime or maximal class ideals of $\Oz$ have nonzero constant-term image and zero intersection with
$\Z$, and which class fields occur as $\Frac(\Oz/P)$ for primes $P\not\supseteq\Pi$? Can the embedded
field $\FF_a$ distinguish such primes?
\end{question}
\emph{Status.} Open. Necessary conditions: \cref{osq:if:cor:primefield,osq:if:prop:reservoir,osq:if:rem:composition}.

\begin{question}[The ambient binomial quotient; 07]\label{osq:q:binomial}
Which normal-form invariants describe the elements and ideals of $\Qa_a$ beyond the Cantor algebra;
which idempotents lie outside it; how are cyclotomic decompositions at different scales related?
\end{question}
\emph{Status.} Open.

\begin{question}[Residue fields of the set-sized models; 08, 09]\label{osq:q:residue}
Which characteristic-zero fields of the threshold cardinality occur as maximal residue fields of
$A^{08}_\kappa$ and of $A^{09}_{\lambda,D}$, and how many maximal ideals are there?
\end{question}
\emph{Status.} Open; existence, cardinality and at least $\lambda$ nonarithmetic maximal ideals
(for 09) are proved.

\begin{question}[Extra structure on the workspace; 08]\label{osq:q:extra}
Can prescribed residue-field behaviour be combined with closure under a surreal derivation or an
exponential, or with an initial embedding, without destroying the finite-coordinate invariant or the
cardinality?
\end{question}
\emph{Status.} Open.

\begin{question}[Intermediate rings; 04, 11]\label{osq:q:intermediate}
Which rings between the finite-support ring and $D+\Pi_k$ (for instance, given by allowed families of
support order types) still have set-sized reflection $D$? Is enough closure under subordinate division
also necessary?
\end{question}
\emph{Status.} Open; \cref{osq:thm:support} settles the case of cardinal support bounds. 09's question on
support bounds below a regular cardinal is addressed for the groups of 06 and 11
(\cref{osq:sub:model06,osq:sub:model11}), where the support sizes of inverses and quotients are checked.

\begin{question}[Beyond ring observations; 04]\label{osq:q:weaker}
Which weaker multiplicative compatibility conditions still force factorization through the residue, and
which structured families of set-sized local models can retain a specified family of scales?
\end{question}
\emph{Status.} Open.

\begin{question}[Invisible class domains and automorphisms; 10, 11]\label{osq:q:invisible}
What are the ideals, class-valued representations and factorization theory of proper-class quotient
domains invisible to all set-sized rings? How much can an automorphism of $\Oz$ do inside $\Pi$ while
fixing $\ct$?
\end{question}
\emph{Status.} Open.

\begin{question}[Normalization; 03]\label{osq:q:normalization}
Give a membership test for $\NN$ on mixed supports; describe the integral closure over a fixed finitely
generated exponent subgroup; decide the descent $\NN\cap F^{03}_\kappa=\NN_\kappa$; find the algebra
generation number of $\NN_\kappa$ between $\cf\kappa$ and $\kappa^{\aleph_0}$; find a class-sized
replacement for the valuative criterion of integrality; describe the infinite set-sized images of $\NN$.
\end{question}
\emph{Status.} Open. 04's question on Diophantine fibers belongs to the sibling report \cite{RepoOdg}.

\section{Verification record and formalization routes}\label{osq:sec:provenance}

\subsection{Finite checks}\label{osq:sub:checks}
Each manuscript ships a program that checks finite identities with exact arithmetic. They are
kept byte-identical in \texttt{code/} with their recorded outputs in \texttt{data/}; all were rerun
for this merge on copies (Python~3.14.4, SymPy~1.14.0) and reproduce the recorded outputs, up to line
endings and the recorded Python version.

\begin{center}\small
\begin{tabular}{@{}l>{\raggedright\arraybackslash}p{0.56\textwidth}rl@{}}
\toprule
Source & What is checked & Checks & Needs\\
\midrule
03 & finite telescopes in a symbolic lattice, the unit series coefficients, lexicographic comparisons & 2{,}362 & stdlib\\
04 & 1{,}950 finite telescopes in $\Q\oplus_{\mathrm{lex}}\Q$ with remainders and positive supports, 600 constant-term and 900 augmentation products, 3 support failures and 3 distinguishing examples & 5{,}406 & stdlib\\
06 & geometric remainders and supports, constant terms, boundary cases, Euler derivation, real complex-structure blocks & 17{,}586 & stdlib\\
07 & weighted geometric identities, binomial remainders, first idempotents, dyadic refinement maps & 621 & SymPy\\
08 & telescopes, constant terms, Hahn Leibniz identities, binomial powers, root corrections & 1{,}277 & stdlib\\
09 & telescopes, truncated comaximality, coefficient maps, finite-support evaluation, Gaussian derivation, gcd identities & 597 & stdlib\\
10 & divided differences, Newton reconstruction, digit indicators and lifts, telescopes, rational differences & 171{,}395 & SymPy\\
11 & finite division identities, positive quotient supports, augmentation and conjugation identities, Gaussian derivation & 1{,}858 & stdlib\\
\bottomrule
\end{tabular}
\end{center}

These checks catch sign, index and coefficient errors in finite special cases. None of them
verifies an infinite support, a class-size argument, a cardinal computation, the existence of a
maximal ideal, an $\Ext$ computation, a primality statement, or novelty; those rest on the written
proofs (all sources). The merge results were not machine-checked.

\subsection{Formalization routes}\label{osq:sub:formalization}
No statement of this report is formalized, and the repository has no omnific module. The sources
propose compatible routes, recorded here as proposals only. First formalize the scaled-field collision
lemma and the separation lemma, which need no surreals (06, 04, 07). Then the countable telescope in a
set-sized Hahn workspace, with its two support inequalities, proved coefficientwise and not by
convergence of partial sums (04, 06, 07, 08, 10). The set-sized criteria
(\cref{osq:thm:criterion,osq:thm:coinitial}) avoid proper-class syntax and are good first targets (08). The full-class theorem needs an explicit universe interface in which the target and the collision
index are small and the exponent class supplies smaller scales below every small support; a
universe-polymorphic theorem about an arbitrary field does not supply them (04, 08, 10, 11). For the
constructions, the finite-coordinate invariant of 08 must be preserved through field generation,
real closure and truncation, and the module bound at a cardinal must use cyclic quotients or the
direct collision, not an endomorphism ring (08, 06). The homological part comes last, over set-sized
rings only (06). For 07's quotient, the support characterization of $\Oz[\omega^{-a}]$ and the exact
binomial kernel should be kept separate from the finite Chinese-remainder algebra. The repository's
small-normal-form modules (\texttt{SmallNormalForm.lean}, \texttt{SmallNormalFormTruncation.lean},
\texttt{SmallNormalFormEvaluation.lean}, \texttt{SignSequenceValuationBalls.lean}, see
\cite{RepoBridge}) are possible prerequisites (04); they are not a certificate for any result here.

\begingroup
\raggedright
\small
\begin{thebibliography}{99}

\bibitem{Abramov}
D.~Abramov, \emph{Conway-refinement}, public Lean development,
\url{https://github.com/gaearon/conway-refinement}, inspected by 07 on 22 September 2026 for its
proof claim and stated review caveats only; nothing from it is used.

\bibitem{BagayokoHoeven}
V.~Bagayoko and J.~van der Hoeven, Surreal substructures, \emph{Fundamenta Mathematicae}
\textbf{266} (2024), 25--96, \href{https://doi.org/10.4064/fm231020-23-2}{doi:10.4064/fm231020-23-2};
arXiv:2305.02001.

\bibitem{BGR}
V.~Barucci, S.~Gabelli and M.~Roitman, Complete integral closure and strongly divisorial prime
ideals, arXiv:math/0302223, 2003. Cited for terminology only.

\bibitem{Berarducci}
A.~Berarducci, \emph{Surreal numbers, exponentiation and derivations}, arXiv:2008.06878, 2020
(Sections 3 and 8--10 for cuts, normal forms and the Hahn support calculus).

\bibitem{BM}
A.~Berarducci and V.~Mantova, Surreal numbers, derivations and transseries, arXiv:1503.00315, 2015.

\bibitem{CC}
P.-J.~Cahen and J.-L.~Chabert, \emph{Integer-Valued Polynomials}, Mathematical Surveys and
Monographs \textbf{48}, American Mathematical Society, 1997.

\bibitem{CGG}
P.~Cegielski, S.~Grigorieff and I.~Guessarian, Characterizing congruence preserving functions
$\Z/n\Z\to\Z/m\Z$ via rational polynomials, arXiv:1506.00133, 2015.

\bibitem{Conway}
J.~H.~Conway, \emph{On Numbers and Games}, Academic Press, 1976; second edition, A K Peters, 2001.

\bibitem{ElliottBinomial}
J.~Elliott, Binomial rings, integer-valued polynomials, and $\lambda$-rings, \emph{Journal of Pure
and Applied Algebra} \textbf{207} (2006), 165--185,
\href{https://doi.org/10.1016/j.jpaa.2005.09.003}{doi:10.1016/j.jpaa.2005.09.003}.

\bibitem{MO188430}
J.~Elliott, Transcendence degree of the surreals over the subfield generated by the ordinals,
MathOverflow question 188430, posted 30 November 2014,
\url{https://mathoverflow.net/questions/188430}. Its final addition asks the quotient question
answered in \cref{osq:thm:ordinal}; the answer by E.~Wofsey concerns transcendence degree.

\bibitem{Gonshor}
H.~Gonshor, \emph{An Introduction to the Theory of Surreal Numbers}, London Mathematical Society
Lecture Note Series \textbf{110}, Cambridge University Press, 1986,
\href{https://doi.org/10.1017/CBO9780511629143}{doi:10.1017/CBO9780511629143}.

\bibitem{Kuhlmann}
F.-V.~Kuhlmann, Dense subfields of henselian fields, and integer parts, arXiv:1003.5681, 2010.
Context only.

\bibitem{LM}
S.~L'Innocente and V.~Mantova, A factorisation theory for generalised power series and omnific
integers, \emph{Advances in Mathematics} \textbf{442} (2024), article 109513,
\href{https://doi.org/10.1016/j.aim.2024.109513}{doi:10.1016/j.aim.2024.109513}; arXiv:1710.07304v5,
22 January 2024. Cited items: Fact~2.1.1 (closedness of Hahn fields), Remark~3.4.4 (the geometric
product), Fact~8.0.1 (Gonshor's Theorem~8.6) with the example on p.~46, Proposition~8.2.1
(truncation criterion for divisibility), Theorem~B (primality of $\omega^{\sqrt2}+\omega+1$).

\bibitem{Poonen}
B.~Poonen, Maximally complete fields, \emph{L'Enseignement Math\'ematique} \textbf{39} (1993),
87--106 (Corollary~4).

\bibitem{StacksFlat}
The Stacks Project Authors, \emph{The Stacks Project}, Section 10.39, Tag 00H9, Lemma 10.39.3,
\url{https://stacks.math.columbia.edu/tag/00H9}.

\bibitem{StacksVal}
The Stacks Project Authors, \emph{The Stacks Project}, Section 10.50, Tag 00I8, Lemmas 10.50.3 and
10.50.11, \url{https://stacks.math.columbia.edu/tag/00I8}.

\bibitem{RepoFound}
\emph{Surreal} repository, foundations report,
\texttt{docs/foundations-and-computation/foundations/} (label \nolinkurl{found:eq:omnific}, the
definition $\Oz=\Pi\oplus\Z$, and the fine-topology size obstructions).

\bibitem{RepoE3}
\emph{Surreal} repository, \emph{Euclidean Three-Space over the Surreal Numbers},
\texttt{docs/surreal/euclidean-three-space/}, Part~III (labels \nolinkurl{e3:cut:thm:smallquotient},
\nolinkurl{e3:cut:cor:onlyimage}) and \texttt{10-rotation-quotients-SOURCE\_AUDIT.md}.

\bibitem{RepoTrig}
\emph{Surreal} repository, trigonometry report, \texttt{docs/surcomplex/trigonometry/}
(label \nolinkurl{trigonometry:eq:split}).

\bibitem{RepoFKC}
\emph{Surreal} repository, \emph{The First-$\kappa$ Coefficients},
\texttt{docs/surcomplex/first-kappa-coefficients/} (label \nolinkurl{fkc:thm:main}).

\bibitem{RepoOdg}
\emph{Surreal} repository, sibling report on the Diophantine geometry of the omnific integers,
\texttt{docs/surreal/omnific-diophantine-geometry/}, merged from manuscripts 01, 02 and 05 of the
same batch; written concurrently with this report.

\bibitem{RepoNotation}
\emph{Surreal} repository, \texttt{docs/NOTATION.md}.

\bibitem{RepoBridge}
\emph{Surreal} repository, \texttt{docs/NORMAL\_FORM\_BRIDGE.md}.

\bibitem{RepoGammaZeta}
\emph{Surreal} repository, gamma and zeta functions report,
\texttt{docs/surcomplex/gamma-and-zeta-functions/}.

\bibitem{RepoDiffEq}
\emph{Surreal} repository, differential equations report,
\texttt{docs/surcomplex/differential-equations/} (label \nolinkurl{diff:warn:ct}).

\end{thebibliography}
\endgroup
\end{document}
